\documentclass[12pt]{article}
\usepackage{fullpage}
\usepackage[T1]{fontenc}
\usepackage{amsmath,amssymb,amsthm,mathrsfs}
\usepackage{xurl}
\usepackage[hidelinks,hypertexnames=false]{hyperref}
\newtheorem{theorem}{Theorem}[section]
\newtheorem{proposition}[theorem]{Proposition}
\newtheorem{lemma}[theorem]{Lemma}
\newtheorem{corollary}[theorem]{Corollary}
\newcommand{\Rea}{\operatorname{Re}}
\newcommand{\N}{\operatorname{N}}
\newcommand{\pp}[1]{(#1)_+}
\usepackage{verbatim}
\allowdisplaybreaks[2]

\title{A certified uniform descent for all Dirichlet characters}
\author{}
\date{9 October 2026}
\begin{document}
\maketitle

\section{Problem statement and interpretation}
The requested all-Dirichlet threshold is
\[
 \frac{3499999}{4000000}=0.87499975.
\]
We seek a strict improvement of this incumbent, not merely of \(7/8\).
``Smallest'' is interpreted as the best explicit rational established
by the developed argument, not a claimed absolute optimum among
all possible analytic methods.  The restricted optimality statement
below names its precise exponent program.  No assumed zero-free
theorem, finite-height verification, or workflow approval is used
as a premise.  All essential method-specific reductions used for the
stated theorem are included in this document.

\section{Result, attribution, and scope}\label{sec:result}
The explicit rational bound proved in this document is
\begin{equation}\label{eq:theta}
 \boxed{\theta=\frac{874957019421}{1000000000000}
                   =0.874957019421}.
\end{equation}
% QRH_UNIFORM_BOUND: 874957019421/1000000000000
Its improvement over the stage incumbent is
\[
 \frac{3499999}{4000000}-\theta
       =\frac{42730579}{1000000000000}>0.
\]
The number is not a numerical zero-search bound.

\begin{theorem}\label{thm:main}
For every integer \(q\ge1\) and every Dirichlet character \(\chi\)
modulo \(q\), \(L(s,\chi)\) has no zero in the strict half-plane
\(\Rea s>\theta\), with \(\theta\) exactly as in \eqref{eq:theta}.
The assertion includes \(\zeta\), all real zeros, all heights,
principal characters, and imprimitive characters.  A principal
pole at \(1\) is allowed.  Moreover, every finite-order Hecke
\(L\)-function over \(F=\mathbb Q(\sqrt{-3})\) is zero-free in
this same half-plane, with its principal pole allowed.
\end{theorem}

Write \(\mathcal O=\mathbb Z[\omega]\), \(\omega=e^{2\pi i/3}\),
and \(q_a=\N a\).  We develop the supplied manuscript
\cite{Seed}; its arithmetic construction and much of its dependency
proof are retained, with attribution.  Its two AI workflow reports
are not mathematical premises.  References to \cite{QRH} mean
only revision
\[
 \texttt{fd4aeeb2ee4fc729c18d98444fed42fd0529eeeb}.
\]
We assume no zero-free conclusion from that paper or its companion.
The optional Landau--Siegel manuscript is not used.
Nor is the incumbent treated as an axiom.

The unconditional arithmetic inputs are completed cubic-theta
reflection, a masked quadratic/cubic hybrid norm, an additive Gram
bound, and two different finite moment inductions.  They are proved
or reduced to precisely identified primary formulas in
Appendices~\ref{app:conventions}--\ref{app:plain}.
The physical probe, both exact representations, the buffered
detector, and both moment inductions are written out below,
including the transformed summation sets and their inherited zeros.
Pinned source locations record attribution, not omitted proof steps.
The plain moment is strengthened to \(13/18\le\kappa\le1\);
its positive-slot premise is \(B\le(1+\kappa)/2\) when
\(\kappa<1\), and there is no such premise at \(\kappa=1\).
Here \(B\) is a supremum, not an assumed zero-free constant.

The improvements are joint geometry optimization, retention of the
actual moment parameter below \(3/4\), and an exact certificate for
every point of the continuous exponent domain.
Section~\ref{sec:optimization} gives a lower bound for a precisely
restricted exponent program, including variations with \(M+\ell\ne1\).
It is not a lower bound for all possible developments of the method.
The decimal \(0.8749570194200989\ldots\) appearing there is an
algebraic optimization threshold, not the rational number in
Theorem~\ref{thm:main}.

A public-source check on 9 October 2026 located the independent
heuristic exploration \cite{Exploration} and the Argonaut Math
preprint \cite{Argonaut}, whose announced rational is the same
stage incumbent, \(3499999/4000000\).  Its release is dated
8 October 2026; the TeX in release v0.1.8 is dated 7 October 2026.
Neither publication nor any verification claim is imported as a
mathematical premise.  No priority claim is made.
The actual attachment hashes and the distinction between present
byte checks and older reported downloads appear in
Appendix~\ref{app:checks}.

A further primary-source check located Liu's
\cite[Theorem 21.1]{Liu} all-Hecke and all-Dirichlet statement at
\[
 B_{\rm Liu}=\frac{1507-2\sqrt{921}}{1653}
              =0.874957069799\ldots .
\]
We pin arXiv version 1, submitted on 8 October 2026, rather than
an unversioned mirror.  This is the fixed-\(\kappa_0=3/4\)
threshold discussed in Section~\ref{sec:optimization}.
The present rational satisfies
\[
 \frac{5037}{10^{11}}<B_{\rm Liu}-\theta
                         <\frac{5038}{10^{11}};
\]
the exact certificate verifies these comparisons by squaring
positive rational quantities.  Liu's conclusion and its reported
formalization are not premises here.  The distinction between
joint parameter optimization already present in that literature
and the extended plain-moment interval used here is essential;
no claim of priority is made.

\section{Geometry, constraints, and continuation}\label{sec:setup}
Fix the following rational numbers, in the displayed order:
\begin{equation}\label{eq:new-geometry}
\begin{gathered}
 \ell=\frac{11}{3}-4\theta=\frac{125128941737}{750000000000},
 \qquad \kappa_0=2\theta-1=\frac56-\frac{\ell}{2},\\
 b=-\frac{(3\ell+1)(36\ell^2-75\ell+19)}
                {3(114\ell^2-159\ell-7)},\qquad
 M=1-\ell,\quad l_x=\frac{1-\ell-b}{2},\quad
 l_y=\frac{1-\ell+b}{2},\\
 h=1-l_x+\ell=\frac{1+3\ell+b}{2},\qquad
 z_0=\frac{17}{50},\quad d_{\min}=\frac1{100},
 \quad\zeta=10^{-14}.
\end{gathered}
\end{equation}
In particular \(b=0.123405685424\ldots\) is not fixed at \(1/8\).
The exact identities replacing the old geometry are
\begin{equation}\label{eq:geometry}
 C(s)=s+\frac{l_x}{2}-1+\frac h6=s-\frac23-\frac b6,
 \qquad
 \theta=1-\frac h6+\frac b{12}=\frac{11}{12}-\frac\ell4,
 \qquad C(\theta)=\frac{l_x}{2}+\frac b{12}.
\end{equation}
The coefficient of \(s\) in \(C\) is still one.

Here is a ledger of the real constraints used in the proof.
All inequalities in the right column follow by substitution of
\eqref{eq:new-geometry}; the exact verification is reproduced in
Appendix~\ref{app:certificate-code}.  Strict margins, not decimal
approximations, are used.
\[
\begin{array}{c|c}
 \text{purpose or parameter}&\text{certified range}\\ \hline
 \text{low energy}&1/6<\ell<167/1000<1/5\\
 \text{asymmetry and Gram scale}&123/1000<b<124/1000\\
 \text{rescaled row range}&M-2\ell>49/100\\
 \text{rescaled first window}&l_x-\ell>9/50\\
 \text{rescaled Gram remainder}&l_y-\ell-11b/6>2/25\\
 \text{physical lengths}&35/100<l_x<36/100,\quad47/100<l_y<49/100\\
 \text{row endpoint and prime supply}&81/100<h<82/100,\quad
                                  \ell/(h+\zeta)>1/5\\
 \text{principal local region}&43/50<\theta<3499999/4000000\\
 \text{reference plain parameter}&13/18<37/50<\kappa_0<3/4\\
 \text{detector and amplitudes}&1/50\le\delta\le1,\quad0\le q\le\delta/2\\
 \text{adaptive witness}&1\le t\le3/2,\quad \alpha=5/6 .
\end{array}
\]
The floor \(\delta=1/50\) has no required witness and is counted
separately.  The refined count is needed only for \(d\ge1/2\);
small rows have \(d\le d_{\min}\), and large rows have \(d>h+\zeta\).
Local Euler regions, moment length inequalities, height restrictions,
and the order of all losses are recorded where they are used.
They are not omitted from the optimization.

Let \(\mathcal F\) consist of all primitive finite-order Hecke
characters over \(F\), including the principal character, and put
\begin{equation}\label{eq:B-definition}
 B=\sup\bigl(\{1/2\}\cup
 \{\Rea\rho:\Rea\rho\ge1/2,\ L_F(\rho,\eta)=0,\
                                     \eta\in\mathcal F\}\bigr).
\end{equation}
Poles are excluded; real zeros and every height are included.
Euler convergence gives \(1/2\le B\le1\), not \(B\le11/12\).
Finite-order products stay in this family after primitive induction;
norm twists only translate the variable vertically.

\begin{lemma}[Continuation criterion]\label{lem:criterion}
Suppose \(B>\theta\).  For every target suppose there is a
finite set \(S_\eta\), a holomorphic \(H_\eta\) with
\(|H_\eta-1|\le1/2\) on \(\Rea s>\theta\), and a function
\(J_\eta\), such that for common \(0<\omega<B-\theta\)
and \(\sigma>0\),
\[
 |J_\eta(Z)|\ll_\eta Z^{C(\theta)+\omega},\qquad
 |J_\eta(Z)-f_\eta(Z)|\ll_\eta Z^{C(B)-\sigma},
\]
where
\begin{equation}\label{eq:signal}
 f_\eta(Z)=\frac1{2\pi i}\int_{(2)}
 Z^{C(s)}e^{(s-5/6)^2}
 \frac{H_\eta(s)}{L_F^{S_\eta}(s,\eta)}\,ds.
\end{equation}
Starting scales may depend on the target.  These assumptions are
contradictory.
\end{lemma}
\begin{proof}
Set \(e_*=\min(B-\theta-\omega,\sigma)>0\) and
\(b_*=B-e_*\), so \(\theta<b_*<B\).
The two estimates give \(f_\eta(Z)\ll_\eta Z^{C(b_*)}\) at infinity.
The reciprocal Euler product is bounded on \(\Re s\ge2\).
Shifting \eqref{eq:signal} right to any fixed \(A>2\), using
Gaussian decay, gives \(f_\eta(Z)\ll_{\eta,A}Z^{C(A)}\) near zero.
Consequently
\[
 {\cal M}_\eta(s)=\int_0^\infty f_\eta(Z)Z^{-C(s)}\,\frac{dZ}{Z}
\]
is holomorphic on \(\Re s>b_*\), with locally uniform convergence
also after differentiation.  Under \(Z=e^u\), the defining integral
is an inverse Fourier transform on \(\Re s=2\).
More explicitly its inverse is
\(u\mapsto e^{-C(2)u}f_\eta(e^u)\).
This weighted function and its transform are integrable, by
the endpoint estimates and Gaussian decay, respectively.  Fourier inversion
and then the meromorphic identity theorem give
\[
 {\cal M}_\eta(s)L_F^{S_\eta}(s,\eta)
       =e^{(s-5/6)^2}H_\eta(s)\quad(\Re s>b_*).
\]
The right side never vanishes, so there is no zero of the
\(L\)-function there.  Deleted Euler factors are nonzero for
\(\Re s>0\).  Applying this to every target gives \(B\le b_*<B\).
A principal pole is harmless; no supremum need be attained.
\end{proof}

\paragraph{Hecke induction of characters.}
If a finite-order Hecke character \(\eta\) modulo \(\mathfrak f\)
is induced from a primitive \(\eta^*\) of conductor \(\mathfrak f^*\),
Euler comparison, followed by continuation, gives
\[
 L_F(s,\eta)=L_F(s,\eta^*)
 \prod_{\substack{\mathfrak p\mid\mathfrak f\\
                    \mathfrak p\nmid\mathfrak f^*}}
       (1-\eta^*(\mathfrak p)(\N\mathfrak p)^{-s}).
\]
Every displayed factor is a holomorphic unit on \(\Re s>0\).
Thus a zero-free conclusion for the primitive Hecke family gives
exactly the same conclusion for all finite-order Hecke characters;
a principal pole is preserved, not counted as a zero.

\paragraph{Transfer and principal characters.}
We use meromorphic continuation as in \cite{DLMF,Poonen}.
For a Dirichlet character induced from the primitive \(\chi^*\)
of conductor \(f\mid q\), Euler comparison and continuation give
\begin{equation}\label{eq:imprimitive}
 L(s,\chi)=L(s,\chi^*)\!
 \prod_{\substack{p\mid q\\p\nmid f}}(1-\chi^*(p)p^{-s}).
\end{equation}
Each extra factor is a holomorphic unit on \(\Re s>0\), since
\(|\chi^*(p)p^{-s}|=p^{-\Re s}<1\).  Orders, including zero
multiplicities and poles, agree there.  For the principal character
this is \(\zeta(s)\prod_{p\mid q}(1-p^{-s})\), whose residue at
one is \(\prod_{p\mid q}(1-p^{-1})>0\); \(q=1\) uses the empty product.

For completeness the Hecke-to-Dirichlet implication
\cite[Proposition 11.3]{QRH} also includes a possible cancellation
at one.  Lift a primitive \(\chi\) by
\(\eta(\mathfrak a)=\chi(\N\mathfrak a)\); norm congruences show
this is a finite-order ray character.  Superscripts below delete
primes dividing \(3q\), and their lifts.  With \(\varepsilon_3\) the nonprincipal character modulo three,
split and inert Euler factors give
\[
 L_F^S(s,\eta)=L^{3q}(s,\chi)L^{3q}(s,\chi\varepsilon_3).
\]
If \(\eta^*\) induces the norm lift, then
\(L_F^S(s,\eta)=L_F^S(s,\eta^*)\): all its possibly additional
local factors have already been deleted.  Enlarging \(S\) if
necessary, or applying the same holomorphic-unit argument as
in \eqref{eq:imprimitive}, proves this equality and preserves all
zero orders.
At an inert prime the identity is
\((1-a)^{-1}(1+a)^{-1}=(1-a^2)^{-1}\).
Away from one neither factor has a pole to cancel a zero.
At one the only nonprincipal primitive target with principal
companion is \(\chi=\varepsilon_3\), and
\[
 L(1,\varepsilon_3)
 =\sum_{n\ge0}\frac1{(3n+1)(3n+2)}>0 .
\]
Pairing the convergent Dirichlet series justifies this absolutely
convergent expression.  Thus a Hecke-family conclusion transfers
to every Dirichlet character, preserving strictness, all heights,
and real exceptional zeros.

\section{The common physical probe and its two exact representations}\label{sec:common-probe}
The following construction and finite calculations are adapted from
\cite[Sections 6, 7, 12, and 16]{QRH}.  Every defining sum is
included here.  Arithmetic notation, the finite phases \(\mathcal R,G\),
and the central completion are fixed in
Appendices~\ref{app:conventions}--\ref{app:reflection}.
A one-argument \(G(a)\) is a finite ray function; the later
two-argument \(G(a,k)\) is a normalized Gauss sum.
\subsection{Fixed data and the base sum}
Take \(T\) to be the ray class group modulo \((36)\).
On the unique primary representatives, the functions
\(\mathcal R\) and \(G\) in Appendix~\ref{app:conventions}
factor through this group: their displayed residue formulas depend
only on the primary generator modulo \(4\), with the fixed cubic
supplement at \(2\).  The primary convention is fixed modulo \(3\).
Thus \(T\) is independent of the target.  The target-dependent
fixed-numerator characters below need not factor through $T$.  For a
finite-order Hecke character $\eta$, write
\[
 \Theta=\langle\eta,\widehat T\rangle.
\]
Use one excluded set $S$ for the finite family
$\{\eta\theta:\theta\in\widehat T\}$.  It contains the primes over
$6$, the prime supports of the defining moduli of the fixed
zero-extended presentations of $\eta$ and the characters in
$\widehat T$, and every prime of norm at most a fixed $P_0\ge10^{10}$.
We allow $P_0$ to be enlarged in terms of the later chosen buffer $e$;
that choice is made before the target data and every scale.  These
full presentations belong to the fixed arithmetic data.  They are fixed
before $X,Y,Z$ and the varying row and prime parameters are chosen.
Enlarging $S$ will not change $T$.

The finite transform should exclude frequencies meeting $S$, including
zero, and should leave coefficients that factor prime by prime. The
following character and phase corrections arrange these two properties.

Let $b_*$ generate the squarefree product of the primes in $S$.
Choose a residue character $\xi$ modulo $b_*$ whose restriction to
each prime factor is nonprincipal and whose order divides six.  Such
a choice exists: at a prime of odd residue characteristic one may
use the quadratic character, and the residue field at the prime over
$2$ has order four and has a character of order three.  The Chinese
remainder theorem makes $\xi$ primitive modulo $b_*$.  All residue
characters below are extended by zero on nonunits.  Put
\[
 \begin{gathered}
 g_\psi(c,k)=\sum_{d\bmod c}\psi(d)e(kd/c),\qquad
 \tau=q_{b_*}^{-1/2}g_\xi(b_*,1),\qquad
 \Xi(a)=\xi(a)\chi_a(b_*),\\
 \Psi_{m,s,\eta}(a)=
       \eta(a)\Xi(a)^{-1}\mathcal R(a,s)\chi_a(m)
       \qquad ((a,S)=1).
 \end{gathered}
\]
Finite character orthogonality at each prime gives $|\tau|=1$.
By the fixed-numerator ray argument in Appendix~\ref{app:conventions}, the fixed-numerator symbol
$a\mapsto\chi_a(b_*)$ is a finite ray character whose conductor is
supported on $S$.  Thus $\Xi$ is a fixed finite character for this
target, although it need not factor through $T$.

The primitive character $\xi$ at the primes in $S$ will force the
Poisson frequency to be prime to $S$. The accompanying normalizations
cancel the unit factors introduced by the auxiliary modulus $b_*$.
For the completed index $A=cn^3$, the factor $\mathcal R(A,s)$
cancels the cross-prime reciprocity phases between $A$ and $s$, while
$\overline{G(A)}$ cancels the remaining pair phases within $A$.
We verify these cancellations coefficientwise in
Section~\ref{sec:local-euler}.

Define the corrected row at spectral parameter $t+1/2$ by
\[
 T_{m,s,\eta}(t)=
 \sum_{\substack{c\ {\rm sf}\\n}}
 \gamma_2(c)\overline{\alpha(cn^3)}\Psi_{m,s,\eta}(cn^3)
 \overline{G(cn^3)}q_c^{-1/2-t}q_n^{-1-3t}.
\]
Here and below the ideal variables avoid $S$.  In particular every
factor in $\Xi(a)^{-1}$ is evaluated on a unit.  The value
$G(cn^3)$ is the finite-ray extension in
Equation~\eqref{eq:G-quadratic-refinement}, also when $cn^3$ is not
squarefree.

Fix nonnegative, nonzero functions $W_{\rm b,0},W_{\rm b,1}\in C_c^\infty(0,\infty)$
and put $\Phi(t)=e^{t^2}$.  Define
\begin{equation}\label{eq:general-probe}
\begin{aligned}
 I_\eta(X,Y,Z)={}&\frac1Y\sum_s
 \frac{W_{\rm b,1}(q_s/Y)\chi_s(b_*)}
      {\tau\xi(s)\sqrt{q_{b_*}q_sX}}\\
 &\quad\cdot\sum_{m\in\mathcal O}
       \xi(m)q_s^{-1/2}g_{\chi_s}(s,-m)
       W_{\rm b,0}\!\left(\frac{q_m}{q_{b_*}q_sX}\right)
       \frac1{2\pi i}\int_{(4)}
             Z^t\Phi(t)T_{m,s,\eta}(t)\,dt.
\end{aligned}
\end{equation}
The masks in this definition remove nonunits at $S$.  All sums in
Equation~\eqref{eq:general-probe} are absolutely convergent on the
displayed line; the $m$ and $s$ sums are finite because the weights
are annular.

For a squarefree product \(D\) of good primes, the marked probe
\(I_{\eta;D}\) is exactly \eqref{eq:general-probe} with
\(1_{D\mid cn^3}\) inserted in \(T_{m,s,\eta}(t)\).
For every multiplicative zero-extended \(\psi\), put
\begin{equation}\label{own:completed-marked-series}
 T_D(s_0,\psi)=\sum_{\substack{c\ {\rm sf}\\n}}
 \gamma_2(c)\overline{\alpha(cn^3)}\psi(cn^3)
 1_{D\mid cn^3}q_c^{-s_0}q_n^{-3s_0+1/2}.
\end{equation}
Here \(T=T_1\); \(c,n\) may share primes.  This agrees with the
completion in Appendix~\ref{app:reflection} by direct multiplication
of its two series.
The fixed finite Fourier expansion of the correction is
\[
 \overline{G(A)}=\sum_{\theta\in\widehat T}a_\theta\theta(A),
 \qquad
 a_\theta=\frac1{|T|}\sum_{v\in T}
                 \overline{G(v)}\,\overline{\theta(v)}.
\]
Parseval and Cauchy--Schwarz give
$\sum_\theta|a_\theta|\le |T|^{1/2}$.  For a ray class
$\sigma\in T$ choose a representative $s_\sigma$ and put
\[
 \begin{aligned}
 \nu_\sigma(a)&=\eta(a)\Xi(a)^{-1}\mathcal R(a,s_\sigma)
       \quad ((a,S)=1),\\
 B_{m,\sigma}(Z)&=\sum_{\theta\in\widehat T}a_\theta
 \frac1{2\pi i}\int_{(4)}Z^t\Phi(t)
       T(t+1/2,\nu_\sigma\theta\chi_\bullet(m))\,dt.
 \end{aligned}
\]
Extend $\nu_\sigma$ by zero away from the primary elements prime to $S$, without
evaluating $\Xi(a)^{-1}$ there.  The function $T$ here is the explicit series
\eqref{own:completed-marked-series} with $D=1$, with all its zeros.
Since $\mathcal R(A,s)=\mathcal R(A,s_\sigma)$ for $s\in\sigma$,
the displayed Mellin integral of the corrected row equals
$B_{m,\sigma}(Z)$ throughout that class.

Write $Q=q_{b_*}XY$.  Choose a smooth compactly supported function
$\Omega$ on $(0,\infty)$ equal to one on every value of $q_m/Q$
that can occur on the support of
$W_{\rm b,1}(q_s/Y)W_{\rm b,0}(q_m/(q_{b_*}q_sX))$.  For $v\in\mathbb R$, set
\[
 A_{m,\sigma,v}(Y)=\frac1Y\sum_{s\in\sigma}
 \frac{W_{\rm b,1}(q_s/Y)\chi_s(b_*)}{\tau\xi(s)}
 (q_s/Y)^{-1/2+iv}q_s^{-1/2}g_{\chi_s}(s,-m).
\]
Our Mellin convention is
\[
 \widehat W_{\rm b,0}(iv)=\int_0^\infty W_{\rm b,0}(r)r^{iv}\,\frac{dr}{r},
 \qquad
 \frac{q_m}{q_{b_*}q_sX}=\frac{q_m/Q}{q_s/Y}.
\]
Mellin inversion gives the exact identity
\begin{equation}\label{eq:low-separated}
 I_\eta(X,Y,Z)=\frac{Q^{-1/2}}{2\pi}
 \sum_{\sigma\in T}\int_{\mathbb R}\widehat W_{\rm b,0}(iv)
 \sum_{m\ne0}\Omega(q_m/Q)\xi(m)(q_m/Q)^{-iv}
                 A_{m,\sigma,v}(Y)B_{m,\sigma}(Z)\,dv.
\end{equation}
The omitted $m=0$ term vanishes by the annular support of $W_{\rm b,0}$.
The factor $(q_m/Q)^{-iv}$ has absolute value one.  In $A$, the
arithmetic coefficient is fixed and bounded independently of $m,v$;
all dependence on $v$ occurs in a norm power of the annular variable.
\subsection{Physical slots and whole-index marks}
Fix positive slot lengths $\ell_1,\ldots,\ell_K$ of total length
$\ell$, and put $P_i=Z^{\ell_i}$.  Each physical slot has a
nonnegative, nonzero smooth annular weight $W_i(q_p/P_i)$.  Its allowed
prime set is
\[
 \mathcal P_i(Z)=\{p\text{ prime}:p\notin S,\ p\in1_T,
                         \ q_p/P_i\in\operatorname{supp}W_i\}.
\]
The underlying window sets
$\{p\text{ prime}:q_p/P_i\in\operatorname{supp}W_i\}$ for distinct slots
are required to be disjoint before imposing the ray and $S$ restrictions.
In particular, the sets $\mathcal P_i(Z)$ are disjoint, and every tuple in
$\prod_i\mathcal P_i(Z)$ consists of distinct primes.
The number and lengths of the slots will be chosen later, but they are
fixed independently of $Z$.  A sum over $p\in1_T$ for slot $i$ below
always retains this same exclusion and annular support.

For a squarefree product $D$ of slot primes, let
$I_{\eta;D}(X,Y,Z)$ denote Equation~\eqref{eq:general-probe} with
the indicator $1_{D\mid cn^3}$ inserted in its completed row.
Thus $I_{\eta;1}=I_\eta$.  For a tuple
$(p_1,\ldots,p_K)\in\prod_i\mathcal P_i(Z)$ and
$J\subseteq\{1,\ldots,K\}$ write $p_J=\prod_{i\in J}p_i$, with
$p_\varnothing=1$.  The subset $J$ indexes the rescaled slots and $J^c$
the marked slots; this subset notation is distinct from the normalized
probe $J_\eta$ defined in Section~\ref{sec:closure}.  Define the modified probe by the finite identity
\begin{equation}\label{eq:probe}
\begin{split}
 I_{\eta,\mathrm{modified}}(Z)
 =\sum_{(p_i)\in\prod_i\mathcal P_i(Z)}\prod_{i=1}^K W_i(q_{p_i}/P_i)
 \sum_{J\subseteq\{1,\ldots,K\}}&(-1)^{|J|}q_{p_J}^{-3/2}
          \overline{\eta(p_{J^c})}\\
 &\cdot I_{\eta;p_{J^c}}
       (X/q_{p_J},Y/q_{p_J},Zq_{p_{J^c}}).
\end{split}
\end{equation}
Equivalently, at each prime $p$ the operation is the marked term
$\overline{\eta(p)}I_{\eta;p}(X,Y,Zq_p)$ minus the rescaled term
$q_p^{-3/2}I_\eta(X/q_p,Y/q_p,Z)$.  Formula
\eqref{eq:probe} specifies their composition:
every slot window stays at its original scale $P_i$, including when
another slot changes the completed scale.  On $1_T$ one has
$\theta(p)=1$ for every $\theta\in\widehat T$.  Thus marking commutes
with the Fourier decomposition of $\overline G$ and preserves the
same excluded set and zero masks in every summand.
For fixed $J$ and a fixed rescaled tuple $p_J$, put
$X'=X/q_{p_J}$, $Y'=Y/q_{p_J}$, and let
\[
\begin{split}
 B^J_{m,\sigma}(Z)=
 \sum_{p_i\in\mathcal P_i(Z),\ i\notin J}
 \prod_{i\notin J}\overline{\eta(p_i)}W_i(q_{p_i}/P_i)
 \sum_{\theta\in\widehat T}a_\theta\frac1{2\pi i}
 \int_{(4)} &(Zq_{p_{J^c}})^t\Phi(t)\\
 &\cdot T_{p_{J^c}}(t+1/2,
             \nu_\sigma\theta\chi_\bullet(m))\,dt.
\end{split}
\]
Thus $B^J$ is exactly the completed-row factor for the $J$ summand
of Equation~\eqref{eq:probe} after the
marked slots have been summed.  The fixed tuple $p_J$ affects
only $X',Y'$ in its low separation.

Writing \(Q=q_{b_*}X'Y'\), the contribution of this rescaled
tuple after summing the marked slots is exactly
\begin{equation}\label{own:marked-low-identity}
\begin{split}
 &\frac{(-1)^{|J|}q_{p_J}^{-3/2}
               \prod_{i\in J}W_i(q_{p_i}/P_i)}
      {2\pi\sqrt Q}
 \sum_{\sigma\in T}\int_{\mathbb R}\widehat W_{\rm b,0}(iv)\\
 &\qquad{}\times
 \sum_{m\ne0}\Omega(q_m/Q)\xi(m)(q_m/Q)^{-iv}
       A_{m,\sigma,v}(Y')B^J_{m,\sigma}(Z)\,dv .
\end{split}
\end{equation}
There is no additional mask, spectral factor, or row-dependent
slot coefficient.  This follows by substituting the defining sum,
expanding \(\overline G\), and applying the same Mellin inversion
as in \eqref{eq:low-separated}; all exchanges begin in absolute
convergence.  Summing \eqref{own:marked-low-identity} over the
rescaled tuples and all \(J\) is the exact low representation of
\eqref{eq:probe}.
\subsection{Poisson on the very same probe}

The correction $\overline{G(A)}$ is independent of $m$, and the
varying character in the completed row is precisely $\chi_A(m)$,
including its zeros.  Consequently Poisson summation uses the full
modulus $sb_*A$, even if $s$ and $A$ share primes.  Expand the
$s$-Gauss sum and first sum the lifts of a residue modulo $b_*A$.
The lifts impose $H\equiv b_*Ad\pmod s$ and give the coefficient
\begin{equation}\label{eq:local-F}
 F(s,A,H)=\sum_{\substack{d\bmod s\\H\equiv b_*Ad\pmod s}}
 \chi_s(d)g_{\xi\chi_A}\!\left(b_*A,\frac{H-b_*Ad}{s}\right).
\end{equation}
Indeed, finite Fourier transformation modulo \(sb_*A\) sends
\[
 m\longmapsto \xi(m)\chi_A(m)q_s^{-1/2}g_{\chi_s}(s,-m)
\]
to \(H\longmapsto q_s^{1/2}F(s,A,H)\).
The Poisson prefactor for the row scale
$q_{b_*}q_sX$ is $X/q_A$.  After the outer square-root normalization in
Equation~\eqref{eq:general-probe}, leaving its separate $Y^{-1}$ outside,
their product is
$X^{1/2}/(q_{b_*}^{1/2}q_A)$.  This also verifies all factors of
$q_s$ in the transformation.

At the primes of $b_*$ the primitive Gauss sum vanishes unless
$(H,S)=1$.  It therefore also removes $H=0$.  Every remaining
element has a unique expression
$H=ua^6$, where $u$ is a sixth-power-free element, including its
unit factor, and $a$ is the primary generator of an ideal.  Write
$\sum_u^{(6)}$ for the sum over such nonzero $u$ with $(u,S)=1$.

Let $\widetilde W_{\rm b,0}(|y|^2)$ be the planar Fourier transform of
$W_{\rm b,0}(|y|^2)$ for the self-dual measure and character fixed in
Appendix~\ref{app:conventions}.  The radial Fourier argument after
Poisson is $Xq_H/q_A$.  Put
\[
 m_0(z)=\int_0^\infty\widetilde W_{\rm b,0}(r)r^{z-1}\,dr,
 \qquad
 m_1(w)=\int_0^\infty W_{\rm b,1}(r)r^{w-1}\,dr.
\]
Smoothness at zero and Schwartz decay show that $m_0$ is holomorphic
for $\Re z>0$.  On every compact positive real strip it has
arbitrary polynomial decay in $\Im z$: every derivative of
$e^{(\Re z)v}\widetilde W_{\rm b,0}(e^v)$ is integrable in $v$, uniformly
on that strip, so repeated integration by parts applies.  The same
argument gives arbitrary polynomial decay for $m_1$ on
every fixed real strip.

We will need the strict positivity
\begin{equation}\label{eq:radial-mellin-positive}
 m_0(1/6)>0.
\end{equation}
Here is a proof that does not require the Fourier transform itself
to be nonnegative.  For $0<\sigma<1$,
\[
 |y|^{2\sigma-2}=\frac1{\Gamma(1-\sigma)}
             \int_0^\infty t^{-\sigma}e^{-t|y|^2}\,dt.
\]
The pairing of the left side with $\widetilde W_{\rm b,0}(|y|^2)$ is
absolutely integrable.  Fubini and Fourier duality express it as
\[
 \frac1{\Gamma(1-\sigma)}\int_0^\infty t^{-\sigma}
 \int_{\mathbb C}W_{\rm b,0}(|x|^2)
       \frac{2\pi}{\sqrt3\,t}
       \exp\!\left(-\frac{4\pi^2|x|^2}{3t}\right)
                  d\mu(x)\,dt>0.
\]
The inner integral is strictly positive for every $t>0$, since
$W_{\rm b,0}$ is nonnegative and nonzero.  Absolute integrability follows
either from the displayed representation on the Fourier side or
from the annular support of $W_{\rm b,0}$ on this side.  Polar integration
identifies the original pairing with $(2\pi/\sqrt3)m_0(\sigma)$,
proving Equation~\eqref{eq:radial-mellin-positive}.

Set $x=t+1-z$ after Mellin inversion of the two weights.  For a
nonzero sixth-power-free $u$ with $(u,S)=1$, define
\begin{equation}\label{eq:local-full-series}
 \mathcal F_{\eta,u}(x,w,z)=
 \sum_{\substack{c\ {\rm sf}\\n,s,a}}
 K_\eta(c,n,s,a;u)\,
 q_c^{-1/2-x}q_n^{-1-3x}q_s^{-w}q_a^{-6z},
\end{equation}
where $A=cn^3$ and the coefficient, with all normalizations retained,
is
\begin{equation}\label{eq:local-full-coefficient}
 K_\eta(c,n,s,a;u)=
 \frac{F(s,A,ua^6)\chi_s(b_*)}
 {\sqrt{q_{b_*}}\tau\xi(s)\overline{\xi(u)}}
 \gamma_2(c)\overline{\alpha(A)G(A)}
 \eta(A)\Xi(A)^{-1}\mathcal R(A,s).
\end{equation}
For convenience define the common Mellin weight
\[
 \mathcal W(X,Y,Z;x,w,z)=
 X^{1/2-z}Z^{x+z-1}Y^{w-1}
 \Phi(x+z-1)m_0(z)m_1(w).
\]
For example, the lines $\Re x=3$, $\Re w=3$, $\Re z=2$ correspond
to the original line $\Re t=4$ and lie in absolute convergence.
The complete high identity is
\begin{equation}\label{eq:probe-poisson-identity}
 \begin{split}
 I_\eta(X,Y,Z)=\frac1{(2\pi i)^3}
 \int_{(3)}\int_{(3)}\int_{(2)}
 &\mathcal W(X,Y,Z;x,w,z)\\
 &\cdot\sum_u^{(6)}q_u^{-z}\overline{\xi(u)}
             \mathcal F_{\eta,u}(x,w,z)\,dz\,dw\,dx.
 \end{split}
\end{equation}
One may justify all interchanges on these lines by the elementary
bound $|F(s,A,H)|\le q_sq_{b_*A}$ and by the annular weights.
In particular, Equations~\eqref{eq:low-separated} and
\eqref{eq:probe-poisson-identity} are identities for the same
probe, not estimates for separately chosen test expressions.

\subsection{The scalar Euler identity}
\label{sec:local-euler}

We next evaluate the coefficient in
Equation~\eqref{eq:local-full-coefficient}.  Our objective is to prove
that the high series is a scalar Euler product for the target $\eta$.
The calculation also records separately the terms with positive valuation
of the completed index $A$ at a given prime.  We keep every ray phase until
it has been cancelled or assigned to a local factor.

Fix a nonzero sixth-power-free row $u$ with $(u,S)=1$ and a prime
$p\notin S$.  Set
\begin{equation}\label{eq:local-parameters}
 \begin{gathered}
 Q=q_p,\qquad j=v_p(u)\in\{0,\ldots,5\},\qquad
 \rho=\chi_p(u/p^j),\qquad \omega_p=\chi_p(-1),\\
 a_p=\overline{\alpha(p)}^3\eta(p)^3,\qquad
 b_p=\overline{\alpha(p)}^2\eta(p)^2\overline{\chi_p(4)},
 \qquad b_p^3=a_p^2,\\
 V=Q^{-6z},\qquad R=a_p^2Q^{4-6x-6z},\qquad
 W_{\rm loc}=\rho Q^{-w}.
 \end{gathered}
\end{equation}
The values $\rho,a_p,b_p$ have absolute value one, and
$\omega_p\in\{1,-1\}$.  No ray-class restriction on $p$ is imposed.

\paragraph{The finite scalar and its unit factors.}
Write $G_r=g_{\chi_p^r}(p,1)$, with $G_0=-1$, and
$\gamma_r=Q^{-1/2}G_r$ for $1\le r\le5$.  For $t\ge1$ and
$j'\ge0$, summing the lifts of a residue modulo $p$ gives
\[
 g_{\chi_p^r}(p^t,p^{j'})=
 \begin{cases}
 Q^{t-1}G_r1_{j'=t-1},&6\nmid r,\\
 Q^t1_{j'\ge t}-Q^{t-1}1_{j'\ge t-1},&6\mid r.
 \end{cases}
\]
The second line is the Ramanujan sum for the principal character
extended by zero.  Orthogonality also gives
$G_1G_{-1}=\omega_pQ$.

Let $e_0=v_p(c)\in\{0,1\}$, $l=v_p(n)$, $t=e_0+3l$,
$k=v_p(s)$, and $m=v_p(a)$, so $v_p(H)=j'=j+6m$.
The part of Equation~\eqref{eq:local-F} at $p$, before its unit
factors are restored, is
\begin{equation}\label{eq:local-scalar-definition}
 C_p(t,k,j')=
 \sum_{\substack{d\bmod p^k\\p^{j'}\equiv p^td\pmod{p^k}}}
 \chi_p^k(d)
 g_{\chi_p^t}\!\left(p^t,\frac{p^{j'}-p^td}{p^k}\right).
\end{equation}
At modulus one both factors in this formula mean one.  Its complete
evaluation is
\begin{equation}\label{eq:local-scalar-values}
 C_p(t,k,j')=
 \begin{cases}
 1,&t=k=0,\\
 1_{j'=0},&t=0,\ k>0,\\
 g_{\chi_p^t}(p^t,p^{j'}),&t>0,\ k=0,\\
 \omega_p^{-1}G_1g_{\chi_p^{t-1}}(p^t,p^{j'-1}),
                      &t>0,\ k=1,\ j'\ge1,\\
 0,&\text{otherwise}.
 \end{cases}
\end{equation}
For $t=0$ and $k>0$, the congruence fixes $d=p^{j'}$ modulo
$p^k$, whose zero-extended character is nonzero exactly when
$j'=0$.  For the fourth line the congruence first forces $j'\ge1$.
Expanding the inner Gauss sum, the sum over $d\bmod p$ is
\[
 \sum_{d\bmod p}\chi_p(d)e(-vd/p)
       =\omega_p^{-1}G_1\chi_p(v)^{-1}
       \qquad(p\nmid v),
\]
which changes the Gauss character from $\chi_p^t$ to
$\chi_p^{t-1}$.  If $t>0$ and $k\ge2$, the permitted lifts
$d=d_0+p^{k-1}v$ do not change $\chi_p^k(d)$, but their Gauss phase
is $e(-yv/p)$ with the Gauss variable $y$ a unit.  Their sum is
zero.  This proves Equation~\eqref{eq:local-scalar-values}, including
the cases where the Gauss character is principal and zero-extended.

For the original units write $h_p=H/p^{j'}$, $s_p=s/p^k$, and
$A_p^\circ=A/p^t$.  Chinese remaindering and the substitution
$d=h_p(b_*A_p^\circ)^{-1}d'$ give the additional unit factor
\begin{equation}\label{eq:local-unit-factor}
 U_p=\chi_p^{k-t}(h_p)\chi_p^t(s_p)
       \chi_p^{t-k}(A_p^\circ)\chi_p^{t-k}(b_*).
\end{equation}
More explicitly, the outer residue character gives
$\chi_p^k(h_p)\chi_p^{-k}(b_*A_p^\circ)$, while the Chinese
remainder factor and the rescaling of the Gauss argument give
$\chi_p^t(b_*A_p^\circ)\chi_p^{-t}(h_p/s_p)$.  Their product is
Equation~\eqref{eq:local-unit-factor}.  Since
$h_p=(u/p^j)(a/p^m)^6$, one has $\chi_p(h_p)=\rho$.

At the primes of $b_*$, one has
$(H-b_*Ad)/s\equiv H/s\pmod {b_*}$.  The local factor in the
complete Gauss sum is
$\sqrt{q_{b_*}}\tau\xi(A)\overline{\xi(H/s)}$.
Here $\overline{\xi(H/s)}=\overline{\xi(u)}\xi(s)$ because
$\xi(a)^6=1$.  This cancels every displayed $\xi$ factor in
Equation~\eqref{eq:local-full-coefficient}.  The product of the last
factors in Equation~\eqref{eq:local-unit-factor} cancels
$\chi_s(b_*)\chi_A(b_*)^{-1}$.  This cancellation is coefficientwise.
For any fixed enlargement of $S$, the same calculation uses the same
ray group $T$.

\paragraph{Cancellation of the remaining pair phases.}
With $t_p=v_p(A)$ and $k_p=v_p(s)$, reciprocity at each pair of
distinct primes gives
\[
 \mathcal R(A,s)\prod_p
       \chi_p^{t_p}(s_p)\chi_p^{-k_p}(A_p^\circ)
       =\prod_p\mathcal R(p,p)^{t_pk_p}.
\]
Indeed, for $p\ne r$ the two orientations have opposite symbol
exponents, and their quotient is $\mathcal R(p,r)$; this cancels
the corresponding pair in the bicharacter $\mathcal R(A,s)$.
Only its diagonal remains.

Equation~\eqref{eq:signal-checked} gives
$\gamma_2(c)=\mu(c)\alpha(c)G(c)/\gamma_1(c)$.
The Chinese remainder formula
\[
 \gamma_1(c)=\prod_{p\mid c}\gamma_1(p)
             \prod_{p<r,\ p,r\mid c}\chi_p(r)\chi_r(p)
\]
removes the squarefree part of the remaining pair product.  If
$e_p=v_p(c)$ and $l_p=v_p(n)$, its exponent at a pair $p<r$ is
$t_pt_r-e_pe_r=3(e_pl_r+e_rl_p+3l_pl_r)$.  The cube of
$\chi_p(r)\chi_r(p)$ is $\mathcal R(p,r)$, whose square is one.
Thus the pair product left after this division is
\[
 \prod_{p<r}\mathcal R(p,r)^{e_pl_r+e_rl_p+l_pl_r}
 =\mathcal R(c,n)G(n^3)
   \prod_p\overline{G(p^3)}^{\,l_p}
            \mathcal R(p,p)^{-e_pl_p-l_p(l_p-1)/2}.
\]
The equality follows by expanding the quadratic refinement
Equation~\eqref{eq:G-quadratic-refinement} on the prime factors of
$n^3$.  The same refinement gives
$G(c)\mathcal R(c,n)G(n^3)=G(cn^3)$, which is cancelled by the
inserted $\overline{G(cn^3)}$.  Finally
Equation~\eqref{eq:fixed-gauss-phases} gives
$G(p^3)=\gamma_3(p)$ and $\mathcal R(p,p)=\omega_p$.
These identities account for all pair phases, with no restriction
on the ray class of $p$.

The coefficient in Equation~\eqref{eq:local-full-coefficient}
has therefore separated into prime factors.  In absolute convergence,
its factor at $p$ is the series
\begin{equation}\label{eq:local-P}
\begin{aligned}
 P_p={}&\sum_{\substack{e_0=0,1\\l,k,m\ge0}}
 (-1)^{e_0}\gamma_1^{-e_0}\eta(p)^{e_0}
 (a_p\overline{\gamma_3})^l\rho^{k-t}
 \omega_p^{tk-e_0l-l(l-1)/2}C_p(t,k,j+6m)\\
 &\hspace{33mm}\cdot
 Q^{-(x+1/2)e_0-(1+3x)l-wk-6zm},\qquad t=e_0+3l.
\end{aligned}
\end{equation}
Let $P_p^*$ be the same sum restricted to $t>0$, equivalently to
positive valuation of the completed index $A$.

\begin{lemma}[The complete local identity]\label{lem:local-euler}
Fix a sixth-power-free row $u$ prime to $S$ and a good prime $p$.
The complete local factors are
\begin{equation}\label{old-eq:5.13a}
\begin{aligned}
 P_p^*&=\frac1{1-R}\left\{
 \frac{R(1-Q^{-1})-\eta(p)(Q-1)
 Q^{-x-w}V^{1_{j\le1}}}{1-V}+J_j\right\},\\
 P_p&=\frac1{1-V}+1_{j=0}\frac{W_{\rm loc}}{1-W_{\rm loc}}+P_p^*,
\end{aligned}
\end{equation}
where the six values of $J_j$ are
\[
\begin{array}{c|l}
 j&J_j\\\hline
 0&-\eta(p)\rho^{-1}Q^{-x}+W_{\rm loc}R\\
 1&\eta(p)Q^{-x-w}\\
 2&a_p\rho^{-3}Q^{3/2-3x}\\
 3&-\eta(p)b_p\rho^{-2}Q^{2-3x-w}
       +b_p^2\rho^{-4}Q^{2-4x}\\
 4&-\eta(p)a_p\rho^{-3}Q^{5/2-4x-w}\\
 5&-a_p^2Q^{3-6x}.
\end{array}
\]
Put $W=\chi_p(u)Q^{-w}$ and
$D=\eta(p)\overline{\chi_p(u)}Q^{-x}$, with the original zero
extension at $p\mid u$, and define
\begin{equation}\label{eq:local-H-definition}
 H_p=P_p\frac{(1-V)(1-W)}{1-D},\qquad
 \mathcal H_{\eta,u}=\prod_{p\notin S}H_p.
\end{equation}
Then the series in Equation~\eqref{eq:local-full-series} has the
scalar factorization
\begin{equation}\label{eq:probe-high}
 \mathcal F_{\eta,u}(x,w,z)=
 \frac{\zeta_F^S(6z)L^S(w,\chi_\bullet(u))}
 {L^S(x,\eta\overline{\chi_\bullet(u)})}
 \mathcal H_{\eta,u}(x,w,z).
\end{equation}
Write $x_r=\Re x$, $w_r=\Re w$, and $z_r=\Re z$.  For every
fixed $\epsilon_0>0$, the correction product converges normally
on a neighborhood of every point in each of the following regions,
and hence defines a holomorphic function there:
\begin{align}
 x_r\ge51/100,\quad z_r\ge17/50,\quad w_r\ge-1/100,\quad
 x_r+w_r\ge1+\epsilon_0;\label{eq:local-region-one}\\
 x_r\ge7/8,\quad z_r\ge33/200,\quad w_r\ge19/20.
 \label{eq:local-region-two}
\end{align}
Uniformly in imaginary parts and unit phases,
$\mathcal H_{\eta,u}\ll_\epsilon q_u^\epsilon$, with the constant
also depending on $\epsilon_0$ in the first region.  For $u=1$ in the
second region,
$\mathcal H_{\eta,1}=1+O(P_0^{-c_H})$ with, for example,
$c_H=4/5$.  Equation~\eqref{eq:probe-high} begins in absolute
convergence and supplies a meromorphic continuation through these
regions.
\end{lemma}

The first region will contain the buffered contours for nonprincipal
rows, including the reflected numerator line. The second contains the
principal residue point $w=1$, $z=1/6$ and the contours used to reach it.

\begin{proof}
There are four families with $t>0$:
$(e_0,l)=(0,2r+2),(1,2r),(0,2r+1),(1,2r+1)$, where $r\ge0$.
In each family, increasing $r$ by one and the first permitted $m$
by one multiplies the summand of Equation~\eqref{eq:local-P} by
$R$.  In fact the phase ratio is $a_p^2$: the remaining ratio is
one because $\gamma_3^2=\omega_p$, $\omega_p^2=1$, and $\rho^6=1$.
Increasing $m$ further multiplies a summand by $V$.

For transparency, the following table lists every nonzero summand
after division by $R^r$.  In its first row $j$ is arbitrary; all
other conditions are as displayed.  Every omitted case is zero by
Equation~\eqref{eq:local-scalar-values}.
\begin{equation}\label{eq:local-family-table}
\begin{array}{c|c|l|l}
 (e_0,l)&k&\text{condition on }(j,m)&R^{-r}\text{ times summand}\\\hline
 (0,2r+2)&0&m\ge r+1&R(1-Q^{-1})V^{m-r-1}\\
 (0,2r+2)&0&j=5,\ m=r&-a_p^2Q^{3-6x}\\
 (0,2r+2)&1&j=0,\ m=r+1&W_{\rm loc}R\\
 (1,2r)&0&j=0,\ m=r&-\eta(p)\rho^{-1}Q^{-x}\\
 (1,2r)&1&j=1,\ m=r&\eta(p)Q^{-x-w}\\
 (1,2r)&1&m\ge r+1_{j\le1}&-\eta(p)(Q-1)Q^{-x-w}V^{m-r}\\
 (0,2r+1)&0&j=2,\ m=r&a_p\rho^{-3}Q^{3/2-3x}\\
 (0,2r+1)&1&j=3,\ m=r&-\eta(p)b_p\rho^{-2}Q^{2-3x-w}\\
 (1,2r+1)&0&j=3,\ m=r&b_p^2\rho^{-4}Q^{2-4x}\\
 (1,2r+1)&1&j=4,\ m=r&-\eta(p)a_p\rho^{-3}Q^{5/2-4x-w}.
\end{array}
\end{equation}
Here is a direct check of its entries.  In the first family $t=6r+6$.
For $k=0$, the principal Gauss lift is
$Q^t(1-Q^{-1})$ for $m\ge r+1$; its extra boundary value at
$(j,m)=(5,r)$ is $-Q^{t-1}$.  For $k=1$, the product
$\omega_p^{-1}G_1G_{-1}=Q$ allows only $(j,m)=(0,r+1)$.
These give the first three rows.  In the second family $t=6r+1$.
For $k=0$ the nonprincipal equality $j'=t-1$ gives the fourth row.
For $k=1$, the principal Gauss lift has its negative boundary at
$j'=t$, giving the fifth row, and its value $Q^{t-1}(Q-1)$ for
$j'>t$, giving the sixth row.  The latter condition is exactly
$m\ge r+1_{j\le1}$.

For the two odd families $t=6r+3$ and $t=6r+4$.  Their $k=0$
lines require $j'=t-1$ and their $k=1$ lines require $j'=t$.
The phase reductions are
\[
 \gamma_1\gamma_2
       =-\alpha(p)\overline{\chi_p(4)}\gamma_3,
 \qquad
 \frac{\gamma_4}{\gamma_1}
       =-\frac{\overline{\alpha(p)}}{G(p)}.
\]
The first is Equation~\eqref{eq:signal-checked} at $p$; the second follows
from the first and $\gamma_2\gamma_4=1$.  Substitution gives the
last four rows, including both $j=3$ terms.  Thus the table covers
all five ramified valuations as well as $j=0$.

Summing $r\ge0$ gives $(1-R)^{-1}$, and the two unbounded $m$
ranges give $(1-V)^{-1}$.  These sums are precisely the formula for
$P_p^*$ in Equation~\eqref{old-eq:5.13a}.  If $t=0$, then
$e_0=l=0$.  The $k=0$ terms give $(1-V)^{-1}$, while the $k>0$
terms require $j=m=0$ and give $W_{\rm loc}/(1-W_{\rm loc})$.
This proves the
formula for $P_p$ as well.

It remains to justify the analytic assertions, especially where
$w_r$ is negative.  Put $\mathcal E_p=P_p^*+D$.  The contribution
from $t=0$ is $(1-V)^{-1}+W/(1-W)$, so an exact simplification gives
\begin{equation}\label{eq:local-H-defect}
 H_p-1=
 \frac{D(V+W-VW)-VW+(1-V)(1-W)\mathcal E_p}{1-D}.
\end{equation}
This expression has no $1-W$ denominator.  At $p\mid u$, where
$D=W=0$, it reduces to $(1-V)\mathcal E_p$.  In both stated regions
$|R|<1$, $|V|<1$, and $|D|<1$, uniformly away from one, so these
are holomorphic local expressions.

For $p\nmid u$ and $\vartheta=(-w_r)_+$, the $j=0$ row gives
\[
 |\mathcal E_p|\ll
 Q^{4-6x_r-6z_r+\vartheta}+Q^{1-x_r-w_r-6z_r}.
\]
In the first region $4-6x_r-6z_r\le-11/10$ and
$\vartheta\le1/100$.  The extra factor $1-W$ in
Equation~\eqref{eq:local-H-defect} therefore leaves this contribution
at most $Q^{-27/25}$.  The term $DW$ is at most
$Q^{-1-\epsilon_0}$; the other terms in that equation are smaller.
For $p\mid u$, the strict second-family term has exponent
$1-x_r-w_r\le-\epsilon_0$.  The closest other exponents in the
table are
\[
 3/2-3x_r\le-3/100,\qquad
 2-3x_r-w_r\le-1/50-\epsilon_0,
\]
and all remaining ones are no larger than $-3/100$.
Consequently, with $\epsilon_H=\min(\epsilon_0,1/50)$,
\[
 H_p-1=O(Q^{-1-\epsilon_H})\quad(p\nmid u),\qquad
 H_p-1=O(Q^{-\epsilon_H})\quad(p\mid u).
\]
In the second region the same comparison gives the stronger bounds
\[
 H_p-1=O(Q^{-363/200})\quad(p\nmid u),\qquad
 H_p-1=O(Q^{-33/40})\quad(p\mid u).
\]
The closest good-prime term here is
$Q^{1-x_r-w_r-6z_r}$, and the closest ramified term is
$Q^{1-x_r-w_r}$.

The sum over primes not dividing $u$ converges normally, and the
finite product over primes of $u$ is $O_\epsilon(q_u^\epsilon)$
by the divisor-product bound.  The estimates are uniform in all
imaginary parts and unit phases.  When $u=1$, the good-prime tail
above $P_0$ is $O(P_0^{-163/200})$ by the ideal count, which implies
the asserted $O(P_0^{-4/5})$ estimate for the product.
Finally extracting $(1-V)^{-1}(1-W)^{-1}(1-D)$ at every prime
gives Equation~\eqref{eq:probe-high}.  The normal convergence of
the remaining product proves its claimed continuation.
\end{proof}
\subsection{The full holomorphic correction}

For a fixed slot prime $p$, restricting the completed index to
$p\mid A$ replaces $P_p$ by $P_p^*$ in the high series.  Changing
$Z$ to $ZQ$ multiplies its Mellin weight by $Q^{x+z-1}$.  The
marked term of Equation~\eqref{eq:probe}
therefore has, at points where $P_p$ is defined and nonzero, multiplier
$\overline{\eta(p)}Q^{x+z-1}P_p^*/P_p$.  The rescaled term has
multiplier
\[
 Q^{-3/2}Q^{-(1/2-z)}Q^{-(w-1)}=Q^{z-w-1}.
\]
Thus, on the same locus, the exact local multiplier of the two-term operation is
\begin{equation}\label{old-eq:5.13}
 Q^{z-1}\mathcal B_p,\qquad
 \mathcal B_p=\overline{\eta(p)}Q^xP_p^*/P_p-Q^{-w}.
\end{equation}
Define the combined local replacement by
\begin{equation}\label{old-eq:5.13c}
 G_p=
 \frac{(\overline{\eta(p)}Q^x-Q^{-w})(1-V)(1-W)P_p^*
       -Q^{-w}(1-VW)}{1-D}.
\end{equation}
At points in the Euler regions where the raw quotient in
Equation~\eqref{old-eq:5.13} is defined, one has $G_p=H_p\mathcal B_p$.
The displayed formula defines
$G_p$ also where that quotient is not defined.  There is neither a $P_p$
nor a $1-W$ denominator on the right.  The bounds $|R|,|V|,|D|<1$ in
Lemma~\ref{lem:local-euler}, together with its formula for $P_p^*$, show
that $G_p$ is holomorphic in both stated Euler regions, including at zeros
of $P_p$ or $H_p$.

At points in the Euler regions where the selected quotients are defined,
write the full slot
multiplier
\[
 \mathcal B_i(u;x,w,z)=\sum_{p\in\mathcal P_i(Z)}
           W_i(q_p/P_i)q_p^{z-1}\mathcal B_p.
\]
The correction used for contour moves is defined without these quotients:
\begin{equation}\label{eq:holomorphic-selected-tuple}
 \mathfrak H_{\eta,u,Z}(x,w,z)=
 \sum_{(p_i)\in\prod_i\mathcal P_i(Z)}
 \prod_{i=1}^K\bigl[W_i(q_{p_i}/P_i)q_{p_i}^{z-1}G_{p_i}\bigr]
 \prod_{\substack{p\notin S\\p\notin\{p_1,\ldots,p_K\}}}H_p.
\end{equation}
This is the full correction after the scalar quotient in
Equation~\eqref{eq:probe-high} has been extracted.  To verify this assertion,
fix an allowed tuple and put $\mathcal P=\{p_1,\ldots,p_K\}$.  On the
absolute starting lines the selected operation replaces $P_p$ by
\begin{equation}\label{eq:selected-local-operation}
 \overline{\eta(p)}Q^{x+z-1}P_p^*-Q^{z-w-1}P_p
 =\frac{1-D}{(1-V)(1-W)}Q^{z-1}G_p.
\end{equation}
The equality follows by substituting
$P_p-P_p^*=(1-V)^{-1}+W/(1-W)$ into
Equation~\eqref{old-eq:5.13c}.  It extracts the same scalar local factor
$(1-V)^{-1}(1-W)^{-1}(1-D)$ as at an unselected prime.  Each selected
prime occurs exactly once, because the slot
supports are disjoint.  The complete factor for this tuple is therefore
the scalar quotient in Equation~\eqref{eq:probe-high} times its summand in
Equation~\eqref{eq:holomorphic-selected-tuple}.  This reasoning involves
no division by $P_p$ or $H_p$.  At points in the Euler regions where all
raw selected quotients are defined, the same finite sum also factors as
$\mathcal H_{\eta,u}\prod_i\mathcal B_i$.

The unselected product in each summand converges normally in both Euler
regions.  More quantitatively, in any fixed subregion of
Equation~\eqref{eq:local-region-one}, put
$\epsilon_H=\min(\epsilon_0,1/50)>0$.  The primewise defect bounds in
Lemma~\ref{lem:local-euler} give, uniformly in the selected tuple,
\begin{equation}\label{eq:unselected-local-product}
 \prod_{\substack{p\notin S\\p\notin\mathcal P}}|H_p|
 \le
 \prod_{\substack{p\notin S\\p\nmid u}}
       (1+Cq_p^{-1-\epsilon_H})
 \prod_{p\mid u}(1+Cq_p^{-\epsilon_H})
 \ll_\epsilon q_u^\epsilon.
\end{equation}
The first positive product converges by the ideal count, and the second
satisfies the divisor-product bound.  Omitting selected factors only removes
factors from these positive majorants.  In the second Euler region the same
argument uses the respective defect exponents $-363/200$ and $-33/40$.
Thus Equation~\eqref{eq:unselected-local-product} holds there as well,
with the corresponding positive majorants.  No nonvanishing of $H_p$ is
asserted by these upper bounds.

For each fixed $Z$, Equation~\eqref{eq:holomorphic-selected-tuple} is a
finite sum of products of holomorphic selected factors and normally
convergent unselected products.  It is therefore holomorphic on a
neighborhood of every point in both Euler regions.  Its all-height majorant follows directly from the same expressions.
Indeed, on a fixed real box in either region, the formula for $P_p^*$ and
the denominators bounded away from zero give $|G_p|\ll Q^B$ for some
fixed $B$, uniformly in all imaginary parts and unit phases.  There are
$O(P_i)$ ideals in a slot and $q_p\asymp P_i$ there.  The finite tuple
sum and Equation~\eqref{eq:unselected-local-product} consequently give
\begin{equation}\label{eq:compensated-all-height-majorant}
 |\mathfrak H_{\eta,u,Z}(x,w,z)|
 \ll_\epsilon q_u^\epsilon\prod_iP_i^{z_r+B}
 \ll_\epsilon q_u^\epsilon Z^{B'}.
\end{equation}
for a fixed $B'$ on that box.  The height exponent in this estimate is zero.  All constants here
are fixed before any later order of integration by parts.

Combining the exact finite operation with
Equations~\eqref{eq:probe-poisson-identity} and \eqref{eq:probe-high}
now gives
\begin{equation}\label{eq:compensated-poisson-identity}
\begin{split}
 I_{\eta,\mathrm{modified}}(Z)
 ={}&\frac1{(2\pi i)^3}\int_{(3)}\int_{(3)}\int_{(2)}
       \mathcal W(X,Y,Z;x,w,z)\\
 &\quad\cdot\sum_u^{(6)}q_u^{-z}\overline{\xi(u)}
       \frac{\zeta_F^S(6z)L^S(w,\chi_\bullet(u))}
       {L^S(x,\eta\overline{\chi_\bullet(u)})}
       \mathfrak H_{\eta,u,Z}(x,w,z)\,dz\,dw\,dx.
\end{split}
\end{equation}
This identity is absolutely convergent on the displayed lines, because for
fixed $Z$ it is obtained from finitely many absolutely convergent rescaled
base-probe identities.  It has no additional outside Euler factor.
This is the exact high representation of \eqref{eq:probe}, with
the same fixed data as its low representation.  It holds at all
positive \(X,Y,Z\); no inequality among the lengths has been used.
Outside this section we rename its Mellin variable \(x\) to \(s\).
All contour moves use the quotient-free correction
\eqref{eq:holomorphic-selected-tuple}, not an assumed globally
holomorphic product of individual quotients.

\section{The low estimate in the optimized geometry}\label{sec:low}
Use the fixed physical slots and exactly the sum \eqref{eq:probe}
of Section~\ref{sec:common-probe}, at \(X=Z^{l_x},Y=Z^{l_y}\).
Its whole-index marks, fixed data and all zero extensions are
those defined there.  The vector of lengths and all windows will
be chosen independently of the target.

The following elementary inequality is the changed part of the
energy calculation.
\begin{lemma}\label{lem:energy}
Let \(0\le d\le\ell\le1/5\), \(M'=1-\ell-2d\),
\(\ell'=\ell-d\), and \(r\ge0\).  Suppose
\[
\begin{gathered}
 A,N,S,G,z\ge0,\quad O,v,t,e\ge-r,\quad|\vartheta|\le r,\\
 H\le M'-O+r,\quad 2A\le O+r,\quad N\le A,\quad z\le\ell',\\
 T=2H+2A+2z-1-\ell'-\vartheta-N-3G,\\
 y=v+3t+e\le T+r,\qquad u=\min\{v,z,(v+z)/3\}.
\end{gathered}
\]
Then
\begin{equation}\label{eq:energy}
 O/2+\max(H,v+t)-S-G+z-u-t-2e/3-(T-y)_+/2
 \le M'+(5\ell+d-1)_+/4+11r.
\end{equation}
\end{lemma}
\begin{proof}
Put \(\Gamma=M'-O-H\ge-r\).  Using \(M'+\ell'-1=-3d\),
\[
 T-(H-3d)=H-M'+2A+2(z-\ell')-N-3G-\vartheta\le3r.
\]
Thus \(v+t=y-2t-e\le H+7r\).
If \(u=z\), dropping nonpositive terms bounds the left side by
\(O/2+H+26r/3\le M'+61r/6\).
Otherwise \(u\ge v/2-r/2\): for \(u=v\) use \(v\ge-r\),
and for \(u=(v+z)/3\le z\) use \(2z\ge v\).
Hence \(u+t+2e/3\ge y/4-17r/12\).
The inequality \(y/4+(T-y)_+/2\ge T/4\) holds for all real
\(T,y\).  Now use the exact identity
\[
 O/2+\Gamma+S+G+T/4
 =\{2M'+2z-1-\ell'+2\Gamma-\vartheta+2A-N+G+4S\}/4.
\]
Its last four terms have nonnegative sum.  The resulting bound is
\(M'+(1+3\ell'-2M')/4+55r/6\).
Since \(1+3\ell'-2M'=5\ell+d-1\), taking the larger branch gives
\eqref{eq:energy}.
\end{proof}

\begin{proposition}[Low bound]\label{prop:low}
For every \(\epsilon>0\),
\[
 |I_{\eta,\mathrm{modified}}(Z)|
       \ll_{\mathrm{fixed},\epsilon}Z^{C(\theta)+\epsilon}.
\]
\end{proposition}
\begin{proof}
Fix \(J\) and its rescaled tuple, and put
\(d=\sum_{i\in J}\ell_i\), \(\ell'=\ell-d\),
\(M'=M-2d\), \(r_J=q_{p_J}/Z^d\).
The fixed supports bound \(r_J\) above and below.  The exact low
identity \eqref{own:marked-low-identity} is a finite sum of Mellin integrals of
\[
 Q^{-1/2}\sum_{m\ne0}\Omega(q_m/Q)\xi(m)(q_m/Q)^{-i\nu}
                     A_{m,\sigma,\nu}(Y')B^J_{m,\sigma}(Z),
 \quad
 Q=q_{b_*}Z^{M'}/r_J^2,\quad Y'=Z^{l_y-d}/r_J.
\]
Here \(\Omega\) is the original compact row cutoff.  \emph{Every
mean square below is over \(0<q_m\ll Q\)}, with its fixed
comparison constant; the displayed bounded row multipliers are
retained.  Also \(P_a=Y'^2/Q=q_{b_*}^{-1}Z^b\).
For \(0\le d\le\ell\),
\[
 M'\ge M-2\ell>49/100,\quad l_x-d>9/50,\quad
 l_y-d-11b/6>2/25.
\]
Thus \(Q,Y',P_a\ge1\) eventually and \(P_a^2/Y'\ll P_a^{1/6}\).
The Gram bound proved in
Appendix~\ref{app:gram} gives
\(\sum|A_m|^2\ll (1+|\nu|)^{J_{\rm Gr}}(Q/Y')P_a^{1/6}Z^\epsilon\).

For the other factor, apply the general reflected-energy estimate
\eqref{eq:full-energy-proved}, not a previously specialized low bound.  On a Gaussian annulus
the joint profile is
\[
 \chi(\log x)V_{\rm G}\!\left(e^kx/
          \prod_{i\notin J}\varrho_i\right)
             \prod_{i\notin J}W_i(\varrho_i),
 \qquad\varrho_i=q_{p_i}/Z^{\ell_i}.
\]
Separate all variables jointly before evaluating rows or slots.
At each Fourier mode slot coefficients are individual norm powers
and cutoffs times \(\overline{\eta(p_i)}\), independent of the row.
Gaussian annular decay gives summable fixed seminorms and negligible
tails outside \(|k|\le r\log Z+O(1)\); see
Appendix~\ref{app:smooth}.  The completed scale is
\(N_*=1+\ell'+\vartheta\), \(|\vartheta|\le r\), with its original
central normalization.

Freeze the powerful and supported parts on the actual row annulus.
In the notation of Appendix~\ref{app:general-energy}, its exponent is the left side
of \eqref{eq:energy}, under
\((A,N,S,G,z,t,e)=(A_0,N_0,S_0,B_0,z_a,\ell_b,e_\lambda)\).
Its support conditions are \(H\le M'-O+r\), \(2A_0\le O+r\),
\(N_0\le A_0\), \(z_a\le\ell'\), and retention
\(v+3\ell_b+e_\lambda\le T+r\).
All fixed logarithmic offsets fit in \(r\) at a sufficiently large
threshold.  Dual support, including shared primes, is unchanged.
The kernel saving is extracted before positivity; the energy input
already counts the powerful parts.  Lemma~\ref{lem:energy}, followed
by choosing \(r\) sufficiently small, therefore gives
\[
 \sum_{0<q_m\ll Q}|B^J_m|^2
       \ll Z^{M'+(5\ell+d-1)_+/4+\epsilon}.
\]
This explains the entire changed use of the general energy theorem.

Cauchy--Schwarz and Mellin integration bound one unscaled tuple by
\[
 Z^{(l_x-d)/2+b/12+(5\ell+d-1)_+/8+O(\epsilon)}.
\]
There are \(O(Z^{d+\epsilon})\) rescaled tuples, weighted by
\(O(Z^{-3d/2})\).  The resulting excess over
\(C(\theta)\) is
\[
 -d+(5\ell+d-1)_+/8\le-7d/8\le0,
\]
because \(\ell\le1/5\).  Allocate the arbitrarily small losses
and sum the finitely many subsets.
\end{proof}

\section{Local factors and the principal comparison}\label{sec:principal}
Use the quotient-free correction
\(\mathfrak H_{\eta,u,Z}\) in \eqref{eq:holomorphic-selected-tuple}
and the exact identity \eqref{eq:compensated-poisson-identity}.
Its selected factors \(G_p\), unselected factors \(H_p\), and
tests \(m_0,m_1\) have all been defined and derived in
Section~\ref{sec:common-probe}.  Its common Mellin weight is
\[
 X^{1/2-z}Z^{s+z-1}Y^{w-1}
 e^{(s+z-1)^2}m_0(z)m_1(w).
\]
The tests decay rapidly on every fixed vertical strip
(\(\Re z>0\)), and \(m_0(1/6),m_1(1)>0\), as proved there.

\begin{lemma}\label{lem:principal-local}
Include all primes of norm at most \(10^{10}\) in \(S\), and
retain the fixed, nonzero, nonnegative slots.  On
\[
 \Re s\ge43/50,\qquad \Re w\ge19/20,\qquad \Re z\ge33/200
 \tag{P}\label{eq:Preg}
\]
the principal correction satisfies, on each fixed bounded
box \(\mathcal K\) of real parts contained in (P),
\[
 |\mathfrak H_{\eta,1,Z}|
       \ll_{\mathrm{fixed},\mathcal K} Z^{\ell\Re z},
\]
uniformly in all imaginary parts.
There are \(H_\eta\), holomorphic on \(\Re s>43/50\), and
\(A_T(Z)\ne0\) eventually, such that
\[
 \begin{gathered}
 |H_\eta-1|\le189/500000<1/2,\qquad |A_T|^{-1}\ll_\epsilon Z^\epsilon,\\
 \mathfrak H_{\eta,1,Z}(s,1,1/6)
 =H_\eta(s)Z^{\ell/6}A_T(Z)(1+O_{\rm fixed}(Z^{-\mu})),
 \end{gathered}
\]
uniformly in height, where \(\mu=(43/50)\min_i\ell_i>0\).
\end{lemma}
\begin{proof}
Here is the required extension below \(7/8\), directly from the
rational factors, without a local division at an unknown zero.
For a good prime of norm \(Q\), write
\(V=Q^{-6z}\), \(R=vQ^{4-6s-6z}\), \(W=\rho Q^{-w}\),
\(D=u\rho^{-1}Q^{-s}\), \(T=u(Q-1)Q^{-s-w}\), with unit phases.
The identities \eqref{old-eq:5.13a} and \eqref{eq:local-H-defect}
give, with \(\mathcal E=P_p^*+D\),
\[
 \begin{aligned}
 \mathcal E&=
 \frac{R((1-Q^{-1})/(1-V)+W-D)-TV/(1-V)}{1-R},\\
 H_p-1&=\frac{D(V+W-VW)-VW+(1-V)(1-W)\mathcal E}{1-D}.
 \end{aligned}
\]
On (P), \(|D|,|V|,|W|<3/4\), \(|R|<1/4\).
The decay exponents of \(R,TV,DV,DW,VW\) are respectively at least
\(43/20,9/5,37/20,181/100,97/50\).
Hence \(|\mathcal E|\le20Q^{-9/5}\) and
\(|H_p-1|\le336Q^{-9/5}\).

The Eisenstein ring is norm-Euclidean (coordinate rounding leaves
norm at most \(3/4\)), so every ideal is principal.
Counting \(a+b\omega\) using
\(a^2-ab+b^2\ge(a^2+b^2)/2\) bounds the ideal count by \(25X\).
Partial summation gives
\(\sum_{q_p>P_0}q_p^{-9/5}\le(225/4)P_0^{-4/5}\).
Thus the sum of the local defects above \(10^{10}\) is at most
\(189/10^6\).  Expanding finite products and using
\(e^t-1\le2t\) for \(t\le1/4\) proves the asserted contraction.
On the larger open region
\(\Re s>17/20,\Re w>47/50,\Re z>4/25\), the same calculation
gives \(336Q^{-7/4}\); this supplies normal convergence on a
neighborhood of (P).

For the principal row \(\rho=1\), formula \eqref{old-eq:5.13c} becomes
\[
 G_p=\frac{(D^{-1}-W)(1-V)(1-W)P_p^*-W(1-VW)}{1-D}.
\]
Its exact cancellation is
\[
 G_p+H_p=(1-W)\left(H_p-1+\frac{V-D}{1-D}\right)
       +\frac{(1-V)(1-W)\mathcal E}{D(1-D)}.
\]
Now \(|\mathcal E/D|\le12Q^{-129/100}+8Q^{-47/50}\);
the last display is at most \(1008Q^{-43/50}\).
Since \(|H_p|\ge1/2\) above the cutoff,
\[
 |G_p/H_p+1|\le2048Q^{-43/50}.
\]
The positive majorant for the unselected product is less than two, \(G_p=O(1)\),
and ideal counting in each annulus proves the tuple bound on
\(\mathcal K\).  The dependence on \(\mathcal K\) includes
\(\sup_{y\in\operatorname{supp}W_i}y^{\Re z-1}\);
there is no claim of a uniform constant as \(\Re z\to+\infty\).

Put \(S_i(Z)=\sum_{p\in\mathcal P_i(Z)}
 W_i(q_p/Z^{\ell_i})q_p^{-5/6}\) and
\(A_T=(-1)^KZ^{-\ell/6}\prod_iS_i\).
The fixed-ray prime theorem
\cite[Lemma 13.1]{QRH}, by partial summation from
\cite[Theorem 1.1]{TZ} for a fixed extension, gives
\[
 S_i(Z)\sim
 \frac{Z^{\ell_i/6}}{|T|\ell_i\log Z}
             \int_0^\infty W_i(y)y^{-5/6}\,dy>0.
\]
This uses a fixed extension, not a moving-conductor error estimate;
a possible exceptional zero can affect the threshold only.
Thus \(|A_T|\asymp_{\rm fixed}(\log Z)^{-K}\).
Disjoint windows give the exact product of the \(K\) selected
sums times \(H_\eta=\prod_{p\notin S}H_p(s,1,1/6)\).
Nonnegative weights make each sum
\(-S_i(1+O(Z^{-43\ell_i/50}))\), proving the result.
\end{proof}

\begin{proposition}\label{prop:principal}
Let \(B>\theta\), \(0<e<10^{-3}\), and let
\(\mathcal P_\eta\) denote the \(u=1\) term of the exact high
identity.  With the preceding data and \(f_\eta\) from
\eqref{eq:signal},
\[
 \frac{\mathcal P_\eta}{c_SA_T}-f_\eta
 \ll_{\rm fixed,e,\epsilon}
 Z^{C(B)+(1+h)e-l_y/20+\epsilon}
 +Z^{C(B)+e-h/600+\epsilon}
 +Z^{C(B)+e-\mu+\epsilon},
\]
where \(c_S=m_1(1)m_0(1/6)
(\operatorname{Res}_{v=1}\zeta_F^S(v))^2/6>0\).
This includes \(B=1\) and principal targets.
\end{proposition}
\begin{proof}
The principal triple integral has scalar quotient
\(\zeta_F^S(6z)\zeta_F^S(w)/L_F^S(s,\eta)\).
Move its initial lines \((3,3,2)\) to
\((B+e,1+e,1/6+e)\).  The reciprocal is holomorphic for
\(\Re s>B\).  Move \(w\) to \(19/20\), and in its residue move
\(z\) to \(33/200\).  All paths lie in (P); only \(w=1,z=1/6\)
are crossed.  The raw power after the tuple bound is
\(l_x/2+\Re s-1+h\Re z+l_y(\Re w-1)\).
This gives the first two exponents.  The double residue, after
normalization, is
\[
 \frac1{2\pi i}\int_{(B+e)}
 Z^{C(s)}e^{(s-5/6)^2}\frac{H_\eta(s)}{L_F^S(s,\eta)}
                         (1+O(Z^{-\mu}))\,ds.
\]
Estimate its error on this line; move only its main part to \(2\).
This gives the third exponent and \eqref{eq:signal}.

To justify the contours, the functional equation gives polynomial
strip growth for each numerator.  On closed substrips of
\(\Re s>B\), a zero-free logarithm, Borel--Carath\'eodory, and
the Euler value at \(2+it\) give polynomial growth for the
reciprocal, by Lemmas~\ref{lem:hecke-strip-growth}--\ref{lem:deleted-euler-factors}.
For a principal target use
\((s-1)\zeta_F(s)/(s+1)\); its pole creates a reciprocal zero.
Deleted factors are nonzero on positive real strips.
The tests decay in the independent coordinates
\((\Im(s+z),\Im z,\Im w)\), dominating every fixed polynomial.
Horizontal joins therefore vanish.  Division by \(A_T\) costs
only the displayed arbitrarily small power.
\end{proof}

\section{The buffered detector and simultaneous saturated witnesses}\label{sec:detector-complete}
Write \(L(s,\psi)=L_F(s,\psi)\) and \(\beta_*:=B\) in this section and in the moment appendices.  All sums over ideals
use the good primary generators of Appendix~\ref{app:conventions}.
For any base \(V>1\) and nonnegative log-lengths, define
\begin{equation}\label{own:central-polynomials}
 \begin{split}
 M_\psi(r;W)&=V^{-r/2}\sum_n\mu(n)\psi(n)W(q_n/V^r),\\
 S_\psi(m;W)&=V^{-m/2}\sum_n\psi(n)W(q_n/V^m).
 \end{split}
\end{equation}
The full zero-extended presentation, not only its primitive character,
is used in both sums.  We give the detector proof of \cite[Section 8]{QRH}
at the range \(B\le1\), including its height allocation.
The analytical preliminaries are included to specify exactly the
uniformity in conductor that the detector uses.
\subsection{Uniform growth, logarithmic control, and Euler deletions}

The disk estimate will be applied only after zeros have been excluded from
that disk.  The following elementary growth bound provides the input for the
complex-analytic argument and makes clear which constants are uniform in the
conductor.  When $\psi$ is a Hecke character, $L(s,\psi)$ abbreviates
$L_F(s,\psi)$.  In this subsection $\Gamma(s)$ denotes Euler's gamma function;
the finite quadratic function $\Gamma(c)$ was confined to the arithmetic
identities above.

\begin{lemma}[Hecke strip growth]\label{lem:hecke-strip-growth}
Let $\psi$ be a primitive nonprincipal finite-order Hecke character of $F$,
with conductor norm $Q$.  Then
\begin{equation}\label{eq:hecke-growth}
 |L(\sigma+it,\psi)|\ll Q^{3/5}(3+|t|)^2,
 \qquad-1/10\le\sigma\le11/10.
\end{equation}
For the principal character, the same bound with $Q=1$ holds for
$(s-1)\zeta_F(s)/(s+1)$, with the removable value used at $s=1$.
\end{lemma}

\begin{proof}
A finite-order character has trivial infinite type, since $\mathbb C^\times$
is connected.  The primitive Hecke functional equation in this case is
\[
 \Lambda(s,\psi)=(3Q)^{s/2}(2\pi)^{-s}\Gamma(s)L(s,\psi),\qquad
 \Lambda(s,\psi)=\varepsilon(\psi)\Lambda(1-s,\overline\psi),
 \quad |\varepsilon(\psi)|=1.
\]
This is the finite-part normalization of
\cite[Section 4.7 and Lemma 4.20, complex-place calculation; Section 5.11,
Theorem 5.22]{Poonen}.  The complex factor is
$\Gamma_{\mathbb C}(s)=2(2\pi)^{-s}\Gamma(s)$, and
$|D_F|=3$; its constant factor $2$ cancels on the two sides.
A nonprincipal $L(s,\psi)$ is entire.  Absolute Euler convergence bounds $L(11/10+it,\psi)$ uniformly in
$Q,t$.  Applying the functional equation and Stirling's formula on
$\Re s=-1/10$ gives
\[
 |L(-1/10+it,\psi)|\ll Q^{3/5}(3+|t|)^{6/5}.
\]
The estimate for bounded $t$ follows by continuity of the gamma quotient on
that line, so the constant is uniform there as well.

For completeness, the growth hypothesis needed to use the strip principle
can be obtained directly from a lattice theta integral.  Let
$\mathfrak f=(f)$ be the conductor and define the periodic residue character
$\phi(a)=\psi((a))$ when $(a,\mathfrak f)=1$, with value zero otherwise.
For the principal conductor $(1)$ put $\phi(a)=1$ for all $a$, including
$a=0$.  Since the ideal class group is trivial, conductor one has no
nonprincipal character.  Periodicity modulo $\mathfrak f$ follows from the ray character
property; $\phi$ is trivial on the six units because $(u)=(1)$ for a unit.
Thus, for $\Re s>1$,
\[
 6\pi^{-s}\Gamma(s)L(s,\psi)
 =\int_0^\infty\bigl(\Theta_\phi(v)-\phi(0)\bigr)v^{s-1}\,dv,
 \qquad
 \Theta_\phi(v)=\sum_{a\in\mathcal O}\phi(a)e^{-\pi vq_a}.
\]
The factor six counts the generators of each ideal.  If
$\widehat\phi(h)=Q^{-1}\sum_{a\bmod f}\phi(a)e(-ha/f)$,
finite Fourier inversion and Gaussian Poisson summation give
\[
 \Theta_\phi(v)=\frac{2}{\sqrt3v}
 \sum_{h\bmod f}\widehat\phi(h)
 \sum_{b\in\mathcal O}
    \exp\left(-\frac{4\pi|b-h/f|^2}{3v}\right).
\]
Changing the generator $f$ only reindexes the finite sum.  For each fixed
conductor, the nonzero dual vectors have a positive minimum length.  With
$A_\phi=2\widehat\phi(0)/\sqrt3$, this proves
\[
 \Theta_\phi(v)=A_\phi v^{-1}+O_\phi(v^{-1}e^{-c_\phi/v})\quad(0<v\le1),
 \qquad
 \Theta_\phi(v)=\phi(0)+O_\phi(e^{-c_\phi v})\quad(v\ge1)
\]
for some $c_\phi>0$.  Splitting the Mellin integral at one consequently gives
\[
 \begin{split}
 6\pi^{-s}\Gamma(s)L(s,\psi)
 ={}&\frac{A_\phi}{s-1}-\frac{\phi(0)}s
 +\int_0^1\bigl(\Theta_\phi(v)-A_\phi/v\bigr)v^{s-1}\,dv\\
 &+\int_1^\infty\bigl(\Theta_\phi(v)-\phi(0)\bigr)v^{s-1}\,dv.
 \end{split}
\]
Both integrals are entire and bounded in height on each bounded real strip,
with constants that may depend on the fixed conductor.  Stirling's formula
for $1/\Gamma(s)$ therefore gives at most exponential height growth for
$L(s,\psi)$ on that strip.  In the principal case the same conclusion holds
after multiplication by $(s-1)/(s+1)$.  Only this qualitative growth, for
each fixed conductor, is used in the strip principle.

To see explicitly that the fixed-conductor growth constants do not enter
the uniform strip bound, put
\[
 H_Q(s)=Q^{-3/5}\frac{L(s,\psi)}{(s+2)^2}.
\]
It is holomorphic in the closed strip and is bounded by one absolute
constant $C$ on both vertical boundary lines, by the bounds already proved.
For $\delta>0$, apply the maximum principle on the rectangle of height $T$
to $H_Q(s)\exp(\delta(s-1/2)^2)$.  On the vertical sides the exponential
has modulus at most $e^{9\delta/25}$, because
$|\Re s-1/2|\le3/5$.  On the horizontal sides its modulus is at most
$e^{\delta(9/25-T^2)}$, which tends to zero faster than the qualitative
fixed-$Q$ exponential growth as $T\to\infty$.  First let $T\to\infty$ for
this fixed $Q,\delta$, and then let $\delta\downarrow0$.  The result is
$|H_Q(s)|\le C$ throughout the strip, with the same constant for every $Q$.
Multiplication by $Q^{3/5}|s+2|^2$ proves
Equation~\eqref{eq:hecke-growth}.

For the principal function, the factor $(s-1)/(s+1)$ removes its only pole
in the strip and is bounded on the boundary.  Its functional equation gives
the same boundary bounds with $Q=1$.  This equation also follows from the
preceding Poisson formula with $\phi=1$: after writing $v=2t/\sqrt3$, the
theta relation is
$\Theta_\phi(2t/\sqrt3)=t^{-1}\Theta_\phi(2/(\sqrt3t))$, whose split Mellin
integral is invariant under $s\mapsto1-s$.  Apply the same damped-rectangle
argument to $(s-1)\zeta_F(s)/((s+1)(s+2)^2)$.
\end{proof}

\begin{lemma}[Logarithmic control]\label{lem:logarithmic-control}
Let $\psi$ be a primitive nonprincipal finite-order Hecke character of
conductor norm $Q$, and put $\mathcal L(s)=L(s,\psi)$.  For the principal
character put $Q=1$ and
$\mathcal L(s)=(s-1)\zeta_F(s)/(s+1)$, with its removable value at one.
Fix $a\in[1/2,1]$ and $0<e<10^{-3}$.  Suppose $\mathcal L$ has no zero in
the open disk of radius $2-a-2e$ centered at $2+it$.  Uniformly on the
closed concentric disk of radius $2-a-6e$,
\begin{equation}\label{eq:disk-control}
 |\mathcal L(s)|+|\mathcal L(s)^{-1}|
 \ll_{e,\epsilon}\{2Q(3+|t|)^2\}^{\epsilon}.
\end{equation}
On the disk of radius $2-a-8e$ one also has
\[
 |\mathcal L'(s)/\mathcal L(s)|\ll_e\log\{2Q(3+|t|)^2\}.
\]
In particular, for every fixed $b\ge\beta_*$ and $v>0$, these bounds hold
on $\Re s\ge b+v$, with constants depending on $v,\epsilon$ (and with the
height of $s$ in place of $t$).  In the principal case the reciprocal bound
passes to $1/\zeta_F$; the upper and logarithmic-derivative bounds are for
the regularized function $\mathcal L$.
\end{lemma}

\begin{proof}
Write $R_j=2-a-je$ and $\mathcal C=2Q(3+|t|)^2$.  All the disks used here
lie in $\Re s>1/2$, where the principal regularizer is holomorphic.  On the
zero-free disk choose a holomorphic logarithm $g=\log\mathcal L$ whose value
near the center is the Euler logarithm, together with the logarithm of the
regularizing factor in the principal case.  The value $g(2+it)$ is uniformly
bounded.  Lemma~\ref{lem:hecke-strip-growth} and Euler convergence to the
right of $11/10$ give
\[
 \Re g(s)=\log|\mathcal L(s)|\ll\log\mathcal C
       \qquad(|s-(2+it)|\le R_3).
\]
The possible heights differ from $t$ by at most $3/2$, which is included in
$\mathcal C$.  Borel--Carath\'eodory on radii $R_3,R_4$, whose difference is
$e$, now gives $|g|\ll_e\log\mathcal C$ on the disk of radius $R_4$.

On the disk of fixed radius $r_0=49/100$, one has $\Re s\ge151/100>1$.
The absolutely convergent Euler logarithm is uniformly bounded there.  For
the principal function the logarithm of $(s-1)/(s+1)$ is also bounded on
this disk, using its branch in $\Re s>1$.  Since
$r_0<R_6<R_4$, Hadamard's three-circles theorem applied to $g$ gives
\[
 \max_{|s-(2+it)|\le R_6}|g(s)|
 \ll_e(\log\mathcal C)^\theta,
 \qquad
 \theta=\frac{\log(R_6/r_0)}{\log(R_4/r_0)}<1.
\]
For fixed $e$, continuity on the compact interval $1/2\le a\le1$ makes
$\theta$ uniformly smaller than one.  For every $\epsilon>0$,
$\exp(C_e(\log\mathcal C)^\theta)\ll_{e,\epsilon}\mathcal C^\epsilon$.
Applying this to both $e^g$ and $e^{-g}$ proves
Equation~\eqref{eq:disk-control}.  Cauchy's estimate between radii $R_6$ and
$R_8$ gives the stated bound for $g'=\mathcal L'/\mathcal L$.

For the global assertion first suppose $\beta_*\le b\le1$.  Choose
$0<e<\min(10^{-3},v/8)$ and put $a=b$.  The disk of radius $R_2$ is contained
in $\Re s>b+2e>\beta_*$, so it is zero-free by the definition of $\beta_*$
and absolute Euler convergence beyond one.  The disk of radius $R_8$ covers
the points $b+v\le\Re s\le2$ at its center's height; Euler convergence
covers $\Re s\ge2$.  The preceding constants are uniform in $b$.  If $b>1$,
Euler convergence on $\Re s\ge1+v$ suffices.  Finally
\[
 \frac1{\zeta_F(s)}=\frac{s-1}{s+1}\mathcal L(s)^{-1},
 \qquad \left|\frac{s-1}{s+1}\right|\le1\quad(\Re s\ge0),
\]
which proves the principal reciprocal assertion.
\end{proof}

\begin{lemma}[Deleted Euler factors]\label{lem:deleted-euler-factors}
Let $R$ be a squarefree ideal and let $|a_p|\le1$ for $p\mid R$.  Put
$D_R(s)=\prod_{p\mid R}(1-a_pq_p^{-s})$.  For fixed $\sigma_0>0$ and every
$\epsilon>0$,
\[
 |D_R(s)|+|D_R(s)^{-1}|\ll_{\sigma_0,\epsilon}q_R^\epsilon
       \qquad(\Re s\ge\sigma_0).
\]
On a bounded real strip one has, uniformly in the height,
\[
 |D_R(s)|\ll_\epsilon q_R^{(-\Re s)_++\epsilon},
 \qquad
 \left|\frac{D_R'(s)}{D_R(s)}\right|
 \ll_{\sigma_0}\log(2q_R)\quad(\Re s\ge\sigma_0).
\]
In particular a support $q_R\le Z^B$ with bounded $B$ costs any prescribed
positive power of $Z$ on a fixed positive real half-plane.  For an
imprimitive function $L(s,\psi^*)D_R(s)$, both the primitive conductor and
the deletion radical $R$ must be included when applying logarithmic control.
\end{lemma}

\begin{proof}
Every factor is nonzero on $\Re s>0$.  For $\Re s\ge\sigma_0$, both its
absolute value and the absolute value of its inverse are at most
$(1-q_p^{-\sigma_0})^{-1}$.  For sufficiently large $q_p$, depending on
$\sigma_0,\epsilon$,
$-\log(1-q_p^{-\sigma_0})\le\epsilon\log q_p$.  The finitely many smaller
primes contribute a fixed constant.  This proves the first estimate.
For any real $\sigma$,
\[
 \prod_{p\mid R}(1+q_p^{-\sigma})
 \le q_R^{(-\sigma)_+}2^{\#\{p:p\mid R\}}
 \ll_\epsilon q_R^{(-\sigma)_++\epsilon}.
\]
The last inequality follows by separating the finitely many primes with
$q_p<2^{1/\epsilon}$.  Finally, logarithmic differentiation gives
\[
 \left|\frac{D_R'}{D_R}(s)\right|
 \le\sum_{p\mid R}\frac{\log q_p}{q_p^{\sigma_0}-1}
 \ll_{\sigma_0}\log(2q_R).
\]
If $q_R\le Z^B$, choose the exponent in the first estimate smaller than the
desired exponent of $Z$ divided by the bounded value of $B$; the case $B=0$
is immediate.  This proves the stated interpretation for deleted factors.
\end{proof}
\subsection{The row family and its buffer}\label{sec:detector}

The high expansion contains a sum over sixth-power-free rows \(u\).
We associate each nonprincipal row with a buffered zero-free rectangle
for the finite family of character presentations that it determines.
Whenever the resulting bin lies above a fixed floor, a zero produces
an inverse polynomial and a plain polynomial with simultaneous lower
bounds.  These witnesses can then be counted using different mean
estimates.  The analytic inputs here are Lemma~\ref{lem:logarithmic-control},
the smooth calculus of Appendix~\ref{app:smooth}, and the global
growth estimate in Equation~\eqref{eq:hecke-growth}, together with the
deleted Euler-factor bounds of Lemma~\ref{lem:deleted-euler-factors}.

Throughout this section, assume \(\beta_*>51/100\), where \(\beta_*\)
is the global supremum in Equation~\eqref{eq:B-definition}.  This setting
does not select either of the two boundaries in the continuation
criterion.

\subsection{Buffered rectangles and pointwise bounds}

Fix the arithmetic data \(\mathcal A\) independently of \(Z\), as in
Appendix~\ref{app:conventions}.  In particular, the finite group
\[
 \Theta=\langle\eta,\widehat T\rangle
\]
contains the target and the fixed ray twists used in the high-row
family, and every conductor prime of a member of \(\Theta\) belongs
to \(S\).  For a
sixth-power-free element \(u\), define the finite collection of
presentations
\[
 \mathcal X_u=
 \bigl\{\psi_{u,\nu,\varsigma}(n)=\nu(n)\chi_n(u)^\varsigma:
       \nu\in\Theta,\ \varsigma\in\{-1,1\}\bigr\}.
\]
At a nonunit the value of \(\chi_n(u)^{-1}\), like every other
power of the symbol, is zero.  The group property of \(\Theta\)
shows that \(\mathcal X_u\) is closed under conjugation.  It contains
the numerator character \(\chi_\bullet(u)\), the denominator
character \(\eta\overline{\chi_\bullet(u)}\), and every zero-extended
presentation obtained by multiplying either orientation of the sextic
symbol by a member of \(\Theta\).

Let \(\psi^*\) be the primitive character inducing a presentation
\(\psi\in\mathcal X_u\), and let \(Q_\psi\) be its conductor norm.
In \eqref{own:row-fixed-ray}, every good local factor is a
character modulo \(p\), and every other ray modulus is fixed.
Multiplying by a member of \(\Theta\) therefore gives
\[
 Q_\psi\ll_{\mathcal A}q_u,\qquad
 L_{\rm orig}(s,\psi)
   =L(s,\psi^*)\prod_{p\in E_u}(1-\psi^*(p)q_p^{-s}),
 \qquad
 E_u\subset S\cup\{p:p\mid u\}.
\]
The radical of the deleted product has norm \(O_S(q_u)\).
These identities retain the original zero extensions; in particular,
they do not replace a row-dependent mask by an independently chosen
mask.

\phantomsection\label{par:physical-row-ramification}
If \(v_p(u)=j\in\{1,\ldots,5\}\) at a prime \(p\notin S\), the
local character of \(\chi_\bullet(u)\) on units at \(p\) has order
\(6/\gcd(6,j)>1\).  A character in \(\Theta\) is unramified there,
so it cannot cancel this local character.  Consequently a
presentation in \(\mathcal X_u\) can induce the principal character
only when \(u\) is supported on \(S\).  There are finitely many such
sixth-power-free rows, including unit factors.  Remove them for the
present detector; each application must estimate these bounded physical
rows separately.

Fix \(0<d_{\min}<d_{\max}<\infty\) and write \(U=Z^d\), where
\(d_{\min}\le d\le d_{\max}\).  The rows in the present dyad satisfy
\(q_u\asymp U\), with fixed comparison constants.  Let
\[
 0<\tau\le d_{\min}/100,\qquad T_1=Z^\tau>2,\qquad
 0<e<10^{-3},\qquad I=\lceil1/e\rceil+2.
\]
The value of \(\tau\) will be chosen after the fixed height orders
are known.  For now it is fixed independently of \(Z\).  In
particular \(T_1\le U^{1/100}\).

\begin{lemma}[Buffered zero-free bins]\label{lem:buffered-bins}
For every retained row \(u\), there are an index
\(i\in\{1,\ldots,I-1\}\) and a grid point
\[
 a\in(51/100+e\mathbb Z_{\ge0})\cap[51/100,1]
\]
with the following properties.  For \(j=1,\ldots,I\), set
\[
 M_j(u)=\max\left(\{51/100\}\cup
 \left\{\Re\rho:
 \begin{array}{l}
  \psi\in\mathcal X_u,\quad L(\rho,\psi^*)=0,\\
  \Re\rho\ge51/100,\quad |\Im\rho|\le3jT_1
 \end{array}\right\}\right).
\]
Then
\[
 a\le M_i(u)<a+e,\qquad M_{i+1}(u)<a+2e.
\]
If \(a>51/100\), some \(\psi\in\mathcal X_u\) has a zero
\(\rho=\sigma+i\gamma\) with
\[
 a\le\sigma<a+e,\qquad |\gamma|\le3iT_1.
\]
For every \(\epsilon_1>0\), every \(\psi\in\mathcal X_u\), and
sufficiently large \(Z\), uniformly in the retained row,
\[
 |L_{\rm orig}(s,\psi)|+|L_{\rm orig}(s,\psi)^{-1}|
 \ll_{\mathcal A,e,\epsilon_1} U^{\epsilon_1}
 \quad
 \left(\Re s\ge a+6e,\ |\Im s|\le(3i+2)T_1\right).
\]
On the reflected line in the same height range,
\[
 |L_{\rm orig}(1-a-6e+it,\psi)|
 \ll_{\mathcal A,e,\epsilon_1}
 U^{a-1/2+12e+\epsilon_1}(3+T_1)^C,
\]
where \(C\) depends only on the fixed real strip and the field.
There are \(O_e(1)\) possible pairs \((i,a)\), independently of
the row and its conductor.
\end{lemma}

\begin{proof}
Each maximum exists: the collection is finite, the zeros of each
nonprincipal primitive \(L\)-function are discrete in compact
rectangles, and absolute Euler convergence excludes zeros with real
part greater than one.  The sequence \(M_j(u)\) is nondecreasing
and lies in \([51/100,1]\).  If every one of its \(I-1\) increments
exceeded \(e\), then
\[
 M_I(u)-M_1(u)>(I-1)e>1,
\]
which is impossible.  Choose an index with
\(M_{i+1}(u)-M_i(u)\le e\), and round \(M_i(u)\) down on the stated
grid.  This gives the two inequalities.  When \(a>51/100\), the
maximum \(M_i(u)\) is attained by an actual zero, giving \(\rho\).
This reasoning allows a zero on the line one.

For \(|t|\le(3i+2)T_1\), the closed disk centered at \(2+it\)
with radius \(2-a-2e\) has real part at least \(a+2e\).
Its radius is less than \(3/2<T_1\), so every point in it has
imaginary part of absolute value less than \(3(i+1)T_1\).
The bound \(M_{i+1}(u)<a+2e\) therefore excludes every zero of every
primitive function in this disk.  Lemma~\ref{lem:logarithmic-control}
controls the concentric disk of radius \(2-a-6e\).  Apply it with
center \(2+i\Im s\), and use absolute Euler convergence farther
right.  Since \(Q_\psi\ll_{\mathcal A}U\) and
\((3+|t|)^2\ll_e U^{1/50}\), the arbitrarily small power in that
lemma can be chosen so that its contribution is \(U^{\epsilon_1/2}\).
The estimates for polynomial-size deleted Euler products following
that lemma supply the other \(U^{\epsilon_1/2}\).

For the last assertion, the primitive functional equation recalled
in the proof of Lemma~\ref{lem:hecke-strip-growth}
\cite[Theorem 5.22 and Section 5.11]{Poonen} relates the value at \(1-a-6e+it\) to
the conjugate primitive function at \(a+6e-it\).  The conjugate
presentation is in \(\mathcal X_u\).  The conductor and gamma
quotient contribute at most
\[
 Q_\psi^{a-1/2+6e}(3+|t|)^C.
\]
The reflected primitive value has an arbitrarily small \(U\)-power
by the preceding disk bound.  Finally, \(1-a-6e\ge-6e\), so the
upper bound for the deleted product is \(U^{6e+\epsilon_1}\), after
reducing the preliminary losses.  This proves the reflected
estimate.  The ranges of \(i\) and \(a\) give \(O_e(1)\) pairs.
\end{proof}

We call \((i,a)\) the bin of the row and put
\[
 \delta=2a-1,\qquad 1/50\le\delta\le1.
\]
Every primitive character inducing a presentation in \(\mathcal X_u\)
is a finite-order Hecke character.  The definition of \(\beta_*\),
together with \(\beta_*>51/100\), therefore gives \(M_i(u)\le\beta_*\).
In particular, the bin satisfies the exact inequalities
\[
 a\le\beta_*,\qquad \delta\le2\beta_*-1.
\]
The bin \(a=51/100\) will be bounded by the trivial row count.
For \(a>51/100\), the zero in Lemma~\ref{lem:buffered-bins} produces
large polynomials to which the moments can be applied.

\begin{lemma}[Pointwise dyadic estimates]\label{lem:detector-dyads}
Fix bounded nonnegative ranges for \(r,m\), and let
\[
 M_r=M_\psi(r;W_M),\qquad S_m=S_\psi(m;W_S)
\]
denote the two polynomials in \eqref{own:central-polynomials},
now at base \(U\).  Suppose
\(\psi\in\mathcal X_u\), the untwisted profiles have uniformly
bounded smooth seminorms on a fixed annulus, and any pure norm twist
has height at most \((3i+1)T_1\).  All additional Mellin frequencies
entering the same \(L\)-argument are required to have total absolute
value at most \(T_1/2\).

For every \(\epsilon>0\), there are \(e_0>0\), depending only on
\(\epsilon\) and the bounded length ranges, and a finite height
order \(A_{\mathcal A}\), uniform in the moving rows and profiles,
such that, when
\[
 0<e<e_0,\qquad (1+T_1)^{A_{\mathcal A}}\le U^{\epsilon/10},
\]
the following estimates hold for sufficiently large \(Z\):
\begin{equation}
 |M_r|^2\ll_{\mathcal A,\epsilon}U^{\delta r+\epsilon},
 \label{old-eq:6.2}
\end{equation}
\begin{equation}
 |S_m|^2\ll_{\mathcal A,\epsilon}
       U^{\delta\min(m,1-m)+\epsilon}.
 \label{old-eq:6.3}
\end{equation}
The external tail order may be chosen after the positive number
\(\tau\).  It does not change \(A_{\mathcal A}\).
\end{lemma}

\begin{proof}
We spell out the height restriction because a full infinite contour
shift would not be justified by the bin.  For a fixed annular \(W\),
a bounded real number \(\sigma\), and a pure twist \(y^{i\omega}\),
Mellin inversion on a line \(c+\sigma>1\) gives
\[
 \begin{split}
 &U^{-r/2}\sum_n\mu(n)\psi(n)
       W(q_n/U^r)(q_n/U^r)^{-\sigma+i\omega}\\
 &\quad=\frac1{2\pi i}\int_{(c)}
  \mathcal MW(s)\,
  U^{r(s+\sigma-i\omega-1/2)}
  L_{\rm orig}(s+\sigma-i\omega,\psi)^{-1}\,ds.
 \end{split}
\]
Here \(\mathcal MW\) is the Mellin transform defined in
the smooth calculus of Appendix~\ref{app:smooth}.  The twist occurs in the
\(L\)-argument, not in the transform whose tails are estimated.
Shift only the portion with added frequency at most the assigned
fraction of \(T_1\) to \(\Re(s+\sigma)=a+6e\).
The horizontal joins and the shifted portion remain in the
zero-free rectangle of Lemma~\ref{lem:buffered-bins}.  Its bound
gives \(U^{(a-1/2+6e)r+\epsilon_1}\) for the central portion.

Leave the remaining tails on the absolute line, which we may take
to be \(\Re(s+\sigma)=2\).  There the reciprocal Euler product is
absolutely bounded.  On the inverse joins the reciprocal remains
inside the buffered rectangle.  Since the real join interval and
the length range are bounded, the remaining arithmetic and scale
factors on a join are \(O(U^B(1+|\Im s|)^J)\) for fixed \(B,J\)
chosen before the external tail order.  The pointwise Mellin bound
in Equation~\eqref{eq:pointwise-mellin-tail}, applied on that compact
real interval with derivative order \(N+\lceil J\rceil+1\), bounds
each horizontal join by \(O(U^B T_1^{-N})\).  The separate vertical
tails have the same bound by Equation~\eqref{eq:late-fourier-tail}.
Both estimates apply to the untwisted profile, uniformly in the
stated family.  Because \(T_1=Z^\tau\), \(\tau>0\), and all lengths
lie in a fixed bounded interval, a sufficiently large fixed \(N\)
makes these errors smaller than the claimed bound.
Taking \(e\) and \(\epsilon_1\) small in the prescribed
\(\epsilon\) proves Equation~\eqref{old-eq:6.2}.

For the plain polynomial the same calculation uses
\(L_{\rm orig}\) instead of its reciprocal.  Shifting to
\(\Re(s+\sigma)=a+6e\) gives
\[
 |S_m|\ll U^{(a-1/2+6e)m+\epsilon_1}
          +O(U^B T_1^{-N}).
\]
Alternatively shift its central portion to
\(\Re(s+\sigma)=1-a-6e\).  The function is entire because the
row is nonprincipal.  The reflected estimate of
Lemma~\ref{lem:buffered-bins} gives
\[
 |S_m|\ll
 U^{a-1/2+12e+\epsilon_1}
 U^{(1/2-a-6e)m}(3+T_1)^C
 +O(U^B T_1^{-N}).
\]
The joins here need only the global upper strip bound and the
negative-strip bound for the deleted product, since no reciprocal
is present.  These give fixed \(B,J\) as above, so the pointwise
Mellin estimate bounds the joins and the integrated Fourier estimate
bounds the absolute-line tails.  Taking the better central bound,
choosing \(N\) to
dominate also the possibly negative exponent when \(m>1\), and
using the displayed height hypothesis proves
Equation~\eqref{old-eq:6.3}.  A primitive conductor smaller than
\(U\) reduces the reflected conductor factor.  The original deleted
Euler factors have already been included in
Lemma~\ref{lem:buffered-bins}.

There are only finitely many separated variables in any one
\(L\)-argument.  Assign each a fixed fraction of the single
\(T_1/2\) allowance; do not assign that allowance anew at successive
shifts.  The profile windows and the finite height orders used in the
row estimates are fixed before \(N\).  The smooth calculus of Appendix~\ref{app:smooth} differentiates only the
untwisted separating profiles when increasing \(N\), which proves
the last assertion.
\end{proof}

\subsection{Simultaneous saturated witnesses}

At a selected zero, the truncated-inverse construction below produces
a product with squared size at least $U^{\delta(r+m)-\epsilon}$.
The pointwise estimates bound its square by
\[
 U^{\delta\{r+\min(m,1-m)\}+\epsilon}.
\]
Since $\delta\ge1/50$,
these two bounds force $m\le1/2+O(\epsilon)$. Each factor then
attains its own pointwise exponent up to the prescribed loss; this
is the saturation asserted in the next proposition.

\begin{proposition}[Two saturated witnesses]\label{prop:detector-witness}
Let \(u\) have a bin \((i,a)\) with \(a>51/100\), and assume the
loss and height hypotheses of Lemma~\ref{lem:detector-dyads}.
For every \(t\in[1,3/2]\),
and every prescribed \(\epsilon>0\),
after reducing its preliminary losses, there are
\(\psi\in\mathcal X_u\), a zero \(\rho=\sigma+i\gamma\) as in
Lemma~\ref{lem:buffered-bins}, dyadic lengths \(D=U^r\), \(N=U^m\),
and two polynomials \(M_r,S_m\) for the common row character
\(\psi\) such that
\begin{equation}\label{eq:detector-note}
 r\le t+O(1/\log U),\qquad
 r+m\ge t-O(1/\log U),\qquad
 |M_rS_m|^2\gg_{\mathcal A,\epsilon}
 U^{\delta(r+m)-\epsilon}.
\end{equation}
Their lengths and individual values also satisfy
\begin{equation}\label{eq:saturated-witness}
 \begin{gathered}
 t-\tfrac12-O(\epsilon)\le r\le t+O(1/\log U),
 \qquad 0\le m\le\tfrac12+O(\epsilon),\\
 |M_r|^2\gg_{\mathcal A,\epsilon}U^{\delta r-\epsilon},
 \qquad
 |S_m|^2\gg_{\mathcal A,\epsilon}U^{\delta m-\epsilon}.
 \end{gathered}
\end{equation}
The constants in the \(O(\epsilon)\) terms are absolute on the
stated parameter ranges.  Both polynomials have the same twist
height \(\gamma-\nu\), where \(|\nu|\le cT_1\) for a fixed
\(c>0\) chosen within the cumulative frequency allowance.
Their untwisted profiles form a uniformly smooth annular family.
For each row, both witnesses use one presentation
\(\psi_{u,\vartheta,\varsigma}\), selected by a label
\((\vartheta,\varsigma)\) in the fixed finite set
\(\Theta\times\{-1,1\}\).  The presentation itself varies with \(u\).
The dyadic pair can likewise
be selected separately for each row from \(O((\log U)^2)\) possibilities.
\end{proposition}

\begin{proof}
Here is the complete truncated-inverse construction.  The buffered
row-dependent rectangle and both simultaneous lower bounds are retained.
Choose the zero \(\rho\) supplied by the bin and set
\[
 D_*=U^t,\qquad Y_*=U^{20}.
\]
Let \(V_{\le}\) be a fixed smooth function equal to one on
\([0,1]\) and zero on \([2,\infty)\).  Define
\[
 C_\psi(s)=\sum_l\mu(l)\psi(l)V_{\le}(q_l/D_*)q_l^{-s},
 \qquad
 J=\frac1{2\pi i}\int_{(2)}
       Y_*^z\Gamma(z)L_{\rm orig}(\rho+z,\psi)
       C_\psi(\rho+z)\,dz.
\]
The row character is nonprincipal, so \(L_{\rm orig}\) is entire.
Moving the line to \(\Re z=-1/4\) crosses only the pole of
\(\Gamma(z)\) at zero, and its residue vanishes because
\(L_{\rm orig}(\rho,\psi)=0\).  The global strip bound
in Equation~\eqref{eq:hecke-growth}, together with the deleted
Euler factors, bounds the new line by
\[
 U^{-5}U^2D_*^{\,1-\sigma+1/4+o(1)}T_1^2.
\]
Indeed \(\Re(\rho-1/4)\ge26/100>0\), so the deleted product has
an arbitrarily small \(U\)-power; the deliberately weaker
\(U^2\) includes the primitive conductor bound.  The finite
polynomial is bounded by the displayed power of \(D_*\) using the
ideal count.  Finally, the fixed polynomial in the added height is
integrable against \(\Gamma(-1/4+iv)\), since
\((1+|\gamma+v|)^2\le(1+|\gamma|)^2(1+|v|)^2\).
Uniformly for \(1\le t\le3/2\),
\[
 -3+t(5/4-\sigma)
 \le-3+\frac32\left(\frac54-\frac{51}{100}\right)
 =-\frac{189}{100}.
\]
Thus \(J=o(1)\).  This Gamma integration uses only global upper
bounds and consumes no reciprocal height allowance.

On the original line, absolute convergence and the Mellin formula
for the exponential give
\[
 J=\sum_{l,k}\mu(l)\psi(lk)V_{\le}(q_l/D_*)
       (q_lq_k)^{-\rho}e^{-q_lq_k/Y_*}.
\]
For \(1<q_n\le D_*\), the coefficient of \(\psi(n)q_n^{-\rho}\)
is \(\sum_{l\mid n}\mu(l)=0\).  For \(n=1\) it is one.
When \(q_n>D_*\), \(V_{\le}(2q_n/D_*)=0\).
Consequently insertion of \(1-V_{\le}(2q_lq_k/D_*)\) removes
exactly the unit contribution \(e^{-1/Y_*}=1+o(1)\).
The resulting tail has absolute value \(1+o(1)\).
Multiplying it by a fixed smooth terminal cutoff equal to one for
\(q_lq_k\le U^{21}\) and zero for \(q_lq_k\ge2U^{21}\) changes it
by \(o(1)\); this follows from \(Y_*=U^{20}\), exponential decay,
and the ideal count.

Partition \(q_l\asymp D\), \(q_k\asymp N\) by a fixed smooth
dyadic partition.  There are \(O((\log U)^2)\) relevant pairs,
and their support satisfies
\[
 D\ll D_*,\qquad DN\gg D_*,\qquad DN\ll U^{21}.
\]
With \(x=q_l/D\), \(y=q_k/N\), the profile for a pair is
\[
 W_1(x)V_{\le}(Dx/D_*)\,W_2(y)\,H_{D,N}(xy),
\]
where \(\Omega\) is a fixed annular cutoff equal to one on the
product of the supports of \(W_1,W_2\), and
\[
 H_{D,N}(s)=\Omega(s)
   [1-V_{\le}(2DNs/D_*)]e^{-DNs/Y_*}
   V_{\rm term}(DNs/U^{21}).
\]
Every fixed logarithmic derivative is bounded uniformly in
\(D,N,U\).  On a cutoff transition its argument is bounded, and a
logarithmic derivative of the exponential is a polynomial in its
argument times that exponential.  Logarithmic Fourier inversion
therefore gives
\[
 \begin{split}
 H_{D,N}(xy)&=\frac1{2\pi}\int_{\mathbb R}
      \widehat H_{D,N}(\nu)x^{i\nu}y^{i\nu}\,d\nu,\\
 \int_{|\nu|>cT_1}|\widehat H_{D,N}(\nu)|
       (1+|\nu|)^J\,d\nu&\ll_{J,N_0,c}T_1^{-N_0}
 \end{split}
\]
for every fixed \(J,N_0\).  The coefficient \(L^1\)-norm on the
full line is also uniformly bounded.  Trivial bounds for the two
finite polynomials have a fixed \(U\)-power because their lengths
are bounded by \(U^{21+o(1)}\).  After \(\tau>0\) is fixed, choose
\(N_0\) so that the discarded Fourier tail is \(o(1)\).

The absolute value of the remaining sum of integrals is bounded
below by a positive constant.  The number of dyadic pairs and the
uniform \(L^1\)-norm show that, for one pair and one
\(|\nu|\le cT_1\), the product of the two unnormalized blocks has
absolute value \(\gg(\log U)^{-2}\).  Both profiles have the
same height \(\gamma-\nu\):
\[
 \begin{split}
 W_M(x)&=W_1(x)V_{\le}(Dx/D_*)x^{-\sigma-i(\gamma-\nu)},\\
 W_S(y)&=W_2(y)y^{-\sigma-i(\gamma-\nu)}.
 \end{split}
\]
The first extra cutoff remains part of the inverse annular profile
and creates no second frequency.  Removing the factors
\(D^{1/2-\sigma}\) and \(N^{1/2-\sigma}\) by central normalization
gives
\[
 |M_rS_m|^2\gg(\log U)^{-4}(DN)^{2\sigma-1}
 \ge U^{\delta(r+m)-\epsilon}
\]
for sufficiently large \(U\).  The support inequalities give the
two length bounds in Equation~\eqref{eq:detector-note}.

Apply Lemma~\ref{lem:detector-dyads} to these profiles, choosing its
loss smaller than the present \(\epsilon\).  Comparison of the
product lower bound with the two upper bounds gives
\[
 \delta\{m-\min(m,1-m)\}
   =2\delta(m-1/2)_+\ll\epsilon.
\]
Since \(\delta\ge1/50\), this implies
\(m\le1/2+O(\epsilon)\) with an absolute constant.
The support inequality then gives
\(r\ge t-1/2-O(\epsilon)\).
Dividing the product lower bound in turn by each individual upper
bound gives the two individual lower bounds in
Equation~\eqref{eq:saturated-witness}.  Only the dyadic support
errors are \(O(1/\log U)\); all other losses are arbitrarily small
fixed powers.  The smooth calculus of Appendix~\ref{app:smooth} permits the stated
rowwise choices with a fixed polynomial height factor.
\end{proof}
\subsection{Prime amplitudes}\label{sec:amplitudes}

For a main physical slot of scale \(P_i=Z^{\ell_i}\), write its
central factor as
\[
 Q_i(u;z)=P_i^{-1/2}\sum_{p\in1_T}
   \overline{\chi_p(u)}W_i(q_p/P_i)(q_p/P_i)^{z-1}.
\]
The main part \(\mathcal Q_i\) in Proposition~\ref{prop:probe-errors} is exactly
\[
 \mathcal Q_i(u;z)=-\sum_{p\in1_T}\overline{\chi_p(u)}q_p^{z-1}W_i(q_p/P_i)
 =-P_i^{z-1/2}Q_i(u;z).
\]
As in Section~\ref{sec:common-probe}, each sum over \(p\in1_T\) retains
the same allowed set \(\mathcal P_i(Z)\), including its exclusion of \(S\).
The original zero extension is retained, so a prime dividing \(u\)
does not contribute to this main slot.  The real part of \(z\) is
in a fixed bounded range and its imaginary part is one of the
external frequencies restricted by the common height allowance.

\begin{lemma}[Prime bound in a bin]\label{lem:bin-prime-bound}
Under the hypotheses of Lemma~\ref{lem:detector-dyads}, with the
prime-annulus Mellin frequency included in its cumulative allowance,
for every \(\epsilon_1>0\)
\[
 |Q_i(u;z)|\ll_{\mathcal A,e,\epsilon_1}
 U^{\epsilon_1}P_i^{a-1/2+O(e)}.
\]
The implied \(O(e)\) is uniform for \(51/100\le a\le1\) and the
fixed real ranges.  It holds for every main slot in the row.
\end{lemma}

\begin{proof}
Expand \(1_{p\in1_T}=|T|^{-1}\sum_{\theta\in\widehat T}\theta(p)\).
Every resulting prime character belongs to \(\mathcal X_u\).
Lemma~\ref{lem:logarithmic-control}, applied to the buffered disks,
bounds its logarithmic derivative on \(\Re s=a+8e\) by
\(O_{\mathcal A,e}(\log U)\); the logarithmic derivative of the
deleted product has the same bound.  Mellin inversion of the
corresponding smooth von Mangoldt annulus, shifted only in the
retained central range, gives
\(U^{\epsilon_1}P_i^{a-1/2+O(e)}\) after normalization.
The pure twist \((q_p/P_i)^{i\Im z}\) is translated into the
logarithmic-derivative argument before the added Mellin frequency
is truncated.  The joins have bounded real length.  On them the
logarithmic derivative is \(O_{\mathcal A,e}(\log U)\) inside its
allocated buffer, while on the starting line \(\Re s=2\) it is
absolutely bounded.  Apply the pointwise estimate in
Equation~\eqref{eq:pointwise-mellin-tail} to the untwisted transform
for the joins, and Equation~\eqref{eq:late-fourier-tail} for the
absolute-line tails.  Choosing the external order after their fixed
polynomial scale and height bounds gives the asserted estimate.

To pass from von Mangoldt coefficients to primes, divide the
annular weight by \(\log(P_i y)\).  On the fixed annulus,
\[
 (y\partial_y)^j\frac1{\log(P_i y)}
 =\frac{(-1)^j j!}{\{\log(P_i y)\}^{j+1}},
\]
so this preserves every fixed smooth seminorm for sufficiently
large \(P_i\), without introducing a power of \(P_i\) depending on
the derivative order.  Bounded \(P_i\) are handled by absolute
counting.  Prime powers
contribute \(P_i^{o(1)}\) after central normalization, since their
number in a norm annulus is \(O(P_i^{1/2+o(1)})\).
The polynomial-size punctures and their zero extensions are already
included in the logarithmic derivative.  This proves the bound.
\end{proof}

Here is a precise amplitude subdivision.  Choose a fixed bin width
\(\vartheta>0\).  First reduce \(e\) and \(\epsilon_1\), after the
slot lengths have been fixed, so that the preceding bound is at
most \(P_i^{\delta/2+\vartheta}\) for all sufficiently large \(Z\).
This is possible because
\(\log U/\log P_i=d/\ell_i\) stays in a fixed bounded range for
each fixed mesh.  For a main slot with \(Q_i\ne0\), put
\[
 g_i=\min\left\{\frac{\delta}{2},
   \max\left(0,\vartheta
     \left\lfloor\frac{\log|Q_i|}{\vartheta\log P_i}\right\rfloor
   \right)\right\}.
\]
Put \(g_i=0\) if \(Q_i=0\), and also put \(g_i=0\) for an error
slot from Equation~\eqref{old-eq:5.13e}.  Then, for a main slot,
\[
 |Q_i|\le P_i^{g_i+\vartheta},\qquad
 g_i>0\ \Longrightarrow\ |Q_i|\ge P_i^{g_i}.
\]
No lower bound is asserted for a slot with \(g_i=0\).
The finitely many possible vectors \((g_i)\), with the possible
endpoint value \(\delta/2\), partition the rows into amplitude
bins.  Define their length-weighted mean by
\begin{equation}
 q=\frac{\sum_i\ell_i g_i}{\ell},\qquad
 0\le q\le\delta/2,\qquad \ell=\sum_i\ell_i.
 \label{old-eq:6.6}
\end{equation}
This \(q\) records prime amplitude; it is not the conductor
exponent denoted \(q\) in the auxiliary fourth moment.

\section{Row envelopes with the actual plain parameter}\label{sec:counts}
Throughout the contradiction argument suppose \(B>\theta\), and put
\[
 \Delta=B-\theta>0,\qquad \kappa=2B-1=\kappa_0+2\Delta\in(\kappa_0,1].
\]
Appendix~\ref{app:plain} proves the plain moment on
\(13/18\le\kappa\le1\).  If \(\kappa<1\), its premise is
satisfied by equality; if \(B=\kappa=1\), its absolute prime
bound applies.  This is not an appeal to the incumbent.

The detector proved in Section~\ref{sec:detector-complete}
applies for \(B\le1\).  At base \(U=Z^d\), a bin has
\(a\le B\), \(\delta=2a-1\in[1/50,1]\).
Only \(\delta>1/50\) has a witnessing zero.
For any \(t\in[1,3/2]\), its saturated inverse/plain lengths obey
\[
\begin{gathered}
 r\le t+o(1),\quad r+m\ge t-o(1),\quad
 r\ge t-1/2-O(\epsilon),\quad m\le1/2+O(\epsilon),\\
 |M_r|^2\gg U^{\delta r-\epsilon},\qquad
 |S_m|^2\gg U^{\delta m-\epsilon}.
\end{gathered}
\]
The presentations are \(\nu\chi_\bullet(u)^{\pm1}\),
\(\nu\in\Theta=\langle\eta,\widehat T\rangle\).
The detector uses only \(a\le1\); its truncated-inverse error is
\(O(U^{-189/100+o(1)}T_1^2)=o(1)\) for \(T_1\le U^{1/100}\).
It requires no preliminary zero-free half-plane.

\paragraph{Orientation and slot compatibility.}
Both witnesses may be put in the negative orientation \emph{together}.
If their presentation is \(\nu\chi_\bullet(u)\), conjugate the entire
inverse polynomial and the entire plain polynomial, including their
tests and norm twists.  The coefficients \(\mu(n)\) are real,
\(\overline\nu\in\Theta\), and absolute values and lengths are unchanged.
Conjugation sends the selected zero to its conjugate and stays in
the same conjugation-closed row family.  Thus write the resulting
common presentation as \(\psi_u(n)=\nu(n)\overline{\chi_n(u)}\).
At every allowed slot prime, which is outside \(S\), the exact identity
\[
 \psi_u(p)\,\overline{\nu(p)}1_{p\in1_T}
       =\overline{\chi_p(u)}1_{p\in1_T}
\]
holds, including when \(p\mid u\), when both sides vanish.
The coefficient on the left is fixed independently of the row and
is bounded.  Moreover
\[
 \overline{\nu(p)}1_{p\in1_T}
   =|T|^{-1}\sum_{\vartheta\in\widehat T}
                    (\overline\nu\vartheta)(p)
\]
has all its finite characters in \(\Theta\).  It is therefore admissible
in the marked inverse estimate.  Conjugating the \emph{whole selected
product}, including its slots, gives the positive orientation required
by the plain estimate, with coefficients
\(\nu(p)1_{p\in1_T}\), again in the stated class.
All physical profiles, and their conjugates, remain in the smooth
families used in the moment proofs.  In particular, there is no
row-dependent modification of a prime coefficient in this passage.

Fix an error subset and the physical external parameters.
Give error slots amplitude zero and main slots amplitudes
\(0\le g_i\le\delta/2\), with
\[
 |Q_i|\le P_i^{g_i+\vartheta},\qquad
 g_i>0\ \Longrightarrow\ |Q_i|\ge P_i^{g_i},\qquad
 q=\ell^{-1}\sum_i\ell_i g_i\in[0,\delta/2],\quad P_i=Z^{\ell_i}.
\]
These are exactly the bins constructed after Lemma~\ref{lem:bin-prime-bound}.
Their number is fixed.  At base \(U\), put \(w_i=\ell_i/d\).
Order positive slots by decreasing amplitude and take a fractional
initial segment of requested total length \(z\le\ell/d\);
if the positive length is smaller, take all positive slots.
Its weighted amplitude is at least \(qz\).
Dropping the fractional last slot loses at most
\((\delta/2)\max w_i\).  The squared spike is therefore at least
\(U^{2qz-\delta\max w_i}\).

For the adaptive count use the following reference functions:
\begin{equation}\label{eq:adaptive-functions}
\begin{gathered}
 c=\frac1{3\kappa_0},\quad x=\frac q\delta,\quad
 D=3-(1+2c)x,\quad P=(2-2cx)(1-x),\\
 {\cal J}=(\alpha-\delta)D+\delta P,\qquad
 R_*=1-\delta+\frac{(\alpha-\delta)\delta P}{2{\cal J}},
 \qquad \alpha=\frac56 .
\end{gathered}
\end{equation}
For \(0\le\delta\le\alpha,\ 0\le x\le1/2\),
\[
 \frac49<c<\frac6{13},\qquad
 \frac{10}{13}\le P\le D\le3,\qquad
 \frac58<\frac{25}{39}\le{\cal J}\le\frac52,\qquad
 1-\delta\le R_*\le1-\frac\delta2.
\]
Indeed \(D-P=1+x-2cx^2\ge1\), and
\({\cal J}\ge(5/6)\min(P,D)\).
These bounds include both endpoints in \(\delta\).

\begin{proposition}\label{prop:counts}
For every prescribed \(\epsilon>0\), there is a target-independent
slot mesh fine enough that all corresponding pointwise amplitude sets satisfy
\[
 \#\mathcal B\ll_{\rm fixed,\epsilon}
 U^{R+\epsilon}(1+T_1)^{A_\eta}.
\]
The mesh may depend on \(\epsilon\) and the bounded real parameter
ranges, but not on the target or \(Z\); a single fixed positive mesh
is not asserted to work for arbitrarily small later losses.
For \(1/50<\delta\le5/6\), with supply \(\ell/d>1/5\)
by a fixed margin, one may take
\[
 R=R_*+\Delta/4.
\]
For \(5/6\le\delta\le1\), one may take \(R=1-\delta\)
without slots.  For \(1/50<\delta\le5/6\), one may alternatively
take \(R=1-2\delta/3\) without slots.  At the floor take \(R=1\).
All internal height orders precede the external tail order.
\end{proposition}
\begin{proof}
Subdivide by witness presentation and the
\(O((\log U)^2)\) dyadic length pairs.
For a negative-orientation witness the physical prime coefficient
is \(\overline{\nu(p)}1_{p\in1_T}\), a fixed finite combination
of characters from \(\Theta\), independent of the row.
For the plain moment conjugate the whole product; the coefficient
then is \(\nu(p)1_{p\in1_T}\).
All masks, zeros, and windows remain unchanged.
Every sufficiently large physical sixth-power-free row has a good
prime of valuation \(1,\ldots,5\), where its local character
cannot be canceled by \(\Theta\).
Thus its primitive inducing character lies outside \(\Theta\).
Bounded exceptions are covered by the small-row estimate.
Sobolev in the witness parameters costs a fixed height polynomial,
not row-dependent prime coefficients.

For \(r<1,m<1/2\), the inverse and plain capacities are
\[
 z_M=(1-r)/2,\qquad z_P=(1-2m)/(6\kappa).
\]
Their count exponents are \(1-\delta r-2qz_M\) and
\(1-2\delta m-2qz_P\).
Use a fixed positive decrement \(\nu_0\) in inverse capacity.
Without slots, including \(r=1\), Appendix~\ref{app:inverse} gives
\begin{equation}\label{eq:amplification}
 e(r)-\delta r,\qquad e(r)=\max\{1,(1+5r)/6\}.
\end{equation}
For \(\delta\le\alpha\) and \(r\ge1\), this is at most
\[
 L(t)=1-\delta+(\alpha-\delta)(t-1)+O(\epsilon).
\]
For \(m\ge1/2\), the unselected plain count is at most \(1-\delta\).

In the remaining range first compute capacities at \(\kappa_0\).
The two affine exponents are
\[
 A_I(r)=1-\delta\{x+(1-x)r\},\qquad
 S_t(r)=1-\delta\{cx+(2-2cx)(t-r)\}.
\]
They meet at
\[
 r_*=\frac{(2-2cx)t+(c-1)x}{D}.
\]
As \(1\le t\le3/2\) and \(0\le x\le1/2\),
\[
 \frac{23}{37}\le r_*\le1,\qquad
 \frac13\le t-r_*\le\frac12.
\]
For verification, both expressions are affine in \(t\);
at \(t=3/2\) they are \(1\) and \(1/2\).
At \(t=1\), the minimum of \(r_*\) is
\((3-c)/(5-2c)\ge23/37\), and the minimum of \(1-r_*\)
is \(1/3\).
Use the inverse side for \(r\ge r_*\).
It needs at most \(7/37\) of supply, and its strict margins are
\[
 1-r-2z\ge2\nu_0,\qquad 3-2r-8z\ge9/37.
\]
On the plain side \(m\ge1/3-O(\epsilon)\) and the reference
capacity is at most \(1/(18\kappa_0)+O(\epsilon)\le1/13+O(\epsilon)\).
Both capacities are smaller than the supplied \(1/5\).

Replacing the reference plain capacity by the actual one
increases its count exponent by
\[
 \frac{q(1-2m)}3\left(\frac1{\kappa_0}-\frac1\kappa\right)
 =\frac{2q(1-2m)}{3\kappa_0\kappa}\Delta
 \le\frac{30}{169}\Delta+O(\epsilon)
 <\frac\Delta4+O(\epsilon).
\]
Here \(q\le5/12,\ m\ge1/3-O(\epsilon)\), and
\(\kappa,\kappa_0\ge13/18\).
This controls the whole interval \(B\in(\theta,1]\), not a grid.

The reference crossing count is
\[
 R_{\rm short}(t)=1-\delta+\frac{\delta P}{D}(3/2-t).
\]
Taking \(t=1+\delta P/(2{\cal J})\in[1,3/2]\) equates it
to \(L(t)\), with common value \(R_*\).
If \(z_M\le\nu_0\), \eqref{eq:amplification} gives
\(1-\delta+2\delta\nu_0+O(\epsilon)\).
If \(z_P\le\nu_0\), the zero-slot plain theorem gives
\(1-\delta+6\delta\nu_0+O(\epsilon)\).
Thus no limiting strict-moment constant is used.
Decrement and rounding losses fit the prescribed \(\epsilon\).

For \(\delta\ge\alpha\), take \(t=3/2\).
Then \(r\ge1-O(\epsilon)\), and for \(r\ge1\)
\[
 e(r)-\delta r-(1-\delta)=(\alpha-\delta)(r-1)\le0.
\]
The adjacent \(r<1\) interval costs \(O(\epsilon)\).
For the unselected alternative take \(t=1\):
the counts \(1-\delta r\) and \(1-2\delta(1-r)\)
cross at \(r=2/3\), giving \(1-2\delta/3\).
The floor is counted by ideals, not by this witness argument.
Mesh and height uniformity follow from the finite inductions
in Appendices~\ref{app:inverse}--\ref{app:plain}.
\end{proof}
\section{The full dynamic local-error estimate}\label{sec:dynamic-errors}
\subsection{Complete local-error table and simultaneous conductor allocation}\label{own:local-errors}

At an identity-ray prime every fixed phase from $T$ is one, although
$\eta(p)$ may be any unit.  Suppose first that $p\nmid u$ and put
$v=\chi_p(u)$.  Then $D=\eta(p)v^{-1}Q^{-x}$, and
$\overline{\eta(p)}Q^xD=v^{-1}$.  Using
$v^{-1}W=Q^{-w}$ and $\mathcal E_p=P_p^*+D$ gives the exact
normalized cancellation
\begin{equation}\label{old-eq:5.13b}
\begin{split}
 G_p+v^{-1}H_p={}&
 \frac{(1-V)(1-W)}{1-D}
 \Bigl[(v^{-1}-Q^{-w})\{V/(1-V)-D\}\\
 &\hspace{17mm}+(\overline{\eta(p)}Q^x+v^{-1}-Q^{-w})
                    \mathcal E_p\Bigr].
\end{split}
\end{equation}
This is an identity of holomorphic local expressions in the Euler regions.
Within these regions, on the raw quotient locus its left side is
$H_p(\mathcal B_p+v^{-1})$.
To see the cancellation directly on that locus, substitute
$P_p^*=-D+\mathcal E_p$ and
$P_p=(1-V)^{-1}+W/(1-W)+P_p^*$ into the left side before
normalization.  The constants $-v^{-1}$, $v^{-1}-Q^{-w}$, and
$(v^{-1}-Q^{-w})W/(1-W)=Q^{-w}$ sum to zero.
The resulting rational identity has only the denominators $1-R$, $1-V$,
and $1-D$, which are nonzero in the Euler regions; hence it gives the
displayed holomorphic identity there.
For $p\mid u$, the main term is zero by the original zero extension,
and Equation~\eqref{old-eq:5.13c} instead becomes
\[
 G_p=
 \overline{\eta(p)}Q^x(1-V)\mathcal E_p-Q^{-w}H_p.
\]

We now give the bound on the remaining error slots.  Its analytic
input is stated explicitly: the reflected primitive numerator must
be small at the retained height.  Instantiate Section~\ref{sec:detector}
with the current fixed data and $\Theta=\langle\eta,\widehat T\rangle$,
using its zero-extended presentations $\nu(n)\chi_n(u)^\varsigma$ for
$\nu\in\Theta$ and $\varsigma\in\{-1,1\}$.  The
numerator $\chi_\bullet(u)$ and its conjugate belong to $\mathcal X_u$
with fixed multiplier $1\in\Theta$, while the denominator
$\eta\overline{\chi_\bullet(u)}$ belongs to $\mathcal X_u$ with fixed
multiplier $\eta\in\Theta$.  The buffered estimates there, together with
the inverse deleted-factor bound of Lemma~\ref{lem:deleted-euler-factors},
therefore supply the reflected-numerator hypothesis below at the retained
heights: for the points used below, $\Re(1-w)=a+6e$ and
$|\Im(1-w)|\le T_1$, within the buffered rectangle.  The identity-ray
restriction makes every fixed $T$ phase equal to one at a
slot prime, without changing the physical row or any nonunit zero.

\begin{proposition}[Dynamic local errors and conductor allocation]
\label{prop:probe-errors}
Let $51/100\le a\le1$, $0<e\le10^{-3}$, and take
\[
 x_r=a+16e,\qquad w_r=1-a-6e,\qquad z_r=17/50.
\]
Choose $P_0$ sufficiently large in terms of $e$ and the fixed data.
For $U,T_1\ge1$ and a sixth-power-free row $u$ with $(u,S)=1$
and $q_u\asymp U$, let $\psi_u^*$
be the primitive character inducing $\chi_\bullet(u)$.  Let
$\mathcal W_u$ be any subset of
$\{w:\Re w=w_r,\ |\Im w|\le T_1\}$.  Suppose $\psi_u^*$ is
nonprincipal and, uniformly for $w\in\mathcal W_u$,
\[
 |L(1-w,\overline{\psi_u^*})|\ll U^\epsilon
\]
with any prescribed small power.
For this row on the displayed dynamic region, define \(\mathcal B_i\) by
its preceding slot sum, interpreting each \(\mathcal B_p\) as \(G_p/H_p\),
with \(G_p\) from Equation~\eqref{old-eq:5.13c}.  The stated choice of
\(P_0\) makes \(H_p\ne0\) there for every slot prime, as proved at the start
of the proof.  This statement asserts individual continuation only on this
dynamic region.
Define the main and error parts of a slot by
\[
 \mathcal Q_i(u;z)=
       -\sum_{p\in\mathcal P_i(Z)}W_i(q_p/P_i)q_p^{z-1}\overline{\chi_p(u)},
 \qquad
 \mathcal D_i(u)=\mathcal B_i(u;x,w,z)-\mathcal Q_i(u;z).
\]
Then, for every subset $I$ of the slots and the same points $w\in\mathcal W_u$,
\begin{equation}\label{old-eq:5.13e}
 \left|L^S(w,\chi_\bullet(u))\prod_{i\in I}\mathcal D_i(u)\right|
 \ll U^{a-1/2+12e+\epsilon}(3+T_1)^C
       \prod_{i\in I}P_i^{z_r-1/2}.
\end{equation}
The main parts retain the physical row, all zero masks, and the
fixed identity-ray restriction.  In particular, their central
normalization is
\[
 \mathcal Q_i(u;z)=-P_i^{z-1/2}
       \left(P_i^{-1/2}\sum_{p\in\mathcal P_i(Z)}
          \overline{\chi_p(u)}W_i(q_p/P_i)(q_p/P_i)^{z-1}\right).
\]
At the same points $w\in\mathcal W_u$, and on this dynamic region only, the
full correction has the finite decomposition
\begin{equation}\label{eq:dynamic-compensated-decomposition}
 \mathfrak H_{\eta,u,Z}
 =\mathcal H_{\eta,u}\prod_{i=1}^K(\mathcal Q_i+\mathcal D_i)
 =\mathcal H_{\eta,u}
   \sum_{I\subseteq\{1,\ldots,K\}}
       \prod_{i\in I}\mathcal D_i\prod_{i\notin I}\mathcal Q_i.
\end{equation}
\end{proposition}

\begin{proof}
In this region $\vartheta=(-w_r)_+\le6e$.  The estimates in the
proof of Lemma~\ref{lem:local-euler} sharpen to
$H_p-1=O(Q^{-1-10e})$ for $p\nmid u$ and
$H_p-1=O(Q^{-10e})$ for $p\mid u$.  Increase the fixed $P_0$ so
that $|H_p-1|<1/2$ for all these primes.  The holomorphic
factor $G_p$ in Equation~\eqref{old-eq:5.13c}, divided by $H_p$,
now defines $\mathcal B_p$ throughout this region.  It is $H_p$,
not the possibly singular raw factor $P_p$, that is bounded away
from zero when $w_r<0$.

Consequently the tuple sum in Equation~\eqref{eq:holomorphic-selected-tuple}
factors as $\mathcal H_{\eta,u}\prod_i\mathcal B_i$ on this dynamic
region.  Substituting $\mathcal B_i=\mathcal Q_i+\mathcal D_i$ proves
Equation~\eqref{eq:dynamic-compensated-decomposition}.  This factorization
is not used to continue the correction outside the present region.

For $p\nmid u$, Equation~\eqref{old-eq:5.13b} and the estimate for
$\mathcal E_p$ give
\[
 |H_p(\mathcal B_p+\overline{\chi_p(u)})|
 \ll Q^{-x_r+2\vartheta}+Q^{-6z_r+2\vartheta}
       +Q^{4-5x_r-6z_r+2\vartheta}
       +Q^{1-w_r-6z_r+\vartheta}
 \ll Q^{-51/100}.
\]
The four exponents are bounded respectively by
$-a-4e$, $-51/25+12e$, $49/25-5a-68e$, and
$a-51/25+12e$.  Division by the bounded $H_p$ preserves this
estimate.  The number of primes in a slot is $O(P_i)$ by the ideal
count; hence its total contribution from $p\nmid u$, including
$Q^{z-1}$, is $O(P_i^{z_r-51/100+\epsilon})$.

For $p\mid u$ there are only divisor-many possible labels.  A
monomial in $\mathcal E_p$ is multiplied by $Q^{z-1+x}$ in
$Q^{z-1}\mathcal B_p$, so after separating $Q^z$ its exponent is
$x_r-1$ plus the exponent of that monomial.  For the non-tail
boundary terms of the six-valuation table, these exponents at
$e=0$ are
\[
\begin{array}{c|ccccc}
 j&1&2&3&4&5\\\hline
 &a-2&1/2-2a&-a,\ 1-3a&1/2-2a&2-5a.
\end{array}
\]
At positive $e$ their changes are respectively
$+6e,-32e,(-26e,-48e),-42e,-80e$.  They are all strictly below
$-1/2$ in the stated range.  The common $R$ term is smaller still:
its exponent after separating $Q^z$ is at most $49/25-1-5a-80e<-1/2$.
Additional valuation pairs and $m$ values have ratios
$|R|\le Q^{-11/10}$ and $|V|=Q^{-51/25}$, so their geometric sums
preserve these bounds.  The strict second-family term with an
additional $V$ is also smaller than $Q^{-1/2}$ after this
normalization.

The sole remaining term is the strict
$(e_0,l,k,m)=(1,0,1,0)$ term, which occurs for $j\ge2$ and is
$-\eta(p)(Q-1)Q^{-x-w}$.  Its exponent after separating $Q^z$ is $-w_r$.
The explicit rescaling term $-Q^{-w}$ has exponent after separating $Q^z$
$-1-w_r<-1/2$, so it needs no further estimate.

Let $\mathfrak f_u$ be the conductor of $\psi_u^*$.  Write
$u=\epsilon_u\prod_{p\mid u}p^{j_p}$, with $\epsilon_u$ a unit.
For primary $n$ coprime to $uS$, reciprocity gives
\[
 \chi_n(u)=\epsilon_u^{(q_n-1)/6}
       \prod_{p\mid u}\mathcal R(p,n)^{j_p}\chi_p(n)^{j_p}.
\]
The unit exponent is read modulo six.  The prime formula extends to
composite $n$: when $A\equiv B\equiv1\pmod6$,
\[
 \frac{AB-1}{6}\equiv\frac{A-1}{6}+\frac{B-1}{6}\pmod6.
\]
Thus the unit factor depends only on $q_n$ modulo $36$, and the
$\mathcal R$ product depends only on $n$ modulo $4$.  Together with
the multiplicativity of $\mathcal R$ and the local symbols, the
displayed norm congruence makes the right side multiplicative on
integral ideals prime to $6u$; extend it to their fractional ideal
group.  A generator congruent to one modulo
$(36)\operatorname{rad}(u)$ is already primary and makes every
displayed factor one.  This character therefore has a ray
presentation with that fixed $2,3$ modulus times
$\operatorname{rad}(u)$.  It agrees with the Kummer character
$\chi_\bullet(u)$ of the fixed-numerator ray argument in Appendix~\ref{app:conventions} on ideals
avoiding $S$.  These characters induce
the same primitive character: in the principal ideal ring
$\mathcal O$, the Chinese remainder theorem supplies a representative
avoiding the additional finite set $S$ in every ray class of a
common modulus.  At a good prime
$p\mid u$, hold the fixed class and all other good residues at one
and vary the residue modulo $p$ by the Chinese remainder theorem.
The remaining local character $\chi_p^{j_p}$ has exact order
$6/\gcd(6,j_p)>1$.  Its conductor exponent at $p$ is therefore
exactly one: the displayed ray modulus has only the first power of
$p$, and the character cannot descend to a modulus omitting $p$.
This is a local assertion; the ideal character may still need the
fixed normalization modulus at $2,3$.

It follows that, for any set $J_0$ of distinct selected ramified
labels,
\begin{equation}\label{old-eq:5.13d}
 q_{\mathfrak f_u}\ll_S q_{\operatorname{rad}(u)}
       \le q_u\prod_{p\in J_0}q_p^{-(j_p-1)}.
\end{equation}
The functional equation for the primitive numerator
\cite[Theorem 5.22 and Section 5.11]{Poonen}, the assumed reflected bound, and Stirling
give a cost
$q_{\mathfrak f_u}^{A_*}U^\epsilon(3+T_1)^C$, where
\[
 A_*=1/2-w_r=a-1/2+6e>0.
\]
At a selected $p\mid u$ the primitive character already has value
zero, so restoring the original local factor there multiplies by
exactly one.  The remaining deleted Euler factors have radical
$O_S(U)$ and cost at most $U^{6e+\epsilon}$, because $w_r\ge-6e$.

Expand a product of error slots into their coprime-prime terms,
their already bounded ramified terms, and their strict ramified
terms.  For each fixed tuple in this triangle expansion, let $J_0$
be precisely its strict ramified labels.  They are distinct because
the slot supports are disjoint.  Apply the single inequality
Equation~\eqref{old-eq:5.13d} to this entire $J_0$ before summing
the labels.  The numerator together with these local factors,
including $q_p^{z-1}$ at each selected prime, is then bounded by
\[
 U^{A_*+6e+\epsilon}(3+T_1)^C
 \prod_{p\in J_0}q_p^{z_r-w_r-(j_p-1)A_*}
 \le U^{a-1/2+12e+\epsilon}(3+T_1)^C
       \prod_{p\in J_0}q_p^{z_r-1/2},
\]
because $j_p\ge2$ and $-w_r-A_*=-1/2$.  Thus different selected
labels use different factors of the conductor deficit, and the
bounds hold simultaneously.  Summing the divisor-many ramified labels costs
$U^\epsilon$, while the coprime labels have the stronger exponent
$z_r-51/100$.  This proves Equation~\eqref{old-eq:5.13e}.

This argument uses triangle only on error labels.  It never changes
the physical row or the coefficients and masks in any main slot.
The tuple-independent factor $\mathcal H_{\eta,u}$ separately costs
$U^\epsilon$ by Lemma~\ref{lem:local-euler}.  Relative to the
central scale $P_i^{z_r-1/2}$ of a main slot, an error slot thus
has no positive amplitude exponent in the later row count.
\end{proof}

\section{The nonprincipal comparison below seven eighths}\label{sec:nonprincipal}
This section proves the analytic connection to the actual probe,
using the row envelopes just proved.
The full correction is holomorphic in
\[
 {\cal D}_1(e_0)=
 \{\Re s\ge51/100,\ \Re z\ge17/50,\ \Re w\ge-1/100,\
                       \Re(s+w)\ge1+e_0\}.
\]
Lemma~\ref{lem:local-euler} and the quotient-free construction
\eqref{eq:compensated-all-height-majorant} prove this with
\(|\mathfrak H_{\eta,u,Z}|\ll_{\rm fixed,\epsilon}q_u^\epsilon Z^K\)
on fixed real boxes, uniformly in all heights.  The exponent \(K\)
is fixed before any external decay order.  This proof uses the
quotient-free tuple formula and is independent of our lengths.

For a fixed dynamic bin \(q_u\asymp U=Z^d\), first move
\(z:2\to z_0=17/50\), \(s:3\to B+20e\), with \(w=3\);
then move \(w\) to \(1-a-6e\).  Since \(a\le B\le1\),
\(\Re(s+w)\ge1+14e\) and \(\Re w\ge-6e\).
Every retained numerator is nonprincipal.  All these paths are
in \({\cal D}_1(10e)\), with global reciprocal bounds.
Restrict heights to fixed fractions of \(T_1\), and move only
the retained \(s\)-segment to \(a+16e\).
The buffered disk excludes all denominator zeros on this rectangle,
and \(\Re(s+w)\ge1+10e\).
This proves the extension of \cite[Lemmas 10.3--10.4]{QRH}:
their printed restriction \(\sigma_0\ge7/8\) is unused here.
Only \(B>\theta>43/50\) and \(a\le B\le1\) were needed.

All discarded portions and joins have bound
\(O_{\eta,N}(Z^{K_\eta}T_1^{-N})\), for \(K_\eta\) fixed before
\(N\).  Indeed their arithmetic factors have polynomial bounds;
the tests decay rapidly in
\((\Im(s+z),\Im z,\Im w)\).
On an \(s\)-join only its fixed \(s\)-height uses the buffered
estimate; extended \(w,z\) axes use global bounds.
An external trace tail follows by integrating
\((1+|t|)^{J-L}\) on the fixed-height face, with \(L>N+J+3\).
This also proves the integrated-tail assertion, without
differentiating arithmetic factors.

On the retained lines use Proposition~\ref{prop:probe-errors},
proved in Section~\ref{sec:dynamic-errors} for
\(51/100\le a\le1\), and needs no change.
Here are the relevant dependency checks.  Off \(u\), the selected
error is \(O(Q^{-51/100})\); at a ramified prime the only larger
term is the strict second-family term of \eqref{eq:local-family-table},
namely the \((e_0,l,k,m)=(1,0,1,0)\) term for \(j=v_p(u)\ge2\).
The numerator conductor satisfies, simultaneously for all distinct
strict labels \(J_0\),
\[
 q_{\mathfrak f_u}\ll_Sq_u
                  \prod_{p\in J_0}q_p^{-(j_p-1)} .
\]
At a good prime its local conductor exponent is exactly one.
The functional equation uses \(A_*=1/2-\Re w=a-1/2+6e\);
each strict label then has exponent
\(-\Re w-(j_p-1)A_*\le-1/2\).
Thus one numerator allowance pays for the entire error subset.
No main-slot mask or coefficient is changed.  The complete table, all ramified boundary cases, and both geometric
tails were checked explicitly in that proposition; no case is
assigned an unspecified exponent loss.

Consequently, on a pointwise amplitude set with count exponent \(R\),
\[
 \sum_u|L_F^S(w,\chi_\bullet(u))\mathfrak H_{\eta,u,Z}^{(I)}|
 \ll U^{R+\delta/2+12e+\epsilon_c'}
 Z^{\ell(z_0-1/2)+q\ell+\epsilon_c'}(1+T_1)^{A_\eta},
\]
where \(\epsilon_c'>0\) is freely prescribable after the slots
are fixed.  The \(12e\) is not a freely prescribable loss at fixed
\(e\): \(6e\) comes from \(A_*\), and at most another \(6e\)
comes from restoring Euler deletions on \(\Re w\ge-6e\).
The non-strict local terms already save a fixed positive power,
and each strict selected term is exactly paid to exponent
\(z_0-1/2\).  Thus no additional \(e\)-loss per slot is needed.
This is precisely the aggregate bound of Proposition~\ref{prop:probe-errors}.  Multiply by the buffered reciprocal and the
outside Mellin powers.  Relative to \(C(\theta)\), the
ideal exponent is exactly
\begin{equation}\label{eq:E}
 E(d;\delta,q,R)=a-\theta+
 h(z_0-1/6)-a l_y-\ell/2+q\ell
                       +d(R+\delta/2-z_0).
\end{equation}
The other real losses are
\[
 (16-6l_y+12d)e+(1+d)\epsilon_c'
                          +d\epsilon_d+\epsilon_p.
\]
The last term absorbs logarithmic subdivisions.
Subdivide only to bound the retained integrand, \emph{after}
moving the original fixed-bin sum; pointwise pieces are not
separately continued or integrated.
Taking the supremum of \eqref{eq:E} over the occurring pieces
therefore bounds the original integral.

\paragraph{Small rows.}
For \(1\le U\le Z^{1/100}\), \(u\ne1\), use the global lines
\((\Re s,\Re w,\Re z)=(B+e,1/2,z_0)\).
They lie in \({\cal D}_1(7/20)\).
The complete factors \eqref{old-eq:5.13a} and \eqref{old-eq:5.13c}
give the absolute tuple bound on these lines: off \(u\) the four selected-error exponents are
\(-s_r,-6z_0,4-5s_r-6z_0,1/2-6z_0\), all negative for
\(s_r\ge43/50\).
On \(u\), after multiplication by \(Q^{s_r}\), the boundary
exponents in order \(j=1,\ldots,5\) are
\[
 -1/2,\quad3/2-2s_r,\quad(3/2-2s_r,2-3s_r),
             \quad2-3s_r,\quad3-5s_r .
\]
All are at most \(1/2\), and the strict term is at most \(Q^{1/2}\).
The geometric ratios have modulus less than one uniformly.
Thus \(G_p=O(1)\) off \(u\), \(G_p\ll Q^{1/2}\) on \(u\);
divisor counting of the latter labels proves the positive tuple
bound \(O_\epsilon(U^\epsilon Z^{\ell z_0})\).
The unselected product has the same \(U^\epsilon\) majorant.
Lemma~\ref{lem:hecke-strip-growth}, Lemma~\ref{lem:logarithmic-control},
and Lemma~\ref{lem:deleted-euler-factors} give a
\(U^{3/5+\epsilon}\) numerator bound and a subpower reciprocal,
times fixed height polynomials.
Counting \(O(U)\) rows, including the nontrivial unit rows, gives
\[
 \sum_{U\le Z^{1/100}}|\mathcal R_\eta(U;Z)|
 \ll Z^{C(B)+h(z_0-1/6)-l_y/2+2d_{\min}+O(e+\epsilon)}.
\]
Here \(h(z_0-1/6)-l_y/2+2d_{\min}<-7/100\), by exact
substitution.  We enlarged the row power \(63/50\) to two.
The sole principal numerator is \(u=1\): a good prime of a
nonunit ramifies, and a nonidentity unit is not a sixth power in \(F\).
Indeed, if \(a^6\) is a unit, every finite valuation of \(a\)
is zero, so \(a\in\mathcal O^\times\) and \(a^6=1\).
The fixed-numerator Artin description then gives a nontrivial
character for each nonidentity unit.
A bounded row's denominator may be principal; its reciprocal
remains holomorphic on the global line.

\paragraph{Unbounded large rows.}
On \((s_r,w_r,z_r)=(2,2,z_\infty)\), \(z_\infty>2\),
the same table gives \(G_p=O(1)\) for every selected prime and
\(|\mathfrak H|\ll_\epsilon U^\epsilon Z^{\ell z_\infty}\).
Absolute Euler bounds and dyadic summation therefore give, with
\(B_0=l_x/2+1+l_y\),
\[
 \sum_{U>Z^{h+\zeta}}|\mathcal R_\eta(U;Z)|
 \ll Z^{B_0+(h+\zeta)(1+\epsilon)-\zeta z_\infty+\epsilon}.
\]
Each finite-dyad contour move is absolute and the final tail
is summable.  Take the fixed \(\zeta=10^{-14}\) of
\eqref{eq:new-geometry} and \(z_\infty=4\cdot10^{14}\).
For \(0<\epsilon\le10^{-14}\), the displayed exponent is less
than \(-1\): \(B_0<42/25,\ h<41/50\), and
\(\zeta z_\infty=4\).
These fixed large real parts change constants and test orders,
not the order of the later external height choice.

\section{A continuous rational certificate and uniform closure}\label{sec:closure}
All functions in this section use the exact rational geometry
\eqref{eq:new-geometry}, not rounded decimals.
Define
\[
 E^*(\delta,x)=E(h;\delta,x\delta,R_*),\qquad
 F(\delta,y)=-{\cal J}(\delta,1/2-y)E^*(\delta,1/2-y).
\]
By \eqref{eq:E} and \eqref{eq:adaptive-functions}, \(F\) is a
polynomial of degree at most two in \(\delta\) and three in \(y\).
Write \(F=\sum_{i=0}^3 F_i(\delta)y^i\).
We prove
\begin{equation}\label{eq:endpoint-margin}
 E^*(\delta,x)<-\frac{2}{10^{12}}
 \quad(0\le\delta\le5/6,\ 0\le x\le1/2).
\end{equation}

\paragraph{The near-active part \(1/3\le\delta\le5/6\).}
Put
\[
\begin{gathered}
 p(L)=657L^3-954L^2+21L+20,\qquad
 \delta_v=\frac{5-9\ell}{6(3\ell+1)},\\
 A_v=\frac{(3\ell+1)(414\ell^3-1191\ell^2+878\ell-29)}
             {4(3\ell-5)(114\ell^2-159\ell-7)},\\
 D_v=-\frac{(153\ell^2-201\ell-20)p(\ell)}
             {36(3\ell-5)(3\ell+1)(114\ell^2-159\ell-7)} .
\end{gathered}
\]
Direct polynomial division gives the exact identity
\begin{equation}\label{eq:square-cert}
 F_0(\delta)=A_v(\delta-\delta_v)^2+D_v.
\end{equation}
The formula for \(b\) in \eqref{eq:new-geometry} was chosen to
make this completion of the square exact.
In the entire interval \(1/6\le\ell\le167/1000\),
the displayed \(A_v\) is greater than \(1/5\), and
\[
 \frac{D_v}{p(\ell)}
 =\frac{-153\ell^2+201\ell+20}
       {36(5-3\ell)(3\ell+1)(-114\ell^2+159\ell+7)}
 >\frac1{160}.
\]
For example, a lower numerator and upper denominator for the latter
ratio are \(49232983/10^6\) and \(3694400793/500000\);
their ratio exceeds \(1/160\).
For \(A_v\), a lower numerator and upper positive denominator
are \(258102603/2000000\) and \(273477/500\);
their ratio exceeds \(1/5\).  These are interval product bounds
valid throughout the displayed \(\ell\)-range, not a grid.
For our rational \(\ell\),
\[
 p(\ell)=
 \frac{40963882847150664198386369}{46875\cdot10^{30}}
       >\frac{8}{10^{10}}.
\]
Consequently \(D_v>5/10^{12}\).

For \(u=\delta-1/3\), the coefficients of \(F_i(1/3+u)\),
in increasing powers \(1,u,u^2\), have the following
coefficientwise rational lower bounds:
\[
\begin{array}{c|rrr}
 & [u^0]&[u^1]&[u^2]\\ \hline
 10^3F_1(1/3+u)&3&112&605\\
 10^3F_2(1/3+u)&41&269&434\\
 10^3F_3(1/3+u)&16&98&148
\end{array}
\]
These are inequalities between rational coefficients, not sampled
function values.  They follow by substituting
\eqref{eq:new-geometry} in the polynomial defining \(F\);
Appendix~\ref{app:certificate-code} reproduces each comparison.
Thus \(F\ge D_v>5/10^{12}\) on this part of the rectangle.

\paragraph{The remainder \(0\le\delta\le1/3\).}
For each of the intervals \(I=[0,1/4]\) and \(I=[1/4,1/3]\),
put \(\delta=a+(b-a)t,\ y=v/2\), where \(I=[a,b]\).
Write
\[
 F(a+(b-a)t,v/2)
   =\sum_{i=0}^2\sum_{j=0}^3
        \beta_{ij}\binom2i t^i(1-t)^{2-i}
                       \binom3j v^j(1-v)^{3-j}.
\]
The two following matrices are coefficientwise integer lower
bounds for \(10^6(\beta_{ij})\), with row index \(i\):
\[
 \begin{pmatrix}
 35772&41252&46731&52210\\
 12757&12166&11370&10371\\
 4549&4195&5696&10209
 \end{pmatrix},\qquad
 \begin{pmatrix}
 4549&4195&5696&10209\\
 1813&1539&3804&10155\\
 723&1228&5211&14732
 \end{pmatrix}.
\]
For full reproducibility, if
\(F(a+(b-a)t,v/2)=\sum p_{kl}t^kv^l\), the exact conversion is
\[
 \beta_{ij}=\sum_{k\le i,\ l\le j}
       p_{kl}\frac{\binom ik}{\binom2k}
                    \frac{\binom jl}{\binom3l}.
\]
Each comparison is a finite integer cross-multiplication.
The Bernstein basis is nonnegative and has sum one on
\([0,1]^2\).  Hence \(F\ge723/10^6\) on the remainder.
Together with \({\cal J}\le5/2\), this proves
\eqref{eq:endpoint-margin} on the entire continuous domain.
It includes \(x=0,1/2\), \(\delta=0,5/6\), and the common
interval endpoints.

\paragraph{All other real ranges.}
In \eqref{eq:E}, replace \(q\) by its upper bound \(\delta/2\)
when using an unselected count.  Elementary substitution gives
\[
\begin{array}{c|c|c}
 \text{range}&R&\text{endpoint expression}\\ \hline
 \delta=1/50&1&
 E(h)=(25b+192\ell-36)/150\\
 0\le\delta\le5/6,\ d\le1/2&76/75-2\delta/3&
 E(1/2)=
 \dfrac{-75b\delta-49b+225\delta\ell+50\delta+78\ell-73}{300}\\
 5/6\le\delta\le1&1-\delta&
 E(h)=(-6b\delta+2b+15\ell-3)/12 .
\end{array}
\]
The first is below \(-1/200\).
The second increases in \(\delta,\ell\) and decreases in \(b\);
the ledger bounds give at most \(-133/56250<-1/500\).
Its \(d\)-slope is \(101/150-\delta/6>0\).
The third is at most \(-9/125<-7/100\);
its \(d\)-slope is \(33/50-\delta/2\ge4/25\).
The enlarged short count includes the floor, since its value
there is one.  Thus no witnessing-zero formula is used at the floor.

For \(1/2\le d\le h,\ 1/50<\delta\le5/6\), the adaptive count
is \(R_*+\Delta/4\).  Its \(d\)-slope is positive.
Relative to \(C(B)\), its endpoint exponent is at most
\[
 E^*(\delta,x)+h\Delta/4-\Delta
 <-2\cdot10^{-12}-(1-h/4)\Delta
 <-2\cdot10^{-12}.
\]
Every chosen slope is less than two, including \(B=1\).
Extending to \(h<d\le h+\zeta\) costs at most
\(2\zeta=2\cdot10^{-14}\), leaving more than \(10^{-12}\).
Prime supply stays greater than \(1/5\).
The small-row exponent was below \(-7/100\), and the unbounded
large-row exponent was below \(-1\).
This covers every physical row.

\paragraph{Order of choices.}
Reserve the target-independent nonprincipal margin
\(m_{\rm np}=10^{-12}\).
Choose the moment losses and the capacity decrement first, then
the common mesh from Appendix~\ref{app:plain}, within fixed small
fractions of \(m_{\rm np}\).
Choose an even integer \(K\), independent of the target, large
enough that \(2\ell/K\) is below the mesh, the allowed rounding
loss, and \(1/185\).
Set \(\ell_i=\ell/K\), and use nonzero nonnegative smooth
bumps in the disjoint intervals
\[
 \left(1+\frac{2i-1}{2K+1},
             1+\frac{2i}{2K+1}\right),\qquad1\le i\le K.
\]
Only \(d\ge1/2\) uses selection, so \(w_i\le2\ell/K\).
The internal amplifier pool has exponent strictly above the
allowed mesh; it is disjoint from the physical windows eventually.
All constants depending on the fixed slot count are now fixed.

The principal margins \(l_y/20,h/600,\mu\), with
\(\mu=(43/50)\ell/K\), are positive and target-independent.
Choose amplitude width, \(e\), the normalizer loss, and the
remaining small power losses after the slots.
They can leave a nonprincipal margin \(m_{\rm np}/4\) and half each
principal margin.  The real ranges are bounded, all denominators
\({\cal J}\) exceed \(5/8\), and the inverse margins are fixed.
Let \(m>0\) be a smaller common remaining high margin.
A large but finite \(K\), or a very small \(\mu\), changes no
logical quantifier.

Now fix a target and its arithmetic data, including the same
excluded set \(S\) in every identity.  It contains all primes of norm at most
\[
 P_0\ge \max\{10^{10},(2C_e)^{1/(10e)}\},
\]
where \(C_e\) is a uniform constant in the two local defect bounds
at the start of Proposition~\ref{prop:probe-errors}.
This is a choice made after \(e\), not a requirement on a
previously fixed target modulus.  Further fixed defining-modulus
primes of the target are included as well.
This changes neither the fixed ray group \(T\) nor the windows.
Put \(J_\eta=I_{\eta,\mathrm{modified}}/(c_SA_T)\).
The low estimate holds with \(\omega=(B-\theta)/2\), after
sharing this loss with \(A_T^{-1}\).
All internal seminorm and Sobolev orders are fixed before the
external tail order.  Let \(A_\eta\) dominate their height powers
and \(K_\eta\) the scale powers on discarded pieces.

List each external coordinate entering a buffered argument once.
For arguments \(\gamma_j+\sum_k a_{jk}v_k\), restrict
\[
 |v_k|\le c_kT_1,\qquad
 c_k=\{2n(1+\max_j|a_{jk}|)\}^{-1},
\]
where \(n\) is the total number of coordinates.
The added height is at most \(T_1/2\); it is not renewed
at successive applications.
For the smallest detector height loss \(\epsilon_{\rm ht}>0\),
choose
\[
 0<\tau\le\min\left\{\frac1{10000},
 \frac{d_{\min}\epsilon_{\rm ht}}{20(A_\eta+1)},
 \frac{m}{4(A_\eta+1)}\right\},\qquad T_1=Z^\tau .
\]
For all sufficiently large \(Z\), every required height factor
fits \(U^{\epsilon_{\rm ht}/10}\).
External tail differentiation acts only on untwisted smooth
profiles, not on arithmetic factors.  The estimates give
\[
 |J_\eta-f_\eta|
 \ll_{\eta,N}Z^{C(B)-m}(1+T_1)^{A_\eta}
                     +Z^{K_\eta}T_1^{-N}.
\]
Choose \(N\) last with
\(K_\eta-N\tau<C(B)-m/2\), and enlarge the starting scale.
The high saving is the common \(\sigma=m/2\).
Neither \(J_\eta\) nor \(f_\eta\) contains the auxiliary
\(T_1,\tau,N\).
Lemma~\ref{lem:criterion} contradicts \(B>\theta\).
The norm lift and \eqref{eq:imprimitive} complete the proof of
Theorem~\ref{thm:main}.  No part of the proof assumes a nonzero
imaginary part or excludes a real exceptional zero.

\section{Joint optimization, exact barriers, and failed extrapolations}
\label{sec:optimization}

This section separates the proved theorem from optimization of its
exponent bounds.  It does not assume that numerical tuning exhausts
the arithmetic method.

\subsection{\texorpdfstring{What happens if \(M+\ell\ne1\)}{Allowing general total length}}
To inspect the normalization rather than presuppose it, allow
\(M=l_x+l_y\), \(b=l_y-l_x\), and
\(h=1-(M-b)/2+\ell\).
At zero slack the reflected-energy variables of
\eqref{eq:full-energy-proved}, with \(S_0=0\), obey
\[
\begin{gathered}
 O,H,A,N,G,z,v,t,e\ge0,\quad O+H\le M,\quad
 2A\le O,\quad N\le A,\quad z\le\ell,\\
 v+3t+e+N+3G-2H-2A-2z\le-1-\ell .
\end{gathered}
\]
Dropping \(S_0\) only enlarges the exponent.
For \(0\le\ell\le1/2\), this relaxation has exact maximum
\begin{equation}\label{eq:general-low-max}
 {\cal E}_{\max}(M,\ell)=
 \max\left\{M,\ \frac M2+\frac14+\frac{3\ell}4,\
                         2M+\ell-1\right\}
\end{equation}
whenever it is nonempty.

Here is an exact linear-program certificate, without a grid or a
floating-point optimizer.  List the five constraints in the order
displayed above, with right sides \(M,0,0,\ell,-1-\ell\).
In the energy exponent choose each of
\(\max(H,v+t)\) and \(\min(v,z,(v+z)/3)\).
The six nonnegative dual multiplier vectors and resulting bounds are
\[
\begin{array}{c|c|c}
 \text{branch}&\text{multipliers}&\text{upper bound}\\ \hline
 H,\ v&(1/2,0,1/2,1/2,1/4)&M/2+1/4+3\ell/4\\
 H,\ z&(1,0,0,0,1/2)&M\\
 H,\ (v+z)/3&(1/2,0,1/2,1/6,1/4)&M/2+1/4+5\ell/12\\
 v+t,\ v&(1/2,0,0,1,1/2)&M/2+\ell\\
 v+t,\ z&(2,1,0,2,3/2)&2M+\ell-1\\
 v+t,\ (v+z)/3&(4/3,2/3,0,2,7/6)&4(M+\ell)/3-2/3 .
\end{array}
\]
Subtracting each branch from the corresponding linear combination
leaves a nonnegative combination of the nine variables.
In this calculation each branch has the additional constant
\((1+\ell)/2\); the last column includes it.  The script checks
both coefficient domination and every full right-hand-side identity.
The third and fourth bounds are below the second expression in
\eqref{eq:general-low-max}.
The sixth is below \(2M+\ell-1\) if \(2M-\ell-1\ge0\);
otherwise it is below \(M/2+1/4+3\ell/4\), since \(\ell\le1/2\).

Nonemptiness implies \(2M+\ell\ge1\).
The second expression in \eqref{eq:general-low-max} is attained
at \(H=(1-\ell)/2,\ O=M-H,\ z=\ell\), all other variables zero.
When either of the other two expressions is maximal, take
\(H=M,\ O=A=N=G=t=e=0,\ z=\ell,\ v=2M+\ell-1\).
Its required minimum branch is \(z\); the necessary \(v\ge2\ell\)
follows from maximality and \(\ell\le1/2\).
This proves exactness for the relaxation, not existence of arithmetic
rows realizing every relaxed point.

For a rescaled subset of length \(d\), replace \(M,\ell\) by
\(M-2d,\ell-d\).
After the square root and the tuple weight, its excess over the
nominal low exponent before the unrescaled energy correction is
\[
 -d+\frac12\max\left\{0,\
              \frac{1+3\ell-2M+d}{4},\
              M+\ell-1-3d\right\}.
\]
At \(d=0\) this includes the unrescaled energy correction;
the excess over the \emph{actual} unrescaled exponent is obtained
by subtracting that \(d=0\) value.  Each displayed branch is
maximized at \(d=0\), so the conclusion about the worst subset is
unchanged.
Fixed support slack changes these six linear bounds by \(O(r)\),
with an absolute constant; it can be absorbed before taking the
arbitrarily small analytic losses.
Thus the low exponent furnished by this energy/Gram calculation
corresponds to the boundary
\begin{equation}\label{eq:three-low-planes}
 {\cal T}(M,\ell)=
 \max\left\{
 T_1=\frac56+\frac M{12}-\frac\ell6,\quad
 T_2=\frac{23}{24}-\frac M6+\frac{5\ell}{24},\quad
 T_3=\frac13+\frac{7M}{12}+\frac\ell3
 \right\}.
\end{equation}
This remains only the relevant low-side calculation unless the
positive scale, Gram, local, and moment hypotheses also hold.
They have all been verified for \eqref{eq:new-geometry}.
For that geometry \(T_1=T_3=\theta\) and \(T_2<\theta\).
This explains why varying \(M\) was part of the optimization.

\subsection{A proved lower bound for the restricted exponent program}
Define the restricted program as follows: use the six-branch
reflected-energy bound just certified, the Gram exponent \(b/12\),
the selected-prime main/error exponents in \eqref{eq:E}, the
amplification slope \(5/6\), and the crossing count
\eqref{eq:adaptive-functions} with reference parameter
\(\kappa_0=2\vartheta-1\), where \(\vartheta\) is the proposed
boundary.  Allow all positive geometries \(h>0\),
\(0\le\ell\le1/2\), and \(M,b\), subject to their analytic
hypotheses.  Require \(\vartheta\ge T_i\) and require the
endpoint high envelope to be nonpositive as \(B\downarrow\vartheta\).
Allow arbitrarily small positive losses, arbitrarily fine slots,
and any order of choices respecting the analytic hypotheses.

Let \(L_*\) be the unique root of
\[
 p(L)=657L^3-954L^2+21L+20=0
       \quad\text{in }[1/6,1/5],
 \qquad \vartheta_*=\frac{11}{12}-\frac{L_*}{4}.
\]
Uniqueness follows from \(p(1/6)=1/24>0\),
\(p(1/5)=-1088/125<0\), and \(p'<0\) throughout this interval.
Exact rational endpoint substitution gives
\begin{equation}\label{eq:root-isolation}
 \frac{8749570194200989}{10^{16}}
 <\vartheta_*<
 \frac{8749570194200990}{10^{16}}.
\end{equation}

\begin{proposition}\label{prop:restricted-barrier}
No feasible boundary of this restricted program is smaller
than \(\vartheta_*\).
The statement allows \(M+\ell\ne1\) and is not merely a local
numerical obstruction.
\end{proposition}
\begin{proof}
The convex combination
\((17/30)T_1+(2/5)T_2+(1/30)T_3=13/15\)
first gives \(\vartheta\ge13/15\).
Only \(\vartheta<7/8\) needs discussion.
Set
\[
 L=\frac{11}{3}-4\vartheta\in[1/6,1/5],\qquad
 \delta_0=\frac{5-9L}{6(3L+1)}
          =\frac{9\vartheta-7}{18(1-\vartheta)}.
\]
Then \(1/3\le\delta_0\le7/18\); this is an admissible witness
bin, below \(2\vartheta-1\), not the detector floor.
Take \(x=1/2\).

For general \(M,\ell,b\), the endpoint high exponent at count
\(R=2/3\) is
\[
 H_0-\vartheta,\qquad
 H_0=(1+\delta_0)(1-M/2)+\delta_0\ell .
\]
There is no \(b\)-term.
The two inequalities \(T_1,T_3\le\vartheta\), with positive
dual multipliers \(6\delta_0+4/3\) and \(2/3\), imply
\[
 H_0\ge -(6\delta_0+2)\vartheta+6\delta_0+7/3
       =\vartheta.
\]
For the actual reference count, exact substitution gives
\[
 R_*(\delta_0,1/2)-\frac23
 =\frac{p(L)}
        {2(3L+1)(153L^2-201L-20)}.
\]
If \(\vartheta<\vartheta_*\), then \(L>L_*\) and \(p(L)<0\).
The denominator is negative, so this difference is positive.
The actual endpoint error is
\(H_0-\vartheta+h(R_*-2/3)>0\), since \(h>0\).
This contradicts the program's high constraint.
\end{proof}

The proved rational \(\theta\) is less than \(10^{-12}\) above
this restricted lower bound.  At the \emph{algebraic barrier geometry}
\(L=L_*\), the active bin has
\(\delta=(5-9L_*)/\{6(3L_*+1)\}=0.388583354\ldots\),
\(x=1/2\), and \(R_*=2/3\).
This equality is not asserted for the submitted rational geometry.
For that geometry the same exact identity gives
\[
 R_*(\delta_v,1/2)-\frac23
 =\frac{p(\ell)}
 {2(3\ell+1)(153\ell^2-201\ell-20)}<0,
\]
because \(p(\ell)>0\) and the denominator is negative.
Numerically the difference is approximately
\(-5.90956800257\cdot10^{-12}\); this decimal is only diagnostic.
The continuous certificate uses the exact rational count.
Neither the floor, local holomorphy, prime supply, nor a
small-row tail is active at the algebraic barrier.
The equation \(M+\ell=1\) and the rational expression for \(b\)
are consequences of balancing these constraints, not guesses
from one favorable inequality.
No claim of optimality beyond the stated exponent program is made.

\subsection{Conceptual improvements and limitations}
The extension \(13/18\le\kappa<3/4\) is an actual strengthening
of an intermediate lemma, proved in Appendix~\ref{app:plain}.
Its use gives the self-consistent parameter \(\kappa_0=2\theta-1\).
Keeping instead \(\kappa_0=3/4\) leads to the larger restricted
threshold
\[
 \frac{11}{12}
  -\frac14\left(\frac{11}{551}+\frac{8\sqrt{921}}{1653}\right)
  =0.874957069799\ldots .
\]
The difference is small but it is analytic, not a rounding change.

A prior heuristic exploration \cite{Exploration} reports values
near \(0.8745\) after changing the long-row slope.
Here that slope is derived from \(ua^6\):
\[
 e(r)=\max\{1,(1+5r)/6\},\qquad \alpha=5/6.
\]
A previously proved bound on \(B\) changes the plain parameter
\(\kappa\); it does not change \(\alpha\).
At \(r=3/2\), replacing \(5/6\) by \(3/4\) would require
a new saving \(1/24\) in the squared inverse-moment exponent.
No such estimate is supplied by a zero-free premise.

Nor does adding slots to the same sixth-power amplification
automatically improve the long witness.  Even in the favorable
case with amplification height \(r+2z\), the proposed count is
\[
 \frac16+\left(\frac56-\delta\right)r
                     +\left(\frac53-2q\right)z .
\]
For \(q\le5/12\) its last coefficient is at least \(5/6>0\).
If the other marked length inequality forces a still larger
height, the bound only worsens.
This rules out that particular shortcut, not other moment arguments.

The work ledger is therefore: the all-character rational
\eqref{eq:theta} is proved; \eqref{eq:root-isolation} is a certified
lower bound for the specified exponent program; a finite-grid
optimum is not a theorem; and the smaller extrapolated constants
would need genuinely stronger inverse or mixed moments.
Improving the arithmetic energy itself, altering the completion,
or exploiting correlations discarded by the two separate moments
remains outside that program.
\appendix
\section{Arithmetic dependencies and conventions}\label{app:conventions}

These appendices supply the arithmetic inputs to the common physical
probe and both moment bounds.  The construction is attributed to the
pinned manuscript, but the essential reductions are included here.
Source equation numbers elsewhere are provenance only; local labels
identify the arguments used in this document.

\begin{proposition}[Arithmetic package]\label{prop:arithmetic}
The general reflected-energy estimate \eqref{eq:full-energy-proved},
the Gram estimate \eqref{eq:Gram-proved}, the marked inverse
moment of Lemma~\ref{lem:marked}, and its sixth-power amplification
in Appendix~\ref{app:inverse} hold with the coefficient and zero
conventions specified in their proofs.  They are unconditional.
The plain moment of Lemma~\ref{lem:plain} holds on
\(13/18\le\kappa\le1\), under \(B\le(1+\kappa)/2\) for
positive slots when \(\kappa<1\), and without that premise at
\(\kappa=1\) or in the zero-slot case.
Its mesh is independent of the fixed arithmetic datum and fixed
slot count.  Each estimate has finite-seminorm and fixed-polynomial
height uniformity on bounded length ranges.
\end{proposition}

The proof occupies the following appendices.  For comparison with
\cite{QRH}, its lemma numbers give this acyclic dependency map:
\[
\begin{gathered}
 \text{theta/cusp formulas of \cite{DR}}
 \ \Longrightarrow\ \text{completed and marked reflection},\\
 \text{quadratic sieve \cite{GL}, cubic sieve \cite{HB}}
 \ \Longrightarrow\ \text{masked hybrid norm},\\
 \text{these two lines}\ \Longrightarrow\ 14.3
 \ \Longrightarrow\ 17.2\ \Longrightarrow\ 17.1
 \ \Longrightarrow\ 17.6,\\
 \text{finite correlation and lattice Poisson}\ \Longrightarrow\ 15.2,\\
 \text{finite correlation, Poisson, Hecke functional equation,
 and logarithmic control}\ \Longrightarrow\ 18.1 .
\end{gathered}
\]
In particular, 18.1 does not use 17.1.  The use of \(B\) in its
prime estimate occurs only to the right of the supremum defining \(B\);
it does not assume any fixed zero-free half-plane below \(B\).

\subsection{Residue symbols, finite phases, and fixed data}

In these appendices write \(\lambda_3=1+2\omega=\sqrt{-3}\).
This denotes the fixed ramified element, not a varying real length.
An ideal coprime to \(3\) has a unique primary generator
\(a\equiv1\pmod3\).  Put \(\alpha(a)=a/|a|\).  Fix a finite set
\(S\) containing the primes above \(6\) and all the defining
modulus primes of the \emph{fixed} character presentations.
For a prime \(p\notin S\), let
\(\chi_p(x)=(x/p)_6\), extended by zero at \(p\mid x\).
Every power, including zero and negative powers, retains this zero.
For a primary good ideal \(a\), put
\[
 \chi_a(x)=\prod_{p\mid a}\chi_p(x)^{v_p(a)},\qquad
 \gamma_j(a)=q_a^{-1/2}\sum_{x\bmod a}\chi_a(x)^j e(x/a)
 \quad(a\ {\rm squarefree}),
\]
where
\[
 e(z)=\exp(2\pi i\operatorname{Tr}(z/\lambda_3)),\qquad
 \check e(z)=\exp(2\pi i\operatorname{Tr}z).
\]
The additive pairing on \(\mathcal O\) is self-dual for
\(d\mu=(2/\sqrt3)\,dx\,dy\).
Indeed \(e(y)=e(\omega y)=1\) implies \(y\in\mathcal O\).
Thus the finite pairings used in Poisson summation below are perfect.

We record the elementary finite phases that occur in the transformations.
For odd \(c=a+b\omega\), Gaussian Poisson summation gives
\[
 \Gamma(c):=|c|^{-1}\sum_{x\bmod c}e(x^2/c)
 =\frac12\sum_{y\bmod2\mathcal O}e(-cy^2/4)
 =\frac{1+i^{-b}+i^a+i^{b-a}}2.
 \tag{A1}\label{eq:four-phase}
\]
One direct derivation is to apply Poisson to
\(e(z^2/c)e^{-\pi\varepsilon|z|^2}\) and let
\(\varepsilon\downarrow0\), after division by
\(2/(\sqrt3\varepsilon)\).
Its transform is
\[
 \frac{|c|}{2\sqrt{r_\varepsilon}}\,
 e^{-\pi\varepsilon q_c|y|^2/(4r_\varepsilon)}
 e(-cy^2/(4r_\varepsilon)),\qquad
 r_\varepsilon=1+3q_c\varepsilon^2/16.
\]
Splitting the two sides into classes modulo \(c\) and \(2\)
gives \eqref{eq:four-phase}; replacing \(r_\varepsilon\) by one
has total error \(O_c(1)\), negligible after this division.
For example its phase error is bounded by
\(O_c(\varepsilon^2\sum_y|y|^2e^{-c'\varepsilon|y|^2})=O_c(1)\).
This proves the formula also for nonsquarefree odd \(c\).

The square-class representatives modulo \(4\) are
\(1,-1,\lambda_3,-\lambda_3\), with \(\Gamma\)-values
\(1,1,i,-i\).  Thus
\[
 \mathcal R((-1)^e\lambda_3^f,(-1)^g\lambda_3^h)
       =(-1)^{eh+fg+fh}
\]
is the symmetric bicharacter
\(\Gamma(ab)/(\Gamma(a)\Gamma(b))\).
Its arguments in this table are residue square classes, not
ideal classes of the unit ideals.
Ordinary cubic reciprocity \cite[(1.4)]{DR}, together with the
quadratic Gauss CRT formula, gives
\(\chi_b(a)=\mathcal R(a,b)\chi_a(b)\) on coprime primary pairs.
We evaluate \(\mathcal R\) by the displayed finite table on all
pairs, never by dividing two zero residue symbols.

For clarity the signal identities can also be derived without an
unproved Gauss phase.  At a good prime \(p\), write \(H=\chi_p\).
The finite Gauss--Jacobi identity gives
\(\tau(A)\tau(B)=J(A,B)\tau(AB)\).
For \(A=B=H^2\), the Jacobi sum is divisible by \(p\), has
absolute value \(\sqrt{q_p}\), and is \(-1\pmod3\).
Divisibility follows by summing the polynomial
\(x^{(q_p-1)/3}(1-x)^{(q_p-1)/3}\) over the residue field.
The congruence modulo \(3\) follows from
\(\prod_{x\ne0,1}x(1-x)=1\) and
\(\omega^j\equiv1+j(\omega-1)\pmod3\).
The primary unit is therefore fixed: \(J(H^2,H^2)=-p\).
Also the count of solutions of \(4x(1-x)=y\) gives
\(J(H,H^3)=H(4)J(H,H)\).  Consequently
\[
 \gamma_2(p)^3=-\alpha(p),\qquad
 \gamma_1(p)\gamma_2(p)
       =\overline{H(4)}\gamma_3(p)\gamma_2(p)^3 .
\]
Define \(G(a)=\overline{\chi_a(4)}\Gamma(a)\) on every primary
good ideal, not just squarefree ones.
Cubic reciprocity at the primary prime \(-2\) shows that
\(\chi_a(4)=(a/(-2))_3\); hence \(G\) factors through a fixed
ray group at \(2,3\).  CRT now gives, for squarefree \(a\),
\[
 \gamma_2(a)^3=\mu(a)\alpha(a),\qquad
 \gamma_1(a)\gamma_2(a)=\mu(a)\alpha(a)G(a).
 \tag{A2}\label{eq:signal-checked}
\]
These are the finite identities used in the low separation and in
the two Poisson transformations of the inverse moment.
They agree with \cite[Lemmas 4.2--4.4]{QRH}.

A numerator supported on \(S\), with \(u\) a unit, evaluated by
\(A\mapsto\chi_A(u\prod_{\mathfrak p\in S}
\pi_{\mathfrak p}^{e_{\mathfrak p}})\) on ideals outside \(S\),
belongs to a fixed finite ray family: reduce the exponents modulo
six and apply the Artin map of
\(F((u\prod\pi_{\mathfrak p}^{e_{\mathfrak p}})^{1/6})/F\).
This extension is abelian and unramified outside \(6\) and the numerator.
It gives the asserted symbol at unramified primes by Frobenius.
The ray factorization is the power-residue consequence of
Artin reciprocity, \cite[Chapter VIII, (5.3) and (5.5)]{MilneCFT}.
There are only finitely many reduced numerators.
A moving good prime is never put into this fixed family.

The same finite formulas also give, on all primary good ideals,
\begin{equation}\label{eq:G-quadratic-refinement}
 G(ab)=G(a)G(b)\mathcal R(a,b),
\end{equation}
and at a good prime
\begin{equation}\label{eq:fixed-gauss-phases}
 G(p^3)=\gamma_3(p),\qquad \mathcal R(p,p)=\chi_p(-1).
\end{equation}
Indeed \(\chi_a(4)\) is multiplicative of order dividing three,
and the four-class table gives the quadratic refinement of
\(\Gamma\), its invariance under multiplying its argument by a
square, and \(\Gamma(p)^2=\chi_p(-1)\).  The defining quadratic
Gauss sum gives \(\Gamma(p)=\gamma_3(p)\).
For any nonzero row write \(k=u k_S k_{\rm good}\), with
\(u\) a unit and \(k_S\) supported on \(S\).  On primary good \(n\),
\begin{equation}\label{own:row-fixed-ray}
 \chi_n(k)=\chi_n(uk_S)\mathcal R(n,k_{\rm good})
                 \prod_{p\mid k_{\rm good}}\chi_p(n)^{v_p(k)}.
\end{equation}
The \(S\)-numerator part has only finitely many choices after
reducing its valuations modulo six.  None of the moving good
prime factors is put into that fixed family.

\subsection{Smooth separation and exponent uniformity}\label{app:smooth}

Let \(W(e^{\mathbf x})\) have fixed compact logarithmic support.
For each fixed \(J\), integration by parts in its Fourier transform gives
\[
 \int_{\mathbb R^d}|\widehat W(\mathbf t)|
               (1+|\mathbf t|)^J\,d\mathbf t
 \ll \sum_{|\beta|\le J+2d+2}\|\partial^\beta W_{\log}\|_\infty .
 \tag{A3}\label{eq:joint-fourier}
\]
The right side is homogeneous in the single profile.
If \(K\) has uniform Euler-derivative bounds and
\(m(\mathbf y)=\prod y_i^{a_i}\), then
\[
 y_i\partial_{y_i}K(Ym(\mathbf y))
     =a_i(x\partial_xK)(x)\big|_{x=Ym(\mathbf y)} .
\]
There is no factor \(Y^J\) from taking \(J\) such derivatives.
The coordinates and the cutoff must be chosen before evaluating
any row or live label.  Applying \eqref{eq:joint-fourier} then gives
one common Fourier density for the whole Hilbert-space norm.
Minkowski is applied to that density, not to a separate density for
each row.  A pure twist \(y^{it_0}\) merely translates its Fourier
transform and has a fixed polynomial cost in \(1+|t_0|\).
A supremum over polynomially many scale boxes or bounded-width
height boxes follows from the elementary parameter Sobolev inequality
and derivatives of these same profiles.

These observations prove the smooth rules of
\cite[Lemma 4.5 and Section 17.3]{QRH} in the forms used here.
For a Gaussian logarithmic profile, a smooth partition into unit
logarithmic annuli has derivative bounds
\(O_J((1+|k|)^J e^{-c k^2})\); summation over \(k\) proves
the Gaussian extension \cite[Lemma 4.6]{QRH}.
In particular there is no hidden loss of the form \(Z^{J\xi}\)
when a derivative order is increased after a length tolerance \(\xi\)
has been fixed.

Finally, on any finite acyclic reduction with uniform weighted
Fourier densities as in \eqref{eq:joint-fourier}, all required
seminorm and height orders can be chosen backwards, starting at
its terminal estimates.  If a child costs
\((1+|\mathbf t|+|\boldsymbol\xi|)^H\), choose the preceding
Fourier moment at that finite order.  Repeating this over the
finite depth proves a fixed polynomial in the initial heights.
This elementary induction is the content of
\cite[Lemmas 17.3--17.4]{QRH}.  It requires verification of the
common densities and finite depth; those verifications for the
two actual recursions are given below.

\paragraph{Conventions for analytic orders.}
For a profile \(W\) on a fixed logarithmic box, \(p_j(W)\) is
the sum of the suprema of its Euler derivatives of total order at
most \(j\).  For a tuple it means \(1+\sum p_j(W_i)\).
For a fixed real strip \(a\le\Re s\le b\), integration by parts
also gives
\begin{equation}\label{eq:pointwise-mellin-tail}
 |\mathcal MW(\sigma+it)|
       \le C_{a,b,N}p_N(W)(1+|t|)^{-N}.
\end{equation}
Consequently
\begin{equation}\label{eq:late-fourier-tail}
 \int_{|t|>T}|\mathcal MW(\sigma+it)|(1+|t|)^J\,dt
       \le C_{a,b,J,N}p_{N+J+2}(W)T^{-N}.
\end{equation}
The first, pointwise estimate is used for a horizontal trace;
the second is used only for an integrated tail.
For a pure norm twist the Mellin transform is translated, so these
bounds are applied to the untwisted profile.  Gaussian logarithmic
profiles satisfy the analogous estimates by their absolutely
convergent annular partition.

\begin{lemma}[Normalized radial kernels]\label{lem:kernel-seminorms}
If \(\Phi(q_z)\) is a fixed radial Schwartz function, its transform
\(K(q_y)\) for the self-dual measure satisfies, for every \(A,j\),
\[
 \sup_{x>0}(1+x)^A|(x\partial_x)^jK(x)|<\infty.
\]
The bound uses finitely many Schwartz seminorms of \(\Phi\).
For any fixed norm monomial \(m(\mathbf y)\) and compact
logarithmic cutoff \(\chi\),
\begin{equation}\label{eq:normalized-kernel-profile}
 p_j\bigl(\chi(\mathbf y)K(Rm(\mathbf y))\bigr)
       \ll_{A,j}\mathfrak S_{A,j}(\Phi)\min(1,R^{-A})
       \qquad(R>0).
\end{equation}
The same assertion with the extra small-argument factor
\(R^{1/4}\) holds for the reflected kernels of
\eqref{eq:uniform-kernel-proof}.
\end{lemma}
\begin{proof}
The planar Fourier transform is a continuous map of the Schwartz
space to itself, by differentiation under its integral and integration
by parts.  A smooth radial Schwartz function is a smooth function
of \(|y|^2\), with bounded Euler derivatives at zero.
This proves the first assertion.  For the second,
\(y_i\partial_{y_i}K(Rm(\mathbf y))\) is the exponent of \(y_i\)
in \(m\), times \(x\partial_xK(x)\) at \(x=Rm(\mathbf y)\).
The product rule and the upper and lower bounds for \(m\) on the
box finish the proof; there is no derivative-order power of \(R\).
The same computation uses the three bounds for the reflected kernel.
\end{proof}

The preceding finite-depth propagation argument also allows positive
discrete label measures: first divide their weights by their already
recorded total counting power.  The remaining mass has only the
assigned small-power loss.  If children have height orders \(H_j\)
and occur to powers \(\theta_j\ge0\), use the common Fourier
density at the finite order \(\sum_j\theta_jH_j\), enlarged for
the child-profile derivatives.  Inserting their bounds gives
exactly the weighted sum of their norm exponents, with this discrete
mass counted once.  Reverse induction on the finite depth proves
this indexed version.  A label averaged inside a child is never
also a label of the external measure.

\section{Completed reflection, with every frequency branch}
\label{app:reflection}

This appendix establishes the specialized input
\cite[Proposition 5.1 and Corollary 14.1]{QRH}.  Its only
automorphic imports are the following unconditional formulas of
Dunn--Radziwi{\l}{\l} \cite{DR}, arXiv version 3:
the hyperbolic action (5.1), Kubota multiplier (5.4), theta
automorphy and infinity coefficients (5.7)--(5.8), cusp expansions
(5.9), (5.12)--(5.16), and Appendix A, rows \(1,10,19\).
Their PDF locations are pp.\ 16--19 and 78.  Their prime
asymptotics, which have a GRH premise, are not used.  Nor do we
apply their twisted-reflection statement with its unit-frequency
restriction.

\subsection{The source coefficients and the primal mask}

Put \(L(t)=\left(\begin{smallmatrix}1&0\\t&1\end{smallmatrix}\right)\),
\(\Gamma_1(3)=\{g\in{\rm SL}_2(\mathcal O):g\equiv I\pmod3\}\),
and \(\Gamma_2=\langle\Gamma_1(3),{\rm SL}_2(\mathbb Z)\rangle\).
The imported smooth theta function satisfies
\[
 \Theta(gw)=\kappa(g)\Theta(w),\qquad
 \kappa\begin{pmatrix}a&b\\c&\delta'\end{pmatrix}
                  =(c/a)_3\quad(g\in\Gamma_1(3),\ c\ne0).
\]
The multiplier is one on \({\rm SL}_2(\mathbb Z)\) and on the
upper or lower translations in \(\Gamma_1(3)\).
It follows that translations with parameter in
\(\mathbb Z+3\mathcal O\) are also invariances.

The three cusp functions needed are
\(\Theta(w),\Theta(L(\omega^2)w),\Theta(L(\omega)w)\).
They are \(F_1,F_{19},F_{10}\) of \cite{DR}.
Indeed
\(\lambda_3-\omega^2=2+3\omega\) and
\(-\lambda_3-\omega=-1-3\omega\).
Also, if \(u=s+t\omega\pmod3\), \(0\le s,t<3\), then
\[
 \begin{pmatrix}u&-1\\1&0\end{pmatrix}
       =\begin{pmatrix}0&-1\\1&0\end{pmatrix}L(-u),\qquad
 -u-t_t\in\mathbb Z+3\mathcal O,\quad
 (t_0,t_1,t_2)=(0,\omega^2,\omega).
\]
Thus these matrices give the same triple.

Write their conjugate expansions as
\[
 \overline{\Theta(H(z,v))}
   =C_H(v)+\sum_{\mu\ne0}d_H(\mu)vK_{1/3}(4\pi|\mu|v)
                                                 \check e(\mu z).
 \tag{B1}\label{eq:theta-cusp}
\]
If \(d_j\) is a coefficient in \cite{DR}, the coefficient after
conjugation is \(\overline{d_j(-\mu)}\), not
\(\overline{d_j(\mu)}\).  In particular the three functions are
\[
 \overline{\tau(-\mu)},\quad
 \omega\overline{\tau_1(-\omega^2\mu)}\check e(\mu),\quad
 \omega^2\overline{\tau_2(-\omega\mu)}\check e(\mu).
\]
The source formulas imply
\[
 \operatorname{supp}d_H\subset
 \{u\lambda_3^k n b^3:k\ge-4,\ n,b\equiv1\pmod3,\ n\ {\rm sf}\},
 \qquad |d_H(\mu)|\le27\,3^{k/6}|b|.
 \tag{B2}\label{eq:theta-support}
\]
To check the constant, the two infinity-coefficient magnitudes are
\(3^{k/6+8/3}|b|\) and \(3^{k/6+3}|b|\);
the other two families have magnitude \(9|b|\) and \(k=-4\).
The normalized squarefree cubic Gauss sum has modulus one by
finite-field orthogonality and CRT.  This includes the primary
prime \(-2\), where one uses the ordinary cubic character, not
a nonexistent good-prime sextic character.  No condition
\((n,b)=1\), or extra \(S\)-mask on these dual indices, occurs in
the source.  Prime valuations make the displayed \(n,b\) unique
when that representation exists.

At infinity the condition \(\lambda_3^3\mu\equiv1\pmod3\)
selects precisely \(\mu=\lambda_3^{-3}nb^3\).
All other ramified powers are divisible by \(\lambda_3\), and
the negative sign gives residue \(-1\).
At this index
\[
 d_I(\lambda_3^{-3}nb^3)
   =3^{5/2}|b|\,\widetilde g_3(n),\qquad
 \widetilde g_3(n)
 =q_n^{-1/2}\sum_{x\bmod n}(x/n)_3\check e(x/n).
 \tag{B3}
\]
On the primal domain \((n,S)=1\), the substitution
\(y=\lambda_3x\) gives
\(\widetilde g_3(n)=\chi_n(\lambda_3)^{-2}\gamma_2(n)\).

Fix a finite family of ray characters \(\Psi_0\), extended by
zero off a fixed modulus \(\mathfrak E\) supported exactly on \(S\).
Choose one primary-period modulus \(L_0\in\mathcal O\), supported
on \(S\), divisible by \(18\) and \(\mathfrak E\), for which
\[
 \phi(x)=
 \begin{cases}
  \chi_x(\lambda_3)^2\Psi_0(x),&x\equiv1\pmod3,\ (x,\mathfrak E)=1,\\
  0,&\text{otherwise}
 \end{cases}
 \tag{B4}
\]
is periodic.  The cubic supplementary laws make
\((\lambda_3/x)_3\) periodic modulo \(9\) on primary \(x\);
ray periodicity supplies the other factors.  The zero branch is
evaluated first.

Let \(\mathcal P\) be any finite set of good primary primes and
\(j_p\in\{0,\ldots,5\}\).
Put \(\Psi=\Psi_0\prod_{p\in\mathcal P}\chi_p^{j_p}\) and
\[
 T(s,\Psi)=
 \left(\sum_n^{\rm sf}
       \gamma_2(n)\overline{\alpha(n)}\Psi(n)q_n^{-s}\right)
 \left(\sum_b\overline{\alpha(b)}^3\Psi(b)^3q_b^{-3s+1/2}\right),
 \tag{B5}\label{eq:completion-product}
\]
where the two primal sums avoid \(\mathfrak E\).
The first mask in (B4) acts on the whole index \(nb^3\).
On the retained domain its multiplier is
\(\chi_n(\lambda_3)^2\Psi(n)\Psi(b)^3\).
The \(b\)-factor in the supplementary symbol has exponent six.
This assertion includes shared moving primes because the zero
powers retain their zeros.  It never evaluates a sextic
denominator on an inadmissible \(S\)-index.

\subsection{Finite Fourier inversion and rational cusp sectors}

Let
\[
 \begin{aligned}
 \widehat\phi(h_0)&=q_{L_0}^{-1}
                       \sum_{x\bmod L_0}\phi(x)e(-h_0x/L_0),\\
 \tau_{p,j}^{\pm}&=q_p^{-1/2}\sum_{x\bmod p}
                    \chi_p(x)^j e(\pm x/p)
                       \qquad(j\not\equiv0\pmod6).
 \end{aligned}
\]
Finite Fourier inversion gives the complete local table
\[
 C_{p,j}(h):=q_p^{-1}\sum_{x\bmod p}\chi_p(x)^j e(-hx/p)
 =\begin{cases}
 q_p^{-1/2}\tau_{p,j}^-\chi_p(h)^{-j},&j\ne0,\ h\ne0,\\
 0,&j\ne0,\ h=0,\\
 -q_p^{-1},&j=0,\ h\ne0,\\
 1-q_p^{-1},&j=0,\ h=0 .
 \end{cases}
 \tag{B6}\label{eq:all-frequencies}
\]
Here \(|\tau_{p,j}^\pm|=1\).  Define
\[
 {\cal F}(z,v)=\sum_{h_0,\,(h_p)}
 \widehat\phi(h_0)\prod_pC_{p,j_p}(h_p)
             \overline{\Theta(z+z_h,v)},\qquad
 z_h=\lambda_3^2\left(h_0/L_0+\sum_p h_p/p\right).
\]
The coefficient multiplier is exactly (B4) times the moving
masks, since
\[
 \check e(\mu z_h)=e\!\left(\lambda_3^3\mu
                                   (h_0/L_0+\sum_p h_p/p)\right).
\]
In particular every \(h_0\), including nonunits, is present.

Call \(p\) active if \(h_p\ne0\), and put \(r=\prod_{\rm active}p\).
Reduce \(\lambda_3^2h_0/L_0=a_F/c_F\).
Normalize \(a_F\) primary if \(\lambda_3\mid c_F\), and
\(c_F\) primary otherwise.  This is possible since the six units
exhaust \((\mathcal O/3\mathcal O)^\times\), and includes
\(h_0=0,(a_F,c_F)=(0,1)\).
Then
\[
 c=c_Fr,\qquad
 a=a_Fr+\lambda_3^2c_F\sum_{p\mid r}(r/p)h_p,\qquad z_h=a/c.
 \tag{B7}
\]
At primes of \(c_F\), the numerator is \(a_Fr\), a unit; at an
active \(p\) it is \(\lambda_3^2c_F(r/p)h_p\), again a unit.
Thus the fraction is reduced in every case.

Set \(M_{\rm ref}=\lambda_3^{12}L_0^4\).
Lift each \(h_p\) to zero modulo \(M_{\rm ref}^2\), and fix the
class of \(r\) modulo \(M_{\rm ref}^2\).
Then \(a\bmod M_{\rm ref}c_F\) is fixed.
Choose \(\delta'\) by CRT as follows:
\[
 \begin{array}{ll}
 \delta'\equiv a^{-1}\pmod{M_{\rm ref}c_F}
       &\text{at prime-power factors supported on }c_F,\\
 \delta'\equiv0\pmod{M_{\rm ref}}
       &\text{at the other primes of }M_{\rm ref},\\
 \delta'\equiv a^{-1}\pmod p& (p\mid r).
 \end{array}
\]
Put \(b=(a\delta'-1)/c\).
The matrix \(g=\left(\begin{smallmatrix}a&b\\c&\delta'\end{smallmatrix}\right)\)
is integral of determinant one.
At fixed primes the residues of all four entries are determined:
\(b\equiv0\pmod{M_{\rm ref}}\) on \(c_F\), and
\(b\equiv-c^{-1}\pmod{M_{\rm ref}}\) on its complement.
If a symbol would involve a zero entry, translating the cusp by
\(\lambda_3^2M_{\rm ref}^2\mathcal O\) and changing \(\delta'\)
by its congruence modulus avoids the finitely many forbidden
values, without changing any theta function or residue data.

There is a factorization \(g=G H\), \(G\in\Gamma_1(3)\),
with \(H\) one of the preceding cusp representatives.  Its
multiplier is
\[
 \kappa(G)=\kappa_F(a/r)_3
          =\kappa_F\prod_{p\mid r}\chi_p(a)^2 ,
 \tag{B8}\label{eq:fixed-cusp-multiplier}
\]
where \(\kappa_F\) depends only on \(h_0,r\bmod M_{\rm ref}^2\).
Here are all three cases, including their calculation.
\begin{enumerate}
\item If \(3\mid c\), take \(H=I\).
Then \(g\equiv I\pmod3\) and
\(\kappa_F=(c_F/a)_3\), by reciprocity of \(a\) with the
primary factors of \(r\).
\item If \(v_{\lambda_3}(c)=1\), take
\(H=L(u)\), \(u=\lambda_3\) or \(-\lambda_3\) chosen so that
\(u\equiv c\pmod3\).  Let \(A=a-ub\).
The determinant gives \(a(c-u\delta')=cA-u\) and
\(bc\equiv-1\pmod a\).
The primary elements \(a,A\) are coprime.  Reciprocity therefore gives
\[
 \kappa(G)=(-u/A)_3(a/A)_3^{-1}
          =(-u/A)_3((c_F/u)/a)_3(a/r)_3.
\]
For the second equality use
\((A/a)_3=(u/a)_3(c/a)_3^{-1}\).
\item If \(\lambda_3\nmid c\), choose \(u\equiv a\pmod3\) and
\(H=\left(\begin{smallmatrix}u&-1\\1&0\end{smallmatrix}\right)\).
The first column of \(gH^{-1}\) is \((-b,-\delta')\).
The primary elements \(-b,c\) satisfy
\(-bc=1-a\delta'\equiv1\pmod9\).
Factor \(\delta'\) into a unit, a \(\lambda_3\)-power, and a
primary factor.  Supplementary laws and reciprocity give
\((\delta'/(-bc))_3=1\).  Since \(a\delta'\equiv1\pmod c\),
\[
 \kappa(G)=(\delta'/(-b))_3=(a/c)_3
                          =(a/c_F)_3(a/r)_3 .
\]
\end{enumerate}
In the first two cases the supplementary symbols depend only on
the fixed residues modulo \(9\), and all other fixed-prime symbols
depend only on the residues at \(c_F\).
In the last case the denominator \(c_F\) is fixed.
This proves the asserted sector dependence of \(\kappa_F\) and
of the chosen cusp function, not just a bound for their magnitudes.

For \(x=\lambda_3^4\mu\), put
\(D_F=\lambda_3^3c_F\),
\(\sigma_p=\lambda_3^2c/p\), and
\(\epsilon_p=-((\lambda_3^3c/p)\sigma_p)^{-1}\pmod p\).
The same CRT gives
\[
 \check e(-\delta'\mu/c)
   =\vartheta(x)\prod_{p\mid r}e(\epsilon_p h_p^{-1}x/p),
 \qquad
 \vartheta(x)=e(-\delta'_F r^{-1}x/D_F).
 \tag{B9}
\]
The congruences already fix \(\delta'_F\bmod D_F\); hence
\(\vartheta\) has one fixed modulus and is fixed in the sector.

Conjugating (B8), and then using (B6), leaves a sum in \(h_p\ne0\).
The substitution \(y=h_p^{-1}\) makes its character exponent \(j_p+2\).
It is a nontrivial Gauss sum unless \(j_p=4\), when it is the
Ramanujan sum \(-1+q_p1_{p\mid x}\).
For \(j_p=0\) its coefficient is the separate \(-q_p^{-1}\)
in (B6).  Consequently the local dual factors are exactly
\[
 B_p(x)=
 \begin{cases}
  \chi_p(x)^{-j_p-2},&j_p\ne0,4,\\
  q_p^{-1/2}(-1+q_p1_{p\mid x}),&j_p=4,\\
  q_p^{-1/2}\chi_p(x)^{-2},&j_p=0\text{ active}.
 \end{cases}
 \tag{B10}
\]
Their scalar phases are
\[
 \omega_{p,j}=
 \begin{cases}
 \tau_{p,j}^-\tau_{p,j+2}^+\chi_p(\epsilon_p)^{-j-2},
                         &j\ne0,4,\\
 \tau_{p,4}^-,&j=4,\\
 -\tau_{p,2}^+\chi_p(\epsilon_p)^{-2},&j=0 .
 \end{cases}
\]
The calculation is valid also at \(p\mid x\).
An inactive zero exponent contributes only \(1-q_p^{-1}\).
An absent prime and an inactive prime have the same active
radical \(r\), so they have the same \(c\), cusp function,
\(\vartheta\), \(\kappa_F\), other local phases, and kernel.
This simultaneous compatibility holds for all \(h_0\), including
nonunit frequencies.

\subsection{Mellin normalization, entire continuation, and kernel bounds}

Define
\[
 J(s)=\int_0^\infty
      \partial_{\bar z}{\cal F}(0,v)v^{2s-1}\,dv.
\]
The horizontal derivative kills each cusp constant term.
The Bessel Mellin identity used in \cite[Lemma 5.3]{DR}
\[
 \int_0^\infty K_{1/3}(y)y^{w-1}\,dy
     =2^{w-2}\Gamma((w-1/3)/2)\Gamma((w+1/3)/2)
\]
and (B3)--(B5) give, on \(\Re s>1\),
\[
 J(s)=\frac{3^{5/2}}4
       \left(\frac{27}{(2\pi)^2}\right)^s
       \Gamma(s+1/3)\Gamma(s+2/3)T(s,\Psi).
 \tag{B11}\label{eq:direct-mellin-proof}
\]
Here \(q_{\lambda_3^{-3}}=1/27\) and
\(\overline{\alpha(\lambda_3^{-3})}=-i\).

For one translated cusp, the hyperbolic action gives exactly
\[
 g^{-1}(z+a/c,v)
 =\left(-\frac{\delta'}c-\frac{\bar z}{c^2(v^2+|z|^2)},
             \frac{v}{q_c(v^2+|z|^2)}\right).
\]
At \(z=0\), the \(\bar z\)-derivative of its first coordinate
is \(-1/(c^2v^2)\).  The conjugate-coordinate and vertical
derivatives are zero.
Applying the chain rule in \eqref{eq:theta-cusp}, followed by
\(v=1/(q_cu)\), gives on \(\Re s<0\) the contribution
\[
 -\frac i4(2\pi)^{2s-2}\bar\alpha(c)^2q_c^{1-2s}
 \bar\kappa(G)\Gamma(4/3-s)\Gamma(5/3-s)
 \sum_{\mu\ne0}d_H(\mu)\alpha(\mu)
            \check e(-\delta'\mu/c)q_\mu^{s-1}.
 \tag{B12}
\]
For a check of the scale, before the Bessel integral the scalar is
\(-2\pi i c^{-2}q_c^{2-2s}\), and that integral is
\(\int_0^\infty u^{2-2s}K_{1/3}(4\pi|\mu|u)\,du\).
This verifies both the sign and the conductor exponent.

The dual series converges absolutely for \(\Re s<0\), by
\eqref{eq:theta-support}: the independent majorants are
\[
 \sum_n^{\rm sf}q_n^{-(1-\Re s)},\quad
 \sum_b q_b^{1/2-3(1-\Re s)},\quad
 \sum_{k\ge-4}3^{k(1/6-(1-\Re s))}.
\]
All converge.  The differentiated theta and every logarithmic
\(v\)-derivative decay exponentially at infinity by the Bessel
expansion.  The cusp chain rule proves the same at zero.
Thus \(J\) is entire and rapidly decreasing on closed vertical
strips.  Dividing \eqref{eq:direct-mellin-proof} by its gamma
factors continues the \emph{angularly weighted} product \(T\)
entirely.  The direct and dual absolutely convergent series
bound \(T\) polynomially on outer lines \(\Re s>1\) and \(\Re s<0\).
Within a strip the crude bound is
\(O((1+|t|)^C e^{\pi|t|})\).
Divide by a sufficiently high power of \(A+s\), multiply by
\(e^{\varepsilon s^2}\), apply the maximum principle on growing
rectangles, and let \(\varepsilon\downarrow0\).
This proves polynomial growth throughout the strip.
There is no assertion here that an unweighted theta series is entire.

For a rapidly decreasing smooth \(V\), set
\[
 \widehat V(t)=\int_0^\infty V(x)x^t\,\frac{dx}{x},\quad
 {\cal K}=(2\pi)^4/27,\quad
 R(t)=\frac{\Gamma(1+t-1/6)\Gamma(1+t+1/6)}
            {\Gamma(1-t-1/6)\Gamma(1-t+1/6)},
\]
\[
 V^\sharp(x)=\frac1{2\pi i}
                  \int_{(\sigma)}\widehat V(-t)R(t)({\cal K}x)^{-t}\,dt,
 \qquad \sigma\ge0 .
\]
Moving the \(s\)-integral for \(T\) to the dual absolutely
convergent line and putting \(s=1/2-t\), the ratio of (B12)
to \eqref{eq:direct-mellin-proof} is
\[
 -\frac i{81}\bar\alpha(c)^2q_\mu^{-1/2}
 R(t)\left(\frac{{\cal K}q_\mu X}{q_c^2}\right)^{-t}.
\]
Summing the already computed local phases proves
\[
 \begin{aligned}
 &\sum_{n,b}\frac{\gamma_2(n)\overline{\alpha(nb^3)}\Psi(nb^3)}
                   {\sqrt{q_n}\,q_b}V(q_nq_b^3/X)\\
 &\quad=\sum_{\rm branches}\zeta
       \sum_{\mu\ne0}\frac{d_H(\mu)\alpha(\mu)
                         \vartheta(\lambda_3^4\mu)}{\sqrt{q_\mu}}\,
       \prod_{\rm active}B_p(\lambda_3^4\mu)
                         V^\sharp(q_\mu X/q_c^2).
 \end{aligned}
 \tag{B13}\label{eq:complete-reflection-proof}
\]
where
\[
 \zeta=-\frac i{81}\bar\alpha(c)^2\widehat\phi(h_0)\bar\kappa_F
           \prod_{\rm inactive}(1-q_p^{-1})
           \prod_{\rm active}\chi_p(\sigma_p)^{-2}\omega_{p,j_p}.
\]
Every dual sum is absolute and \(|\zeta|\le1/81\).
There are at most \(q_{L_0}2^{|\mathcal P|}\) branches, and at
most \(q_{L_0}4^{|\mathcal P|}\) after splitting the Ramanujan
terms.

Here are the uniform smooth estimates needed for the row norm.
Put
\[
 {\mathfrak M}_{A,B}(V)=
 \sup_{-A\le\eta\le1/4}\int_{\mathbb R}
                  (1+|u|)^B|\widehat V(\eta+iu)|\,du.
\]
The first numerator pole of \(R\) is \(-5/6\), and its
denominator reciprocals have no poles.
Stirling gives \(|R(\sigma+iu)|\ll_A(1+|u|)^{4\sigma}\)
on \(-1/4\le\sigma\le A\).
The three kernel lines \(-1/4,0,A\) therefore prove
\[
 |(x\partial_x)^jV^\sharp(x)|
 \ll_{A,j}\min(x^{1/4},(1+x)^{-A})
                      {\mathfrak M}_{A,\lceil4A\rceil+j+2}(V).
 \tag{B14}\label{eq:uniform-kernel-proof}
\]
All shifts are justified by Mellin integration by parts on fixed
strips.  The rapid large-\(x\) bound first proves absolute
convergence of (B13), including when the kernel is subsequently
represented on its zero line.  No conditionally convergent dual
Dirichlet series is shifted termwise.

Logarithmic integration by parts bounds
\({\mathfrak M}_{A,B}\) by
\[
 C_{A,B}\int_0^\infty
 (1+x^{-A}+x^{1/4})
       (|V(x)|+|(x\partial_x)^{B+2}V(x)|)\,\frac{dx}{x}.
\]
For a fixed cutoff in several logarithmic variables,
\eqref{eq:joint-fourier} and \eqref{eq:uniform-kernel-proof}
give a common Fourier density for
\(W(\mathbf y)V^\sharp(Ym(\mathbf y))\), with weighted mass
\[
 \ll \min(Y^{1/4},(1+Y)^{-A})
                 \times\text{a fixed finite seminorm of }V,W.
 \tag{B15}\label{eq:common-reflected-profile}
\]
The statement is uniform for \(Y>0\) and all the frozen labels.
It is homogeneous in this single profile.
Extract its displayed scalar \emph{before} taking a row norm;
squaring then squares that scalar.
A later positive row enlargement acts only on the separated
formula, not on the original kernel.
This proves the uniform assertion needed from
\cite[Lemmas 5.3--5.4]{QRH}, as well as Proposition 5.1.

\subsection{Whole-index marks and phase separation}

Insert a product of marks \(1_{p\mid nb^3}\) with primes
disjoint from the base local set.
The pointwise identity \(1_{p\mid A}=1-\chi_p(A)^0\) and
the simultaneous branch compatibility in (B6)--(B10) show that
an inactive mark contributes \(q_p^{-1}\), while an active
mark contributes \(q_p^{-1/2}\chi_p(\lambda_3^4\mu)^{-2}\)
with scalar \(-\omega_{p,0}\).
An inactive mark is absent from \(c\) and its kernel scale.
If a mark is a base prime, the primal term vanishes and is
removed first.  This proves the whole statement of
\cite[Corollary 14.1]{QRH}, including shared primes of \(n,b\).

Finally \(\epsilon_p=-\lambda_3^{-5}(c/p)^{-2}\) gives
\[
 \chi_p(\sigma_p)^{-2}\omega_{p,j_p}
       =\xi_{p,j_p}\chi_p(c/p)^{2j_p+2},\qquad|\xi_{p,j_p}|=1 ,
\]
where \(\xi_{p,j_p}\) has no other moving-prime or column
dependence.  Two active primes have pair phase
\[
 \chi_p(q)^{2j_p+2}\chi_q(p)^{2j_q+2}
                      =(q/p)_3^{j_p+j_q+2}.
 \tag{B16}\label{eq:pair-phase-proof}
\]
For a residual row prime \(j_p=1\) and an active mark \(j_q=0\)
it is one.  For \(j_p=1,j_q=4\) it need not be one.
Thus \emph{all} nonresidual base primes, especially \(j=4\)
primes, are frozen before the inner norm.

Writing the active radical as \(F_{\rm act}RP\), fix separate
classes of \(R,P\bmod M_{\rm ref}^2\).
This makes the cusp and additive coefficients common without
introducing a row--tuple restriction.
Row--row phases are row scalars; mark--mark phases are tuple
scalars; interactions with \(F_{\rm act}\) separate in the same
way.  After extraction of the frozen dual factors the moving
columns are, including their zeros,
\[
 \chi_R(nb^3)^3=\chi_R(nb)^3,\qquad
 \chi_P(nb^3)^{-2}=\chi_P(n)^{-2}1_{(P,b)=1}.
\]
The remaining pair masks are exactly \((P,R)=1\) and \((P,b)=1\).
The tuple coefficient is independent of \(R,n,b\), and the
dual coefficient is independent of \(R,P\).
This is the coefficient-independence assertion used in the next appendix.

\section{The reflected energy and the additive Gram estimate}\label{app:energy}

\subsection{The imported large sieves and the two collision masks}

The quadratic input is \cite[Definition 1 and Theorem 1.1]{GL}.
In the present field it gives
\[
 \sum_{\substack{a\ {\rm sf,good}\\q_a\le K}}
   \left|\sum_{\substack{b\ {\rm sf,good}\\q_b\le N}}
                    c_b\chi_a(b)^3\right|^2
 \ll_\epsilon(KN)^\epsilon(K+N)\sum_b|c_b|^2 .
 \tag{C1}\label{eq:quad-sieve-proof}
\]
Here is the family-hypothesis verification, including primitivity.
For squarefree \(a\) prime to \(6\), let
\[
 e(a)=\frac{q_a-1}{2}\pmod2,\qquad
 \epsilon_3:(\mathcal O/(\lambda_3))^\times\longrightarrow\{1,-1\}
\]
be the nontrivial character.  For an element \(x\) prime to
\(a\lambda_3\), consider
\[
 \widetilde\chi_a(x)=(x/a)_2\,\epsilon_3(x)^{e(a)} .
\]
For a unit \(u=(-1)^j\omega^k\), the first factor is
\((-1)^{j e(a)}\): powers of \(\omega\) are squares, and the
quadratic symbol of \(-1\) is \((-1)^{(q_a-1)/2}\).
The second factor has the same value, since \(\omega\equiv1
\pmod{\lambda_3}\).  Thus \(\widetilde\chi_a\) is invariant under
all six units and defines an ideal character of trivial infinite type.
Its conductor is exactly \(a\lambda_3^{e(a)}\): the local factors at
the primes of \(a\) are nonprincipal quadratic characters, and the
factor at \(\lambda_3\) is present precisely when \(e(a)=1\).
The Chinese remainder theorem tests each local factor independently.
When \(e(a)=0\), use the primitive character modulo \(a\);
the extra deletion at \(\lambda_3\) is fixed and outside all indices.
For a primary generator \(x\equiv1\pmod3\), one has
\(\epsilon_3(x)=1\), so its value is exactly \(\chi_a(x)^3\).

The primary residue square class modulo \(4\) is a homomorphism
from ideals prime to \(6\) to a fixed finite group.  In each such
class, \(e(a)\) is constant, as is seen from the representatives
\(1,-1,\lambda_3,-\lambda_3\) in \eqref{eq:four-phase}.
For coprime indices in the same class the two factors at
\(\lambda_3\) cancel, and the product has primitive conductor \(ab\).
The reciprocity factor is the finite bicharacter
\(\mathcal R(a,b)\).  These verify the three hypotheses of
\cite[Definition 1]{GL}; its Theorem~1.1, followed by finite class
partition and Hilbert-space duality, proves
\eqref{eq:quad-sieve-proof}.
Fixed \(S\)-parts of a squarefree column are separated as bounded
row scalars before this application.

The cubic input is \cite[Theorem 2]{HB}:
\[
 \sum_{\substack{a\ {\rm primary,sf}\\q_a\le N}}
  \left|\sum_{\substack{b\ {\rm primary,sf}\\q_b\le L}}
                      c_b(b/a)_3\right|^2
 \ll_\epsilon(NL)^\epsilon
       \{N+L+(NL)^{2/3}\}\sum_b|c_b|^2 .
 \tag{C2}\label{eq:cubic-sieve-proof}
\]
The stars in \cite[Theorem 2, pp.\ 3--4]{HB} restrict both
indices to squarefree elements of \(\mathcal O\) congruent to one
modulo three.  The same inequality and indexing are restated in
\cite[Theorem 9.1]{DR}.  Squarefree rational norms and the exclusion
of rational prime factors are not required.
Only the displayed versions of these primary theorems are used.
The family-hypothesis verification above is required independently
of the source's stated large-sieve bound.

We prove the masked hybrid consequence
\[
 \begin{split}
 &\sum_{\substack{k\ {\rm sf,good}\\q_k\ll K}}
 \left|\sum_{\substack{P\ {\rm sf,good}\\q_P\asymp L}}
 \frac{a(P)}{\sqrt{q_P}}1_{(P,k)=1}
 \sum_{\substack{n\ {\rm sf},\,b\\q_n\asymp N,\ q_b\asymp B_1}}
      \beta(n,b)\chi_k(nb)^3\chi_P(n)^{-2}1_{(P,b)=1}\right|^2\\
 &\hspace{15mm}\ll_\epsilon
       (KNB_1L)^\epsilon(K+NB_1)B_1\{N+L+(NL)^{2/3}\},
 \end{split}
 \tag{C3}\label{eq:hybrid-proof}
\]
where \(|a|,|\beta|\le1\), \(a\) is independent of \(k,n,b\),
and \(\beta\) is independent of \(k,P\).
The source dual \(n,b\) may contain \(S\)-primes other than
\(\lambda_3\).  Fixed support restrictions preserving those
independences are allowed.

First consider just the quadratic sum over \(n,b\).
Write \(b=gt^2\), \(g\) squarefree,
\(c=(n,g)\), \(n=cm\), \(g=ch\).
The three ideals \(c,m,h\) are squarefree and pairwise coprime,
but \(g,t\) need not be coprime.
Exactly, with zeros,
\[
 \chi_k(nb)^3=\chi_k(mh)^3\,1_{(k,ct)=1}.
\]
On dyads \(q_c\asymp C,q_g\asymp G,q_t\asymp T\), with
\(B_1\asymp GT^2\), Cauchy in \(t\) costs \(T\).
Freeze \(t\), and expand the entire remaining mask
\(1_{(k,c)=1}=\sum_{r\mid(k,c)}\mu(r)\).
Divisor Cauchy permits a positive sum over \(r\), after which
\(k=rk'\), \(c=rc'\).  There is no remaining \((k',c')\) mask.
The factor \(\chi_r(mh)^3\) is column-only.
Separate the finitely many squarefree \(S\)-parts of \(m,h\),
group their good product as a squarefree column, and apply
\eqref{eq:quad-sieve-proof}.  Cauchy in the \(O(C/q_r)\)
choices of \(c'\) and the divisor-many factorizations of \(mh\)
gives, for arbitrary coefficients \(\beta\),
\[
 (KNB_1)^\epsilon T\sum_{q_t\asymp T}\sum_r
  \left(\frac K{q_r}+\frac{NG}{C^2}\right)\frac C{q_r}
  \sum_{c',m,h}|\beta(rc'm,rc'ht^2)|^2 .
 \tag{C4}\label{eq:completed-quad-proof}
\]
This is \cite[Lemma 5.5]{QRH} in its coefficient-sensitive form.
No positivity argument has deleted a pair-dependent mask.

For \eqref{eq:hybrid-proof}, first expand
\(1_{(P,k)=1}=\sum_{D\mid(P,k)}\mu(D)\).
Divisor Cauchy and \(P=DP'\), \(k=Dk'\) give a positive outer
sum with factor \(q_D^{-1}\).
Put
\[
 a_D(P')=a(DP')1_{(D,P')=1},\quad
 \beta_D(n,b)=\beta(n,b)\chi_D(nb)^3\chi_D(n)^{-2}1_{(D,b)=1}.
\]
They have the requisite independence.  Again there is no residual
\((P',k')\) mask.  Apply \eqref{eq:completed-quad-proof} with
row length \(K/q_D\).  Its positive output contains the square
of
\[
 \sum_{P'}\frac{a_D(P')\chi_{P'}(r)^{-2}}{\sqrt{q_{P'}}}
           \chi_{P'}(c'm)^{-2}1_{(P',ht)=1}.
\]
This factorization keeps all zeros at \(r,c',m\); it follows
directly by multiplicativity from the original
\(1_{(P',rc'ht^2)=1}\).
Only at this positive stage discard \(|\beta_D|\le1\).
For fixed \(h,t,D,r\), the displayed coefficients are independent
of \(j=c'm=n/r\), which is squarefree.
Grouping by \(j\) costs a divisor factor.  Its possible
\(S\)-primes are allowed in \eqref{eq:cubic-sieve-proof}.
The cubic sieve bounds this positive \(j\)-sum by
\[
 (KNB_1L)^\epsilon
 \left\{\frac N{r_0}+\frac L{D_0}
                    +\left(\frac{NL}{r_0D_0}\right)^{2/3}\right\}
\]
on \(q_r\asymp r_0,q_D\asymp D_0\), because
\(\sum_{P'}|a_D(P')|^2/q_{P'}\ll (KNB_1L)^\epsilon\).
The \(h\)-count is \(O(G/C)\).
The factors \(T\sum_t\), \(C/q_r\), and that count together
cost \(B_1/r_0\); there are \(O(r_0)\) choices of \(r\) and
\(\sum_{q_D\asymp D_0}q_D^{-1}=O(1)\).
Increasing both sieve factors by replacing \(r_0,D_0\) by one,
and using \(NG/C^2\ll NB_1\), proves (C3).
All dyadic and divisor factors fit an arbitrary small power.
A nonempty quotient range is bounded below by a fixed constant;
subunit ranges below that are empty.  This proves precisely the
masked norm \cite[Lemma 14.2]{QRH}.

\subsection{The general reflected energy, including growing labels}
\label{app:general-energy}

Let \(k\) be an element row of norm \(O(Z^M)\).
Write its ideal as a powerful part of norm \(\asymp Z^O\),
a squarefree supported part dividing \(\operatorname{rad}(Sf\mathfrak r_\rho)\),
and a squarefree residual part \(R\) of norm \(\asymp Z^H\).
Here \(f\) is a fixed squarefree fourth-power label for this norm,
and \(\rho=1_{(\cdot,\mathfrak r_\rho)=1}\) is a fixed puncture.
Both may grow polynomially in \(Z\), but neither becomes part of
the fixed arithmetic ray modulus.  The original row support gives
\[
 H\le M-O+O(1/\log Z);
 \tag{C5}\label{own:actual-row-dyad}
\]
equality is neither assumed nor needed.

Freeze the first two parts and all nonresidual local primes in
\(\rho(n)\nu(n)\chi_n(k)\chi_n(f)^4\).
Call a good prime active when it occurs in the squarefree moving
part of the reflected denominator.  Combine the exponents of
\(\rho(n)\chi_n(k)\chi_n(f)^4\) modulo six, retaining every zero.
The residual row primes have exponent one.  Define
\(A_0,S_0,N_0,B_0\) as the logarithmic norms of the following
actual prime products, respectively:
active non-slot radical length; small Ramanujan or active
zero-mask length; \(j=4\) divisibility length assigned to the
squarefree dual variable; and its length assigned instead to
the cube variable.
The assignment is the exact disjoint identity
\[
 1_{p\mid nb^3}=1_{p\mid n}+1_{p\nmid n}1_{p\mid b}.
\]
Write \(z_a\) for the nominal length of active slots.
After extraction put \(q_n\asymp Z^v,q_b\asymp Z^t\) and
\(q_{\lambda_3^k}=Z^{e_\lambda}\), \(k\ge-4\).
The actual residual-row annulus is retained.
For completed scale \(X=Z^{N_*}\), define
\[
 T_d=2H+2A_0+2z_a-N_*-N_0-3B_0.
\]
\begin{equation}\label{own:retained-energy-dyad}
 v+3t+e_\lambda\le T_d+\tau_{\rm ref}
            \quad\hbox{for a retained dyad}.
\end{equation}

The central coefficient per active tuple has exponent
\[
 -v/2-t-e_\lambda/3-(S_0+B_0+z_a)/2.
 \tag{C6}\label{eq:energy-coefficient}
\]
Indeed a \(j=4\) divisibility factor assigned to \(n\)
has its \(q_p^{1/2}\) canceled by the extracted inverse root.
A factor assigned to \(b\) leaves \(q_p^{-1/2}\).
The small \(j=4\) and active \(j=0\) factors also leave
\(q_p^{-1/2}\).
The source coefficient contributes
\(q_n^{-1/2}q_b^{-1}q_{\lambda_3^k}^{-1/3}\).
These statements allow \(n,b\) to share primes.

Inactive slots have absolute mass
\(\sum_{p\ {\rm in\ slot}}q_p^{-1}\ll_\epsilon Z^\epsilon\)
and are frozen by Minkowski.  Any resulting collision condition is
then a fixed restriction; it does not enter the arithmetic datum.
For active slots, put \(y_R=q_R/Z^H\), and use the separate
normalized norms of \(n,b\) and every slot.
The kernel is exactly of the form
\[
 \frac{q_\mu X}{q_c^2}
   =Yy_ny_b^3y_R^{-2}\prod_i y_i^{-2},\qquad
 Y=q_{c_F}^{-2}Z^{\,v+3t+e_\lambda-T_d}.
 \tag{C7}
\]
Actual extracted frozen norms are used in \(Y\).
All the other coordinates lie in fixed boxes before positivity.
By \eqref{eq:common-reflected-profile}, its common scalar can
therefore be bounded by
\(Z^{-(T_d-v-3t-e_\lambda)_+/4}\) times a uniform seminorm.
It is taken out of the Hilbert-space vector before its square.

At a separated Fourier mode the phase analysis following
\eqref{eq:pair-phase-proof} permits (C3) with
\((K,N,B_1,L)=(Z^H,Z^v,Z^t,Z^{z_a})\).
Tuple aggregation into \(P\) costs only a divisor factor.
Let \(u=\min(v,z_a,(v+z_a)/3)\).
Combining the hybrid bound with (C6), in which its
\(q_P^{-1/2}\) is already included by (C3), gives exponent
\[
 \max(H,v+t)-S_0-B_0+z_a-u-t-2e_\lambda/3.
\]
The squared kernel scalar subtracts
\((T_d-v-3t-e_\lambda)_+/2\).
Counting the powerful parts adds \(O/2\) only once:
a powerful ideal is a square times a squarefree cube, so there
are \(O_\epsilon(Z^{O/2+\epsilon})\) such ideals.
The supported part is a divisor of
\(\operatorname{rad}(Sf\mathfrak r_\rho)\); its count and all
local branches cost arbitrary small powers, uniformly for
polynomial-size \(f,\mathfrak r_\rho\).
This proves the general exponent
\[
 E_{\rm ref}=O/2+\max(H,v+t)-S_0-B_0+z_a-u-t-2e_\lambda/3
                 -(T_d-v-3t-e_\lambda)_+/2.
 \tag{C8}\label{eq:full-energy-proved}
\]
Before freezing the powerful part its fixed-part version omits
\(O/2\).  This is the full estimate \cite[Lemma 14.3]{QRH},
not merely its low-side specialization.

The discarded region \(v+3t+e_\lambda>T_d+\tau_{\rm ref}\)
is negligible by an absolute tail argument before separation.
On the full source support
\(|d_H(\mu)|/\sqrt{q_\mu}\le27\,3^{4/3}\).
All local factors have polynomial bounds.
For either lattice \(\mathcal O,\lambda_3^{-4}\mathcal O\),
shell counting gives, when \(A>1,U\ge1\),
\[
 \sum_{0\ne\mu,\;a q_\mu>U}|V^\sharp(aq_\mu)|
 \ll_A (1+a^{-1})U^{1-A}
                     {\mathfrak M}_{A,\lceil4A\rceil+2}(V).
\]
On discarded dyads \(U\gg Z^{\tau_{\rm ref}/2}\), after
one fixed threshold for the support constants.
Thus their total is
\(O(Z^{B_{\rm tail}-(A-1)\tau_{\rm ref}/2})\),
times a fixed polynomial in the norm-twist heights.
The polynomial counting exponent \(B_{\rm tail}\) is independent
of \(A\).  Choosing \(A\) after the requested internal saving
sums all tail dyads, including unbounded ramified valuations.
No retained-length count is used to justify this tail.

For the terminal inverse application the following width inequality
is also needed.  The coefficients of
\(2A_0-N_0-S_0-4B_0\), prime by prime, are respectively
\[
 \begin{array}{c|ccccc}
 \text{case}&j\text{ odd or }2&j=0\text{ active}
  &j=4\text{ small or in }n&j=4\text{ in }b&j=0\text{ inactive}\\ \hline
 \text{coefficient}&2&1&1&-2&0 .
 \end{array}
\]
They are bounded by the valuations of
\((q_k^2/q_{k_{\rm pow}})q_fq_{\mathfrak r_\rho}\).
A prime belonging only to \(f\) has exponent four and costs one;
a prime belonging only to the puncture has exponent zero and
costs one.  Intersections can only increase the proposed upper
bound.  Consequently
\[
 T_d-S_0-B_0
       \le2M-O+Q+V-N_*+2z_a+O(1/\log Z),
 \quad Q=\log_Zq_{\mathfrak r_\rho},\quad q_f\asymp Z^V .
 \tag{C9}\label{eq:width-verified}
\]
All constants here depend on the original boxes and fixed \(L_0\),
not on the moving labels.
The assertion includes empty, bounded, and zero-length slot
sublists.  Since \(e_\lambda\ge-4\log3/\log Z\), its negative
part is only a fixed logarithmic offset.
This completes the uniformity needed from Lemma 14.3.

\subsection{Full finite Fourier correlations}\label{sec:finite-correlations}\label{app:gram}

We now allow arbitrary prime powers in a modulus.  The two-argument notation
$G(a,k)$ below is a finite Fourier sum; it is distinct from the one-argument
finite-ray function $G(a)$ of the finite phase formulas in Appendix~\ref{app:conventions}.  Define
\begin{equation}
 G(a,k)=q_a^{-1/2}\sum_{x\bmod a}\chi_a(x)e(kx/a),\qquad G(1,k)=1.
 \label{old-eq:2.4}
\end{equation}
All characters in this subsection have the zero extensions specified in
Appendix~\ref{app:conventions},
and we put $v_p(0)=+\infty$.

\begin{lemma}[Prime-power Fourier sums]\label{lem:prime-power-fourier}
For $P=q_p$ and every integer $a\ge1$,
\begin{equation}
 \begin{aligned}
 |G(p^a,k)|&=P^{(a-1)/2}1_{v_p(k)=a-1},&&6\nmid a,\\
 G(p^a,k)&=P^{a/2}1_{p^a\mid k}
             -P^{a/2-1}1_{p^{a-1}\mid k},&&6\mid a.
 \end{aligned}
 \label{eq:gauss-local}
\end{equation}
\end{lemma}

\begin{proof}
Write a residue modulo $p^a$ as $x_0+py$, with $x_0\bmod p$ and
$y\bmod p^{a-1}$.  The sum over $y$ is zero unless $p^{a-1}\mid k$, and
equals $P^{a-1}$ otherwise.  In the latter case write $k=p^{a-1}k_0$.
The remaining normalized sum is
\[
 P^{a/2-1}\sum_{x_0\bmod p}\chi_p(x_0)^a e(k_0x_0/p).
\]
If $6\nmid a$, its inner sum vanishes when $p\mid k_0$ and otherwise has
absolute value $\sqrt P$.  If $6\mid a$, the inner sum is the sum of the
additive character over the units, namely $P1_{p\mid k_0}-1$.  These are
exactly the two cases in the statement.
\end{proof}

\begin{lemma}[Full correlation and common factors]\label{lem:full-correlation}
For primary moduli $u,v$ outside $S$ and $j\in\mathcal O$, put
\begin{equation}
 F(u,v;j)
 =\sum_{\substack{x\bmod u,\ y\bmod v\\
                    vx-uy\equiv j\pmod{uv}}}
       \chi_u(x)\overline{\chi_v(y)}.
 \label{old-eq:2.11}
\end{equation}
Then
\[
 F(u,v;j)=\frac1{\sqrt{q_uq_v}}\sum_{h\bmod uv}
 G(u,h)\overline{G(v,h)}e(-jh/(uv)).
\]
At zero frequency, $F(u,v;0)=0$ unless $u=v$, and
$F(u,u;0)=\varphi(u)$, where $\varphi(u)$ is the number of units modulo $u$.

Let $C=(u,v)$ and write $u=Cn_1$, $v=Cn_2$, where $(n_1,n_2)=1$.
Then $F(u,v;j)=0$ unless $C\mid j$.  For $j=Ck$,
\begin{equation}
 \begin{aligned}
 F(Cn_1,Cn_2;Ck)
 &=\chi_{n_1}(k)\overline{\chi_{n_2}(-k)}
      \mathcal R(n_1,n_2)L_C(n_1,n_2;k),\\
 L_C(n_1,n_2;k)
 &=\sum_{\substack{x,y\bmod C\\n_2x-n_1y\equiv k\pmod C}}
       \chi_C(x)\overline{\chi_C(y)}
   =\prod_{p^c\parallel C}L_{p^c}.
 \end{aligned}
 \label{eq:correlation-lift}
\end{equation}
For $P=q_p$ and $p\nmid n_1n_2$, the local factor is
\begin{equation}
 L_{p^c}=P^{c-1}\chi_p(n_1/n_2)^c
 \begin{cases}
 P-1,&p\mid k,\\
 -1,&p\nmid k,\ 6\nmid c,\\
 P-2,&p\nmid k,\ 6\mid c.
 \end{cases}
 \label{eq:correlation-local}
\end{equation}
If $p$ divides exactly one of $n_1,n_2$, the local factor is
$P^{c-1}(P-1)1_{6\mid c}1_{p\nmid k}$.  On the genuine residual locus
$(n_1,n_2)=1$, $L_C$ means the congruence sum in
Equation~\eqref{eq:correlation-lift}.  When $L_C$ is used on all residual
pairs, it instead denotes the artificial product extension of these local
formulas, with the local value defined to be zero if $p\mid(n_1,n_2)$.
Outside the genuine residual locus this is not the original congruence sum.
The genuine function and this artificial extension both satisfy
$|L_C|\le q_C$, and the extension is periodic modulo
$\operatorname{rad}C$ in each residual column.
\end{lemma}

\begin{proof}
Expanding both Gauss sums in the displayed Fourier transform leaves
\[
 \frac1{q_uq_v}\sum_{x\bmod u,y\bmod v}\chi_u(x)\overline{\chi_v(y)}
 \sum_{h\bmod uv}e\bigl(h(vx-uy-j)/(uv)\bigr).
\]
Additive orthogonality makes the inner sum $q_uq_v$ when the congruence holds
and zero otherwise.  If $j=0$ and a summand is nonzero, $x$ and $y$ are units
modulo $u$ and $v$.  Reducing the congruence modulo $u$ gives $u\mid v$, and
reducing modulo $v$ gives $v\mid u$.  Since the generators are primary,
$u=v$.  The congruence then says $x=y\bmod u$ and gives $\varphi(u)$.  This
argument uses the zero masks also when a local character power is principal.

The divisibility by $C$ is immediate.  After division by $C$, the congruence
for $j=Ck$ is
\[
 n_2x-n_1y\equiv k\pmod{Cn_1n_2}.
\]
Reduction modulo $n_1$ and $n_2$ gives, with zero values retained,
\[
 \chi_{n_1}(x)=\chi_{n_1}(k)\overline{\chi_{n_1}(n_2)},\qquad
 \overline{\chi_{n_2}(y)}
 =\overline{\chi_{n_2}(-k)}\chi_{n_2}(n_1).
\]
The product of the two unit factors is $\mathcal R(n_1,n_2)$.
It remains to identify the multiplicity of lifts of the congruence modulo
$C$.  This can be checked at each prime.  If $c=v_p(C)$ and neither residual
modulus contains $p$, there is no additional lift.  If, say,
$d=v_p(n_1)>0$ and $p\nmid n_2$, then $y\bmod p^c$ is free and the congruence
uniquely determines $x\bmod p^{c+d}$ from it.  Reduction modulo $p^c$ is
exactly the common congruence in Equation~\eqref{eq:correlation-lift}.
The case $p\mid n_2$ is symmetric, and this argument includes $c=0$.
The Chinese remainder theorem thus gives a bijection with the common
solutions used in $L_C$, even when $C$ meets one residual modulus.

For the local calculation put $k_p=\mathcal O/(p)$ and $A=\chi_p^c$, a
character of $k_p^\times$ extended by zero.  Since at least one of
$n_1,n_2$ is a unit at $p$, each
solution modulo $p$ has $P^{c-1}$ lifts modulo $p^c$.  If both are units and
$p\mid k$, the unit variables are proportional and the field sum is
$A(n_1/n_2)(P-1)$.  If $p\nmid n_1n_2k$, put
$y=(k/n_1)t$ and $x=(k/n_2)(1+t)$.  The field sum becomes
\[
 A(n_1/n_2)\sum_{t\ne0,-1}A((1+t)/t)
 =A(n_1/n_2)\sum_{z\in k_p^\times\setminus\{1\}}A(z).
\]
It is $-A(n_1/n_2)$ for nonprincipal $A$ and
$(P-2)A(n_1/n_2)$ for principal $A$.  If $p\mid n_1$ and $p\nmid n_2$, the
field equation forces $x=k/n_2$.  For $p\mid k$ its character is zero.  For
$p\nmid k$ the remaining sum over $y$ is zero unless $A$ is principal, in
which case it is $P-1$.  The other case is symmetric.  This proves all local
formulas on the genuine residual locus.  Extending them by the stipulated
zero when both residuals meet $p$ gives functions of the residual columns
modulo $p$, each of absolute value at most $P^c$.  Their product proves the
final assertions for the artificial extension as well; no congruence-sum
identity is asserted at a newly added pair.
\end{proof}

For some applications it is preferable to remove all primes common to the
two full moduli, with their entire multiplicities.  The next form leaves one
common row character on each remaining product.

\begin{lemma}[Complete-common-support correlation]
\label{lem:complete-support-correlation}
Suppose $u=Da$, $v=Eb$ are primary and outside $S$, with
$(a,b)=1$ and $(ab,DE)=1$.  Then, for every $j\in\mathcal O$,
\begin{equation}
 \begin{split}
 F(Da,Eb;j)={}&F(D,E;j)\mathcal R(a,E)
       \overline{\mathcal R(b,D)}\mathcal R(a,b)\\
 &\hspace{8mm}\cdot\chi_a(j)\overline{\chi_b(-j)}.
 \end{split}
 \label{eq:correlation-child-character}
\end{equation}
The moduli $D,E$ may contain any prime powers, including shared auxiliary
prime factors.
\end{lemma}

\begin{proof}
At primes of $a$ and $b$, solving the congruence in
Equation~\eqref{old-eq:2.11} gives respectively
$\chi_a(j)\overline{\chi_a(Eb)}$ and
$\overline{\chi_b(-j)}\chi_b(Da)$.  These statements remain valid when a
character of $j$ is zero; only the displayed unit factors are inverted.
At primes of $DE$ change variables $x'=bx\bmod D$, $y'=ay\bmod E$.
The congruence becomes $Ex'-Dy'\equiv j\pmod{DE}$ and its character factor
changes by $\overline{\chi_D(b)}\chi_E(a)$.  The remaining unit factor is
\[
 \frac{\chi_E(a)}{\chi_a(E)}
 \frac{\chi_b(D)}{\chi_D(b)}
 \frac{\chi_b(a)}{\chi_a(b)}
 =\mathcal R(a,E)\overline{\mathcal R(b,D)}\mathcal R(a,b).
\]
All denominators here are symbols of units by the hypotheses.  The Chinese
remainder theorem and the definition of $F(D,E;j)$ complete the proof.
\end{proof}

The hypotheses of Equation~\eqref{eq:correlation-child-character} are not
removed by extending $\mathcal R$ as a bicharacter.  To state exactly the
extension used with M\"obius inversion, fix $D,E,j$ and define
$\widetilde F_{D,E}(a,b;j)$ to be the right side of that equation for all
primary $a,b$ outside $S$ with $(ab,DE)=1$, retaining its zero symbols even
when $(a,b)>1$.  For any finite coefficient array $c(a,b)$ on this locus,
\[
 \sum_{(a,b)=1}c(a,b)F(Da,Eb;j)
 =\sum_{a,b}c(a,b)\widetilde F_{D,E}(a,b;j)
                \sum_{t\mid a,\ t\mid b}\mu(t).
\]
Indeed the full inner divisor sum is $1_{(a,b)=1}$, and the two correlation
expressions agree on that locus.  The value of $\widetilde F_{D,E}$ elsewhere
is artificial, not a formula for $F(Da,Eb;j)$.  In particular the full
divisor sum must be inserted before a factorwise estimate separates the two
residual columns; extra common primes introduced by an individual divisor
term do not become primes of the genuine $D,E$ correlation.

\subsection{The additive Gram estimate}
Now \(A_m\) of the low separation has the exact form
\[
 A_m=Y'^{-3/2}\sum_s c_\sigma(s)
       W_{\rm b,1}(q_s/Y')(q_s/Y')^{-1+i\nu}\,
                        g_{\chi_s}(s,-m),
\]
where \(c_\sigma(s)=1_{s\in\sigma}\chi_s(b_*)/(\tau\xi(s))\)
is a fixed ray coefficient, zero off its original domain.
Majorize \(q_m\ll Q\) by a fixed radial Schwartz weight.
Poisson in \(m\) and Lemma~\ref{lem:full-correlation} give prefactor \(Q/Y'^3\)
and first kernel argument comparable to \(q_Cq_k/P_a\),
where \(P_a=Y'^2/Q\).
The frequency \(k=0\) is the diagonal and contributes \(O(Q/Y')\).

For \(k\ne0\), perform the full residual M\"obius inversion just
described and write \(n_i=dm_i\).
Nonzero terms have \((d,Ck)=1\).
After reciprocity the arithmetic coefficient in \(m_1,m_2\)
is periodic modulo
\[
 \mathfrak r=\operatorname{lcm}
  (\mathfrak l_{\rm fixed},\operatorname{rad}C,\operatorname{rad}k),
 \qquad q_{\mathfrak r}\ll_{\rm fixed}q_Cq_k .
\]
The fixed modulus includes the fixed ray phases and the finite
\(S\)-numerator family of Appendix~\ref{app:conventions};
it is independent of \(C,d,k\).
If some prime outside \(C\mathfrak l_{\rm fixed}\) has valuation
in \(k\) not divisible by six, the complete mean in the first
residual variable is zero.  This follows from its nonprincipal
factor \(\chi_p(m_1)^{v_p(k)}\), with all other factors independent
of \(m_1\bmod p\).  The substitution by \(d\) preserves this
because \((d,k)=1\).

On dyads \(q_C\asymp C_0,q_d\asymp D_0\) and
\(q_Cq_k/P_a\asymp R\), including the initial shell \(R=1\),
put \(N=Y'/(C_0D_0)\), \(\mathcal Q=q_{\mathfrak r}\ll RP_a\).
The two-column profile has every required homogeneous seminorm
bounded by
\(O((1+|\nu|)^J R^{-B_2})\), for a fixed large \(B_2\).
For the mean-free frequencies, four-dimensional lattice Poisson gives
\[
 \left|\sum_{\mathbf m}a_{C,d,k}(\mathbf m)
                            W_R(\mathbf m/\sqrt N)\right|
 \ll (1+|\nu|)^J R^{-B_2}q_C N^2
                                  \min\{1,(\mathcal Q/N)^A\},
 \tag{C13}
\]
with \(A>2\).
For \(N\ge\mathcal Q\), the zero dual vector vanishes, the
normalized Fourier coefficients have modulus at most \(q_C\),
and Fourier decay of order \(2A>4\) proves the saving.
For \(N<\mathcal Q\), absolute counting suffices.
A nonempty subunit annulus is bounded below by a fixed positive
scale, so the same \(O(N)\) one-variable count applies.

There are \(O(C_0D_0RP_a/C_0)\) possible \(C,d,k\) in these
dyads.  Multiplying (C13) by this count and \(Q/Y'^3\) gives
\[
 (1+|\nu|)^J R^{-B_2}\frac Q{Y'}
       \frac{RP_a}{C_0D_0}
       \min\left\{1,\left(\frac{RP_aC_0D_0}{Y'}\right)^A\right\}.
\]
For \(\Lambda=Y'/(RP_a)\), the positive dyadic sum is bounded using
\[
 \sum_{i,j\ge0}2^{-i-j}\min\{1,(2^{i+j}/\Lambda)^A\}
       \ll_A\frac{1+\log(2+\Lambda)}{\Lambda}.
\]
Splitting at \(2^{i+j}\asymp\Lambda\) proves this both for
\(\Lambda\ge1\) and \(\Lambda<1\).
It gives the nonexceptional contribution
\(O((1+|\nu|)^J(Q/Y')(P_a^2/Y')Z^\epsilon)\).

The remaining nonzero frequencies have
\((k)=a_1^6e\), where \(e\) is sixth-power-free and supported on
\(C\) and the fixed primes.
There are at most \(6^{\omega(C)+O_S(1)}\ll_\epsilon q_C^\epsilon\)
choices of \(e\).
For \(q_k\ll RP_a/C_0=K_R\), their count is
\(O_\epsilon(K_R^{1/6}q_C^\epsilon)\).
Using \(|L_C|\le q_C\) and absolute column counting now gives,
before the common prefactor,
\[
 (1+|\nu|)^J R^{-B_2}Z^\epsilon
           Y'^2(RP_a)^{1/6}C_0^{-1/6}D_0^{-1}.
\]
Its dyadic sums converge after reducing the subsidiary power
losses.  Together with the diagonal and the mean-free terms,
this proves
\[
 \sum_{q_m\ll Q}|A_m|^2
 \ll (1+|\nu|)^{J_{\rm Gr}}\frac Q{Y'}
            (1+P_a^{1/6}+P_a^2/Y')Z^\epsilon .
 \tag{C14}\label{eq:Gram-proved}
\]
The argument used exactly the fixed ray coefficients in the
source, not arbitrary row-dependent coefficients.
All scale powers are independent of the chosen fixed smooth
orders.  This proves the remaining low-side input
\cite[Proposition 15.2]{QRH}.

\section{The inverse moments: complete transformations and initialization}\label{app:inverse}

\subsection{Statement of the marked estimate}

Let \(\nu\) be a fixed finite-order ray character whose full
zero-extended defining modulus has prime support in \(S\).  This entire
presentation belongs to the fixed arithmetic datum; a locally frozen
moving zero support is instead retained as a puncture whose radical is
included in the explicit norm bound below.
Fix one sign \(\varepsilon_\chi\in\{1,-1\}\), and set
\[
 \psi_u(n)=\nu(n)\chi_n(u)^{\varepsilon_\chi}.
\]
The only masks in this definition are the fixed exclusions and the zero
extension of \(\chi_n(u)\).
For a smooth annular function \(W\), define
\[
 M_u(Z^r;W)=Z^{-r/2}\sum_n\mu(n)\psi_u(n)W(q_n/Z^r).
\]
For a fixed finite set \(I\) of slots, let \(\mathcal P_i\) be disjoint
sets of primes outside \(S\), with \(q_p\asymp Z^{z_i}\) for
\(p\in\mathcal P_i\).  The sets and the coefficients \(a_i(p)\) are
independent of \(u\), and \(|a_i(p)|\le1\).  Given fixed smooth annular
functions \(W_i\), put
\[
 Q_u=\prod_{i\in I}Z^{-z_i/2}
       \sum_{p\in\mathcal P_i}a_i(p)\psi_u(p)W_i(q_p/Z^{z_i}),
 \qquad z=\sum_{i\in I}z_i.
\]
The empty product is one.
A fixed finite sum of whole products \(M_uQ_u\) is also allowed by
triangle inequality, provided each summand uses one common \(\nu\)
in its inverse and all of its slots.

\begin{lemma}[Marked inverse moment]\label{lem:marked}
Fix bounded ranges for the nonnegative parameters \(m,r,z_i\), a bound
for \(|I|\), and positive constants \(c_1,c_2\).  Suppose
\[
 r+2z\le m-c_1,\qquad 2r+8z\le3m-c_2.
\]
Then, for every \(\epsilon>0\),
\begin{equation}\label{old-eq:4.1}
 \sum_{\substack{u\in\mathcal O\\0<q_u\ll Z^m}}
       |M_u(Z^r;W)Q_u|^2\ll Z^{m+\epsilon}.
\end{equation}
The implied constant may depend on the fixed arithmetic data, \(c_1,c_2\),
the bounded parameter ranges, and finitely many smooth seminorms of the
tests, but is uniform in the prime supports and their bounded coefficients.
The assertion holds for either common sign \(\varepsilon_\chi\), and for
every subcollection of the slots.  Norm twists of the tests have a fixed
polynomial cost in their heights.  There is no lower bound on the individual
slot lengths other than the existence of their stated prime supports, and
no mesh condition on those lengths.
\end{lemma}

The proof passes through the canonical family defined next, in which a
fourth-power residue symbol is averaged over an additional squarefree
ideal.  Lemma~\ref{lem:canonical-moment} is proved by a terminal application
of the reflected-energy estimate in Appendix~\ref{app:general-energy} and a recursive step using two masked
Poisson transformations.  After the principal contributions and tails have
been handled, the retained part is bounded in terms of new admissible sums
of that family with a shorter row range.  A separate initialization using
one masked Poisson transformation then deduces Lemma~\ref{lem:marked}, and
the final subsection gives its longer unmarked consequence for
sixth-power-free rows.

\subsection{The canonical estimate}

Use the squarefree coefficient
\(a(n)=\bar\alpha(n)\gamma_2(n)\nu(n)\rho(n)\), where
\(\rho(n)=1_{(n,\mathfrak r_\rho)=1}\).
A product-form mark on the whole index means exactly
\begin{equation}\label{own:mark-definition}
 \mathfrak d(A)=\sum_{(p_i)\in\prod_{i\in I}\mathcal P_i}
           \prod_{i\in I}a_i(p_i)\,1_{p_i\mid A}.
\end{equation}
Here the original lists are disjoint, the coefficients are bounded and
independent of every current row and label, and the nominal slot sum
is at most \(z_0\).  Individual annular weights can be included in
the \(a_i\).  The extra squarefree fourth-power label is averaged below.

For the current \(f,k\), set
\[
 \Psi_k(n)=\nu(n)\rho(n)\chi_n(k)\chi_n(f)^4 .
\]
For \(V_*(y)=y^{1/2}W(y)\), the completion used here is literally
\begin{equation}\label{eq:marked-completion}
 \mathcal T_{\mathfrak d}(X;k)=
 \sum_{\substack{n\ {\rm sf}\\b}}
 \frac{\bar\alpha(n)\gamma_2(n)\Psi_k(n)
           \bar\alpha(b)^3\Psi_k(b)^3}{\sqrt{q_n}q_b}
             \mathfrak d(nb^3)V_*(q_nq_b^3/X).
\end{equation}
Both ideals avoid \(S\) and may share primes.
This is the whole-index completion proved in Appendix~\ref{app:reflection}.
In this appendix \(\ell,r,M,N,V,B,\delta\) are local length
variables; in particular its \(\ell\) is a cube-factor length,
not the fixed physical slot total.
All statements in this appendix are adapted with attribution from
\cite[Section 17]{QRH}; their proofs and transformed sums are included.

The following is the statement closed by the two Poisson transformations.
Its puncture bound is part of the hypothesis, because a frozen moving
modulus cannot be treated as part of the fixed arithmetic data.
The complete marked recursion is proved below; no recursive-moment
conclusion of the source is used as an axiom.

\begin{lemma}[Canonical marked estimate]\label{lem:canonical-moment}
Fix bounded nonnegative ranges for \(M,N,V,z_0\), a bound for the number
of slots, and \(c_*>0\).  Let \(f\) range over squarefree ideals outside
\(S\) with \(q_f\asymp Z^V\).  Let \(w\) be a nonnegative function of
\(f\) alone satisfying \(w(f)\le C d_{\mathcal O}(f)^C\) for one fixed
\(C\ge1\).  Neither \(C\) nor the choice of \(w\) depends on \(Z\),
the current rows, or any other averaged variable.  The implied constant
below may depend on \(C\), but is uniform over all \(w\) satisfying this
fixed bound.  Let \(\nu\) be a fixed
multiplicative finite ray character whose full zero-extended defining
modulus has prime support in \(S\), and let
\(\rho(n)=1_{(n,\mathfrak r_\rho)=1}\) be a fixed puncture.  Both are
independent of \(k,f\).  Let \(\mathfrak d\) have the product form in
Equation~\eqref{own:mark-definition}, with disjoint fixed prime lists of total
length at most \(z_0\) and bounded coefficients independent of \(k,f\).
No other row-dependent column coefficient or puncture is allowed.
There is no additional residual coefficient \(b(n)\), even one
independent of \(k,f\): arbitrary bounded weights are permitted only
as the stated product of individual prime-slot coefficients.

Put \(F_0=N+V\).  Suppose
\begin{equation}\label{eq:invariant}
 \begin{gathered}
 q_{\mathfrak r_\rho}\le Z^{F_0-M-z_0-c_*},\\
 3M+6z_0+c_*\le4F_0.
 \end{gathered}
\end{equation}
For every smooth annular \(W\) and every \(\epsilon>0\), define
\begin{equation}\label{eq:canonical-energy}
 \begin{split}
 \mathcal E={}&Z^{-V}\sum_f w(f)
 \sum_{\substack{k\in\mathcal O\\0<q_k\ll Z^M}}
 \Biggl|Z^{-N/2}\sum_{n\ {\rm sf}}
 \bar\alpha(n)\gamma_2(n)\nu(n)\rho(n)\\
 &\hspace{24mm}\times
 \chi_n(k)\chi_n(f)^4\mathfrak d(n)W(q_n/Z^N)\Biggr|^2.
 \end{split}
\end{equation}
Then \(\mathcal E\ll Z^{F_0+\epsilon}\).  The constant is uniform in
the moving moduli and the frozen outer ideals within the stated ranges.
It depends on finitely many smooth seminorms and has a fixed polynomial
dependence on separated norm-twist heights.  The same conclusion holds
for all subcollections of the slots with the original cap \(z_0\),
including the empty collection.
\end{lemma}
\subsection{Finite Poisson summation with a mask}

We now prove Lemma~\ref{lem:marked}, using the arithmetic conventions of
Appendix~\ref{app:conventions} and the reflected-energy estimate in Appendix~\ref{app:general-energy} for the
terminal case.  The recursive reduction uses two applications of the
following form of Poisson summation, and the initialization uses one.
Its explicit divisor variable retains every zero extension.

\begin{lemma}[Masked primitive Poisson]\label{lem:masked-poisson}
Let \(\psi\) be a primitive finite character modulo an ideal \(m\),
extended by zero on nonunits, and let \(R\) be any ideal.  Let
\(\Phi(q_z)\) be a radial Schwartz function on \(\mathbb C\), and
let \(\mathcal F\Phi(q_y)\) denote its Fourier transform for the
self-dual measure and the kernel \(e(-zy)\).  For \(K>0\),
\begin{equation}\label{eq:masked-poisson}
 \begin{split}
 &\sum_{k\in\mathcal O}\psi(k)1_{(k,R)=1}\Phi(q_k/K)\\
 &\quad=\frac{K\gamma(\psi;m)}{\sqrt{q_m}}
 \sum_{d\mid\operatorname{rad}R}\frac{\mu(d)\psi(d)}{q_d}
 \sum_{h\in\mathcal O}\bar\psi(h)
       \mathcal F\Phi\!\left(\frac{Kq_h}{q_dq_m}\right),
 \end{split}
\end{equation}
where
\[
 \gamma(\psi;m)=q_m^{-1/2}\sum_{x\bmod m}\psi(x)e(x/m).
\]
For a nonprincipal primitive character, \(|\gamma(\psi;m)|=1\) and the
\(h=0\) term is zero.  For the primitive principal character the
conventions are \(m=1,\ \gamma(1;1)=1,\ \psi(h)=1\), including at
\(h=0\).
\end{lemma}

\begin{proof}
Expand the mask as
\(\sum_{d\mid\operatorname{rad}R,\ d\mid k}\mu(d)\) and write \(k=dk'\).
The scale becomes \(K/q_d\) and the character contributes \(\psi(d)\).
Poisson summation in residue classes modulo \(m\) has prefactor
\(K/(q_dq_m)\).  Its finite transform is
\[
 \sum_{x\bmod m}\psi(x)e(hx/m)
   =\sqrt{q_m}\gamma(\psi;m)\bar\psi(h).
\]
For unit \(h\) this follows by changing variables.  For nonunit \(h\)
the transform is zero by primitivity; for \(m=1\) the stated principal
convention applies.  This proves Equation~\eqref{eq:masked-poisson}.
For a nonprincipal primitive character, finite Parseval gives
\(\sum_h|\sum_x\psi(x)e(hx/m)|^2=q_m\sum_x|\psi(x)|^2\).
Writing \(\phi(m)=|(\mathcal O/m)^\times|\), there are \(\phi(m)\)
unit \(h\)'s, all with magnitude
\(\sqrt{q_m}|\gamma(\psi;m)|\), and
\(\sum_x|\psi(x)|^2=\phi(m)\).  Thus \(|\gamma(\psi;m)|=1\).

A fixed ray restriction can first be expanded into finitely many
characters.  Each is then replaced by its primitive inducing character,
with the removed local zero extensions kept in \(R\); the same formula
applies term by term.  For the principal modulus the absolute sum of
the nonzero frequencies is \(O_\Phi(d_{\mathcal O}(\operatorname{rad}R))\).
Indeed, for all \(A>0\),
\[
 A\sum_{h\ne0}|\mathcal F\Phi(Aq_h)|\ll_\Phi1.
\]
Lattice counting proves this when \(A\le1\), and Schwartz decay proves
it when \(A\ge1\).  Sum this bound with \(A=K/q_d\) over \(d\).
At every use below the relevant nonempty row scale has \(K\ge1\), so
this is also \(O(KZ^\epsilon)\).  When principal nonzero frequencies
are added or removed while separating a nonprincipal sum, they will
carry the same outer row mask as that sum; their absolute contribution
is no larger than this full principal bound.
\end{proof}

\subsection{The canonical reduction and its induction}

We now prove Lemma~\ref{lem:canonical-moment}.  The reflected energy
handles the short completion; the remaining blocks undergo two Poisson
transformations to produce admissible canonical sums with shorter rows.

\begin{proof}[Proof of Lemma~\ref{lem:canonical-moment}]
The order of the analytic choices is organized by
the indexed finite-depth argument in Appendix~\ref{app:smooth}.
Its conclusion will be
used on each positive Cauchy-side sum separately; the proof below verifies
the required common coefficient measures, fixed norm losses, and finite depth.

At each node, every joint \((f,k)\) Hilbert space uses the current label
measure \(w(f)\), and that factor is retained exactly once in every
parent \(f\)-sum until Equation~\eqref{eq:inverse-old-label-fibre}
removes the old label.

At each node keep the parameters \(N,V,M\) fixed, and write
\[
 F=N+V,\qquad Q=\log_Zq_{\mathfrak r_\rho}.
\]
Here \(F\) is the real length sum \(N+V\), so \(F=F_0\) at the starting
node.  The two hypotheses at a node with margin \(c_{\rm node}\) are exactly
\begin{equation}\label{eq:inverse-fixed-invariant}
 F-M-Q-z_0\ge c_{\rm node},\qquad
 4F-3M-6z_0\ge c_{\rm node}.
\end{equation}
In particular \(M+z_0+c_{\rm node}\le F\).  The number \(Q\) is the
actual logarithmic norm of the already fixed puncture radical, not a
dyadic center.  The parameter \(M\) is the fixed logarithmic row length;
the first Poisson test has scale \(Z^M\).

Let \(M_{\max}\) bound the starting row lengths.  Choose
\(0<d\le c_*/200\) and put
\[
 D=\left\lceil\frac{M_{\max}+2}{d}\right\rceil+1.
\]
We use positive parameters \(\eta,\tau,\pi,\tau_{\rm ref},\pi_{\rm ref}\),
chosen in the quantified order at the end of the proof.  Here \(\eta\)
is a localization tolerance bounding only the logarithms of fixed annular
ratios, \(\tau\) is the common tolerance in the two Poisson tail comparisons,
and \(\pi\) bounds the aggregate freely chosen local small-power losses.  At depth
\(h\) set
\[
 c_h=c_*-7h\eta,\qquad c_{\rm node}=c_h,
\]
and impose the row cap \(M\le M_{\max}-hd\).  We prove the estimate by
backwards induction on \(h\le D\).  Every retained row parameter below
is nonnegative.  We shall show that a nonterminal passage decreases it
by at least \(d\), loses at most \(7\eta\) in either required margin,
and increases the energy exponent by at most \(40\eta+\tau+\pi\).
No invariant will be transferred from the fixed \(F=N+V\) to an
actual column norm.

\paragraph{The short completion.}
M\"obius inversion of the cube factor in
Equation~\eqref{eq:marked-completion} gives the exact identity
\begin{equation}\label{eq:marked-completion-inversion}
 \begin{split}
 &Z^{-N/2}\sum_{n\ {\rm sf}}a(n)\chi_n(k)\chi_n(f)^4
       \mathfrak d(n)W(q_n/Z^N)\\
 &\quad=
 \sum_h\frac{\mu(h)\bar\alpha(h)^3\Psi_k(h)^3}{q_h}
       \mathcal T_{\mathfrak d_h}(Z^N/q_h^3;k),
 \qquad \mathfrak d_h(A)=\mathfrak d(h^3A).
 \end{split}
\end{equation}
To check both the mark and the normalization, expand the right side and
put \(c=hb\).  The factor \(V_*(y)=y^{1/2}W(y)\) changes the absolute
coefficient to
\[
 Z^{-N/2}\bar\alpha(c)^3\Psi_k(c)^3q_c^{1/2}
       \sum_{h\mid c}\mu(h).
\]
The mark is \(\mathfrak d(nc^3)\), independent of the divisor \(h\).
The divisor sum vanishes unless \(c=1\), proving the identity.
Multiplicativity of \(\Psi_k\), including its punctures, is used here.

Use one fixed smooth dyadic partition with nonnegative ideal centers,
the unit dyad having center zero.  Suppose \(V<d\) and that the center
\(\ell_1\) of the \(h\)-dyad is less than \(d\).  Freeze \(h\), assign
to it the slots it divides, and retain the others on the completed
index.  A surviving slot then has the individual coefficient
\(a_i(p)1_{p\nmid h}\).  Since \(h\) is fixed, this is a bounded
coefficient independent of the current \(k,f\), on the same nominal
annulus and with the same cap.  It is not the unrestricted original
coefficient; the recursive branch below restores its original slot
lists separately by the child zeros.  For an active subcollection define \(z_a\) to be the sum of its
nominal slot lengths.  Assignments only delete slots, so
\(0\le z_a\le z_0\) exactly; an actual tuple norm is never used as this
cap.  Triangle inequality in the joint \((f,k)\) Hilbert space uses
\(\sum_{q_h\asymp Z^{\ell_1}}q_h^{-1}\ll1\) for the fixed annulus,
apart from the separately chosen divisor loss for assignments.  The
bounded factor \(\Psi_k(h)^3\) is a contraction in that space.

Write \(\widehat h=\log_Zq_h\).  The completion is at the actual scale
\(N_*=N-3\widehat h\).  Its fixed support gives
\(\widehat h\le\ell_1+\eta<d+\eta\) once the fixed annular threshold
specified below is imposed.  In the reflected-energy estimate in Appendix~\ref{app:general-energy}, keep
the actual residual-row dyad and choose the same threshold so that
\begin{equation}\label{eq:inverse-terminal-width}
 H\le M-O+\eta,\qquad
 T_d-S_0-B_0\le2M-O+Q+V-N_*+2z_a+\eta.
\end{equation}
These are Equations~\eqref{own:actual-row-dyad} and
\eqref{eq:width-verified}, not estimates for an enlarged row
range.  Substituting \(N_*=N-3\widehat h\) and the first inequality in
Equation~\eqref{eq:inverse-fixed-invariant} gives
\begin{equation}\label{old-eq:4.9}
 T_d-S_0-B_0
 \le M-O+2z_a-z_0-c_{\rm node}+5d+4\eta.
\end{equation}
Indeed, the additional terms are \(2V+3\widehat h+\eta<5d+4\eta\).

Take the threshold also to ensure \(e_\lambda\ge-\eta\), and discard
the whole reflected tail dyads as in the reflected-energy estimate in Appendix~\ref{app:general-energy}.
Thus every retained dyad satisfies
\(v+3\ell_b+e_\lambda\le T_d+\tau_{\rm ref}\).  The row branch of
Equation~\eqref{eq:full-energy-proved}, after dropping its nonpositive terms, obeys
\[
 E_{\rm ref}\le M+z_a+\tfrac53\eta
       \le F-c_{\rm node}+\tfrac53\eta.
\]
For the column branch with \(u=z_a\), the retained dual bound first gives
\(E_{\rm ref}\le O/2+(T_d-S_0-B_0)+5\eta/3+\tau_{\rm ref}\).
Equation~\eqref{old-eq:4.9} and the first invariant then give
\begin{equation}\label{old-eq:4.10}
 E_{\rm ref}\le F-2c_{\rm node}+5d+\tfrac{17}{3}\eta+\tau_{\rm ref}.
\end{equation}
If \(u\ne z_a\), its definition implies \(u\ge v/2\).  The retained
dual bound now first gives
\[
 E_{\rm ref}\le O/2+(T_d-S_0-B_0)/2+z_a
                    +\tfrac76\eta+\tfrac12\tau_{\rm ref}.
\]
Using Equation~\eqref{old-eq:4.9}, \(z_a\le z_0\), and then the second
invariant yields
\begin{equation}\label{old-eq:4.11}
 E_{\rm ref}\le F-M/4-3c_{\rm node}/4+(5/2)d
                    +\tfrac{19}{6}\eta+\tfrac12\tau_{\rm ref}.
\end{equation}
Use the reflected estimate with aggregate local loss
\(\pi_{\rm ref}\), including its freely chosen exponent, the label divisor
bound, and the finitely many logarithmic sums.  If
\[
 c_{\rm node}\ge c_*/2,\qquad
 d\le c_*/200,\qquad \eta,\tau_{\rm ref},\pi_{\rm ref}\le c_*/1000,
\]
all three bounds, after adding \(\pi_{\rm ref}\), are strictly below
\(F\).  In the last bound \(M\ge0\).  For \(z_0=0\) only \(u=z_a=0\) occurs.
The reflected estimate was applied for each fixed \(f\); the weighted
mass bound \(\sum_f w(f)\ll Z^{V+\pi_{\rm ref}}\) shows that summing its
squared bounds over \(f\) costs at most this amount, already included
in the aggregate loss, and cancels the outside \(Z^{-V}\).
Thus no moving fourth-power label entered the hybrid norm, and these
terminal terms satisfy the canonical estimate.

\paragraph{The remaining terms and their marks.}
At every factorization, assign a slot first to an extracted factor
which its prime divides, and otherwise retain it on the residual
column.  The basic identity for the completion is
\begin{equation}\label{old-eq:4.12}
 1_{p\mid nb^3}=1_{p\mid b}+1_{p\nmid b}1_{p\mid n}.
\end{equation}
More generally, for an ordered list of extracted factors
\(D_1,\ldots,D_h\), the exact priority identity is
\begin{equation}\label{eq:inverse-slot-priority}
 1_{p\mid D_1\cdots D_hn}
 =\sum_{j=1}^h1_{p\mid D_j}\prod_{a<j}1_{p\nmid D_a}
   +1_{p\mid n}\prod_{a\le h}1_{p\nmid D_a}.
\end{equation}
It holds even when extracted factors share primes.  Applying it to each
slot and each copy of a square gives a disjoint partition for each tuple
of primes; the two copies of one slot have independent prime choices.
There are at most \(2^{|I|}\) choices at one single-factor assignment,
and at most \((h+1)^{|I|}\) for the displayed ordered list.  The
assigned slots have divisor-bounded coefficients on the extracted
factor.  Because the original coefficient is
\(\prod_i a_i(p_i)\), removing assigned slots leaves exactly a
subcollection with its original product coefficients.  An ideal which
will remain an averaging variable is not fixed merely because a slot
was assigned to it.

If \(V<d\) and \(\ell_1\ge d\), reopen the completion on the right of
Equation~\eqref{eq:marked-completion-inversion}.  If its cube ideal is
\(h'\), first isolate its dyad with center \(\ell_2\ge0\) and the
\(h\)-dyad by triangle inequality in the joint Hilbert space, before
squaring.  Put \(b=hh'\) and define the fixed centers
\[
 \ell=\ell_1+\ell_2,\qquad r=N-3\ell.
\]
Both copies \(b_1,b_2\) in the resulting square have this same nominal
center \(\ell\); their actual norms need not agree.  The mark is
\(\mathfrak d(nb^3)\), independent of the
choice of the divisor \(h\) of \(b\).  Multiplicativity gives a factor
\(\chi_b(k)^3\); the fourth power in \(\Psi_k(b)^3\) is exactly the
mask \(1_{(b,f)=1}\).  The remaining sum over divisors \(h\) of \(b\)
has absolute value at most a divisor function.  After dyadic
decomposition and separation of \(q_nq_b^3/Z^N\), each squared block
is bounded by a small power times
\begin{equation}\tag{A}\label{eq:canonical-block}
 \begin{split}
 &Z^{-r-2\ell-V}\sum_f w(f)\sum_{k\in\mathcal O}\Phi(q_k/Z^M)
 \Biggl|\sum_{q_b\asymp Z^\ell}\beta(b,f)\chi_b(k)^3\\
 &\hspace{26mm}\times
       \sum_{n\ {\rm sf}}a(n)\chi_n(k)\chi_n(f)^4
       \mathfrak d(nb^3)W(q_n/Z^r)\Biggr|^2 .
 \end{split}
\end{equation}
Here \(\Phi\) is a fixed nonnegative radial Schwartz function
majorizing the original row ball, and
\(|\beta(b,f)|\ll Z^\epsilon\) independently of \(k,n\).  Its
\(f\)-dependence includes the original mask \(1_{(b,f)=1}\);
all other original cube masks are retained.  There is no condition
\((b,n)=1\).  In the reopened case, prior to norm-mode separation,
the cube factor is
\[
 \bar\alpha(b)^3\nu(b)^3\rho(b)1_{(b,f)=1}
       (q_b/Z^\ell)^{1/2}
 \sum_{h\mid b}\mu(h)\,\omega_{\ell_1}(q_h)
                         \omega_{\ell_2}(q_{b/h}),
\]
where \(\omega_{\ell_j}\) are the actual dyadic cutoffs.
The separation of \(W((q_n/Z^r)(q_b/Z^\ell)^3)\) merely
multiplies this by a unit norm power and replaces the annular
column test.  This defines the \(\beta(b,f)\) used in the block,
up to the already isolated Fourier mode, and proves its stated
divisor bound and independence of \(k,n\).  The exponent in front is correct because the original
normalization is \(Z^{-N/2}q_b^{1/2}\asymp Z^{-r/2-\ell}\).
If \(V\ge d\), use the same block with \(\ell=0,\ b=1,\ r=N\).
Thus every remaining block has exactly \(V+\ell\ge d\), by the center
classification and \(\ell_2\ge0\).  The added row \(k=0\)
can only contribute in the principal cases counted below.

\paragraph{The two-transform reduction.}
We prove the following conditional estimate for every remaining block
in Equation~\eqref{eq:canonical-block}. Let \(\epsilon_{\rm child}\ge0\).
Suppose every energy \(\mathcal E'\) of the form
\eqref{eq:canonical-energy}, with all coefficient and support hypotheses
of Lemma~\ref{lem:canonical-moment}, the same fixed slot cap, parameters
\[
 0\le M'\le M-d,\qquad 0\le N',V',\qquad F'=N'+V'\le F+11\eta,
\]
and both margins at least \(c_{\rm node}-7\eta\), satisfies
\(\mathcal E'\ll Z^{F'+\epsilon_{\rm child}}\).
The hypothesis is uniform over the bounded ranges used in this induction,
with finite-seminorm and fixed polynomial-height dependence as in that
lemma. For
\(\eta\le d/16\), the remaining block satisfies
\begin{equation}\label{eq:canonical-reduction-contract}
 \text{block in Equation~\eqref{eq:canonical-block}}
 \ll Z^{F+40\eta+\tau+\pi+\epsilon_{\rm child}},
\end{equation}
with the same kind of uniform dependence. The freely chosen local losses
sum to \(\pi\); the principal terms and discarded tails are included
in this estimate. Thus the reduction supplies exactly the implication
needed for backwards induction. Its proof occupies the two transformations
below: the first produces a positive inverse-polynomial norm, and the second
constructs the smaller canonical energies to which the hypothesis applies.

We next specify the support and localization conventions used in both
transformations.  For an actual ideal or nonzero element \(a\), write
\(\widehat a=\log_Zq_a\).  For a named scalar center, a hat will denote
the actual logarithmic norm of its indicated ideal.  All ideal centers
introduced below are nonnegative centers from the fixed dyadic partition.
The already fixed puncture and the ideal \(q_0\) introduced below are used
at their actual logarithmic norms, without independently rounded centers.

Every individual dyad introduced in the two transformations has a fixed
compact normalized support. If \(a\) is its ideal and \(a_{\rm cen}\)
its named logarithmic center, we use
\begin{equation}\label{eq:inverse-ratio-bounds}
 |\log_Zq_a-a_{\rm cen}|\le\eta.
\end{equation}
This convention applies to the original \(n_i,b_i,f\) windows and to
each extracted-ideal dyad when it is introduced. Errors for products
and quotients are the sums of these individual errors; the fresh residual
windows below have the explicitly stated bounds \(4\eta\), \(6\eta\)
and \(4\eta\). These estimates concern fixed compact normalized-ratio
intervals, not annuli of ratio \(Z^\eta\).

For clarity, the intervals are fixed uniformly through the whole finite
depth as follows.  At each factorization choose fresh individual smooth
cutoffs equal to one on the quotient supports, and include the full
supports of these cutoffs, of the larger cutoffs used in Fourier
separation, and of the enlarged label windows.  Include also the bounded
clipping families used below.  Products and quotients of endpoints locate
the initial quotient supports, such as those of \(b=hh'\), the residual
\(n\), and the child column and label.  Iterating this construction
through \(D\) levels gives a finite collection \(\mathcal I_D\) of compact
normalized-ratio intervals.  If \(L_{\rm win}\) is the maximum absolute
logarithm of their endpoints, the threshold
\(\log Z\ge L_{\rm win}/\eta\) implies
Equation~\eqref{eq:inverse-ratio-bounds} and all the fresh-support bounds
specified below.  After a positive sum is enlarged, its newly added
columns or labels need not satisfy an old parent product identity.  Their
norm bounds at this and the next node come directly from the full
independent fresh windows in \(\mathcal I_D\).  Separated norm powers do
not change these supports.

We use the following specialization of the common joint Fourier calculus in
the smooth calculus of Appendix~\ref{app:smooth}.
Every normalized norm occurring in a coupled smooth factor to be separated
from the columns is a coordinate \(y_1,\ldots,y_d\) of one ambient profile.
A derived outer mask or label cutoff retained in the weight remains
outside and needs no profile coordinate.  Keep the current
individual dyad cutoffs and fresh column cutoffs outside the inversion,
until the weighted Cauchy inequality where they are used.  Multiply the
coupled profile by larger individual cutoffs \(\Omega_j\) equal to one
on the full supports of those retained cutoffs.  For the resulting
compact profile \(\mathfrak H\), define before summing any current row or
label
\begin{equation}\label{eq:inverse-log-fourier}
 \begin{split}
 \widehat{\mathfrak H}(\boldsymbol t)
   &=\int_{\mathbb R^d}\mathfrak H(e^{u_1},\ldots,e^{u_d})
             e^{-i\boldsymbol t\cdot\boldsymbol u}\,d\boldsymbol u,\\
 \mathfrak H(\boldsymbol y)
   &=(2\pi)^{-d}\int_{\mathbb R^d}\widehat{\mathfrak H}(\boldsymbol t)
          \prod_{j=1}^dy_j^{it_j}\,d\boldsymbol t.
 \end{split}
\end{equation}
Arithmetic relations among the norms restrict evaluation points of this
identity, not its density.  Outer modes stay in the outer weight; for
two column coordinates the first test receives \(y_1^{it_1}\) and the
second receives \(y_2^{-it_2}\), whose conjugate supplies
\(y_2^{it_2}\).  All normalized real inverse-root powers are included
in the coupled profile, so the final column tests contain no second copy.

This convention also makes the uniformity quantitative.  Write
\(D_j=y_j\partial_{y_j}\).  On a fixed log rectangle, integration by
parts with \((1-\Delta)^{m_0}\), for \(2m_0>J+d\), gives
\begin{equation}\label{eq:inverse-fourier-moment}
 \int_{\mathbb R^d}|\widehat{\mathfrak H}(\boldsymbol t)|
        (1+|\boldsymbol t|)^J\,d\boldsymbol t
 \ll_{J,d,\mathrm{box}}\max_{|\boldsymbol a|\le2m_0}
             \|D^{\boldsymbol a}\mathfrak H\|_\infty.
\end{equation}
For a fixed radial Schwartz Fourier kernel \(K\), every
\(\sup_{x\ge0}|(x\partial_x)^jK(x)|\) is finite, and
\[
 D^{\boldsymbol a}K\!\left(A\prod_jy_j^{e_j}\right)
 =\left(\prod_je_j^{a_j}\right)
   (x\partial_x)^{|\boldsymbol a|}K(x)
       \big|_{x=A\prod_jy_j^{e_j}}.
\]
Thus the separating seminorms are uniform for every scalar \(A>0\),
with no derivative-order power of \(Z\).  Normalized real powers have
bounded Euler derivatives on the fixed boxes, and inherited norm twists
have only a fixed polynomial height cost.  No sharp cutoff in a
column-dependent kernel ratio is put inside such a profile.
For the radial Fourier kernels in the Poisson steps, these are the
normalized-kernel bounds of Lemma~\ref{lem:kernel-seminorms} and
Equation~\eqref{eq:normalized-kernel-profile}.

The raw tail estimate we shall use is, for \(Y\ge1\) and \(A>1\),
\begin{equation}\label{eq:inverse-raw-tail}
 \sum_{h\ne0:\,a q_h>Y}|K(aq_h)|
 \ll_A(1+a^{-1})Y^{1-A}\qquad(a>0).
\end{equation}
Indeed, the shell \(2^jY<a q_h\le2^{j+1}Y\) contains
\(O(1+2^jY/a)\) lattice points and the kernel is
\(O_A((2^jY)^{-A})\) there.  Summing the two geometric series proves the
display; the same argument applies to the fixed lattice
\(\lambda_3^{-4}\mathcal O\).  The proof uses only the large-argument
bound \(|K(x)|\ll_Ax^{-A}\) for \(x\ge1\), so it also applies to the
reflected kernel on such a tail.  Below an actual-ratio inequality is used only
to show that the complement of a sector-fixed outer row ball is contained
in this tail for each supported raw column pair.  That complement is
removed before off-coprime extension and before Fourier absolutization.
The full smooth kernel is retained inside the ball, whose nonsmooth mask
is kept outside Fourier inversion until the relevant weighted Cauchy
inequality.  The tail orders and the crude raw counts are fixed at the end.

\paragraph{The first Poisson transformation.}
Expand the square in Equation~\eqref{eq:canonical-block}.  Put
\[
 \mathfrak B=\operatorname{rad}(b_1b_2),\qquad
 n_i=\mathfrak A_i C u_i\quad(i=1,2),
\]
where \(\mathfrak A_i=(n_i,\mathfrak B)\), and \(C\) is the gcd of
\(n_1/\mathfrak A_1\) and \(n_2/\mathfrak A_2\).  Then
\(\mathfrak A_i\mid\mathfrak B\), while \(u_1,u_2\) are squarefree,
coprime to one another, and prime to \(C\mathfrak B\).  Let \(A_i,B\)
be the fixed centers of \(\mathfrak A_i,C\), and write
\(\widehat A_i=\log_Zq_{\mathfrak A_i}\), \(\widehat B=\log_Zq_C\).
For \(p\mid\mathfrak B\), define
\[
 a_{ip}=1_{p\mid\mathfrak A_i},\qquad
 \pi_p=v_p(b_1b_2)\bmod2,\qquad
 t_p=a_{1p}-a_{2p}+3\pi_p\bmod6,
\]
and put \(R_1=\prod_{t_p\ne0}p\).  Its center is \(R\), and
\(\widehat R=\log_Zq_{R_1}\).
The row character and its remaining zero mask are exactly
\begin{equation}\label{eq:first-masked-data}
 \psi_1=\chi_{u_1}\bar\chi_{u_2}
               \prod_{t_p\ne0}\chi_p^{t_p},\qquad
 m_1=u_1u_2R_1,\qquad
 R_{\rm mask}=C\mathfrak B/R_1.
\end{equation}
Indeed, the exponents at \(u_1,u_2\) are \(1,-1\); the common
off-\(\mathfrak B\) factor \(C\) has exponent zero but retains its
mask; and the exponent at \(p\mid\mathfrak B\) is \(t_p\).  Each
nonzero local power is nonprincipal and primitive modulo \(p\), with
disjoint supports.  Thus \(\psi_1\) is primitive unless \(m_1=1\).
Every original column mask and the two factors
\(\beta(b_1,f)\overline{\beta(b_2,f)}\) remain in the expression.

Apply Lemma~\ref{lem:masked-poisson} with \(K=Z^M\), and denote its
divisor by \(d_k\mid C\mathfrak B/R_1\), with center \(\delta\) and
\(\widehat\delta=\log_Zq_{d_k}\).  On this genuine coprime expression,
the actual conductor length is
\begin{equation}\label{eq:first-poisson-lengths}
 \widehat L_1=\widehat n_1+\widehat n_2-\widehat A_1-\widehat A_2
                    -2\widehat B+\widehat R.
\end{equation}
The actual kernel argument is \(Z^Mq_h/(q_{d_k}q_{m_1})\), and the
Poisson prefactor has exponent \(M-\widehat\delta-\widehat L_1/2\).
We have not yet truncated or Fourier-separated this kernel.

The full nonprincipal expression at this point is
\begin{equation}\label{own:inverse-first-raw}
\begin{split}
 Z^{-r-2\ell-V}
 \sum_{\mathscr R_1}
 &w(f)\beta(b_1,f)\overline{\beta(b_2,f)}
 a(n_1)\overline{a(n_2)}
 \chi_{n_1}(f)^4\overline{\chi_{n_2}(f)^4}\\
 {}\times&
 \mathfrak d(n_1b_1^3)\overline{\mathfrak d(n_2b_2^3)}
 W(q_{n_1}/Z^r)\overline{W(q_{n_2}/Z^r)}\\
 {}\times&
 \frac{Z^M\gamma(\psi_1;m_1)}{\sqrt{q_{m_1}}}
 \sum_{d_k\mid C\mathfrak B/R_1}
       \frac{\mu(d_k)\psi_1(d_k)}{q_{d_k}}
 \sum_{h\ne0}\overline{\psi_1(h)}
       \mathcal F\Phi\!\left(\frac{Z^Mq_h}{q_{d_k}q_{m_1}}\right).
\end{split}
\end{equation}
Here \(\mathscr R_1\) is exactly the set of good \(f\) and
\(b_1,b_2\) in their specified windows and squarefree
\(n_i=\mathfrak A_iCu_i\) in the original column windows, with
\(\mathfrak A_i\mid\mathfrak B\), \((C,\mathfrak B)=1\),
\((u_1,u_2)=1\), \((u_1u_2,C\mathfrak B)=1\), and
\(m_1=u_1u_2R_1\ne1\).  These conditions make the reconstruction
from \(n_1,n_2,b_1,b_2\) unique.  All old zeros in \(a,\beta\)
remain in the displayed factors.  Thus this formula follows directly
by expanding the square and applying the masked Poisson identity
to its exact row character.  It is before frequency truncation,
off-coprime extension, or either Cauchy inequality.

If \(m_1=1\), then \(u_1=u_2=1\) and \(t_p=0\) for every
\(p\mid\mathfrak B\).  The latter condition forces \(\pi_p=0\)
and \(a_{1p}=a_{2p}\).  Hence \(n_1=n_2\) and \(b_1b_2\) is a
square.  Even counting all \(O(Z^{2\ell+\epsilon})\) pairs \(b_1,b_2\),
all \(O(Z^{r+\epsilon})\) equal columns, all \(O(Z^{V+\epsilon})\)
labels, and \(O(Z^M)\) rows, the normalization in
Equation~\eqref{eq:canonical-block} gives \(O(Z^{M+\epsilon})\).
If \(r<0\) but its fixed annular column window is nonempty, then
\(Z^{-r}\) is bounded by its fixed upper endpoint, so its ideal count
is still \(O(Z^r)\) with a fixed constant.
Marks and the retained masks do not increase this bound beyond a
small power.  When the nonzero principal terms are restored below, they
will carry the same outer ball mask as the nonprincipal terms; the last
part of Lemma~\ref{lem:masked-poisson} bounds them at this same cost.

We next identify its column coefficients.  For coprime squarefree primary \(a,b\),
CRT and reciprocity give
\[
 \gamma_2(ab)=\gamma_2(a)\gamma_2(b)
       \chi_a(b)^2\chi_b(a)^2
 =\gamma_2(a)\gamma_2(b)\chi_b(a)^4.
\]
In the first Poisson root, the factors involving \(u_1\), including
the divisor and frequency, are
\[
 \gamma_1(u_1)\overline{\chi_{u_1}(h)}
 \chi_{u_1}(d_kR_1)\prod_{t_p\ne0}\chi_{u_1}(p^{t_p}),
\]
up to a fixed-ray factor.  This follows by writing the CRT factors
between \(u_1\) and \(p\) as
\(\chi_{u_1}(p)\chi_p(u_1)^{t_p}\) and using reciprocity.  On the
conjugated second side \(t_p\) is replaced by \(-t_p\).  The cross
factor between \(u_1,u_2\) is the reciprocity factor
\(\mathcal R(u_1,u_2)\), not a quotient of zero-extended symbols.

Put \(P_T=\prod_{t_p\ne0}p^{t_p}\), regarded in the fixed ray group.
Equation~\eqref{eq:signal-checked} and
Equation~\eqref{eq:G-quadratic-refinement} show that the remaining
two-column ray factor is
\[
 G(u_1u_2^{-1})\mathcal R(u_1u_2^{-1},P_T).
\]
For example, before placing factors inside the conjugated second
polynomial, the raw second Gauss-signal factor is
\(\mu(u_2)\chi_{u_2}(-1)\overline{G(u_2)}\).  Its corresponding
factor written inside that polynomial is
\(\mu(u_2)\chi_{u_2}(-1)G(u_2)\).  Also
\(G(u_2^{-1})=\chi_{u_2}(-1)\overline{G(u_2)}\).
These identities give the displayed quotient.  Fix the ray class of
\(P_T\) and Fourier-expand the displayed function on the finite ray
group.  Its factors on each side are genuine multiplicative ray
characters, denoted by \(\nu_i\).  Their number and coefficient norm
depend only on the fixed ray group.

Define
\[
 E=\prod_{\pi_p=1}p^{a_{1p}+a_{2p}},\qquad
 \widetilde h=hf^2E.
\]
The original fourth power at \(f\) equals
\(\overline{\chi_{u_i}(f^2)}\), including zeros.  Assign the slots
first to \(b_i,\mathfrak A_i,C\).  The remaining local factor at
\(p\mid\mathfrak B\) on side \(i\) has exponent
\[
 4a_{ip}+1_{t_p\ne0}(1\mathbin{\pm}t_p)
          +\pi_p(a_{1p}+a_{2p})\pmod6.
\]
The plus sign belongs to side one.  Direct reduction gives the same
answer on both sides:
\[
 \begin{array}{c|cc|c|c}
 \pi_p&a_{1p}&a_{2p}&t_p&\text{exponent on both sides}\\ \hline
 0&0&0&0&0\\
 0&1&0&1&0\\
 0&0&1&5&0\\
 0&1&1&0&4\\
 1&0&0&3&4\\
 1&1&0&4&4\\
 1&0&1&2&4\\
 1&1&1&3&4
 \end{array}
\]
Define \(\xi(n)\) as the product over \(p\mid\mathfrak B\) of
\(\chi_n(p)\) raised to the corresponding exponent in the last column.
An exponent zero still denotes the puncture at \(p\).  The calculation
proves that \(\xi\) is common to both sides and independent of \(h,f,C\).
The nonprincipal coefficient on side \(i\), after these assignments, is
\begin{equation}\tag{C}\label{eq:first-poisson-column}
 \mu(u_i)\nu_i(u_i)\rho(u_i)
 \overline{\chi_{u_i}(\widetilde h)}
 \chi_{u_i}(C)^4\chi_{u_i}(d_k)\xi(u_i)\mathfrak d_i(u_i).
\end{equation}
It also shows that every character involving the moving
\(b_i,\mathfrak A_i\) is in either \(\xi\) or \(\widetilde h\);
their other factors are outer coefficients.

\paragraph{An explicit outer coefficient for the first transform.}
For a term on the genuine coprime support of
\eqref{own:inverse-first-raw}, write \(c_i=\mathfrak A_iC\),
\(d=d_k\), and
\[
 \xi_R=\prod_{p\mid R_1}\chi_p^{t_p},\qquad
 \Gamma_R=q_{R_1}^{-1/2}
             \sum_{x\bmod R_1}\xi_R(x)e(x/R_1).
\]
When \(R_1=1\), take \(\xi_R=1\), including at zero, and
\(\Gamma_R=1\).  Keep the coefficient
\(a(c)=\bar\alpha(c)\gamma_2(c)\nu(c)\rho(c)\) from the
canonical statement.  The complete column-independent arithmetic
coefficient, before the divisor factor \(\mu(d)\), is
\begin{equation}\label{eq:inverse-first-outer-explicit}
\begin{split}
 \mathcal C^{\rm out}_1={}&w(f)\beta(b_1,f)\overline{\beta(b_2,f)}
       a(c_1)\overline{a(c_2)}
       \chi_{c_1}(f)^4\overline{\chi_{c_2}(f)^4}\\
 &\times\Gamma_R\,\xi_R(d)\overline{\xi_R(h)} .
\end{split}
\end{equation}
This definition uses products only.  It retains all zeros at
\(c_i,f,b_i\) and the primitive frequency zero at \(R_1\).
For a squarefree residual \(u\) set
\[
 \mathcal K(u)=\mu(u)\nu(u)\rho(u)
       \overline{\chi_u(hf^2E)}\chi_u(C)^4\chi_u(d)\xi(u),
 \qquad
 \mathcal G_P(v)=G(v)\mathcal R(v,P_T).
\]
Here \(v\) is a class in the fixed ray group, so its inverse is
always defined.  The exact identity before slot allocation is
\begin{equation}\label{eq:inverse-first-arithmetic-identity}
\begin{split}
 &w(f)\beta(b_1,f)\overline{\beta(b_2,f)}
  a(n_1)\overline{a(n_2)}
  \chi_{n_1}(f)^4\overline{\chi_{n_2}(f)^4}
  \gamma(\psi_1;m_1)\psi_1(d)\overline{\psi_1(h)}\\
 &\hspace{8mm}=
 \mathcal C^{\rm out}_1\,
 \mathcal G_P([u_1][u_2]^{-1})
       \mathcal K(u_1)\overline{\mathcal K(u_2)} .
\end{split}
\end{equation}
The marks, windows, and positive norm denominators have not been
included on either side of this arithmetic identity; they remain
exactly those of \eqref{own:inverse-first-raw}.

Here is a direct verification for all moduli, not a finite-prime test.
CRT, on the pairwise coprime moduli \(u_1,u_2,R_1\), gives
\[
\begin{split}
 \gamma(\psi_1;m_1)
 ={}&\gamma_1(u_1)\gamma_{-1}(u_2)\Gamma_R
       \mathcal R(u_1,u_2)\\
 &\times\chi_{u_1}(R_1)\overline{\chi_{u_2}(R_1)}
             \xi_R(u_1)\xi_R(u_2).
\end{split}
\]
Also
\(a(c_i u_i)=a(c_i)a(u_i)\chi_{u_i}(c_i)^4\).
The factors
\(a(u_1)\gamma_1(u_1)\) and
\(\overline{a(u_2)}\gamma_{-1}(u_2)\) are, respectively,
\[
 \mu(u_1)\nu(u_1)\rho(u_1)G(u_1),\qquad
 \mu(u_2)\overline{\nu(u_2)}\rho(u_2)
                \chi_{u_2}(-1)\overline{G(u_2)}.
\]
The identities remain true if \(\rho\) vanishes.
Reciprocity gives
\(\xi_R(u)=\mathcal R(u,P_T)\chi_u(P_T)\).
Consequently the combined ray factor is exactly
\(\mathcal G_P([u_1][u_2]^{-1})\), by
\eqref{eq:G-quadratic-refinement}.
The factors at \(d,h,f,C\) are those in \(\mathcal K\), with
\(\chi_u(f)^4=\overline{\chi_u(f^2)}\).
At every prime in \(\mathfrak B\), moving
\(\overline{\chi_u(E)}\) into the frequency changes the residual
exponent to
\[
 4a_{ip}+1_{t_p\ne0}(1\mathbin{\pm}t_p)
                       +\pi_p(a_{1p}+a_{2p})\pmod6.
\]
The eight displayed states give exactly \(\xi(u)\) on both sides.
This proves \eqref{eq:inverse-first-arithmetic-identity}.
No power cancellation removed a nonunit zero.

For additional specificity, its ray expansion is
\[
 \mathcal G_P(v)=\sum_{\vartheta\in\widehat T}
          c_\vartheta(P_T)\vartheta(v),\qquad
 c_\vartheta(P_T)=|T|^{-1}\sum_{v\in T}
                 \mathcal G_P(v)\overline{\vartheta(v)} .
\]
Thus in a fixed first ray summand one can take
\(\nu_1=\nu_2=\nu\vartheta\).  The notation \(\nu_i\) merely permits
the two later tests and slot subcollections to differ.  The factor
\(c_\vartheta(P_T)\) stays outer.  On the off-coprime extension used
below, all expressions are defined by these separate factors, not
by the primitive Gauss formula; the full M\"obius indicator then
recovers precisely the genuine coprime domain.

The first transform has produced the inverse-type column coefficient
in Equation~\eqref{eq:first-poisson-column}.  We next collect its
fourth-power factors into one squarefree label before forming the
positive row norm for the second transform.

\paragraph{Retaining the fourth-power label from the cubes.}
Write uniquely \(b_1b_2=q^2s\) with \(s\) squarefree, and define
\[
 J_2=\prod_{\substack{\pi_p=0\\a_{1p}=a_{2p}=1}}p,\qquad
 J=sJ_2,\qquad q=J_2q_0.
\]
Every prime of \(J_2\) has positive even valuation in \(b_1b_2\),
so \(J_2\mid q\).  The ideals \(s,J_2\) are coprime and squarefree;
hence \(J\) is squarefree.  Let \(j\ge0\) be the center of \(J\), and put
\[
 \widehat s=\log_Zq_s,\quad \widehat j_2=\log_Zq_{J_2},\quad
 \widehat j=\log_Zq_J=\widehat s+\widehat j_2,\quad
 \widehat\ell_{\rm pair}=(\widehat b_1+\widehat b_2)/2.
\]
The ideal identity \(b_1b_2=q_0^2J_2^2s\) gives exactly
\begin{equation}\label{eq:inverse-cube-actual}
 \widehat c:=\log_Zq_{q_0}
   =\widehat\ell_{\rm pair}-\widehat s/2-\widehat j_2\ge0,
 \qquad
 4\widehat\ell_{\rm pair}-2\widehat s+\widehat j_2
   =4\widehat c+5\widehat j_2\ge0.
\end{equation}
The table above now proves the exact zero-extended identity
\begin{equation}\label{eq:retained-cube-label}
 \xi(n)=\chi_n(J)^4\,1_{(n,\operatorname{rad}q_0)=1}.
\end{equation}
Indeed, the primes with exponent four are exactly those of \(sJ_2\);
every other prime of \(\mathfrak B\) divides \(q_0\).  A prime in
both \(J\) and \(q_0\) only repeats the zero already supplied by
\(\chi_n(J)^4\).  We will fix \(q_0\) but keep \(J\) in a later
average.  The modulus \(q_0\) contributes a puncture and is not part
of the fixed ray group.
For fixed \(J,q_0\), the choices of \(s,J_2\), the factorizations
\(b_1b_2=q_0^2J_2^2s\), and the ideals \(\mathfrak A_i\) have only
divisor multiplicity.

We now localize the genuine first Poisson expression, before extending its
coprime support or taking absolute Fourier integrals.  The local table gives
the exact actual-norm identity
\[
 \widehat R-\widehat A_1-\widehat A_2+\widehat E
       =\widehat s-2\widehat j_2.
\]
For \(y=hf^2E\), the implication
\(Z^Mq_h/(q_{d_k}q_{m_1})\le Z^\tau\) and
Equation~\eqref{eq:first-poisson-lengths} give
\[
 \begin{split}
 \widehat y
 &\le \widehat n_1+\widehat n_2-2\widehat B-M
       +\widehat s-2\widehat j_2+\widehat\delta+2\widehat f+\tau\\
 &\le \widehat n_1+\widehat n_2-2\widehat B-M
       +4\widehat\ell_{\rm pair}-\widehat j
       +\widehat\delta+2\widehat f+\tau.
 \end{split}
\]
The second inequality is Equation~\eqref{eq:inverse-cube-actual}.
Define the formal center and the enclosing outer row scale by
\begin{equation}\label{eq:enlarged-first-frequency}
 H_c=2r-2B-M+4\ell+2V+\delta-j,\qquad
 H_{\rm use}=H_c+12\eta+\tau.
\end{equation}
Relative to \(H_c\), the last actual upper bound has error
\[
 (\widehat n_1-r)+(\widehat n_2-r)-2(\widehat B-B)
 +2(\widehat b_1-\ell)+2(\widehat b_2-\ell)
 -(\widehat j-j)+(\widehat\delta-\delta)+2(\widehat f-V),
\]
whose positive maximum under Equation~\eqref{eq:inverse-ratio-bounds} is
\(12\eta\).  Thus the actual-ratio inequality implies
\(q_y\le Z^{H_{\rm use}}\) for every supported raw column pair.

On that raw expression insert only the outer mask
\[
 \mathcal B_1(h)=1_{0<q_{hf^2E}\le Z^{H_{\rm use}}}.
\]
For each supported pair the complement lies in the kernel tail
\[
 \frac{Z^Mq_h}{q_{d_k}q_{m_1}}>Z^\tau.
\]
Apply
Equation~\eqref{eq:inverse-raw-tail} to discard that complement, with the
raw counts and order fixed below.  No indicator of the ratio inequality
is inserted.  Inside \(\mathcal B_1=1\) retain the full smooth kernel,
including pairs whose ratio is larger than \(Z^\tau\).  The mask depends
on \(h,f,E\) and the fixed sector, but, after the indicated extraction,
not on \(u_1,u_2\).  If \(H_{\rm use}<0\), its nonzero ball is empty and
the same tail comparison discards every nonzero frequency.  Any principal
nonzero terms now restored to unify the formula carry this identical mask.
Their absolute contribution is bounded by the full principal restoration
already estimated.

Remove the condition \((u_1,u_2)=1\) by
\[
 1_{(u_1,u_2)=1}=\sum_{t'\mid(u_1,u_2)}\mu(t'),\qquad u_i=t'x_i,
 \qquad \widehat t=\log_Zq_{t'}.
\]
Let \(t\) be its fixed nonnegative center.
This inversion is made after the cross phases have been replaced by
the fixed ray functions above.  Those functions, the unchanged outer
mask \(\mathcal B_1\), all remaining
zero-extended local factors, and the formal product
\(q_{u_1}q_{u_2}q_{R_1}\) in the smooth kernel define an expression
also for noncoprime \(u_1,u_2\).  It agrees with Poisson summation on
the coprime support; the displayed divisor identity then recovers
exactly that support.  No primitive Poisson formula is asserted for
the newly introduced noncoprime pairs.  Squarefreeness retains the
puncture \((x_i,t')=1\).  Assigned slots at \(t'\) are outer
coefficients, while the other slots form the mark on \(x_i\).

Insert dyads of \(q_h\) from one fixed partition common to all \(f,E\),
not partitions recentered for those labels.  The mask \(\mathcal B_1\)
implies \(q_h\le q_{hf^2E}\le Z^{H_{\rm use}}\), so there are only
logarithmically many such dyads in the bounded retained scale range.
Keep \(\mathcal B_1\) outside every Fourier inversion.  Define
\[
 s_i=r-A_i-B-t,\qquad
 \kappa_i=M-2r-2\ell-V-\delta+A_i+B-R/2.
\]
The identity \(n_i=\mathfrak A_iCt'x_i\) on the factorized expression
gives \(|\widehat x_i-s_i|\le4\eta\).  The full fresh \(x_i\)-support
is included in \(\mathcal I_D\), so this same bound holds directly on
that support after separation, not only on the original product support.
The reconstructed formal products used in the current prefactor are also
included there.

\paragraph{Separating the first transformed product.}
We record the actual prefactor before applying Cauchy.  It splits exactly
over the two column sides as
\[
 Z^{-r-2\ell-V}Z^{M-\widehat\delta-\widehat L_1/2}
       =Z^{\widehat\kappa_1/2}Z^{\widehat\kappa_2/2},
 \quad
 \widehat\kappa_i=M-r-2\ell-V-\widehat\delta-\widehat n_i
                  +\widehat A_i+\widehat B-\widehat R/2.
\]
Equation~\eqref{eq:inverse-ratio-bounds} gives on each side
\begin{equation}\label{eq:inverse-first-prefactor-error}
 \widehat\kappa_i\le\kappa_i+\tfrac92\eta.
\end{equation}
The five error weights are \(1,1,1,1,1/2\).  We now implement this
normalization with all real inverse roots in one joint profile.
Fix a first dyadic, ray, and slot pattern \(\sigma\), and let \(h_1\)
be the center of its bare-\(h\) dyad.  Use the nine independent
normalized coordinates
\[
 \boldsymbol y=(y_{A_1},y_{A_2},y_C,y_d,y_R,y_t,y_h,y_{x_1},y_{x_2})
 =\left(\frac{q_{\mathfrak A_1}}{Z^{A_1}},
 \frac{q_{\mathfrak A_2}}{Z^{A_2}},\frac{q_C}{Z^B},
 \frac{q_{d_k}}{Z^\delta},\frac{q_{R_1}}{Z^R},\frac{q_{t'}}{Z^t},
 \frac{q_h}{Z^{h_1}},\frac{q_{x_1}}{Z^{s_1}},\frac{q_{x_2}}{Z^{s_2}}\right).
\]
The exact full root and kernel argument on the formal \(u_a=t'x_a\)
expression are
\begin{equation}\label{eq:inverse-first-root-kernel}
 \begin{split}
 \frac{Z^{-r-2\ell-V}Z^M}
 {q_{d_k}\sqrt{q_{R_1}}q_{t'}\sqrt{q_{x_1}q_{x_2}}}
 &=Z^{(\kappa_1+\kappa_2)/2}
 y_d^{-1}y_R^{-1/2}y_t^{-1}(y_{x_1}y_{x_2})^{-1/2},\\
 \frac{Z^Mq_h}{q_{d_k}q_{R_1}q_{t'}^2q_{x_1}q_{x_2}}
 &=A_{{\rm ker},1}\frac{y_h}{y_dy_Ry_t^2y_{x_1}y_{x_2}},\\
 A_{{\rm ker},1}&=Z^{M+h_1-\delta-R-2t-s_1-s_2}.
 \end{split}
\end{equation}
Choose fresh cutoffs \(\omega_{x,a}\) equal to one on the old
\(W\)-support divided by the individual \(\mathfrak A_a,C,t'\)
supports, and retain them in the columns.  Retain every current outer
dyad cutoff outside inversion.  Let \(\Omega_{1,\alpha}\) be larger
cutoffs equal to one on the full supports of these retained cutoffs.
With \(\mathcal K_1=\mathcal F\Phi\), define the single compact profile
\begin{equation}\label{eq:inverse-first-profile}
 \begin{split}
 \mathfrak H_{1,\sigma}(\boldsymbol y)={}&Z^{-9\eta/2}
 \prod_\alpha\Omega_{1,\alpha}(y_\alpha)
 y_d^{-1}y_R^{-1/2}y_t^{-1}(y_{x_1}y_{x_2})^{-1/2}\\
 &\times W(y_{A_1}y_Cy_ty_{x_1})
       \overline{W(y_{A_2}y_Cy_ty_{x_2})}
 \mathcal K_1\!\left(A_{{\rm ker},1}
       \frac{y_h}{y_dy_Ry_t^2y_{x_1}y_{x_2}}\right).
 \end{split}
\end{equation}
The scalar outside this profile is exactly
\(\prod_{a=1}^2Z^{(\kappa_a+9\eta/2)/2}\).
In particular, the profile contains each real root once, also when the
actual column norms differ.  Its transform is defined by
Equation~\eqref{eq:inverse-log-fourier} on the ambient nine-dimensional
box before any actual \(f,h\) or outer ideal is summed.  The norms of
\(b_a,f,J,E,q_0\) occur only in outer coefficients, characters, masks
or individual cutoffs after Equation~\eqref{eq:first-poisson-column};
they require no additional coupled-profile coordinate.

To expose the arithmetic extraction at \(t'\), put \(y=hf^2E\) and
\[
 B_{1,a}(t';y)=\mu(t')\nu_a(t')\rho(t')\overline{\chi_{t'}(y)}
       \chi_{t'}(CJ)^4\chi_{t'}(d_k)1_{(t',\operatorname{rad}q_0)=1}.
\]
Let \(\varepsilon_1=1\) and \(\varepsilon_2=-1\).  For a fixed first
Fourier mode \(\boldsymbol t\), the remaining polynomials are exactly
\begin{equation}\label{eq:inverse-first-polynomial}
 \begin{split}
 P_a(y)={}&\sum_{x\ {\rm sf}}\mu(x)\nu_a(x)\rho(x)
       \overline{\chi_x(y)}\chi_x(CJ)^4\chi_x(d_k)
       1_{(x,\operatorname{rad}q_0t')=1}\\
 &\qquad\times\mathfrak d_{I_a}(x)\omega_{x,a}(q_x/Z^{s_a})
                  (q_x/Z^{s_a})^{\varepsilon_a it_{x_a}}.
 \end{split}
\end{equation}
Multiplicativity is used on squarefree coprime factors, then through
the displayed zero extensions; no character is divided at a zero.
The one additional \(\mu(t')\) from the mutual-gcd inversion is outer.
Apply Equation~\eqref{eq:inverse-slot-priority} with the ordered list
\((b_a,\mathfrak A_a,C,t';x)\).  A surviving prime dividing \(x\)
is already prime to \(b_a,\mathfrak A_a,C,t'\): every prime of
\(\mathfrak B\) is in \(J\operatorname{rad}q_0\), and the displayed
fourth powers and punctures supply these zeros.  Thus
\(\mathfrak d_{I_a}\) has the original individual coefficients and
lists, while assigned prime identities and priority masks stay outer.

Here is the full structure of the first separated component.  Let
\(\Omega_\sigma\) consist of
\[
 (b_1,b_2,\mathfrak A_1,\mathfrak A_2,C,d_k,t',f,h)
\]
and the assigned slot primes, subject to their reconstruction relations and individual
supports.  In particular
\(\mathfrak A_a\mid\mathfrak B\), \((C,\mathfrak B)=1\), and
\(d_k\mid C\mathfrak B/R_1\), while \(R_1,E,J,q_0\) are derived,
not free indices.
Let \(\mathcal A_{1,\sigma}\) be the product of the assigned original
coefficients and priority masks, with conjugation on side two, and let
\(\psi_{1,\sigma}^{\rm out}\) be the product of all retained outer
dyad cutoffs.  For the fixed first ray summand, take explicitly
\[
 e_\sigma(\omega)=
      \mathcal C^{\rm out}_1(\omega)c_\vartheta(P_T),
\]
with \(\mathcal C^{\rm out}_1\) as in
\eqref{eq:inverse-first-outer-explicit}.
Thus it includes the old multiplicity exactly once, both original
\(\beta\) factors and their cube masks, the extracted \(a(c_i)\)
and \(f\)-zeros, and the complete primitive root and frequency factor
at \(R_1\).  The displayed \(\mu(d_k)\) and \(B_{1,a}\) are
excluded.  It is independent of \(x_1,x_2\); no definition uses
division by Equation~\eqref{eq:inverse-first-polynomial}.
Writing \(\mathcal O_1=\{A_1,A_2,C,d,R,t,h\}\), set
\[
 w_{1,\sigma}(\omega;\boldsymbol t)
 =e_\sigma(\omega)\mu(d_k)\mu(t')B_{1,1}(t';y)
       \overline{B_{1,2}(t';y)}\mathcal B_1(h)
       \mathcal A_{1,\sigma}(\omega)\psi_{1,\sigma}^{\rm out}(\omega)
       \prod_{\alpha\in\mathcal O_1}y_\alpha^{it_\alpha}.
\]
For the masked, principal-restored formal first component, Fourier
inversion gives the exact identity
\begin{equation}\label{eq:inverse-first-separated}
 \begin{split}
 \mathcal T_{1,\sigma}=(2\pi)^{-9}\int_{\mathbb R^9}
 &\widehat{\mathfrak H}_{1,\sigma}(\boldsymbol t)
 \sum_{\omega\in\Omega_\sigma}w_{1,\sigma}(\omega;\boldsymbol t)\\
 &\times Z^{(\kappa_1+9\eta/2)/2}P_1(y)
       \overline{Z^{(\kappa_2+9\eta/2)/2}P_2(y)}\,d\boldsymbol t.
 \end{split}
\end{equation}
Expanding the two polynomials restores the two column modes; the seven
outer modes restore the other coordinates of
Equation~\eqref{eq:inverse-first-profile}.  Its larger cutoffs are one
on the retained supports, and the fresh cutoffs are one wherever the
old windows are nonzero.  This verifies the identity term by term.
The original \(\beta\) factors and every outer mask are still in the
complete sum.

\paragraph{The first positive majorant.}
For any finite or absolutely convergent weighted sum we use
\begin{equation}\tag{D}\label{eq:weighted-cauchy}
 \left|\sum_\omega w_\omega U_1(\omega)\overline{U_2(\omega)}\right|
 \le\prod_{i=1}^2
       \left(\sum_\omega|w_\omega||U_i(\omega)|^2\right)^{1/2}.
\end{equation}
Apply this first with \(U_i=Z^{(\kappa_i+9\eta/2)/2}P_i\) on the whole
\(\Omega_\sigma\), for each fixed mode.  Only now bound the outer
factors absolutely.  Their pure arithmetic magnitudes are at most a
separately allocated \(Z^{\pi_\beta}\) times the old multiplicity and
the absolute assigned-slot product; the retained support gives
\(\mathcal B_1(h)1_{(C,J)=1}\).  In particular no factorization of
\(\beta(b,f)\) has been assumed.

For each fixed old outer reconstruction and nonzero \(y=\widetilde h\),
the relation \(y=hf^2E\) implies \(f^2E\mid(y)\).  If its multiplicity
is at most \(C_0d_{\mathcal O}(f)^{C_0}\), then
\begin{equation}\label{eq:inverse-old-label-fibre}
 \sum_{f^2E\mid(y)}w(f)
 \le C_0d_{\mathcal O}((y))^{C_0+1}
 \le C_{\pi_{\rm old}}Z^{\pi_{\rm old}}
 \qquad(0<q_y\le Z^{H_{\rm use}}).
\end{equation}
The last bound is uniform because \(H_{\rm use}\) stays in a bounded
range; \(h=y/(f^2E)\) is then uniquely determined as an element.
This is an old-label fibre bound, not a multiplicity depending on a
later new label.
The polynomial \(P_i(y)\) has no remaining dependence on \(f\) or on
the individual \(b_i,\mathfrak A_i\) beyond \(J,q_0\), the fixed
dyadic length \(A_i\), and the fixed ray sector: this is precisely
Equations~\eqref{eq:first-poisson-column} and
\eqref{eq:retained-cube-label}, with the norm profiles separated.
After the first weighted Cauchy inequality, \(\mathcal B_1\) is simply
\(1_{0<q_y\le Z^{H_{\rm use}}}\).  At this positive-sum stage, and only
now, majorize it by a fixed nonnegative radial Schwartz function
\(\Phi_+(q_y/Z^{H_{\rm use}})\) which is at least one for arguments at
most one.  A nonempty ball has \(H_{\rm use}\ge0\), so the second Poisson
scale is exactly \(Z^{H_{\rm use}}\ge1\).  The majorant is not substituted
as an equality in the original signed expression.  More precisely, let
\(\mathfrak O_\sigma\) consist of
\[
 o=(b_1,b_2,\mathfrak A_1,\mathfrak A_2,C,d_k,t';
       \text{old assigned slot primes}),
\]
retaining their full individual supports, the source reconstruction
relations independent of \(f,h,x\), the assigned priority masks, and
\((C,J)=1\).  Let \(w_o^+\) be the product of the absolute original
coefficients of its assigned primes, and zero off this set.  It is
nonnegative and independent of \(y\).  Each first positive side is at
most \(C Z^{\pi_\beta+\pi_{\rm old}}\), times its scalar
\(Z^{\kappa_i+9\eta/2}\), times the explicit positive majorant
\begin{equation}\label{eq:inverse-first-majorant}
 \mathcal S_{i,\sigma,\boldsymbol t}
 =\sum_{o\in\mathfrak O_\sigma}w_o^+
    \sum_{y\in\mathcal O}\Phi_+(q_y/Z^{H_{\rm use}})|P_i(y)|^2.
\end{equation}
At this positive step the old \(h\) cutoff and mode, and the zeros of
\(B_{1,a}\) depending on \(y\), may be removed by their upper bounds.
The first weighted Cauchy inequality and the old-label fibre bound have
therefore removed the old \(f,h\) from this positive majorant and every
later outer set; their original \(-V\) normalization remains in
\(\kappa_i\).  After the following count, we apply the second Poisson
transformation to its \(y\)-sum.

Let \(\widehat r_{\rm diff}\ge0\) be the actual log-norm of the
even-parity primes with \(a_{1p}\ne a_{2p}\).  The table gives
\(\widehat R=\widehat s+\widehat r_{\rm diff}\), so
\begin{equation}\label{eq:fixed-cube-count}
 \widehat c=\widehat\ell_{\rm pair}+\widehat R/2-\widehat j
                  -\widehat r_{\rm diff}/2
 \le\ell+R/2-j+\tfrac52\eta.
\end{equation}
Ideal counting with this actual upper bound counts the fixed \(q_0\)
while retaining \(J\).  For a diagonal bound which also counts \(J\),
its full window has exponent at most \(j+\eta\), giving the combined
upper count \(Z^{\ell+R/2+7\eta/2+\epsilon}\).  The choices of the
\(\mathfrak A_i\) and of \(d_k\mid C\mathfrak B/R_1\) add only
divisor factors once \(b_1,b_2,C\) are specified.

Combining Equation~\eqref{eq:inverse-first-separated}, weighted Cauchy,
and the old-label fibre bound gives the explicit output of the first
transformation:
\begin{equation}\label{eq:inverse-first-transform-output}
 \begin{split}
 |\mathcal T_{1,\sigma}|\ll{}&Z^{\pi_\beta+\pi_{\rm old}}
 \int_{\mathbb R^9}|\widehat{\mathfrak H}_{1,\sigma}(\boldsymbol t)|\\
 &\quad\times\prod_{i=1}^2
 \left(Z^{\kappa_i+9\eta/2}\mathcal S_{i,\sigma,\boldsymbol t}\right)^{1/2}
 \,d\boldsymbol t.
 \end{split}
\end{equation}
Each \(\mathcal S_{i,\sigma,\boldsymbol t}\) is the positive sum in
Equation~\eqref{eq:inverse-first-majorant}: its row polynomial has the
M\"obius coefficient in Equation~\eqref{eq:inverse-first-polynomial},
the cube-derived label \(J\) is retained, and the old \(f,h\) no longer
occur among the averaged indices. The following transformation is applied
to these positive row norms.

\paragraph{The second Poisson transformation.}
Fix one of the two positive sums obtained from the first Cauchy
inequality.  Its polynomial has the form
\[
 P_i(y)=\sum_x c_i(x)\overline{\chi_x(y)}.
\]
The coefficient \(c_i(x)\) consists of the other factors of
Equation~\eqref{eq:first-poisson-column}, with \(u_i=t'x\), the
puncture at \(t'\), and the separated weight.  It is independent of
\(y\).  In particular \(x\) is squarefree and
\[
 c_i(x)=0\quad\text{unless}\quad
 (x,CJ)=1,\quad (x,\operatorname{rad}q_0t')=1.
\]
The puncture \(\rho\) and the zero of \(\chi_x(d_k)\) are retained
as well.  The \(y\)-sum uses the fixed-shape smooth positive ball at
scale \(Z^{H_{\rm use}}\).

In its expanded square put
\[
 g'=(x_1,x_2),\qquad x_j=g'z_j,\qquad
 \widehat g=\log_Zq_{g'}.
\]
Let \(g\ge0\) be its fixed center.
Then \(z_1,z_2\) are squarefree and coprime, and each is prime to
\(g'\).  The row data for Lemma~\ref{lem:masked-poisson} are
\begin{equation}\label{eq:second-masked-data}
 K=Z^{H_{\rm use}},\qquad
 \psi_2=\bar\chi_{z_1}\chi_{z_2},\qquad
 m_2=z_1z_2,\qquad R_{\rm mask}=g'.
\end{equation}
The common factor \(\bar\chi_{g'}(y)\chi_{g'}(y)\) is exactly this
mask.  The character is primitive away from \(m_2=1\), because its
two nonzero local powers have disjoint supports.  For \(d_2\mid g'\),
let \(\vartheta_c\ge0\) be its center and write
\(\widehat\vartheta_c=\log_Zq_{d_2}\).  The actual conductor length is
\[
 \widehat L_2=\widehat x_1+\widehat x_2-2\widehat g,
 \qquad |\widehat L_2-2(s_i-g)|\le10\eta.
\]
The last inequality uses the full fresh \(x\)-support bound
\(|\widehat x_a-s_i|\le4\eta\) for each copy and
\(|\widehat g-g|\le\eta\).  For \(k_{\rm new}=d_kd_2k''\), the implication
\(Z^{H_{\rm use}}q_{k''}/(q_{d_2}q_{m_2})\le Z^\tau\) gives
\[
 \begin{split}
 \widehat{k_{\rm new}}
 &\le2(s_i-g)-H_{\rm use}+\delta+2\vartheta_c+13\eta+\tau\\
 &=M_c+\eta,
 \end{split}
\]
where
\begin{equation}\label{eq:second-poisson-lengths}
 \begin{split}
 M_c&=2(s_i-g)-H_c+\delta+2\vartheta_c\\
    &=M-4\ell-2A_i-2t-2g+2\vartheta_c-2V+j,\qquad
 M_{\rm ch}=M_c+\eta.
 \end{split}
\end{equation}
The \(13\eta\) consists of \(10\eta\) from the conductor, one from
\(d_k\), and two from \(d_2\).  Subtracting
\(H_{\rm use}=H_c+12\eta+\tau\) cancels the same tolerance \(\tau\)
used in this second Poisson comparison.  Thus no separate row clipping
or unrecorded boundary error is needed.

Separate the original principal contribution for the direct count below.
Before any off-coprime extension or Fourier absolutization, keep on the
genuine nonprincipal second Poisson expression only the outer mask
\[
 \mathcal B_2(k'')=1_{0<q_{d_kd_2k''}\le Z^{M_{\rm ch}}}.
\]
For each supported pair its complement is contained in the actual ratio
tail just considered, so Equation~\eqref{eq:inverse-raw-tail} discards it
with the order fixed below.  Inside the mask retain the full smooth kernel.
The mask depends on \(d_k,d_2,k''\) and the fixed sector, not on the residual
\(z_1,z_2\).  If \(M_{\rm ch}<0\), the nonzero ball is empty and the same
tail comparison discards every nonzero frequency.  In a retained nonempty
passage \(M_{\rm ch}\ge0\).

The actual second prefactor has the exact side split
\[
 Z^{H_{\rm use}-\widehat\vartheta_c-\widehat L_2/2}
       =Z^{\widehat\lambda_1/2}Z^{\widehat\lambda_2/2},
 \qquad \widehat\lambda_a=H_{\rm use}-\widehat\vartheta_c-\widehat x_a+
                           \widehat g.
\]
For \(\lambda_c=H_c-\vartheta_c-(s_i-g)\), the same support bounds give
\begin{equation}\label{eq:inverse-second-prefactor-error}
 \widehat\lambda_a\le\lambda_c+18\eta+\tau.
\end{equation}
Here \(18=12+1+4+1\).  As in the first split, all normalized real
inverse roots will remain in the joint smooth profile until Fourier
inversion; they are not also appended to the final child tests.
All original coefficient masks remain.  In particular
\[
 (g',CJ)=1,\qquad (z_j,g'CJ)=1
\]
on their nonzero support.

The principal case is \(z_1=z_2=1\), equivalently \(x_1=x_2\).
The nominal identity for its count is
\begin{equation}\label{eq:second-principal-exponent}
 \kappa_i+H_c+s_i+B+t+\ell+R/2=F-B-j,
 \qquad F=N+V=r+3\ell+V.
\end{equation}
The actual upper count adds \(9\eta/2\) from the first prefactor,
\(12\eta+\tau\) from the row scale \(H_{\rm use}\), \(4\eta\) from the
diagonal \(x\)-window, \(\eta\) each from \(C,t'\), and \(7\eta/2\) from
the combined \(q_0,J\) count.  Their sum is \(26\eta+\tau\).  With the
aggregate local loss \(\pi\), this principal contribution is therefore
at most \(Z^{F-B-j+26\eta+\tau+\pi}\le Z^{F+26\eta+\tau+\pi}\).
Any nonzero principal frequencies restored to the formal expression carry
the same \(\mathcal B_2\) mask; their absolute sum is bounded by the full
principal restoration in Lemma~\ref{lem:masked-poisson}, since
\(Z^{H_{\rm use}}\ge1\).

For the nonprincipal terms, we now identify the new coefficient
class.  For squarefree \(z\), complex conjugation of a primitive
Gauss sum and Equation~\eqref{eq:signal-checked} give
\[
 \gamma_{-1}(z)=\chi_z(-1)\overline{\gamma_1(z)},\qquad
 \mu(z)\gamma_{-1}(z)
 =\bar\alpha(z)\gamma_2(z)\chi_z(-1)\overline{G(z)}.
\]
The raw second factor is
\(\mu(z)\gamma_1(z)=\overline{\bar\alpha(z)\gamma_2(z)}G(z)\);
the corresponding factor written inside the conjugated second
polynomial has ray factor \(\overline{G(z)}\).  The CRT cross factor of the two primitive
roots is \(\mathcal R(z_1,z_2)\).  Thus their combined ray factor is
\[
 \mathfrak G(a,b)=\chi_a(-1)\overline{G(a)}G(b)\mathcal R(a,b)
       =G(ba^{-1}).
\]
The last equality follows from
Equation~\eqref{eq:G-quadratic-refinement} and
\(\mathcal R(a,a)=\chi_a(-1)\), first on primes by
Equation~\eqref{eq:fixed-gauss-phases} and then multiplicatively.
All quotients in this display are in the fixed finite ray group.
In particular,
\[
 \mathfrak G(v'n_1,v'n_2)=\mathfrak G(n_1,n_2).
\]
The common multiplicative character \(\nu_i(v')\) also cancels
between the two sides.

Insert dyads of \(q_{k''}\) from one partition common to all \(d_k,d_2\).
The mask \(\mathcal B_2\) implies \(q_{k''}\le Z^{M_{\rm ch}}\), so their
number is logarithmic in the bounded retained range.  Let \(h_2\) be
the center of one such bare-frequency dyad.  We now record the exact
component to be separated.  Fix the first mode and selected side \(i\),
write
\[
 H_o=CJ,\qquad r_0=\operatorname{rad}q_0t',\qquad d_0=d_k,
 \qquad a_0(z)=\bar\alpha(z)\gamma_2(z),\qquad
 \mathcal K_2=\mathcal F\Phi_+,
\]
and retain \(K=Z^{H_{\rm use}}\) from
Equation~\eqref{eq:second-masked-data}.  For squarefree arguments put
\begin{equation}\label{eq:inverse-second-outer-coefficients}
 \begin{split}
 B_o(g')&=\mu(g')\nu_i(g')\rho(g')1_{(g',r_0)=1}
              \chi_{g'}(H_o)^4\chi_{g'}(d_0),\\
 L_{o,d_2,k''}(z)&=a_0(z)\nu_i(z)\rho(z)1_{(z,r_0)=1}
       \chi_z(H_o)^4\chi_z(d_0)\overline{\chi_z(d_2)}\chi_z(k'').
 \end{split}
\end{equation}
The old separated test in \(P_i\) is denoted by \(W_i\), so
\(W_i(x)=\omega_{x,i}(x)x^{\varepsilon_iit_{x_i}}\).
Let \(\psi_g,\psi_{d_2},\psi_{k''}\) denote the retained individual
cutoffs on the fixed \(g',d_2,k''\) dyads, evaluated at their normalized
norms.  For this preliminary pattern \(\sigma_0\), the masked
principal-restored nonzero component of
Equation~\eqref{eq:inverse-first-majorant} is exactly
\begin{equation}\label{eq:inverse-second-formal-component}
 \begin{split}
 \mathcal T_{2,\sigma_0,i}={}&
 \sum_{\substack{o\in\mathfrak O_\sigma,\ g'\ {\rm sf}\\d_2\mid g',\ k''\ne0}}
 w_o^+\psi_g\psi_{d_2}\psi_{k''}\mu(d_2)\mathcal B_2(k'')|B_o(g')|^2\\
 &\times\sum_{\substack{z_1,z_2\ {\rm sf}\\(z_1z_2,g')=1}}
 1_{(z_1,z_2)=1}L_{o,d_2,k''}(z_1)\overline{L_{o,d_2,k''}(z_2)}
       G([z_2][z_1]^{-1})\\
 &\times\mathfrak d_{I_i}(g'z_1)\overline{\mathfrak d_{I_i}(g'z_2)}
       W_i(q_{g'}q_{z_1}/Z^{s_i})
       \overline{W_i(q_{g'}q_{z_2}/Z^{s_i})}\\
 &\times\frac{K}{q_{d_2}\sqrt{q_{z_1}q_{z_2}}}
       \mathcal K_2\!\left(\frac{Kq_{k''}}{q_{d_2}q_{z_1}q_{z_2}}\right).
 \end{split}
\end{equation}
Here brackets denote classes in the fixed finite ray group.  Indeed the
arithmetic part of \(c_i(g'z)\) factors into \(B_o(g')\) and its
residual \(\mu(z)\) coefficient on \((g',z)=1\).  The common row factor
is the mask at \(g'\), whose Poisson expansion gives the one divisor
\(d_2\), while the signal calculation above gives the displayed two
\(L\) factors and ray quotient.  The restored principal pair
\(z_1=z_2=1\) agrees by \(G(1)=1\).  Its separate cost, and the raw
complement removed before this formula, have already been bounded.
Thus this is an equality for the specified component of the positive
majorant, not for the original signed block.

Define the summand of Equation~\eqref{eq:inverse-second-formal-component}
on all individually admissible squarefree \(z_1,z_2\) by its displayed
quotient-free \(L\) factors, fixed-ray quotient, formal product norms,
full kernel, and unchanged \(\mathcal B_2\).  It agrees with the
genuine formula on coprime pairs.  Insert the complete identity
\[
 1_{(z_1,z_2)=1}=\sum_{v'\mid(z_1,z_2)}\mu(v'),\qquad z_a=v'n_a,
\]
before any factorwise estimate.  Let \(v\ge0\) be the center of the
squarefree \(v'\)-dyad and \(\widehat v=\log_Zq_{v'}\).  The
\(n_a\) are squarefree, with \((v',n_a)=1\); they need not be mutually
coprime.  No primitive Poisson formula is used on this extension.
For squarefree \(v'\) put
\[
 D_o(v';d_2,k'')=a_0(v')\nu_i(v')\rho(v')1_{(v',r_0)=1}
   \chi_{v'}(H_o)^4\chi_{v'}(d_0)\overline{\chi_{v'}(d_2)}\chi_{v'}(k'').
\]
On squarefree coprime \(v',n\), the exact extraction is
\begin{equation}\label{eq:inverse-second-extraction}
 \begin{split}
 L_{o,d_2,k''}(v'n)={}&D_o(v';d_2,k'')a_0(n)\nu_i(n)\rho(n)1_{(n,r_0)=1}\\
 &\times\chi_n(H_o)^4\chi_n(d_0)\overline{\chi_n(d_2)}
          \chi_n(k'')\chi_n(v')^4.
 \end{split}
\end{equation}
This uses \(a_0(v'n)=a_0(v')a_0(n)\chi_n(v')^4\).  On an overlap the
right side defines the extension to be zero by \(\chi_n(v')^4\),
without evaluating a nonsquarefree Gauss sum.  The two common factors
give \(|D_o(v';d_2,k'')|^2\), including the zero \((v',k'')=1\) and
its fixed-puncture, \(d_0\), and label zeros.  They remain in the outer
weight through Cauchy.  The old gcd also leaves \(1_{(v',g')=1}\).
Since \(D_o\) vanishes on \((v',CJ)>1\), we retain equivalently the
explicit outer factor \(1_{(v',g'CJ)=1}\), which repeats that label zero.

Set \(r_g=g'/d_2\), \(k_{\rm new}=d_kd_2k''\) and
\(f_{\rm new}=JCd_2v'\).  All moving factors on \(n\) satisfy the
complete zero-extended identity
\begin{equation}\label{eq:inverse-moving-characters}
 \chi_n(d_k)\chi_n(k'')\overline{\chi_n(d_2)}
       \chi_n(CJ)^4\chi_n(v')^4
 =\chi_n(k_{\rm new})\chi_n(f_{\rm new})^4.
\end{equation}
It uses \(\overline{\chi_n(d_2)}=\chi_n(d_2)\chi_n(d_2)^4\), also on
nonunits.  The residual gcd with \(g'=d_2r_g\) splits into the repeated
zero at \(d_2\) and the fixed puncture at \(r_g\).  The old punctures
at \(q_0,t'\) remain.  If \(\mathcal C_{\rm ray}\) denotes the fixed
group through which \(G\) factors, put
\[
 \widehat G(\vartheta)=|\mathcal C_{\rm ray}|^{-1}
       \sum_{c\in\mathcal C_{\rm ray}}G(c)\overline{\vartheta(c)}.
\]
Then
\[
 G([n_2][n_1]^{-1})=\sum_{\vartheta\in\widehat{\mathcal C}_{\rm ray}}
   \widehat G(\vartheta)\overline{\vartheta(n_1)}\vartheta(n_2).
\]
Thus both new polynomials in a fixed ray summand use the same
\(\nu_i\overline\vartheta\), and its one bounded coefficient
\(\widehat G(\vartheta)\) remains outer.

\paragraph{The new canonical data and their admissibility.}
The second transformation has returned the moving characters to canonical
form. We now identify which data will be fixed and which will remain
averaged. This distinction also determines the puncture of the child.
Put \(\nu_a'=\nu_i\overline\vartheta\) for both sides and define
\begin{equation}\label{eq:canonical-child-data}
 \begin{gathered}
 \gamma=(q_0,t',r_g),\qquad r_g=g'/d_2,\\
 k_{\rm new}=d_kd_2k'',\qquad f_{\rm new}=JCd_2v',\\
 \rho'(n)=\rho_\gamma(n)=\rho(n)1_{(n,\operatorname{rad}q_0t'r_g)=1}.
 \end{gathered}
\end{equation}
On the nonzero coefficient support, \(J,C,d_2,v'\) are squarefree
and pairwise coprime: \((C,J)=1\) belongs to \(\mathfrak O_\sigma\),
\(|B_o(g')|^2\) forces \((g',CJ)=1\), \(d_2\mid g'\), and the
retained zero of \(D_o\), together with the old gcd restriction, gives
\((v',g'CJ)=1\). Thus \(f_{\rm new}\) is squarefree on that support.
These common factors remain in the outer weight through Cauchy.

The triple \(\gamma\) will be fixed before the child row and label sums.
Its puncture \(\rho_\gamma\) is therefore independent of those two
averaging variables. The ideals \(J,C,d_2,v'\) remain in the new
averaged label; they are not counted as additional fixed labels.

Apply Equation~\eqref{eq:inverse-slot-priority} on each selected-side
copy with the ordered list \((d_2,r_g,v';n_a)\).  Let \(I_a'\) be the
subcollection of the selected side's slots retained on \(n_a\) in copy
\(a\).  Assigned primes and their priority masks remain outer.
A surviving prime dividing \(n\)
is already prime to \(JCd_2v'\) by the fourth-power zero and to
\(\operatorname{rad}q_0,t',r_g\) by the fixed puncture.  These also
supply every earlier survival exclusion because
\(\operatorname{supp}\mathfrak B\subset\operatorname{supp}(Jq_0)\).
The surviving mark is therefore an original product-form subcollection.
An old slot described as assigned to \(J\) is only a regrouping of an
old \(b_a,\mathfrak A_a\) assignment by this derived support, not a
new independent prime choice.

Keep \(\mathcal B_2\) outside the Fourier separation of the full kernel
and old windows. Let \(I_g,I_d,I_v,I_k\) be the full supports of the
current individual dyad cutoffs \(\psi_g,\psi_{d_2},\psi_v,\psi_{k''}\).
Choose fresh cutoffs \(\omega_{n,a}\) equal to one on the support of
\(W_i\) divided by \(I_gI_v\), with fixed full supports in an interval
\([a,b]\), where \(b\ge1\). Choose a nonnegative \(\omega_f\) equal to
one on \(I_JI_CI_dI_v\), with fixed full support \(I_f\). Include these
full supports, and the reconstructed formal product \(g'v'n\) on them,
in the finite window family \(\mathcal I_D\) specified above.
These full fresh child windows give the formal centers and support bounds
\begin{equation}\label{eq:canonical-child-lengths}
 \begin{gathered}
 N_c=r-A_i-B-t-g-v,\qquad V_c=B+\vartheta_c+v+j,\qquad F_c=N_c+V_c,\\
 |\widehat n_{\rm child}-N_c|\le6\eta,\qquad
 |\widehat f_{\rm new}-V_c|\le4\eta.
 \end{gathered}
\end{equation}
On the factorized expression the bounds follow from the exact products
\(n_i=\mathfrak A_iCt'g'v'n_{\rm child}\) and
\(f_{\rm new}=JCd_2v'\).  The full fresh child cutoff and the full enlarged
label window are included in \(\mathcal I_D\), so the same bounds hold
directly on their supports when independent terms are later added by
positivity.  Those added terms are not asserted to arise from a parent
factorization.  During the current separation, the reconstructed formal
product \(x=g'v'n\) on the full fresh cutoff is also included in
\(\mathcal I_D\); thus the \(4\eta\) bound used in the second prefactor
holds on the whole separated side.

Keep the nonnegative label center \(V_c\) without rounding it upward.
Proceed to a child only if the retained row component is nonzero; then
the outer ball has \(M_{\rm ch}\ge0\).  This gate is separate from any
structural outer index set, which may contain zero-weight tuples even
when \(\mathcal B_2\) is identically zero.  Similarly, if a full fresh
child column window contains no squarefree ideal, its polynomial is zero
and that component is omitted.  If \(N_c<0\) and a child column exists, its norm is at least one, and
Equation~\eqref{eq:canonical-child-lengths} gives \(-6\eta\le N_c<0\).
For a retained nonempty child define
\begin{equation}\label{eq:inverse-child-clipping}
 N_{\rm ch}=\max(0,N_c),\qquad \delta_N=N_{\rm ch}-N_c\in[0,6\eta],
 \qquad F_{\rm ch}=N_{\rm ch}+V_c=F_c+\delta_N.
\end{equation}
This is a bounded change of the test, not a \(Z^\eta\)-wide support
enlargement. Use the fixed enclosing support interval \([a,b]\)
chosen above. If
\(N_c<0\), nonemptiness gives \(Z^{\delta_N}\le b\); if \(N_c\ge0\),
then \(\delta_N=0\) and the same bound follows from \(b\ge1\).
Thus \(1\le Z^{\delta_N}\le b\) in both cases.  The fresh annular cutoffs
\(\omega_{n,a}(Z^{\delta_N}y)\) have support in \([a/b,b]\) and uniformly
bounded Euler seminorms. Their scale depends only on the fixed sector
centers, not on a current row or label.  The full union of these clipping
supports is included in \(\mathcal I_D\).

Fix a refinement \(\sigma\) of \(\sigma_0\) by its \(v'\)-dyad,
second ray summand and slot branches.
Let \(\Xi_\sigma\) consist of
\[
 \xi=(o,g'\ {\rm sf},d_2\mid g',v'\ {\rm sf},k''\ne0;
               \text{new assigned slot primes})
\]
with the retained individual supports and slot branches.  It contains
neither the old \(f,h\) nor the current column indices \(n_1,n_2\).
The identities defining \(R_1,E,J,q_0\) are inherited from \(o\),
and \(\gamma,k_{\rm new},f_{\rm new}\) are the displayed functions of
\(\xi\).  The set may be taken before imposing \(\mathcal B_2\) and
\((v',g'CJ)=1\), retaining those masks in the signed outer weight.  Let
\(\mathcal A_{2,\sigma}(\xi)\) be the product of all newly assigned
original slot coefficients and their priority masks, conjugated on
side two.

Let \(\mathfrak X_\sigma^{\rm src}\) be the structural outer set of
this transformed component, before adding any independent child rows
or labels.  It retains
\(o\in\mathfrak O_\sigma\), in particular the exact reconstruction
\(b_1b_2=q_0^2J_2^2s\), the derived \(J=sJ_2,q_0\), and the full
fixed \(b_1,b_2,t'\) dyad supports; it also retains the fixed
\(g',d_2\) supports, \(d_2\mid g'\), and its other outer arithmetic
and slot conditions.  Zero-weight tuples may be retained, and unit
Fourier phases are ignored in defining this structural set.  Define
the set of distinct triples, with no witness multiplicity, by
\begin{equation}\label{eq:inverse-source-triples}
 \Gamma_\sigma=
 \{(q_0,t',r_g=g'/d_2):\xi\in\mathfrak X_\sigma^{\rm src}\}.
\end{equation}
This projection is taken once over all possible current rows and
labels, before fixing any \(k_{\rm new},f_{\rm new}\) or Fourier
mode.  Every nonzero transformed term maps into it.  Every member
has a source witness, which gives
\(\widehat c\le\ell+\eta\), \(\widehat t\le t+\eta\), and
\(\log_Zq_{r_g}\le g-\vartheta_c+2\eta\).  The first of these bounds
will be needed for the child puncture; the containing norm balls
used for counting do not imply it.  The nonempty row and child-column gates stated above
remain separate from this structural projection.

Put
\[
 D_c=\ell+A_i+t+g-\vartheta_c+V.
\]
Since \(r=N-3\ell\) is an exact definition tied to the fixed parent,
substitution in Equations~\eqref{eq:canonical-child-lengths} and
\eqref{eq:second-poisson-lengths} gives the exact formal identities
\begin{equation}\label{eq:canonical-transition}
 \begin{gathered}
 F_c=F-D_c-2\ell+j,\qquad M_c=M-2D_c-2\ell+j,\\
 F_c-M_c=F-M+D_c,\\
 4F_c-3M_c=4F-3M+2(A_i+t+g-\vartheta_c+V)+j.
 \end{gathered}
\end{equation}
These identities involve the same fixed \(F=N+V\) as
Equation~\eqref{eq:inverse-fixed-invariant}, not an actual parent norm.

For every \(\gamma\in\Gamma_\sigma\), choose any source witness.
Its actual divisibility \(d_2\mid g'\) and
Equation~\eqref{eq:inverse-cube-actual} give
\[
 \widehat\vartheta_c\le\widehat g,\qquad
 \widehat j\le2\widehat\ell_{\rm pair},\qquad
 \widehat c\le\widehat\ell_{\rm pair}.
\]
Consequently the only actual-to-center inequalities needed here are
\begin{equation}\label{eq:inverse-center-inequalities}
 g-\vartheta_c\ge-2\eta,\qquad j\le2\ell+3\eta,\qquad
 \widehat c\le\ell+\eta.
\end{equation}
The first two are not asserted with zero error.  In particular the
third inequality holds for every fixed triple because of its source
witness, even though it need not hold throughout the containing count
ball.  The new radical in Equation~\eqref{eq:canonical-child-data}
has actual logarithmic norm \(Q_{\rm new}=Q_\gamma\), satisfying
\begin{equation}\label{eq:inverse-new-puncture}
 \begin{split}
 Q_{\rm new}&\le Q+\widehat c+\widehat t+\widehat g-\widehat\vartheta_c\\
 &\le Q+\ell+t+g-\vartheta_c+4\eta\le Q+D_c+4\eta.
 \end{split}
\end{equation}
This also holds when its factors share primes with the old puncture,
because taking the radical only decreases the norm.  The last inequality
uses \(A_i,V\ge0\).

Use \(F_{\rm ch}=F_c+\delta_N\), \(M_{\rm ch}=M_c+\eta\), and
\(2(A_i+t+g-\vartheta_c+V)+j\ge-4\eta\) in
Equation~\eqref{eq:canonical-transition}.  The two child margins satisfy
\begin{equation}\label{eq:inverse-child-margins}
 \begin{aligned}
 F_{\rm ch}-M_{\rm ch}-Q_{\rm new}-z_0
   &\ge c_{\rm node}+\delta_N-5\eta\ge c_{\rm node}-5\eta,\\
 4F_{\rm ch}-3M_{\rm ch}-6z_0
   &\ge c_{\rm node}+4\delta_N-7\eta\ge c_{\rm node}-7\eta.
 \end{aligned}
\end{equation}
The nominal cap \(z_0\) is unchanged because the surviving mark is a
subcollection.  Moreover \(D_c\ge\ell+V-2\eta\) and
\(2\ell-j\ge-3\eta\), whence
\begin{equation}\label{eq:inverse-row-decrease}
 M-M_{\rm ch}=2D_c+2\ell-j-\eta
       \ge2(\ell+V)-8\eta\ge2d-8\eta.
\end{equation}
For \(\eta\le d/16\) this is at least \(3d/2\), hence at least \(d\).
The same identities give \(F_c\le F-\ell-V+5\eta\), so
\(F_{\rm ch}\le F+11\eta\).  Because its two summands are nonnegative,
this bounds the child \(N_{\rm ch},V_c\) through the finite depth.

The arithmetic construction has thus produced admissible child ranges
with a smaller row length. We next express the transformed sum in their
canonical polynomials. The required Fourier density is common to all
triples \(\gamma\), rows and labels; only after that identity will we
apply Cauchy and dominate the reconstruction multiplicities.

\paragraph{A common density for the second transformed sum.}
Before any actual outer ideal is summed, use the six independent coordinates
\[
 (y_g,y_d,y_v,y_k,y_1,y_2)
 =\left(\frac{q_{g'}}{Z^g},\frac{q_{d_2}}{Z^\vartheta_c},
 \frac{q_{v'}}{Z^v},\frac{q_{k''}}{Z^{h_2}},
 \frac{q_{n_1}}{Z^{N_c}},\frac{q_{n_2}}{Z^{N_c}}\right).
\]
Use the fresh cutoffs \(\omega_{n,a},\omega_f\) already chosen.
Retain the four current dyad cutoffs
and \(\omega_f\) outside inversion, and \(\omega_{n,a}\) in the two
columns.  Let \(\Omega_{2,\alpha}\) be larger cutoffs equal to one on
the full supports of the corresponding four dyad and two column
cutoffs.  These full supports, \(I_f\), and the formal products
\(g'v'n\) on them are the windows already included in \(\mathcal I_D\).

The exact second root and kernel argument are
\begin{equation}\label{eq:inverse-second-root-kernel}
 \begin{split}
 \frac{K}{q_{d_2}q_{v'}\sqrt{q_{n_1}q_{n_2}}}
 &=Z^{\lambda_c+12\eta+\tau}y_d^{-1}y_v^{-1}(y_1y_2)^{-1/2},\\
 \frac{Kq_{k''}}{q_{d_2}q_{v'}^2q_{n_1}q_{n_2}}
 &=A_{{\rm ker},2}\frac{y_k}{y_dy_v^2y_1y_2},\\
 A_{{\rm ker},2}&=Z^{H_{\rm use}+h_2-\vartheta_c-2v-2N_c}.
 \end{split}
\end{equation}
Indeed \(N_c=s_i-g-v\) and
\(\lambda_c=H_c-\vartheta_c-(s_i-g)\).  Define the one ambient profile
\begin{equation}\label{eq:inverse-second-profile}
 \begin{split}
 \mathfrak H_{2,\sigma,\boldsymbol t}(\boldsymbol y)={}&Z^{-6\eta}
 \prod_\alpha\Omega_{2,\alpha}(y_\alpha)
 y_d^{-1}y_v^{-1}(y_1y_2)^{-1/2}\\
 &\times W_i(y_gy_vy_1)\overline{W_i(y_gy_vy_2)}
 \mathcal K_2\!\left(A_{{\rm ker},2}\frac{y_k}{y_dy_v^2y_1y_2}\right).
 \end{split}
\end{equation}
The scalar outside it is exactly \(Z^{\lambda_c+18\eta+\tau}\).
The full \(q_{d_2}\) denominator and both normalized real column roots
are in this profile, and no such root is appended to the final tests.

For \(\boldsymbol\zeta=(\zeta_g,\zeta_d,\zeta_v,\zeta_k,
\zeta_1,\zeta_2)\), put \(z_n=q_n/Z^{N_{\rm ch}}\) and define
\[
 W_a'(z)=\omega_{n,a}(Z^{\delta_N}z)z^{\varepsilon_a i\zeta_a},
 \qquad \nu_a'=\nu_i\overline\vartheta,\qquad
 \mathfrak d_a'=\mathfrak d_{I_a'}.
\]
The polynomial on either new Cauchy side is then exactly
\begin{equation}\label{eq:canonical-child-polynomial}
 \begin{split}
 Q_{a,\gamma,\boldsymbol\zeta}(k_{\rm new},f_{\rm new})={}&
 \sum_{n\ {\rm sf}}\bar\alpha(n)\gamma_2(n)\nu_a'(n)\rho'(n)
 \chi_n(k_{\rm new})\\
 &\times\chi_n(f_{\rm new})^4
 \mathfrak d_a'(n)W_a'(q_n/Z^{N_{\rm ch}}).
 \end{split}
\end{equation}
After \(q_0,t',g'/d_2\) and the finite ray sector are fixed,
\(\nu_a'\) and \(\rho'\) are independent of \(J,C,d_2,v'\) and
the new row.  The only new characters on \(n\) are the explicitly
displayed row and fourth-power factors; all other phases involving
the old \(b_i,\mathfrak A_i\) were placed in
Equation~\eqref{eq:retained-cube-label} or in the old row
\(\widetilde h\).  The separated annular weight \(W_a'\) is also
independent of the new row and label.
The complete remaining outer weight is
\begin{equation}\label{eq:inverse-second-outer-weight}
 \begin{split}
 V_\sigma(\xi;\boldsymbol\zeta)={}&w_o^+\mu(d_2)\mu(v')\widehat G(\vartheta)
 |B_o(g')|^2|D_o(v';d_2,k'')|^2
 1_{(v',g'CJ)=1}\mathcal B_2(k'')\\
 &\times\mathcal A_{2,\sigma}(\xi)
 \psi_g(y_g)\psi_{d_2}(y_d)\psi_v(y_v)\psi_{k''}(y_k)
 \omega_f(q_{f_{\rm new}}/Z^{V_c})\\
 &\times y_g^{i\zeta_g}y_d^{i\zeta_d}y_v^{i\zeta_v}y_k^{i\zeta_k}
       Z^{i\delta_N(\zeta_1+\zeta_2)}.
 \end{split}
\end{equation}
All old pure arithmetic and priority masks are retained in
\(\mathfrak O_\sigma\), with their nonnegative old coefficient weight
\(w_o^+\).  The displayed two common magnitudes retain every other
common coefficient zero, including the one depending on \(k''\).
There is one \(\mu(d_2)\), one \(\mu(v')\), and one ray coefficient.
The label cutoff is inserted as an equality on the original product
supports and remains outer.

Let \(\mathcal T_{\sigma,i}\) be the fixed \(v'\)-dyad, ray, and slot
summand obtained from Equation~\eqref{eq:inverse-second-formal-component}
by the complete M\"obius and slot expansions above.  The exact identity
needed for the child is
\begin{equation}\label{eq:inverse-xi-identity}
 \begin{split}
 \mathcal T_{\sigma,i}={}&Z^{\lambda_c+18\eta+\tau}(2\pi)^{-6}
 \int_{\mathbb R^6}\widehat{\mathfrak H}_{2,\sigma,\boldsymbol t}
       (\boldsymbol\zeta)\\
 &\quad\times\sum_{\xi\in\Xi_\sigma}V_\sigma(\xi;\boldsymbol\zeta)
 Q_{1,\gamma,\boldsymbol\zeta}(k_{\rm new},f_{\rm new})
 \overline{Q_{2,\gamma,\boldsymbol\zeta}(k_{\rm new},f_{\rm new})}
 \,d\boldsymbol\zeta.
 \end{split}
\end{equation}
To check it, expand the two \(Q\)'s.  Since
\(y_a=Z^{\delta_N}z_{n_a}\), their modes together with the last
factor of Equation~\eqref{eq:inverse-second-outer-weight} are exactly
\(y_1^{i\zeta_1}y_2^{i\zeta_2}\).  The other four modes are outer.
Equation~\eqref{eq:inverse-log-fourier} restores the full profile,
whose larger cutoffs are one on the retained supports.  The fresh
column cutoffs are one wherever the old \(W_i\) factors are nonzero,
and \(\omega_f=1\) on every original label product.  The scalar and
profile therefore restore the full root, full kernel, and old windows.
Equations~\eqref{eq:inverse-second-extraction},
\eqref{eq:inverse-moving-characters}, and
\eqref{eq:inverse-slot-priority} restore the arithmetic factors term
by term.  The fixed-box sums are finite and
Equation~\eqref{eq:inverse-fourier-moment} justifies the integral.
Extra combinations in the full fresh windows cancel through the old
windows in this complex Fourier identity before Cauchy; only later are
they kept independently in a positive sum.

The density in Equation~\eqref{eq:inverse-xi-identity} depends on the
fixed sector, \(Z\), the fixed tests and \(\boldsymbol t\), not on
actual \(\gamma,k_{\rm new},f_{\rm new}\).  The coupled kernel uses
the bare \(k''\) coordinate, not the derived row; the latter appears
only in characters and \(\mathcal B_2\).  The derived label appears
only in characters and its outer cutoff.  Relations such as
\(g'=d_2r_g\) restrict evaluation points, not the transform.
For every fixed \(B_{\rm ht},J_{\rm ht}\ge0\), the Euler calculation
in Equation~\eqref{eq:inverse-fourier-moment} gives some fixed \(b\) with
\begin{equation}\label{eq:inverse-joint-fourier-moments}
 \int_{\mathbb R^9}\int_{\mathbb R^6}
 |\widehat{\mathfrak H}_{1,\sigma}(\boldsymbol t)|
 |\widehat{\mathfrak H}_{2,\sigma,\boldsymbol t}(\boldsymbol\zeta)|
 (1+|\boldsymbol t|)^{B_{\rm ht}}(1+|\boldsymbol\zeta|)^{J_{\rm ht}}
 \,d\boldsymbol\zeta\,d\boldsymbol t\ll p_b(W)^2.
\end{equation}
Indeed the inner weighted integral is polynomial in \(\boldsymbol t\),
because \(W_i\) is a fixed cutoff times a unit norm mode; apply the
first transform's bound at that polynomial order plus \(B_{\rm ht}\).
The constants are uniform in every actual ideal and every positive
kernel scalar, with no height-order power of \(Z\).

\paragraph{The second Cauchy inequality and the outer counts.}
Apply Equation~\eqref{eq:weighted-cauchy} to the second transformed
sum with
\(U_a=Z^{(\lambda_c+18\eta+\tau)/2}Q_{a,\gamma,\boldsymbol\zeta}\)
on the complete set \(\Xi_\sigma\), for each fixed Fourier mode.
In particular \(\mathcal B_2\), \(|B_o|^2\), \(|D_o|^2\), and
\(1_{(v',g'CJ)=1}\) are still in its weight. They retain the
squarefreeness of \(f_{\rm new}\) proved above. All other factors in
the weight are bounded by a fixed constant,
including the bounded number of original assigned slot coefficients,
the fixed ray coefficient, the individual cutoffs, and the unit modes.
Consequently, only after Cauchy, the absolute weights obey
\[
 |V_\sigma(\xi;\boldsymbol\zeta)|
 \ll \mathcal B_2(k'')\mu^2(f_{\rm new}).
\]
This is the step which permits the new averaged ideal to remain in
the squarefree class.

We first record the exact reconstruction needed to dominate a fibre.
\begin{equation}\label{eq:poisson-divisor-reconstruction}
 d_k\mid C\mathfrak B/R_1
 \quad\Longrightarrow\quad
 d_k\mid f_{\rm new}\operatorname{rad}q_0.
\end{equation}
To prove this, a prime of \(\mathfrak B/R_1\) has \(t_p=0\).
The local table shows that it has even parity and equal \(a_{ip}\).
If both are one it is in \(J_2\), and if both are zero it is in
\(\operatorname{rad}q_0\).  Primes of \(C\) are already in
\(f_{\rm new}\).  This proves the implication.  Thus \(d_k\) also
has divisor multiplicity after \(f_{\rm new},q_0\) are specified.
Once these factors are fixed, \(k''\) is uniquely determined by
the element \(k_{\rm new}=d_kd_2k''\).

Let \(K_{\rm slot}\) bound the original number of slots.  At fixed
\(\gamma=(q_0,t',r_g)\), squarefree \(f=f_{\rm new}\), and element
\(k=k_{\rm new}\), the number of possible \(\xi\in\Xi_\sigma\) with
nonzero weight is at most a fixed constant times
\(D_{K_{\rm slot}}(f)C_{K_{\rm slot}}(\gamma)\), where
\begin{equation}\label{eq:inverse-new-fibre}
 \begin{split}
 D_{K_{\rm slot}}(f)&=d_{\mathcal O}(f)^{9+4K_{\rm slot}},\\
 C_{K_{\rm slot}}(\gamma)&=
 d_{\mathcal O}(q_0)^{5+2K_{\rm slot}}
 d_{\mathcal O}(t')^{2K_{\rm slot}}
 d_{\mathcal O}(r_g)^{2K_{\rm slot}},
 \end{split}
\end{equation}
Here and below \(d_{\mathcal O}\) denotes the ideal divisor count.
For completeness, the ordered allocation of \(f\) into
\(J,C,d_2,v'\) costs \(4^{\omega(f)}=d_{\mathcal O}(f)^2\), and
splitting \(J=sJ_2\) costs at most \(d_{\mathcal O}(f)\).
The choices of \(b_1,b_2\) with
\(b_1b_2=q_0^2J_2^2s\) cost at most
\(d_{\mathcal O}(q_0)^2d_{\mathcal O}(f)^3\); the two choices
\(\mathfrak A_a\mid\operatorname{rad}(b_1b_2)\) cost at most
\(d_{\mathcal O}(q_0)^2d_{\mathcal O}(f)^2\).  These estimates use
\(d_{\mathcal O}(IJ)\le d_{\mathcal O}(I)d_{\mathcal O}(J)\) and
\(d_{\mathcal O}(I^2)\le d_{\mathcal O}(I)^2\).
Equation~\eqref{eq:poisson-divisor-reconstruction} costs at most
\(d_{\mathcal O}(q_0)d_{\mathcal O}(f)\) choices for \(d_k\).
Then \(g'=d_2r_g\) is fixed, and \(k''=k/(d_kd_2)\) is the unique
element if it is integral.  The at most \(2K_{\rm slot}\) old assigned
primes divide \(q_0ft'\), so their choices cost at most
\([d_{\mathcal O}(q_0)d_{\mathcal O}(f)d_{\mathcal O}(t')]^{2K_{\rm slot}}\).
The at most \(2K_{\rm slot}\) new assigned primes divide \(fr_g\),
and cost at most
\([d_{\mathcal O}(f)d_{\mathcal O}(r_g)]^{2K_{\rm slot}}\).
The other first side's surviving slots disappeared on selecting the
first positive side; the selected surviving slots occur twice in its
new square, exactly as counted here.  Thus the displayed product
dominates the fibre even if its actual size depends on \(k\).
There is no count of the old \(f,h\), no independent count of
\(J,C,d_2,v'\), and no second frequency count.

Equation~\eqref{eq:fixed-cube-count} and the same source witnesses
place \(\Gamma_\sigma\) inside the product of the three norm balls
\[
 q_{q_0}\le Z^{\ell+R/2-j+5\eta/2},\qquad
 q_{t'}\le Z^{t+\eta},\qquad q_{r_g}\le Z^{g-\vartheta_c+2\eta}.
\]
We use these balls only to bound the number of distinct triples, not
as a replacement domain for the child estimate.  Ideal counting gives
\begin{equation}\label{eq:canonical-fixed-count}
 \#\Gamma_\sigma\le C Z^{C_c+11\eta/2+\pi_{\rm count}},\qquad
 C_c=\ell+R/2-j+t+g-\vartheta_c.
\end{equation}
Here \(\pi_{\rm count}\) is part of the separately chosen aggregate
\(\pi\).  In a nonempty sector each upper exponent of these three
balls is nonnegative, since it bounds the norm of an existing ideal;
the constant term in ideal counting therefore adds no boundary power.
In particular \(J\) is not counted here.  The ranges of all three
ideals are bounded, so the adjustable divisor bound gives
\(C_{K_{\rm slot}}(\gamma)\ll Z^{\pi_{\rm fib}}\) uniformly on
the containing balls.  Put \(c_\sigma=C_c+11\eta/2\) and introduce
the genuinely row- and label-independent measure
\begin{equation}\label{eq:inverse-normalized-triples}
 d\nu_\sigma(\gamma)=Z^{-c_\sigma}
 \sum_{\gamma\in\Gamma_\sigma}C_{K_{\rm slot}}(\gamma)\delta_\gamma,
 \qquad \|\nu_\sigma\|\ll Z^{\pi_{\rm count}+\pi_{\rm fib}}.
\end{equation}

\paragraph{The positive canonical children.}
For a fixed Fourier mode define the positive child sums
\begin{equation}\label{eq:inverse-positive-child}
 \begin{split}
 \mathcal H_{a,\gamma}={}&
 \sum_{\substack{f\ {\rm sf}\\q_f/Z^{V_c}\in I_f}}
 D_{K_{\rm slot}}(f)
 \sum_{\substack{k\in\mathcal O\\0<q_k\le Z^{M_{\rm ch}}}}
 |Q_{a,\gamma,\boldsymbol\zeta}(k,f)|^2
 =Z^{F_{\rm ch}}\mathcal E_{a,\gamma}.
 \end{split}
\end{equation}
The equality is the normalization of Equation~\eqref{eq:canonical-energy}
with this same row ball, the label weight \(D_{K_{\rm slot}}\), and
the test and puncture in Equations~\eqref{eq:canonical-child-polynomial}
and \eqref{eq:canonical-child-data}; if a larger fixed row ball is
used there, this equality is instead an inequality in the needed
upper-bound direction.  The weight depends on \(f\) alone.
After reindexing, \(\mathcal B_2\) is exactly the displayed row ball.
Only the now positive outer row and label sums have been enlarged to
their full fixed windows.  The columns use their full fixed fresh
windows, but their puncture, mark, and test have not been deleted or
changed.  The actual fibre domination just proved gives
\[
 \begin{split}
 \left|\sum_{\xi\in\Xi_\sigma}V_\sigma Q_1\overline{Q_2}\right|
 &\le\prod_{a=1}^2
       \left(\sum_{\xi\in\Xi_\sigma}|V_\sigma||Q_a|^2\right)^{1/2}\\
 &\ll Z^{c_\sigma}\prod_{a=1}^2
       \left(\int_{\Gamma_\sigma}\mathcal H_{a,\gamma}
                    \,d\nu_\sigma(\gamma)\right)^{1/2}.
 \end{split}
\]
Here \(Q_a=Q_{a,\gamma,\boldsymbol\zeta}(k_{\rm new},f_{\rm new})\).
The two square roots each supply \(Z^{c_\sigma/2}\), hence there is
one fixed-count factor.  A uniform child bound supplies one total
mass \(\|\nu_\sigma\|\), not its square.  The child bound will be
used only for \(\gamma\in\Gamma_\sigma\), using the hypotheses already verified in
Equation~\eqref{eq:inverse-child-margins}.

\paragraph{The energy exponent.}
The formal centers satisfy the exact identity
\begin{equation}\label{eq:canonical-prefactor}
 \kappa_i+\lambda_c+C_c+2F_c=F.
\end{equation}
This follows by substituting their definitions and
\(H_c=2r-2B-M+4\ell+2V+\delta-j\); all extracted lengths cancel and the
result is \(r+3\ell+V=N+V\).

We now use the child estimate only for
\(\gamma\in\Gamma_\sigma\), for which the preceding margins and row
decrease have been proved.  If it gives
\(\mathcal E_{a,\gamma}\ll C(\boldsymbol t,\boldsymbol\zeta)
Z^{F_{\rm ch}+\epsilon_{\rm child}}\) uniformly there, with its
fixed polynomial height dependence, then
Equation~\eqref{eq:inverse-positive-child} and the final weighted
Cauchy bound above give
\[
 \left|\sum_{\xi\in\Xi_\sigma}V_\sigma Q_1\overline{Q_2}\right|
 \ll C(\boldsymbol t,\boldsymbol\zeta)
 Z^{c_\sigma+2F_{\rm ch}+\epsilon_{\rm child}}\|\nu_\sigma\|.
\]
The normalized mass is the one in
Equation~\eqref{eq:inverse-normalized-triples}; witness multiplicity
is already in its fibre weights.  The common Fourier density is then
integrated once using Equation~\eqref{eq:inverse-joint-fourier-moments}.
It is not part of the discrete measure and is not chosen afresh for
any row or label.  Thus both the genuine outer count and its normalized
mass occur once.  The averaged \(J,C,d_2,v'\) stay inside the positive
child sum and are not also counted.

If a child is bounded by \(Z^{F_{\rm ch}+\epsilon_{\rm child}}\), the
first prefactor error, the second prefactor error, the fixed count, and
the restored normalization give on each positive second-Cauchy side
\begin{equation}\label{eq:inverse-step-loss}
 \begin{split}
 &\kappa_i+\lambda_c+C_c+2F_{\rm ch}
      +(\tfrac92+18+\tfrac{11}2)\eta+\tau+\pi+\epsilon_{\rm child}\\
 &\qquad=F+28\eta+2\delta_N+\tau+\pi+\epsilon_{\rm child}\\
 &\qquad\le F+40\eta+\tau+\pi+\epsilon_{\rm child}.
 \end{split}
\end{equation}
Both copies of \(\delta_N\) are present: one restores the child column
normalization and the other occurs in its asserted exponent.  Each of
the two weighted Cauchy inequalities takes a geometric mean of its
positive sides, so it does not double this loss.  The first principal
bound is \(Z^{M+\pi}\), and the second principal bound above is at most
\(Z^{F+26\eta+\tau+\pi}\); both fit the same allowance.

The parameter \(\pi\) has an independent quantifier.  Let \(J_{\rm step}\)
be the maximum number of uses of a free divisor or sieve exponent and
of logarithmic dyadic-count factors in one fixed two-Poisson, two-Cauchy
factorization tree, maximized over the finitely many assignments with at
most the fixed slot bound.  Once \(L_{\rm pow}\) bounds the log-length of
every scale product in such a use, choose each free input exponent at most
\(\pi/(2J_{\rm step}\max(1,L_{\rm pow}))\).  For sufficiently large \(Z\),
the product of the at most \(J_{\rm step}\) logarithmic factors is at most
\(Z^{\pi/2}\).  These choices make all non-center, non-frequency local
losses at most \(CZ^\pi\).  They include the retained common frequency
dyads and the reconstruction divisor bounds.  In particular the distinct
allocations \(\pi_\beta,\pi_{\rm old},\pi_{\rm count},\pi_{\rm fib}\)
and the finite-union dyadic allocation are parts of this single
aggregate, not repeated allowances of size \(\pi\).  Fixed ray sums and fixed
seminorm constants are constants.  The same convention defines the
aggregate terminal loss \(\pi_{\rm ref}\).  In particular \(\pi\) hides
no multiple of \(\eta\) or \(\tau\).

This proves the conditional reduction estimate
\eqref{eq:canonical-reduction-contract}, including its common-measure and
finite-order uniformity once the choices below are made. It remains to fix
those choices uniformly and apply the reduction through the finite depth.

\paragraph{Order of choices and termination.}
For the requested canonical exponent \(\epsilon_c=\epsilon\), first fix
the starting ranges, tests, slot bound, \(c_*\), and then \(d,D\) as at
the beginning of the proof.  Choose
\begin{equation}\label{eq:inverse-parameter-order}
 \begin{gathered}
 \eta\le\min\{d/16,c_*/(14D),\epsilon_c/(160D),c_*/1000,1/100\},\\
 \tau\le\min\{\epsilon_c/(4D),c_*/1000,1/100\},\qquad
 \pi\le\epsilon_c/(4D),\qquad
 \tau_{\rm ref},\pi_{\rm ref}\le c_*/1000.
 \end{gathered}
\end{equation}
Use this same \(\tau\) in both Poisson comparisons.  At this point \(\pi\)
is a target aggregate loss; its subsidiary input exponents are chosen
only after the bounded scale range is known.

Choose the individual cutoffs through depth \(D\), including every full
larger or fresh support and the bounded clipping families described
above, and form \(\mathcal I_D,L_{\rm win}\).  Include the product
intervals for the active slots and the terminal constants
\(C_{\rm res},C_{\rm width},C_{\rm ker}\), maximized through the depth.
This is finite because there are finitely many factorization types per
passage and finitely many passages.  Put
\[
 F_{\max}=N_{\max}+V_{\max}+11D\eta,\qquad
 L_{\rm all}=100(1+M_{\max}+F_{\max}+z_{0,\max}+c_*).
\]
These bound all retained scale lengths.  Indeed,
Equation~\eqref{eq:inverse-row-decrease} and the bound following it give
\(M\le M_{\max}\) and \(N,V\le F_{\max}\).  A nonempty column window has
\(r\ge-\eta\), so \(\ell\le(F_{\max}+\eta)/3\).  The displayed parent
products bound every extracted divisor length by \(2F_{\max}+6\eta\).
For example \(d_k\mid C\operatorname{rad}(b_1b_2)\) gives
\(\delta\le r+2\ell+4\eta\), and Equation~\eqref{eq:enlarged-first-frequency}
then gives \(H_{\rm use}\le7F_{\max}+18\eta+\tau\).  The terminal active
conductor is supported on the row, \(f\), the puncture, and slots; their
total lengths are at most \(M_{\max}+2F_{\max}+z_{0,\max}\) plus the
already displayed ratio allowances.  Equations~\eqref{eq:inverse-terminal-width}
and \eqref{own:retained-energy-dyad} bound the retained terminal dual lengths by
\(4M_{\max}+4F_{\max}+2z_{0,\max}+3d+12\eta+\tau_{\rm ref}\).
Under Equation~\eqref{eq:inverse-parameter-order}, all these bounds and
the exact conductor product lengths are below \(L_{\rm all}\).
We may now take \(L_{\rm pow}=L_{\rm all}\) and choose the subsidiary
small exponents that realize \(\pi\) and \(\pi_{\rm ref}\).

Define \(B_{\rm crude}\) to be the maximum exponent obtained by replacing
every bounded nonfrequency ideal or element sum in either raw Poisson
expansion, including both column sums, by its lattice count at length
\(L_{\rm all}\), and every explicit norm factor, including \(1+a^{-1}\)
in Equation~\eqref{eq:inverse-raw-tail}, by its absolute upper bound.
Replace divisor-bounded coefficients and marks by their trivial polynomial
norm bounds as well.
Exclude only the frequency kernel itself.  The expansions contain a
fixed finite number of factors, so this defines a finite uniform number.

For the reflected tail define \(B_{\rm ref}\) analogously using its bounded
row, local, and active-product counts and \(1+a^{-1}\) with \(a=X/q_c^2\).
Do not count discarded dual indices at a retained length.  Instead,
Equation~\eqref{eq:theta-support} gives on its support
\[
 \frac{|d(\mu)|}{\sqrt{q_\mu}}
 \le27q_{\lambda_3}^{-k/3}q_n^{-1/2}q_b^{-1}
 \le27q_{\lambda_3}^{4/3}\qquad(k\ge-4).
\]
Every other local factor is bounded by the product of \(q_p^{1/2}\) over
the bounded active primes, independently of \(\mu\) after its indicators
are dropped.  Thus \(B_{\rm ref}\) covers all non-kernel coefficients and
bounded counts, while the remaining \(\mu\)-sum is an unrestricted
lattice sum on \(\lambda_3^{-4}\mathcal O\).

Fix a tail saving \(T>1+B_{\rm crude}+B_{\rm ref}\).  For each Poisson
tail, choose a Schwartz order
\[
 A>1+(B_{\rm crude}+T)/\tau.
\]
Equation~\eqref{eq:inverse-raw-tail} with \(Y=Z^\tau\) then makes the
discarded raw outer-ball complement \(O(Z^{-T})\), up to finitely many
input seminorms and a fixed polynomial in any twist heights.  For the
reflected whole-dyad tail the actual argument is at least
\(Z^{\tau_{\rm ref}/2}\); use the same lattice-shell bound on
\(\lambda_3^{-4}\mathcal O\) and choose
\[
 A_{\rm ref}>1+2(B_{\rm ref}+T)/\tau_{\rm ref}.
\]
This proves the discarded-tail assertion without presuming any bound on
the discarded dual lengths.

Next choose the finitely many Fourier-height and smooth-seminorm orders
backwards through depth \(D\), above the chosen tail orders and the
terminal input orders.  Normalized inverse roots are fixed real powers
on annuli, and Euler differentiation of a full kernel at a norm monomial
introduces no scale power.  Equations~\eqref{eq:inverse-joint-fourier-moments}
and \eqref{eq:inverse-normalized-triples} supply the common coefficient
measure for each positive indexed Cauchy-side sum.  Apply
the indexed finite-depth argument of Appendix~\ref{app:smooth} separately to those sums,
using the child-profile bounds just proved from
the smooth calculus of Appendix~\ref{app:smooth}, and then take the displayed geometric
means.  The two square roots retain one normalized discrete mass, as
already shown, so genuine outer labels are counted only once.
Any external height cutoff in a later application
is chosen after this internal finite propagation, not inserted into it.

Finally choose \(Z\) large enough for \(\log Z\ge L_{\rm win}/\eta\),
the thresholds for \(C_{\rm res},C_{\rm width}\) in
Equation~\eqref{eq:inverse-terminal-width},
\(\log Z\ge2\log C_{\rm ker}/\tau_{\rm ref}\),
\(\log Z\ge4\log q_{\lambda_3}/\eta\), and the fixed \(h\)-annulus threshold
used in the short completion.  Impose also the logarithmic bounds defining
the aggregate local losses and the thresholds of the input lemmas.
All choices precede this final threshold and depend only on the fixed
data.  Bounded smaller \(Z\) are handled by increasing the final constant.

Now \(c_h=c_*-7h\eta\ge c_*/2\) for \(h\le D\).
Equations~\eqref{eq:inverse-child-margins} and
\eqref{eq:inverse-row-decrease} send every retained nonterminal child to
the next depth and its row cap.  At depth \(D\) that cap is negative,
whereas every retained row parameter is nonnegative.  The terminal bounds
are at most \(Z^F\); principals and recursive terms cost at most
\(L_{\rm step}=40\eta+\tau+\pi\) in exponent.  Backwards induction gives
\(Z^{F+(D-h)L_{\rm step}}\) at depth \(h\), and
\[
 D L_{\rm step}\le\epsilon_c/4+\epsilon_c/4+\epsilon_c/4<\epsilon_c.
\]
This proves \(\mathcal E\ll Z^{F_0+\epsilon}\), with the claimed
uniform finite-seminorm and polynomial-height dependence.
\end{proof}

\subsection{Initialization of the marked moment}

We now convert the inverse polynomial in Lemma~\ref{lem:marked} to
the canonical family.  Only one Poisson transformation is required.

\begin{proof}[Proof of Lemma~\ref{lem:marked}]
In the product \(M_u(Z^r;W)Q_u\), let \(P\) be the product of its
slot primes and put \(j=(n,P)\).  Both \(n\) and \(P\) are
squarefree.  Fix the subset of slots occurring in \(j\), and write
\[
 n=jn_0,\qquad P=jP_0,\qquad (n_0,P_0)=1.
\]
The product \(c=n_0P_0\) is squarefree.  The identities
\[
 \mu(n_0)=\mu(c)\mu(P_0),\qquad
 \psi_u(n)\psi_u(P)=\psi_u(j)^2\psi_u(c)
\]
hold with all zero extensions.  The sign \(\mu(P_0)\) is a product
of signs on the surviving slots and can be incorporated into their
bounded coefficients.  The coefficients thus remain of the form
in Equation~\eqref{own:mark-definition}.

For this fixed assigned subset define \(G\) from its nominal slot centers:
\[
 G=\sum_{i\ {\rm assigned}}z_i,\qquad
 z_0=z-G,\qquad D'=r+z-2G.
\]
Then \(G,z_0\ge0\), and \(z_0\) is exactly the surviving nominal cap.
The normalization satisfies
\[
 Z^{-(r+z)/2}=Z^{-G}Z^{-D'/2}.
\]
The actual \(j\) is the product of the assigned primes, so
\(q_j/Z^G\) lies in a fixed product interval.  Ideal counting gives at
most \(CZ^G\) choices, apart from the separately chosen divisor loss
for assignments.  Triangle inequality in the row Hilbert space uses
\(Z^{-G}\) to cancel this count.  The whole factor \(\psi_u(j)^2\) is a
bounded row scalar, including its zeros, and is a contraction in that
space.  The residual column is prime to \(j\), giving the fixed
puncture \(1_{(c,j)=1}\).

Here is the exact identity underlying this overlap reduction.  Let
\(I_0\) be the surviving subset and put
\(y_j=q_j/Z^G\), \(y_c=q_c/Z^{D'}\), and
\(y_i=q_{p_i}/Z^{z_i}\) for every slot.  On the original support,
\[
 \frac{q_n}{Z^r}=\frac{y_jy_c}{\prod_{i\in I_0}y_i}.
\]
For a fixed assigned tuple, its contribution \(\mathcal U_{u,j}\) to
\(M_u(Z^r;W)Q_u\) is exactly
\begin{equation}\label{eq:initial-overlap-polynomial}
 \begin{split}
 \mathcal U_{u,j}={}&\mu(j)
 \prod_{i\notin I_0}a_i(p_i)W_i(y_i)\,
 Z^{-G}\psi_u(j)^2 Z^{-D'/2}
 \sum_{\substack{c\ {\rm sf}\\(c,j)=1}}\mu(c)\psi_u(c)\\
 &\quad\times
 \sum_{\substack{(p_i)\in\prod_{i\in I_0}\mathcal P_i\\P_0\mid c}}
 \mu(P_0)\prod_{i\in I_0}a_i(p_i)W_i(y_i)
 W\!\left(\frac{y_jy_c}{\prod_{i\in I_0}y_i}\right).
 \end{split}
\end{equation}
Indeed a squarefree \(c\) prime to \(j\), together with a surviving
tuple with \(P_0\mid c\), reconstructs uniquely
\(n=j(c/P_0)\) and the prescribed overlap.  Squarefreeness of \(c\)
already gives \((c/P_0,P_0)=1\), so no further condition was lost.
The sign \(\mu(P_0)\) is a product of the individual prime signs.
This also shows that the \(j\) on both sides of every subsequent
square is one fixed ideal: the preceding triangle inequality was
taken before that square.

The identity \(c=nP/j^2\) puts \(y_c\) in a fixed compact interval.
This remains meaningful if \(D'\) is slightly negative in a bounded
nonempty window; \(D'\) is not yet a canonical parameter.  Choose a
fresh individual cutoff \(\omega_c\) equal to one on this full quotient
support.  For fixed \(j\), let \(\Omega_c\) be one on the full support
of \(\omega_c\), and let \(\Omega_i\) be one on each full surviving
individual slot support.  The single joint profile
\[
 \mathfrak H_{{\rm ov},j}(y_c,(y_i)_{i\in I_0})
 =\Omega_c(y_c)\prod_{i\in I_0}\Omega_i(y_i)
 W\!\left(\frac{y_jy_c}{\prod_{i\in I_0}y_i}\right)
\]
is compactly supported in these independent coordinates.  The current
cutoff \(\omega_c\) and the individual factors \(W_i(y_i)\) stay outside
its transform.  Equation~\eqref{eq:inverse-log-fourier}, in dimension
\(1+|I_0|\), therefore expresses Equation~\eqref{eq:initial-overlap-polynomial}
exactly as the same assigned scalar and \(Z^{-G}\psi_u(j)^2\), times
\[
 (2\pi)^{-1-|I_0|}\int_{\mathbb R^{1+|I_0|}}
 \widehat{\mathfrak H}_{{\rm ov},j}(\boldsymbol\upsilon)
 Z^{-D'/2}\sum_{c\ {\rm sf}}\mu(c)\psi_u(c)1_{(c,j)=1}
 \mathfrak d_{\boldsymbol\upsilon}(c)
 \omega_c(y_c)y_c^{i\upsilon_c}\,d\boldsymbol\upsilon,
\]
where \(\mathfrak d_{\boldsymbol\upsilon}\) has the original surviving
lists and individual coefficients
\(\mu(p)a_i(p)W_i(y_i)y_i^{i\upsilon_i}\).  Expanding the mark and
inverting the joint transform proves this equality term by term.
No measure is chosen for an actual surviving tuple.  The measure may
depend on the already fixed \(j\), but Equation~\eqref{eq:inverse-fourier-moment}
is uniform for its \(y_j\) in the fixed product interval.  Minkowski's
inequality is used in the whole row Hilbert space for this one measure.
Include the full support of \(\omega_c\), its larger cutoffs, the
surviving slot supports, and the \(j\)-product interval in the finite
collection \(\mathcal I_D\) before the final threshold is chosen.
With the same \(\eta\) convention as the canonical proof, we then have
\[
 |\widehat c-D'|\le\eta,\qquad |\widehat j-G|\le\eta
\]
on the full supports actually used below, where \(\widehat c=\log_Zq_c\).
For the sign
\(\varepsilon_\chi=1\), conjugate the whole residual polynomial,
including its finite character, mark, and test; its row norm is unchanged.
For the other sign leave it unchanged.  It therefore suffices to estimate
one polynomial of the exact form
\begin{equation}\label{eq:initial-residual-polynomial}
 \mathcal R_0(u)=Z^{-D'/2}\sum_{c\ {\rm sf}}
 \mu(c)\nu_0(c)1_{(c,j)=1}\overline{\chi_c(u)}
 \mathfrak d_0(c)W_0(q_c/Z^{D'}).
\end{equation}
Here \(\nu_0\) is a fixed finite ray character, \(\mathfrak d_0\)
is the product mark of cap \(z_0\) with those individual coefficients
or their conjugates, and \(W_0\) is \(\omega_c(y)y^{i\upsilon_c}\)
or its conjugate.  All are independent of the row.

Majorize the original row ball by a fixed nonnegative radial Schwartz
function \(\Phi_{\rm init}\) at scale \(Z^m\), and expand its square.  Let \(C\) be the gcd
of the two columns and write \(c_i=Cz_i\), with \(z_1,z_2\) squarefree
and coprime.  Put \(\widehat c_i=\log_Zq_{c_i}\), so the preceding
support bound applies to each \(i=1,2\).  Let \(B\ge0\) be the center of \(C\), and write
\(\widehat B=\log_Zq_C\).  The row character has primitive modulus
\(z_1z_2\) and remaining zero mask \(C\).  Apply
Lemma~\ref{lem:masked-poisson}, with \(d'\mid C\).  Let \(\vartheta_c\ge0\)
be its center and \(\widehat\vartheta_c=\log_Zq_{d'}\).  Set
\[
 P_1=B-\vartheta_c\ge-2\eta,
\]
where the inequality follows from \(\widehat\vartheta_c\le\widehat B\) and
the two \(\eta\)-ratio bounds.  The actual conductor length is
\[
 \widehat L_{\rm init}=\widehat c_1+\widehat c_2-2\widehat B,
 \qquad |\widehat L_{\rm init}-2(D'-B)|\le4\eta.
\]
For a positive tolerance \(\tau_{\rm init}\), the implication
\(Z^mq_{h'}/(q_{d'}q_{z_1z_2})\le Z^{\tau_{\rm init}}\) gives
\[
 \widehat{d'h'}\le2D'-m-2P_1+6\eta+\tau_{\rm init}.
\]
Define the fixed formal and enclosing row scales
\begin{equation}\label{eq:inverse-initial-row}
 M_c^{\rm init}=2D'-m-2P_1,\qquad
 M_{\rm init}=M_c^{\rm init}+6\eta+\tau_{\rm init}.
\end{equation}
The \(6\eta\) is \(4\eta\) from the conductor and \(2\eta\) from the
two copies of \(d'\) in the new row bound.

Separate the original principal contribution for the direct count below.
On the genuine nonprincipal coprime Poisson expression, before any
off-coprime extension or Fourier absolutization, keep only the outer mask
\[
 \mathcal B_{\rm init}(h')=1_{0<q_{d'h'}\le Z^{M_{\rm init}}}.
\]
Its complement is contained in the actual ratio tail above for every
supported pair.  Equation~\eqref{eq:inverse-raw-tail}, with the separate
initial raw count and order specified below, discards precisely that
complement.  Retain the full smooth kernel inside the mask, with no
column-dependent ratio selector.  The mask depends on \(d',h'\), the
fixed overlap and sector, but not on the residual columns.  If
\(M_{\rm init}<0\), its nonzero ball is empty and the same tail comparison
discards every nonzero frequency; otherwise the retained row parameter is
nonnegative.

The actual normalization and Poisson prefactor split over the two sides as
\[
 Z^{-D'}Z^{m-\widehat\vartheta_c-\widehat L_{\rm init}/2}
   =Z^{\widehat\kappa_1^{\rm init}/2}
      Z^{\widehat\kappa_2^{\rm init}/2},
 \quad
 \widehat\kappa_i^{\rm init}=m-D'-\widehat\vartheta_c-\widehat c_i+\widehat B.
\]
Thus the actual side exponents obey
\begin{equation}\label{eq:initial-poisson-prefactor}
 \widehat\kappa_i^{\rm init}\le\kappa_c^{\rm init}+3\eta,
 \qquad \kappa_c^{\rm init}=m-2D'+P_1.
\end{equation}
All normalized real inverse roots will be kept in the single joint
profile below, not also in the final child tests.
Principal columns have the direct diagonal bound \(O(Z^{m+\epsilon})\),
also for a bounded nonempty negative \(D'\)-window.  Any principal nonzero
frequencies restored to the formal formula carry the same
\(\mathcal B_{\rm init}\) mask and are bounded by the full principal
restoration in Lemma~\ref{lem:masked-poisson}, since \(Z^m\ge1\).

Insert frequency dyads of \(q_{h'}\) from one partition common to all
\(d'\), not a label-recentered partition.  The outer mask implies
\(q_{h'}\le Z^{M_{\rm init}}\), so there are only logarithmically many
dyads in the bounded retained range.  Let \(H\) be the center of one
such bare-frequency dyad, and let \(\psi_C,\psi_{d'},\psi_{h'}\)
be the retained individual cutoffs on the current \(C,d',h'\) dyads,
evaluated at their normalized norms.

The orientation in Equation~\eqref{eq:initial-residual-polynomial}
is now fixed.  The calculation in the second canonical Poisson
transformation replaces \(\mu(z)\gamma_{-1}(z)\) by
\(\bar\alpha(z)\gamma_2(z)\chi_z(-1)\overline{G(z)}\).
On the second side the factor inside conjugation has ray factor
\(\overline{G(z)}\); together with the CRT cross phase the relative
ray factor is \(G([z_2][z_1]^{-1})\).  Expand it by the fixed group
formula above, and in a fixed ray summand put
\[
 \nu_*=\nu_0\overline\vartheta,\qquad
 A_*(z)=\bar\alpha(z)\gamma_2(z)\nu_*(z),\qquad
 \mathcal K_{\rm init}=\mathcal F\Phi_{\rm init}.
\]
For a preliminary pattern \(\Sigma_0\) fixing the overlap, these
dyads and the ray summand, and with the earlier mode held fixed, the masked
principal-restored nonzero component is exactly
\begin{equation}\label{eq:initial-formal-component}
 \begin{split}
 \mathcal T^{\rm init}_{\Sigma_0}={}&
 \sum_{\substack{C\ {\rm sf},\ d'\mid C\\h'\ne0}}
 \psi_C\psi_{d'}\psi_{h'}\mu(d')\widehat G(\vartheta)
 \mathcal B_{\rm init}(h')1_{(C,j)=1}\\
 &\times\sum_{\substack{z_1,z_2\ {\rm sf}\\(z_1z_2,Cj)=1}}
 1_{(z_1,z_2)=1}
 A_*(z_1)\overline{A_*(z_2)}
 \overline{\chi_{z_1}(d')}\chi_{z_1}(h')
 \chi_{z_2}(d')\overline{\chi_{z_2}(h')}\\
 &\times\mathfrak d_0(Cz_1)\overline{\mathfrak d_0(Cz_2)}
 W_0(q_Cq_{z_1}/Z^{D'})\overline{W_0(q_Cq_{z_2}/Z^{D'})}\\
 &\times\frac{Z^{-D'}Z^m}{q_{d'}\sqrt{q_{z_1}q_{z_2}}}
 \mathcal K_{\rm init}\!\left(
       \frac{Z^mq_{h'}}{q_{d'}q_{z_1}q_{z_2}}\right).
 \end{split}
\end{equation}
Indeed \(|\mu(C)\nu_0(C)|^2=1\) for the squarefree \(C\) outside
\(S\), the common row factor is exactly \(1_{(u,C)=1}\), and its
Poisson expansion contributes the one \(\mu(d')\).  The displayed
two numerator factors are those of the primitive character
\(\overline{\chi_{z_1}}\chi_{z_2}\) and its Fourier transform.
The finite ray sum restores the principal pair by \(G(1)=1\).
The cost of that restoration and the raw complement removed before
this equality were bounded separately above.

Define the summand in Equation~\eqref{eq:initial-formal-component}
on all individually squarefree \(z_1,z_2\) prime to \(Cj\) by the
displayed separate \(A_*\) factors, characters, formal product norms,
full kernel, and unchanged outer ball.  It agrees on the coprime
domain.  Insert the complete identity
\[
 1_{(z_1,z_2)=1}=\sum_{s\mid(z_1,z_2)}\mu(s),\qquad z_a=sn_a,
\]
before any factorwise bound.  Let \(v\ge0\) be the center of a
squarefree \(s\)-dyad, \(\widehat v=\log_Zq_s\), and \(\psi_s\)
its retained individual cutoff.  The \(n_a\) are squarefree and
prime to \(s\), but need not be mutually coprime.  This formal
extension does not apply Poisson summation to a noncoprime conductor.
For squarefree coprime \(s,n\), CRT and the zero-preserving identity
\(\overline{\chi_n(d')}=\chi_n(d')\chi_n(d')^4\) give exactly
\begin{equation}\label{eq:initial-common-extraction}
 \begin{split}
 A_*(sn)\overline{\chi_{sn}(d')}\chi_{sn}(h')
 ={}&A_*(s)\overline{\chi_s(d')}\chi_s(h')\\
 &\times A_*(n)\chi_n(d'h')\chi_n(d's)^4.
 \end{split}
\end{equation}
The fourth power includes the CRT factor \(\chi_n(s)^4\).
On an overlap it defines the residual coefficient to be zero, without
evaluating \(\gamma_2\) on a nonsquarefree ideal.  Since \(|A_*(s)|=1\)
outside \(S\), the two extracted scalars give
\[
 |A_*(s)\overline{\chi_s(d')}\chi_s(h')|^2
 =1_{(s,d'h')=1}=1_{(s,h')=1}
 \quad\text{when }d'\mid C,\ (s,C)=1.
\]
The original residual and overlap exclusions also give
\(1_{(C,j)=1}1_{(s,Cj)=1}\).  Thus, apart from the assigned
coefficients, individual cutoffs, and unit phases, the outer arithmetic
weight is precisely
\begin{equation}\label{eq:initial-outer-arithmetic}
 \mu(d')\mu(s)\widehat G(\vartheta)\mathcal B_{\rm init}(h')
 1_{(C,j)=1}1_{(s,Cj)=1}1_{(s,h')=1}.
\end{equation}
All these factors remain through weighted Cauchy.  Set
\[
 t=C/d',\qquad k=d'h',\qquad f=d's,\qquad
 \rho_{t,j}(n)=1_{(n,\operatorname{rad}(tj))=1}.
\]
The factors on \(n\) in Equation~\eqref{eq:initial-common-extraction},
together with \(\rho_{t,j}\), supply exactly the original exclusions
at \(C,s,j\) and the numerator zeros at \(d',h'\).  On the nonzero
support \(t,d',s,n\) are pairwise coprime.  Apply
Equation~\eqref{eq:inverse-slot-priority} to each mark with this
ordered list.  Assigned primes and their priority masks stay outer.
A surviving prime is already prime to \(d's\) by \(\chi_n(f)^4\)
and to \(tj\) by \(\rho_{t,j}\), so its original individual list and
coefficient can be retained.  Its mark is an original subcollection
of cap \(z_0\), not a row- or label-dependent residual coefficient.

The substitutions \(k=d'h'\), \(f=d's\), and the puncture
\(\rho_{t,j}\) have therefore supplied the canonical character and
coefficient class.  We still need a common separated profile, the
outer multiplicity bound, and admissibility of the resulting ranges.

Keep \(\mathcal B_{\rm init}\) outside the common Fourier separation
of the full kernel and the old windows.  The exact
formal child centers and the full fresh-support bounds are
\begin{equation}\label{eq:initial-canonical-lengths}
 \begin{gathered}
 N_c^{\rm init}=D'-B-v,\qquad V_c^{\rm init}=\vartheta_c+v,\qquad
 F_c^{\rm init}=D'-P_1,\\
 |\widehat n-N_c^{\rm init}|\le3\eta,\qquad
 |\widehat f-V_c^{\rm init}|\le2\eta.
 \end{gathered}
\end{equation}
On factorized terms these bounds follow from \(c=Csn\) and \(f=d's\).
The full fresh column cutoff and enlarged label window are also included
in \(\mathcal I_D\), so the same bounds hold directly for independent
terms later added by positivity.  The reconstructed formal product
\(c=Csn\) on the fresh cutoff is included there as well, so the initial
prefactor bound holds on its entire separated side.

Fix a refinement \(\Sigma\) by the \(s\)-dyad and the slot branches,
with the earlier mode \(\boldsymbol\upsilon\) held fixed.  Before
summing any actual outer ideal, use the six independent coordinates
\[
 (x_C,x_d,x_s,x_h,x_1,x_2)
 =\left(\frac{q_C}{Z^B},\frac{q_{d'}}{Z^\vartheta_c},
 \frac{q_s}{Z^v},\frac{q_{h'}}{Z^H},
 \frac{q_{n_1}}{Z^{N_c^{\rm init}}},
 \frac{q_{n_2}}{Z^{N_c^{\rm init}}}\right).
\]
Let \(I_C,I_d,I_s,I_h\) be the full supports of the four current
individual cutoffs.  Choose \(\omega_{n,a}\) equal to one on the
support of \(W_0\) divided by \(I_CI_s\), with fixed full supports
in an enclosing interval \([a,b]\) whose upper endpoint satisfies
\(b\ge1\).  Choose a nonnegative \(\omega_f\) equal to one on
\(I_dI_s\), with fixed full support \(I_f\).  The current four dyad
cutoffs and \(\omega_f\) remain outside inversion, and
\(\omega_{n,a}\) remain in the columns.  Let
\(\Omega_{{\rm init},\alpha}\) be larger cutoffs equal to one on the
full supports of the corresponding four current and two fresh cutoffs.
These full supports, \(I_f\), and the formal products \(Csn\) on them
are the windows already included in \(\mathcal I_D\).  None of these
supports depends on a current outer tuple or on a Fourier mode.

Keep \(V_c^{\rm init}\ge0\).  Proceed to a child only if the retained
row component is nonzero, which implies \(M_{\rm init}\ge0\).
This is a separate gate from the structural outer set below, which
may include zero-weight tuples even when the row ball is empty.
Independently, if a full fresh child column window contains no
squarefree ideal, its polynomial is zero and that component is omitted.
For a retained nonempty child define
\[
 N_{\rm init}=\max(0,N_c^{\rm init}),\qquad
 \delta_{\rm init}=N_{\rm init}-N_c^{\rm init}\in[0,3\eta],\qquad
 F_{\rm init}=N_{\rm init}+V_c^{\rm init}=F_c^{\rm init}+\delta_{\rm init}.
\]
For a nonempty negative-center window, \(q_n\ge1\) and
Equation~\eqref{eq:initial-canonical-lengths} give
\(\delta_{\rm init}\le3\eta\), while its upper support endpoint gives
\(Z^{\delta_{\rm init}}\le b\).  If \(N_c^{\rm init}\ge0\), then
\(\delta_{\rm init}=0\) and the latter inequality follows from \(b\ge1\).
Thus \(1\le Z^{\delta_{\rm init}}\le b\) in every retained case.
The cutoff \(\omega_{n,a}(Z^{\delta_{\rm init}}y)\) is supported in
\([a/b,b]\) and has uniformly bounded Euler seminorms.  Its rescaling
depends only on fixed centers, not on a current row or label, and its
full clipping family is included in \(\mathcal I_D\).

On the formal factorization \(z_a=sn_a\), the complete root and
kernel argument are exactly
\begin{equation}\label{eq:initial-root-kernel}
 \begin{split}
 \frac{Z^{-D'}Z^m}{q_{d'}q_s\sqrt{q_{n_1}q_{n_2}}}
 &=Z^{\kappa_c^{\rm init}}x_d^{-1}x_s^{-1}(x_1x_2)^{-1/2},\\
 \frac{Z^mq_{h'}}{q_{d'}q_s^2q_{n_1}q_{n_2}}
 &=A_{{\rm ker},{\rm init}}\frac{x_h}{x_dx_s^2x_1x_2},\\
 A_{{\rm ker},{\rm init}}&=Z^{m+H-\vartheta_c-2(D'-B)}.
 \end{split}
\end{equation}
For the first equality use \(N_c^{\rm init}=D'-B-v\) and
\(\kappa_c^{\rm init}=m-2D'+B-\vartheta_c\).  Equivalently its root is
\(x_d^{-1}x_C(x_{c_1}x_{c_2})^{-1/2}\) with
\(x_{c_a}=x_Cx_sx_a\).  Define the one joint profile
\begin{equation}\label{eq:initial-joint-profile}
 \begin{split}
 \mathfrak H_{{\rm init},\Sigma,\boldsymbol\upsilon}(\boldsymbol x)
 ={}&Z^{-3\eta}\prod_\alpha\Omega_{{\rm init},\alpha}(x_\alpha)
 x_d^{-1}x_s^{-1}(x_1x_2)^{-1/2}\\
 &\times W_0(x_Cx_sx_1)\overline{W_0(x_Cx_sx_2)}
 \mathcal K_{\rm init}\!\left(
 A_{{\rm ker},{\rm init}}\frac{x_h}{x_dx_s^2x_1x_2}\right).
 \end{split}
\end{equation}
Its outside scalar is exactly \(Z^{\kappa_c^{\rm init}+3\eta}\).
The old \(c\)-windows are coupled and therefore are in this profile,
as are all normalized real inverse roots and the full kernel.  No
real root is also put in a child test.

For \(\boldsymbol\xi=(\xi_C,\xi_d,\xi_s,\xi_h,\xi_1,\xi_2)\), put
\(z_n=q_n/Z^{N_{\rm init}}\), set again
\(\varepsilon_1=1,\varepsilon_2=-1\), and define
\[
 W_a^{\rm init}(z)=\omega_{n,a}(Z^{\delta_{\rm init}}z)
                       z^{\varepsilon_a i\xi_a}.
\]
If \(\mathfrak d_a'\) is the surviving subcollection on side \(a\),
the precise unnormalized child polynomial is
\begin{equation}\label{eq:initial-child-polynomial}
 Q^{\rm init}_{a,t,j,\boldsymbol\xi}(k,f)
 =\sum_{n\ {\rm sf}}A_*(n)\rho_{t,j}(n)
 \chi_n(k)\chi_n(f)^4\mathfrak d_a'(n)
 W_a^{\rm init}(q_n/Z^{N_{\rm init}}).
\end{equation}
Its fixed finite character, puncture, product mark, and test are
independent of \(k,f\).  The two sides can have different tests and
subcollections, but both have the same canonical coefficient class.

Let \(\Omega_\Sigma\) consist of
\[
 \omega=(C\ {\rm sf},d'\mid C,s\ {\rm sf},h'\ne0;
                      \text{assigned slot primes})
\]
with the retained individual supports and slot branches.  It has no
current column index.  The functions \(t,k,f\) are the products
defined above; the set may be taken before imposing the displayed
outer zero masks, which will be in the weight.  Let
\(\mathcal A_\Sigma^{\rm init}(\omega)\) be the product of all
assigned coefficients of \(\mathfrak d_0\) and their priority masks,
conjugated on side two.  The complete remaining outer weight is
\begin{equation}\label{eq:initial-outer-weight}
 \begin{split}
 V_\Sigma^{\rm init}(\omega;\boldsymbol\xi)={}&
 \mu(d')\mu(s)\widehat G(\vartheta)\mathcal B_{\rm init}(h')
 1_{(C,j)=1}1_{(s,Cj)=1}1_{(s,h')=1}
 \mathcal A_\Sigma^{\rm init}(\omega)\\
 &\times\psi_C(x_C)\psi_{d'}(x_d)\psi_s(x_s)\psi_{h'}(x_h)
 \omega_f(q_f/Z^{V_c^{\rm init}})\\
 &\times x_C^{i\xi_C}x_d^{i\xi_d}x_s^{i\xi_s}x_h^{i\xi_h}
 Z^{i\delta_{\rm init}(\xi_1+\xi_2)}.
 \end{split}
\end{equation}
The cutoff \(\omega_f\) is one on every original label product, so
its insertion is an equality before separation.  There is one copy
of each M\"obius factor and one ray coefficient.  In particular the
common zero at \((s,h')\) is still present.

Let \(\mathcal T_\Sigma^{\rm init}\) be the fixed \(s\)-dyad and
slot summand obtained from Equation~\eqref{eq:initial-formal-component}
by the complete M\"obius and priority expansions.  Its exact separated
identity is
\begin{equation}\label{eq:initial-separated-identity}
 \begin{split}
 \mathcal T_\Sigma^{\rm init}={}&
 Z^{\kappa_c^{\rm init}+3\eta}(2\pi)^{-6}
 \int_{\mathbb R^6}
 \widehat{\mathfrak H}_{{\rm init},\Sigma,\boldsymbol\upsilon}
       (\boldsymbol\xi)\\
 &\quad\times\sum_{\omega\in\Omega_\Sigma}
 V_\Sigma^{\rm init}(\omega;\boldsymbol\xi)
 Q^{\rm init}_{1,t,j,\boldsymbol\xi}(k,f)
 \overline{Q^{\rm init}_{2,t,j,\boldsymbol\xi}(k,f)}
 \,d\boldsymbol\xi.
 \end{split}
\end{equation}
Indeed \(x_a=Z^{\delta_{\rm init}}z_{n_a}\).  The two column modes,
with the last factor of Equation~\eqref{eq:initial-outer-weight},
are exactly \(x_1^{i\xi_1}x_2^{i\xi_2}\), and the other four modes
are outer.  Fourier inversion restores the full profile; its larger
cutoffs are one on the retained supports, and the fresh column
cutoffs are one wherever the old \(W_0\) windows are nonzero.
Equation~\eqref{eq:initial-root-kernel} restores the full prefactor
and kernel.  Equation~\eqref{eq:initial-common-extraction} and the
priority expansion restore every arithmetic factor term by term.
Extra combinations in the full fresh windows cancel through the old
windows in this complex identity before Cauchy.

The transform in Equation~\eqref{eq:initial-separated-identity}
is chosen from the fixed ambient profile before any actual
\(C,d',s,h'\), and hence before any \(t,k,f\), is evaluated.  It
depends on \(\Sigma,Z\), the fixed tests, and the earlier mode
\(\boldsymbol\upsilon\), and may depend on the already frozen \(j\),
but not on any current outer ideal or reindexed row or label.  The bare
\(h'\) norm is a coordinate; the derived row and label occur only
in the outer masks or in the displayed characters.  The relation
\(C=td'\) only restricts evaluation points.  By
Equation~\eqref{eq:inverse-fourier-moment}, for each fixed \(J\)
there are fixed \(L,b_0\) such that
\[
 \int_{\mathbb R^6}
 |\widehat{\mathfrak H}_{{\rm init},\Sigma,\boldsymbol\upsilon}
       (\boldsymbol\xi)|(1+|\boldsymbol\xi|)^J\,d\boldsymbol\xi
 \ll (1+|\boldsymbol\upsilon|)^L,\qquad
 \int_{\mathbb R^{1+|I_0|}}|\widehat{\mathfrak H}_{{\rm ov},j}(\boldsymbol\upsilon)|
       (1+|\boldsymbol\upsilon|)^L\,d\boldsymbol\upsilon
 \ll p_{b_0}(W).
\]
The first bound uses the fixed normalized roots and the uniform Euler
bounds of the full kernel for every positive
\(A_{{\rm ker},{\rm init}}\); \(W_0\) has only polynomial
\(\boldsymbol\upsilon\)-dependence.  The second is uniform over the
fixed \(j\)-product interval.  Taking larger \(L\) if needed also
controls the earlier row-space Minkowski inequality.  No derivative
order introduces a power of \(Z\).

Equation~\eqref{eq:initial-separated-identity} now expresses the component
in polynomials of canonical shape with one common Fourier density chosen
before the current outer ideals are evaluated.  It remains to bound the
positive fibres and source measure, verify child admissibility from source
witnesses, and restore the normalization in Lemma~\ref{lem:marked}.

Apply Equation~\eqref{eq:weighted-cauchy} on the whole
\(\Omega_\Sigma\), including all quotients \(t\), before fixing one
for the child.  Only afterward take absolute weights.  The retained
\((s,C)=1\) and \(d'\mid C\) make \(f=d's\) squarefree, and all
other factors of the weight are bounded by a fixed constant.  Thus
\[
 |V_\Sigma^{\rm init}(\omega;\boldsymbol\xi)|
 \ll\mathcal B_{\rm init}(h')\mu^2(f).
\]
For fixed \(t,f,k\), the identities
\[
 d'\mid f,\qquad s=f/d',\qquad C=td',\qquad h'=k/d'
       \quad\text{if this quotient is an element}
\]
show that the reconstruction has at most \(d_{\mathcal O}(f)\)
choices before slots.  Here \(K_{\rm slot}\) again denotes the fixed
bound for the original number of slots.  The at most \(2K_{\rm slot}\) assigned
residual primes divide \(tf\).  Their choices cost at most
\([d_{\mathcal O}(t)d_{\mathcal O}(f)]^{2K_{\rm slot}}\).
The earlier overlap primes were already frozen and counted in the
overlap triangle.  Hence the fibre, even if its actual size depends
on \(k\), is bounded by
\begin{equation}\label{eq:initial-fibre}
 D_{\rm init}(f)C_{\rm init}(t),\qquad
 D_{\rm init}(f)=d_{\mathcal O}(f)^{1+2K_{\rm slot}},\qquad
 C_{\rm init}(t)=d_{\mathcal O}(t)^{2K_{\rm slot}}.
\end{equation}
There is no independent \(d',s\) count and no second frequency count.

Let \(\Omega_\Sigma^{\rm src}\) be the structural outer set just
used, before independent positive row or label additions and ignoring
unit Fourier phases, and define the set of distinct source quotients
\[
 \mathfrak T_\Sigma=\{t=C/d':\omega\in\Omega_\Sigma^{\rm src}\}.
\]
It is projected once over all current rows and labels, not redefined
for a fixed \(k,f\).  Every member has a source witness on the
\(C,d'\) dyads, and therefore
\(\log_Zq_t=\widehat B-\widehat\vartheta_c\le P_1+2\eta\).
The containing norm ball bounds its cardinality by
\(\#\mathfrak T_\Sigma\ll Z^{P_1+2\eta+\pi_{{\rm count},{\rm init}}}\);
it is not used as a replacement child domain.  In a nonempty sector
\(P_1+2\eta\ge0\), since it bounds the norm of an existing ideal.
The divisor bound on this bounded range gives the normalized measure
\begin{equation}\label{eq:initial-normalized-quotients}
 d\nu_\Sigma^{\rm init}(t)
 =Z^{-P_1-2\eta}\sum_{t\in\mathfrak T_\Sigma}C_{\rm init}(t)\delta_t,
 \qquad
 \|\nu_\Sigma^{\rm init}\|
 \ll Z^{\pi_{{\rm count},{\rm init}}+\pi_{{\rm fib},{\rm init}}}.
\end{equation}
This set and measure ignore the unit Fourier phases and are common
to the whole current row and label sum.  The nonempty gates remain
separate from this structural projection.

For a fixed mode let
\begin{equation}\label{eq:initial-positive-child}
 \mathcal H_{a,t}^{\rm init}
 =\sum_{\substack{f\ {\rm sf}\\q_f/Z^{V_c^{\rm init}}\in I_f}}
 D_{\rm init}(f)\sum_{\substack{k\in\mathcal O\\0<q_k\le Z^{M_{\rm init}}}}
 |Q^{\rm init}_{a,t,j,\boldsymbol\xi}(k,f)|^2
 =Z^{F_{\rm init}}\mathcal E_{a,t}^{\rm init}.
\end{equation}
This is Equation~\eqref{eq:canonical-energy} with the same row ball;
with a larger fixed canonical ball the equality is replaced by the
corresponding upper bound.  The fibre estimate and weighted Cauchy give
\begin{equation}\label{eq:initial-weighted-count}
 \left|\sum_{\omega\in\Omega_\Sigma}V_\Sigma^{\rm init}Q_1\overline{Q_2}\right|
 \ll Z^{P_1+2\eta}
 \prod_{a=1}^2
 \left(\int_{\mathfrak T_\Sigma}\mathcal H_{a,t}^{\rm init}
                         \,d\nu_\Sigma^{\rm init}(t)\right)^{1/2}.
\end{equation}
Here \(Q_a=Q^{\rm init}_{a,t,j,\boldsymbol\xi}(k,f)\).
Only the positive outer \(f,k\) sums were enlarged to their full
fixed windows; the inner \(\rho_{t,j}\), mark, and fresh test are
unchanged.  The weight \(D_{\rm init}\) depends on \(f\) alone.
The new outer ball is exactly \(1_{0<q_k\le Z^{M_{\rm init}}}\).
Thus use of the uniform child bound introduces the displayed quotient-count
exponent once and the normalized mass once.

The fixed puncture is exactly at \(\operatorname{rad}(tj)\).
For each \(t\in\mathfrak T_\Sigma\), its source bound and the
fixed overlap bound give
\begin{equation}\label{old-eq:4.13}
 Q_{\rm init}:=\log_Zq_{\mathfrak r_\rho}
 \le\widehat B-\widehat\vartheta_c+\widehat j\le P_1+G+3\eta.
\end{equation}
Shared puncture primes only decrease this bound.  It continues to hold
when positive new labels or rows are added, because those additions
do not change \(t,j\) or delete their inner puncture.

The fixed formal centers satisfy
\[
 \begin{aligned}
 F_c^{\rm init}-M_c^{\rm init}-(P_1+G)-z_0&=m-r-2z+2G,\\
 4F_c^{\rm init}-3M_c^{\rm init}-6z_0&=3m-2r-8z+10G+2P_1.
 \end{aligned}
\]
Use the marked premises, \(G\ge0\), \(P_1\ge-2\eta\),
Equations~\eqref{eq:inverse-initial-row} and \eqref{old-eq:4.13}, and
the favorable clipping signs.  The actual canonical margins are
\begin{equation}\label{eq:inverse-initial-margins}
 \begin{aligned}
 F_{\rm init}-M_{\rm init}-Q_{\rm init}-z_0
   &\ge c_1+\delta_{\rm init}-9\eta-\tau_{\rm init}
     \ge c_1-9\eta-\tau_{\rm init},\\
 4F_{\rm init}-3M_{\rm init}-6z_0
   &\ge c_2+4\delta_{\rm init}-22\eta-3\tau_{\rm init}
     \ge c_2-22\eta-3\tau_{\rm init}.
 \end{aligned}
\end{equation}
The \(22\eta\) is \(4\eta\) from the possible negative \(2P_1\) and
\(18\eta\) from three copies of the row enclosure.

There is also an explicit energy calculation.  Apply the canonical
bound only for \(t\in\mathfrak T_\Sigma\), where the preceding puncture
and margin checks hold.  Equations~\eqref{eq:initial-positive-child}
and \eqref{eq:initial-weighted-count} then bound the unscaled signed
outer sum by
\[
 C(\boldsymbol\upsilon,\boldsymbol\xi)
 Z^{P_1+2\eta+2F_{\rm init}+\epsilon_c}
 \|\nu_\Sigma^{\rm init}\|.
\]
The one mass is that of Equation~\eqref{eq:initial-normalized-quotients};
the label \(d's\) remains inside the child average.  The common
Fourier density is integrated separately, once, with the uniform
weighted bounds above.  This counts the distinct frozen quotient
once, with its witness multiplicity already in the fibre weights.
The formal identity
\[
 \kappa_c^{\rm init}+P_1+2F_c^{\rm init}=m
\]
then gives, after restoring the child normalization and applying the
canonical bound with loss \(\epsilon_c\), the per-side exponent
\begin{equation}\label{eq:inverse-initial-energy}
 \begin{split}
 &m+3\eta+2\eta+2\delta_{\rm init}+\pi_{\rm init}+\epsilon_c\\
 &\qquad\le m+11\eta+\pi_{\rm init}+\epsilon_c.
 \end{split}
\end{equation}
Here \(\pi_{\rm init}\) is the aggregate freely chosen local divisor,
dyadic, and separation loss, including the assigned-overlap divisor
loss and the distinct allocations
\(\pi_{{\rm count},{\rm init}},\pi_{{\rm fib},{\rm init}}\).
These are parts of one aggregate, not repeated allowances.
The two clipping copies come from the restored normalization and
the child exponent.  There is no \(\tau_{\rm init}\) energy factor: the
row enlargement occurs inside the canonical child and was accounted for
in its margin test.  The initial weighted Cauchy takes a geometric mean,
and the overlap triangle was already canceled by \(Z^{-G}\).

We give the parameter and tail order explicitly.  Let
\(r_{\max},z_{\max},m_{\max}\) bound the marked ranges, let
\(s_{\rm slot}\) bound the number of slots, and put
\(D_{\max}=r_{\max}+z_{\max}\).  Set
\(\bar c=\min(c_1,c_2)\) and choose the canonical margin
\(c_*=\bar c/2\).  For \(\eta,\tau_{\rm init}\le1/100\), valid a priori
starting bounds for that application are
\[
 M_{\max}=2D_{\max}+1,\qquad N_{\max}=V_{\max}=D_{\max}+1.
\]
Indeed, \(P_1\ge-2\eta\) gives
\(M_{\rm init}\le2D_{\max}+10\eta+\tau_{\rm init}\), and
\(F_{\rm init}\le D_{\max}+5\eta\) bounds both nonnegative child
parameters.  Make the choices in Equation~\eqref{eq:inverse-parameter-order}
with \(\epsilon_c=\epsilon/2\), imposing in addition
\[
 \eta,\tau_{\rm init}\le\bar c/100,\qquad
 \eta\le\epsilon/88,\qquad \pi_{\rm init}\le\epsilon/8.
\]
One may take \(\tau_{\rm init}=\tau\) after imposing both sets of bounds.
Equation~\eqref{eq:inverse-initial-margins} then gives both margins at
least \(3\bar c/4>c_*\), so Lemma~\ref{lem:canonical-moment} applies.
Equation~\eqref{eq:inverse-initial-energy} adds at most \(\epsilon/4\)
to its \(\epsilon_c=\epsilon/2\) loss.

With these bounded initial ranges fixed, choose the subsidiary input
exponents realizing the target \(\pi_{\rm init}\), before choosing the
initial tail and Fourier-height orders.

For the separate initial raw tail, the supported columns satisfy
\(q_{c_i}\le C Z^{D_{\max}}\), while \(D'\in[-z_{\max},D_{\max}]\).
On the genuine initial expansion put
\(a=Z^m/(q_{d'}q_{z_1z_2})\).  Since \(m\ge0\), \(d'\mid C\), and
\(c_i=Cz_i\),
\[
 a^{-1}\le q_{d'}q_{z_1z_2}
       \le q_{c_1}q_{c_2}/q_C\le C Z^{2D_{\max}}.
\]
The explicit crude exponent
\[
 B_{\rm crude,init}=10\{1+m_{\max}+(s_{\rm slot}+1)(D_{\max}+1)\}
\]
bounds all raw nonfrequency factors even without cancellation.  To verify
this, the normalization costs at most \(z_{\max}\), the two column counts
cost \(2D_{\max}\), the \(d'\)-divisor count at most \(D_{\max}\), and
\(1+a^{-1}\) at most \(2D_{\max}\).  The two residual marks cost at most
\(2s_{\rm slot}D_{\max}\), using \(d_{\mathcal O}(c)^{s_{\rm slot}}\le q_c^{s_{\rm slot}}\).
Even counting overlap pairs and their assigned coefficients without
their canceling weights costs at most
\(2z_{\max}+2s_{\rm slot}z_{\max}\).  The Poisson prefactor costs at most
\(m_{\max}\).  The sum is at most
\(m_{\max}+(8+4s_{\rm slot})D_{\max}<B_{\rm crude,init}\).

Choose \(T_{\rm init}>1+B_{\rm crude,init}\) and then
\[
 A_{\rm init}>1+(B_{\rm crude,init}+T_{\rm init})/\tau_{\rm init}.
\]
Equation~\eqref{eq:inverse-raw-tail} makes the raw
\(\mathcal B_{\rm init}=0\) complement \(O(Z^{-T_{\rm init}})\), with
fixed seminorm and polynomial-height factors.  This order is chosen after
the marked ranges and \(\tau_{\rm init}\), but before the initial
Fourier-height orders and before the final \(Z\) threshold.  If
\(\tau_{\rm init}=\tau\), take the maximum of the orders required by this
separate crude bound and the canonical tails.  The full kernel retained
inside the outer ball has the uniform Euler bounds already used in the
canonical separation.  The final threshold also enforces the full initial support
ratios in \(\mathcal I_D\).

The principal count has exponent \(m\), and every nonprincipal term has
now been bounded by \(Z^{m+\epsilon}\).  The argument keeps all element
rows and every original mask, applies to subcollections and either common
orientation, and preserves arbitrary bounded row-independent prime
coefficients.  The common separated measures give the stated finite-
seminorm and polynomial-height uniformity.  This proves
Equation~\eqref{old-eq:4.1}.
\end{proof}

\subsection{Sixth-power amplification, including its endpoint}

For sixth-power-free element rows \(u\), let \(D\ge1\),
\(H=\max(2U,D^{1+c})\), \(P=(H/U)^{1/6}\), with fixed \(c>0\).
The zero-slot case of Lemma~\ref{lem:marked}, with fixed strict margins depending
on \(c\), and Sobolev in the scale, gives
\[
 \sum_{0<q_v\ll H}\sup_{0<D'\le D}|M_v(D';W)|^2
                       \ll H(HD)^\epsilon .
\]
The scale derivative replaces \(W(y)\) by
\(-W(y)/2-yW'(y)\); only \(O(\log(2+D))\) boxes are needed.
If a scale \(D'<1\) is nonempty, the upper endpoint of the fixed
annular support bounds \(D'\) below by a fixed positive constant.
It is therefore covered by a bounded dilation of the scale-one
profile; smaller scales are zero.
There are \(\gg_S P\) good primary ideals \(a\) with \(q_a\le P\).
For each such \(a\), splitting the squarefree column at the
primes dividing \(a\) gives the exact identity
\[
 M_u(D;W)=\sum_{d\mid\operatorname{rad}a}
           \mu(d)\psi_u(d)q_d^{-1/2}M_{ua^6}(D/q_d;W).
 \tag{D17}\label{eq:sixth-power-mask-identity}
\]
It remains true when \(a,u\) share primes.
The coefficient mass is divisor-bounded.
The map \((u,a)\mapsto ua^6\) is injective: prime valuations
modulo six recover \(u\), and their quotients by six recover \(a\);
the primary generator then fixes the unit of \(u\).
Averaging \eqref{eq:sixth-power-mask-identity} over \(a\) and using the scale supremum therefore
gives
\[
 \sum_{q_u\asymp U}^{(6)}|M_u(D)|^2
        \ll (H/P)(UD)^\epsilon
        \ll \max\{U,U^{1/6}D^{5(1+c)/6}\}(UD)^\epsilon .
\]
For \(D=U^r\) on a bounded \(r\)-range choose \(c\) once in
terms of the requested loss.  The exponent is
\(\max\{1,(1+5r)/6\}+\epsilon\), including \(r=1\).
Sobolev in the smooth real and imaginary test parameters allows
rowwise witnesses at a fixed polynomial height cost.
It does not allow row-dependent prime coefficients.
This proves \cite[Lemma 17.6]{QRH}, without taking a limit of
strict marked-moment constants.

\section{The extended plain fourth moment: full coefficient induction}\label{app:plain}\label{sec:plain}

The refined row count requires a fourth-moment estimate for two plain
character polynomials and a product of short prime polynomials.  We prove
that estimate here.  The two plain polynomials have no length restriction
when no prime polynomial is present; otherwise the permitted lengths lie
in an affine region.  Two finite Fourier transforms return products of the
same kind at a smaller effective width.  Some transformed rows, however,
induce characters in a fixed
finite family.  Their plain products can have volume-sized main terms.
For the longer inputs we therefore subtract a comparison product with the
same product of scales.  Preserving its common character, mask, and norm
power through the transforms makes those main terms cancel.  The estimate
itself concerns the original, uncentered product.

The two-transform proof below is adapted from \cite[Section 18]{QRH}.
The change is the endpoint \(13/18\), not an invocation of its lemma
outside the printed range \(3/4\le\kappa\le1\).  Every use of this
change is explicit: the centered comparison uses \(6\kappa-1\ge10/3\);
the remaining transformations are independent of \(\kappa\).
We use the conventions in Appendix~\ref{app:conventions}.  In particular, the
rows are elements \(k\in\mathcal O\) with \(0<q_k\ll Z^m\), and
\[
 \psi_k(n)=\tau(n)\chi_n(k),\qquad M=m+q.
\]
The character \(\tau\) is fixed within a row sum.  The union of the
prime supports of all its displayed moving residue-symbol factors,
before canceling factors or reducing exponents modulo six, has norm at
most \(Z^q\).  Every displayed factor retains its zero extension,
including a canceled or six-divisible factor.  A redundant zero may
instead be represented by an additional puncture mask only when an exact
factorization retains every surviving local and fixed-ray phase.  This
mask is the indicator of coprimality to one squarefree ideal of norm at
most \(Z^B\), for a bounded \(B\), fixed within the current row sum;
the same mask is used in both plain factors and in every prime factor.
It may depend on previously frozen labels, but not on the varying row.
In this appendix \(\mathcal S:=S\) is the fixed excluded prime set of
Appendix~\ref{app:conventions}; it is not enlarged by freezing a moving label.
All ideals in the polynomials avoid this same set.

Fix a finite group \(\Theta\) of finite-order ray characters whose
conductor primes belong to \(\mathcal S\).  It contains the fixed
twists and all characters used to separate the fixed reciprocity
phases.  Thus it also contains the supplementary character
\(n\mapsto\chi_n(-1)\): by
Equation~\eqref{eq:fixed-gauss-phases} and the fixed bicharacter
table, this character is the diagonal character
\(n\mapsto\mathcal R(n,n)\).  Membership of an inducing character in
\(\Theta\) means equality with the primitive inducing character of some
member of \(\Theta\), not equality of the chosen zero-extended
presentations.  All zeros of those presentations remain in the polynomials.
Let \(\mathcal R_0\) be the rows for which
\(\psi_k\) induces a nonprincipal character.  For \(z>0\), let
\(\mathcal R_z\) be the rows for which the inducing character of
\(\psi_k\) does not belong to \(\Theta\).  Write
\begin{equation}
 A=n_1+n_2+z.
 \label{old-eq:3.1}
\end{equation}

\begin{lemma}[Fourth moment with short prime factors]\label{lem:plain}
Let \(13/18\le\kappa\le1\).  Let \(n_1,n_2\ge0\), and let
\(Q=\prod_{i\in\mathcal I}Q_{\psi_k,i}\), where
\[
 Q_{\psi_k,i}=P_i^{-1/2}\sum_{p\ {\rm prime}}
       \psi_k(p)\nu_i(p)W_i^{\rm slot}(q_p/P_i),\qquad P_i=Z^{z_i},
\]
be a finite product of prime polynomials of lengths
\(z_i\ge0\) and total length \(z=\sum_{i\in\mathcal I}z_i\).
Their underlying prime supports are pairwise disjoint before common
masks and row zero extensions are imposed.  Their coefficients are
fixed finite linear combinations of finite-ray characters.  For every
positive-length slot the expansion has the form
\[
 \nu_i(n)=\sum_{j=1}^{J_i}c_{ij}\vartheta_{ij}(n),
 \qquad \vartheta_{ij}\in\Theta,
\]
where the finite list and its coefficients are fixed independently of
\(Z\) and of the row.  The row and mask hypotheses are those just stated.

For every \(\epsilon>0\), there is a slot mesh \(\eta>0\) such that,
if \(z_i\le\eta\) for every \(i\), then
\begin{equation}
 \sum_{k\in\mathcal R_z}
 \left|S_{\psi_k}(n_1;W_1)S_{\psi_k}(n_2;W_2)Q\right|^2
 \ll Z^{M+\epsilon}
 \label{old-eq:2.1}
\end{equation}
in either of the following cases:
\begin{enumerate}
\item \(z=0\).  There is no restriction on the bounded nonnegative
      lengths \(n_1,n_2\), and no lower bound on the conductor.
\item \(z>0\) and
\begin{equation}
 n_1+n_2+6\kappa z=A+(6\kappa-1)z\le M.
 \label{old-eq:3.9}
\end{equation}
      If \(\kappa<1\), assume in addition that
      \(\beta_*\le(1+\kappa)/2\).  If \(\kappa=1\), no zero-free
      hypothesis is required.
\end{enumerate}
For an empty slot list, \(Q=1\).  If \(z=0\) and the list contains
zero-length slots, their scales are bounded and absolute counting
absorbs them into the constant; the zero-slot assertion therefore
has the same strength.  All real length parameters range over prescribed bounded
sets.  The mesh depends only on those sets and \(\epsilon\), uniformly
for \(\kappa\in[13/18,1]\).  For each fixed \(Z\)-independent arithmetic
datum \(\mathcal A\) and fixed slot count, the bound uses finitely many
smooth seminorms and a fixed polynomial in the separated norm-twist
heights.  Their orders, and the bound itself, are uniform over all
moving moduli, admissible masks, and frozen outer labels in the stated
ranges.
\end{lemma}

The application to our bound sets \(\kappa=2B-1\), so
its zero-free hypothesis in this lemma holds by equality.  We retain this
dynamic value without assuming any earlier zero-free theorem,
since the refined row count uses the resulting affine
capacity.  The proof below uses the
finite Gauss identities of Section~\ref{sec:finite-correlations},
the smooth calculus of Appendix~\ref{app:smooth}, and Lemma~\ref{lem:logarithmic-control}.
Absorb any zero-length slots by absolute counting at their bounded
scales; henceforth \(z=0\) means that no live slot remains.  We first
make two reductions that preserve the row family.

The later induction is on the effective width \(M=m+q\), with all zero-slot
bands completed before the positive-slot bands.  Within each band the range
\(A\le5M/6\) is proved first.  This threshold comes from the
transformed rows whose inducing characters lie in \(\Theta\): counting
their sixth-power form and
bounding their plain products by volume leaves an excess
\(A-5M/6\), before the common-support savings.  Reflection handles
some longer inputs; the others are reduced to an equal-product-scale
difference, whose cancellation removes that excess.

\subsection{Fixed masks and reflection}

Call a character \emph{natural} when its zeros consist exactly of its
primitive conductor primes, its redundant row and declared moving
radical primes, and the fixed primes in \(\mathcal S\).  Here a
redundant prime is one at which the inducing character is unramified
but the displayed zero-extended product still vanishes.  A displayed
six-divisible power still vanishes on nonunits.  If a redundant factor
of the fixed twist is not included in the declared moving radical,
first factor it exactly into its remaining phase and a coprimality
indicator, and put that indicator in a squarefree ideal
\(\mathfrak R\), with \(q_{\mathfrak R}\le Z^B\), fixed throughout
the row sum.  In particular, a cancellation between the fixed twist
and the varying factor \(\chi_n(k)\) is not reclassified as a
separately chosen extra puncture for each row.

Let \(\psi_k^0\) be the natural character obtained by deleting these
additional punctures, and let \(T_k(X)\) denote its centrally
normalized plain sum at scale \(X\).  Write \(S_{k,\mathfrak R}(X)\)
for the sum with the extra mask.  Multiplicativity, including zero
extensions, gives the exact identities
\begin{align}
 S_{k,\mathfrak R}(X)
 &=\sum_{d\mid\mathfrak R}\mu(d)\psi_k^0(d)q_d^{-1/2}
                  T_k(X/q_d),
 \label{old-eq:2.1a}\\
 Q_{k,i,\mathfrak R}
 &=Q_{k,i}^{\,0}
   -P_i^{-1/2}\sum_{\substack{p\mid\mathfrak R\\p\ {\rm prime}}}
       \psi_k^0(p)\nu_i(p)W_i^{\rm slot}(q_p/P_i),
 \qquad P_i=Z^{z_i}.
 \label{eq:prime-mask-erasure}
\end{align}
For example, the first identity follows by inserting
\(\sum_{d\mid(n,\mathfrak R)}\mu(d)\) and writing \(n=dl\).
The quotient \(l\) is unrestricted at primes of \(d\); the original
natural zero extension remains on it.

Apply these identities simultaneously to both plain factors and all
slots.  If the frozen slots form \(J\subseteq\mathcal I\), and the
plain divisors are \(d_1,d_2\), the scalar depending on the row has
modulus at most one.  The remaining coefficient has absolute value at
most a fixed profile constant times
\[
 q_{d_1d_2}^{-1/2}\prod_{i\in J}P_i^{-1/2}.
\]
The resulting natural product has lengths
\[
 A'=A-\log_Zq_{d_1d_2}-\sum_{i\in J}z_i,\qquad
 z'=z-\sum_{i\in J}z_i.
\]
Its inducing character is unchanged.  If \(z'>0\), the left side of
Equation~\eqref{old-eq:3.9} has decreased; if \(z'=0\), the zero-slot
assertion has no length restriction and
\(\mathcal R_z\subseteq\mathcal R_0\).

For every fixed \(\epsilon_1>0\),
\begin{equation}
 \sum_{d\mid\mathfrak R}q_d^{-1/2}
 \le \prod_{p\mid\mathfrak R}(1-q_p^{-1/2})^{-1}
 \ll_{\epsilon_1,B} Z^{\epsilon_1}.
 \label{old-eq:2.1b}
\end{equation}
This follows from Lemma~\ref{lem:deleted-euler-factors}.  The corresponding
mass for a frozen slot is also \(Z^{\epsilon_1}\): on its annular support,
\(P_i^{-1/2}\ll q_p^{-1/2}\), and the same product bounds the sum over
\(p\mid\mathfrak R\).  The number of slots is fixed.  Minkowski's
inequality therefore reduces Equation~\eqref{old-eq:2.1} to the
natural assertion at the same width, with an arbitrarily small power
loss.  A nonempty annular scale below one is bounded below by a
positive profile-dependent constant and may be rescaled to one.
After this rescaling its new nonnegative length is
\(\max\{n_i-\log_Zq_{d_i},0\}\le n_i\), so the asserted
nonincrease of the affine expression remains valid.

We next record precisely the reflection used for natural characters.
Let \(\psi_k^*\) be the primitive character inducing \(\psi_k^0\),
and let \(\mathfrak R_{0,k}\) be its redundant natural radical.
The primes of \(\mathfrak R_{0,k}\) are disjoint from the primitive
conductor.  If \(Q_k\) is that conductor norm, set
\(C_k=3Q_k/(2\pi)^2\), the conductor scale in the functional
equation.  Tameness at primes outside \(\mathcal S\) gives
\begin{equation}
 C_kq_{\mathfrak R_{0,k}}
 \ll_{\mathcal S,\tau_{\rm fixed}} Z^M.
 \label{old-eq:2.1c}
\end{equation}
Indeed, each good prime contributes at most once, either to the
primitive conductor or to the redundant radical.  Their union is
contained in the union of the row radical and the declared moving
radical, whose norm is at most \(q_kZ^q\).  By
the fixed-numerator ray argument in Appendix~\ref{app:conventions}, the unit and
\(\mathcal S\)-supported parts of a row, with valuations reduced
modulo six when evaluated on primary elements prime to \(\mathcal S\),
range over a fixed finite ray
family.  Thus all primes in \(\mathcal S\) contribute only a fixed
factor.

Apply the primitive Hecke functional equation recorded in the proof of
Lemma~\ref{lem:hecke-strip-growth} to \(\psi_k^*\), using the entireness of
\(L(s,\psi_k^*)\) stated there.  Let \(\varepsilon_k\) be its root number,
so \(|\varepsilon_k|=1\).  Mellin inversion and a contour shift show that
the normalized plain sum of \(\psi_k^*\) at \(X\) equals the root number times the conjugate
character sum at \(C_k/X\), with transformed profile \(W^\sharp\)
defined by
\[
 \mathcal M W^\sharp(s)
 =\mathcal M W(1-s)\frac{\Gamma(s)}{\Gamma(1-s)}.
\]
The Mellin transform convention is the one in
the smooth calculus of Appendix~\ref{app:smooth}.  The quotient of gamma functions has
poles only at the nonpositive integers and zeros at the positive
integers.  Shifting the inverse Mellin contour to the left gives an
expansion in nonnegative integral powers at zero; shifting it to the
right gives arbitrary decay at infinity.  Thus \(W^\sharp\) and its
Euler derivatives are bounded at zero and rapidly decreasing at
infinity.  The height assertion is also quantitative.  For
\(W_\omega(y)=W(y)y^{i\omega}\), the exact identity
\(\mathcal M W_\omega(s)=\mathcal M W(s+i\omega)\) puts all
height dependence in a translate of the rapidly decreasing Mellin
transform.  On each fixed vertical line the gamma quotient and
each fixed number of Euler derivatives have polynomial growth in
the integration height.  Integrating the translated Mellin decay
therefore gives a fixed polynomial in \(1+|\omega|\), with its
degree depending only on the fixed contour and derivative orders.
The pointwise Mellin integration-by-parts estimate on a compact real
strip in the smooth calculus of Appendix~\ref{app:smooth} justifies the horizontal
joins in these shifts; an integrated Fourier tail is not used to
bound a fixed horizontal trace.  This proves the required
finite-seminorm and polynomial-height bounds for reflection.

Deleting the redundant Euler factors before reflection and restoring
them geometrically afterward gives
\[
 T_k(X;W)
 =\varepsilon_k
 \sum_{\substack{d_0\mid\mathfrak R_{0,k}\\
                  \operatorname{rad}(h_0)\mid\mathfrak R_{0,k}}}
 \frac{\mu(d_0)\psi_k^*(d_0)\overline{\psi_k^*(h_0)}}
      {\sqrt{q_{d_0}q_{h_0}}}\,
 T_{\overline{\psi_k^0}}
   \left(\frac{C_kq_{d_0}}{Xq_{h_0}};W^\sharp\right),
 \qquad |\varepsilon_k|=1.
\]
The geometric series is absolutely bounded by
\(\prod_{p\mid\mathfrak R_{0,k}}(1-q_p^{-1/2})^{-1}\).
Consequently its total coefficient mass, together with the \(d_0\)
sum, is \(Z^{\epsilon_1}\).  The scales are
\begin{equation}
 Y=\frac{C_kq_{d_0}}{Xq_{h_0}},\qquad
 d_0\mid\mathfrak R_{0,k},\quad
 \operatorname{rad}(h_0)\mid\mathfrak R_{0,k}.
 \label{old-eq:2.1d}
\end{equation}
By Equation~\eqref{old-eq:2.1c}, there is a fixed \(C_{\mathcal A}\ge1\)
such that
\[
 C_kq_{\mathfrak R_{0,k}}\le C_{\mathcal A}Z^M,\qquad
 Y\le C_{\mathcal A}Z^{M-\log_ZX}.
\]

The profile \(W^\sharp\) can be partitioned into smooth annuli.
Below its main scale, central normalization gives summable
coefficients \(O(2^{-j/2})\) on the annuli of relative scale \(2^{-j}\).
Above that scale, rapid decay gives an arbitrary summable power.
Fix a small length tolerance \(\xi>0\).  Truncate the upper annuli
after an enlargement \(Z^{\xi/2}\); sufficiently many fixed derivatives
make the omitted part negligible by absolute counting.  A main scale
below \(Z^{-\xi}\) is likewise negligible.  If a retained annular
profile has fixed upper support endpoint \(B\), a scale \(S\) that
can contain a nonzero integral ideal satisfies \(SB\ge1\).  On such
a scale,
\[
 \max(S,1)\le\max(1,B)S.
\]
Thus replacing a retained subunit scale by scale one costs this fixed
multiplicative factor.  All these assertions are uniform in the moving labels.

The row-dependent choices \(d_0,h_0\) are handled by their coefficient
mass followed by a rowwise supremum in the resulting scales.
the smooth calculus of Appendix~\ref{app:smooth} bounds a supremum over two
polynomial-range scales using \(O((\log Z)^2)\) unit boxes and
derivatives of the profiles.  The row character and its natural zeros
remain unchanged: after reflection, conjugating that whole factor
inside its absolute value replaces the conjugate character by the
original character and conjugates its profile.  This operation is
valid because the absolute value of a product is unchanged by
conjugating one factor.

Let \(C_{\rm ref}\ge1\) include \(C_{\mathcal A}\), the fixed annular
endpoint multipliers, the logarithmic box multipliers for a reflected
factor and any paired unreflected factor, and \(\max(1,B)\) for the
retained profile types.  This constant is fixed after the data and
profile types, independently of \(Z\) and the moving labels.  For
\(n=\log_ZX\), every retained scale satisfying the preceding support
test, after the permitted clipping and box enlargement, has declared
nonnegative length
\begin{equation}
 n_{\rm ref}\le M-n+\frac\xi2+\frac{\log C_{\rm ref}}{\log Z}
 \le M-n+\xi
 \qquad\left(\frac{\log C_{\rm ref}}{\log Z}\le\frac\xi2\right).
 \label{eq:centered-reflection-length}
\end{equation}
This is a linear bound also when \(M-n\) is slightly negative; if its
right side is negative, no such retained scale exists.  The smaller
upper-annulus range changes only the fixed decay, seminorm, and height
orders used for the discarded tail.

For \(z=0\), reflect each original factor whose length exceeds \(M/2\)
once, and leave the other factor unchanged.  Equation~\eqref{eq:centered-reflection-length}
bounds every retained reflected factor; the fixed box multiplier for
an unreflected factor is included in \(C_{\rm ref}\).  After the preceding
annular truncation and scale-box enlargement, the resulting lengths
satisfy
\begin{equation}
 n_i^*\le M/2+\xi,\qquad
 A^*=n_1^*+n_2^*\le M+2\xi\le M+\delta,
 \label{old-eq:2.1h}
\end{equation}
provided \(2\xi\le\delta\).  We call this the padded zero-slot core.
Here \(\delta\) is the positive padding parameter chosen below after the
width floor and step.
There is no repeated reflection at the boundary \(M/2\).  Reflection
and scale suprema are taken before centering; on a later centered term
whose primitive inducing character lies outside \(\Theta\), one first
applies the triangle inequality to its two rectangles.  No such supremum
is applied to a centered difference whose primitive inducing character
belongs to \(\Theta\).

\subsection{Prime estimates and the induction order}\label{sec:weighted}

For \(\kappa<1\), put \(s_\kappa=(1+\kappa)/2\).  Under the hypothesis
\(s_\kappa\ge\beta_*\), the global part of
Lemma~\ref{lem:logarithmic-control} gives, for every fixed \(e>0\),
\begin{equation}
 \frac{L'}{L}(s_\kappa+e+it,\psi)
 \ll_e \log\!\bigl(2Q_\psi(3+|t|)^2\bigr)
 \label{old-eq:3.4}
\end{equation}
for a primitive nonprincipal inducing character \(\psi\).
To estimate a prime annulus of scale \(P\), apply Mellin inversion to
a smooth von Mangoldt sum and move its contour to
\(\Re s=s_\kappa+e\).  There are no poles in the region of the shift,
and the Mellin transform has arbitrary decay, so
Equation~\eqref{old-eq:3.4} bounds the new integral by
\(P^{s_\kappa+e}\) times a logarithm of the conductor and a fixed
polynomial in the height.  Prime powers of exponent at least two
contribute \(O(P^{1/2+\epsilon_1})\).  For large \(P\), dividing the
annular weight by \(\log(Py)\) replaces the von Mangoldt weight by
the prime weight; its smooth seminorms are bounded on the annulus.
Bounded \(P\) are estimated by absolute counting.  Expand each
positive-length coefficient in the characters \(\vartheta_{ij}\)
from the statement.  If the primitive inducing row character
\(\psi\notin\Theta\) made \(\psi\vartheta_{ij}\) principal, then
\(\psi=\vartheta_{ij}^{-1}\in\Theta\), a contradiction.  Thus
every resulting slot character is nonprincipal on \(\mathcal R_z\).

After extra-mask erasure, fix a positive slot list and put
\(\boldsymbol W=(W_i^{\rm slot})_{i\in\mathcal I}\).  For normalized
twist heights \(\boldsymbol\omega=(\omega_i)_{i\in\mathcal I}\), define
\[
 Q_{\boldsymbol\omega}
 =\prod_{i\in\mathcal I}\left\{P_i^{-1/2}
   \sum_{p\ {\rm prime}}\psi_k^0(p)\nu_i(p)W_i^{\rm slot}(q_p/P_i)
                       (q_p/P_i)^{i\omega_i}\right\}.
\]
For every \(\epsilon_1>0\), central normalization and multiplication over this
fixed list, with \(e\) and all subsidiary losses sufficiently small,
give finite \(j,b,h\ge0\) and a constant \(C\) such that
\begin{equation}
 |Q_{\boldsymbol\omega}|^2
 \le C Z^{\kappa z+\epsilon_1}p_j(\boldsymbol W)^b
                 (1+|\boldsymbol\omega|)^h
 \qquad(k\in\mathcal R_z,\ z>0).
 \label{old-eq:3.5}
\end{equation}
Here \(j\) may be rounded up to an integer and \(2s_\kappa-1=\kappa\).
The orders and \(C\) may depend on the fixed data, slot count, loss,
and any fixed requested internal logarithmic or normalized-twist
derivatives, but not on \(Z\), moving labels, rows, or the numerical
heights.  The exponent \(\kappa z+\epsilon_1\) is independent of
these derivative and height orders.  Indeed, translate each pure twist
in its Mellin variable before integration by parts.  The weighted
integral of the untwisted Mellin transform then contributes a finite
seminorm and a fixed polynomial in \(\boldsymbol\omega\), while the
contour displacement contributes exactly \(2ez\) to the squared
exponent.  Logarithmic derivatives remain annular, and normalized-twist
derivatives insert only powers of the logarithmic profile variable.
The redundant natural radical contains
only \(O(\log Z)\) primes, so deleting it changes a normalized slot
by \(O(P_i^{-1/2}\log Z)\) times its fixed profile factor; this is
within the stated power loss because \(P_i\ge1\).  For \(\kappa=1\),
absolute prime counting gives Equation~\eqref{old-eq:3.5} without a
zero-free assumption and with \(h=0\) for undifferentiated pure twists,
whose modulus is one.  For \(\kappa<1\) the hypothesis remains
\(\beta_*\le(1+\kappa)/2\).  The choice of losses is uniform in
\(\kappa\in[13/18,1]\), since the length ranges are bounded.  In every
later use of Equation~\eqref{old-eq:3.5}, its fixed seminorm and height
factor is retained in the weighted Fourier estimates; it is not included
in the exponent of \(Z\).

We now specify the induction, including the estimates at its smallest
widths.  Choose a width floor \(\rho>0\), a width step \(\sigma>0\),
and then \(\delta>0\), all small in terms of the final
\(\epsilon\), with \(\delta\ll\min(\rho,\sigma)\).  Precise loss
choices will be made after the depth is bounded.  Divide the bounded
range of \(M\) into consecutive bands of length \(\sigma/4\).
Prove all zero-slot bands first, in increasing order of width, and
then all positive-slot bands.  Within a zero-slot band, first prove
the uncentered assertion for \(A\le5M/6\), and then the padded core
in Equation~\eqref{old-eq:2.1h}.  Reflection then supplies all
zero-slot lengths in that band.  Within a positive-slot band, first
prove the uncentered assertion for \(A\le5M/6\) subject to
Equation~\eqref{old-eq:3.9}, and then the remaining part of that
region by centering.  A completed earlier band therefore includes
the unrestricted zero-slot assertion.  Let \(M_{\max}\) bound the
initial width range, and define
\[
 D=2+\left\lceil 2M_{\max}/\sigma\right\rceil.
\]
The strict width decrease proved below will show that at most \(D-2\)
nonterminal calls occur on a branch.

At \(M\le\rho\), the padded zero-slot core follows from absolute
counting: there are \(O(Z^m)\) rows and the squared product is
\(O(Z^{A+\epsilon_1})\).  Its exponent above \(M\) is at most
\(A-q\le\rho+\delta\).  For positive slots, combine the completed
zero-slot estimate with Equation~\eqref{old-eq:3.5}.  In the region
of Equation~\eqref{old-eq:3.9}, \(A\ge z\), whence
\begin{equation}
 z\le\frac{M}{6\kappa}\le\frac{3M}{13},
 \qquad \kappa z\le\frac M6.
 \label{old-eq:3.6}
\end{equation}
Thus the extra terminal exponent for positive slots is at most
\(\rho/6\).  These terminal exponents can be made smaller than the
reserved final loss by choosing \(\rho,\delta\) sufficiently small.

For a width above the floor, the two-transform estimate below will
first prove the uncentered range \(A\le5M/6\).  We explain now why
the remaining inputs can be compared with that range.  Set
\begin{equation}
 L=M/4.
 \label{old-eq:2.2}
\end{equation}
For \(A>5M/6\), the comparison lengths \(L,A-z-L\) are both at
least \(L\).  In the positive-slot case,
\((6\kappa-1)z\le M-A<M/6\), and \(6\kappa-1\ge10/3\); hence
\(z<M/20\).  This also shows \(A-z>2L\).

If one original plain length is \(b<L\), reflect the other factor.
If both are at least \(L\), introduce the comparison with lengths
\(L,A-z-L\) and the same two profiles; reflect its longer factor
only when estimating the comparison separately.  In both cases the
total length after that reflection is at most
\[
 b+\{M-(A-z-b)\}+z+\xi
 \le 3M/2-A+2z+\xi,
\]
where \(b\le L\).  Equation~\eqref{eq:centered-reflection-length}
includes the clipped reflected scale and the paired unreflected box
multiplier in this same \(\xi\); the discarded tails have the stated
fixed-order bounds.  Thus
\begin{equation}
 A_{\rm comp}\le3M/2-A+2z+\xi.
 \label{old-eq:3.11}
\end{equation}
For \(z=0\), this gives
\(5M/6-A_{\rm comp}\ge M/6-\xi\).  For \(z>0\), use
\((6\kappa-1)z\le M-A\), \(A>5M/6\), and \(z<M/20\) to get
\begin{equation}
 \begin{aligned}
 A_{\rm comp}&\le\frac{23}{30}M+\xi,\\
 A_{\rm comp}+(6\kappa-1)z
 &\le\frac52M-2A+2z+\xi\le\frac{14}{15}M+\xi.
 \end{aligned}
 \label{eq:comparison-uncentered-margin}
\end{equation}
Each required boundary has a margin of at least \(M/15\) before
the \(\xi\) term.  Choose \(\xi\le\rho/30\).  After the threshold in
Equation~\eqref{eq:centered-reflection-length}, the unit-box multipliers
are already included in \(\xi\), so every retained reflected
comparison calls the previously proved uncentered assertion at this
same width.  This also completes the original case with a factor
shorter than \(L\).

In the remaining case put \(X_i=Z^{n_i}\),
\(Y_1=Z^L\), and \(Y_2=X_1X_2/Y_1\).  Then
\(X_1X_2=Y_1Y_2\), and all four plain lengths are at least \(L\).
Subtract the unreflected comparison from the original product,
leaving the common product \(Q\) of slots; call the result
\(\Delta_k\).  Multiplicativity and the equal product normalization
give the explicit formula
\[
 \begin{split}
 \Delta_k
 =\frac{Q}{\sqrt{X_1X_2}}\sum_{l_1,l_2}\psi_k(l_1l_2)
 \bigl\{
 &W_1(q_{l_1}/X_1)W_2(q_{l_2}/X_2)\\
 -{}&W_1(q_{l_1}/Y_1)W_2(q_{l_2}/Y_2)
 \bigr\}.
 \end{split}
\]
The comparison is bounded on the original permissible
rows.  The squared norm of the whole \(\Delta_k\) is nonnegative,
so it can then be enlarged to all rows in the smooth row ball.  In
the uncentered case, enlarge the squared norm of the original
product instead.  The Poisson transforms below are therefore never
applied to an indicator selecting exceptional or nonexceptional rows.

\subsection{The centered coefficient and its support}

We keep the subtraction as one coefficient while transforming its row
norm.  For fixed ideals
\(\mathbf b=(\mathfrak b_1,\mathfrak b_2)\), define
\begin{equation}
 D_{\mathbf b}(l_1,l_2)
 =
 \prod_{i=1}^2W_i(q_{\mathfrak b_i}q_{l_i}/X_i)
 -
 \prod_{i=1}^2W_i(q_{\mathfrak b_i}q_{l_i}/Y_i),
 \qquad X_1X_2=Y_1Y_2.
 \label{old-eq:2.18a}
\end{equation}
An allocation \(\mathbf b\) records prime powers already extracted
from the two plain variables.  If its conditions are impossible for
one rectangle, that rectangle's profile is zero on the corresponding
sum; the formal difference in Equation~\eqref{old-eq:2.18a} is
nevertheless retained.

The reason for retaining these common data is already visible in the
exceptional case.  For a fixed \(\vartheta\in\Theta\), a common mask
\(\mathfrak R_*\), and a common norm power \(q_l^{it}\), the lattice
estimate proved below has leading term
\[
 \sum_l\vartheta(l)1_{(l,\mathfrak R_*)=1}q_l^{it}W_i(q_l/T)
 =c_{\vartheta,\mathfrak R_*}T^{1+it}I_i(t)+\text{error},
\]
where \(c_{\vartheta,\mathfrak R_*}\) does not depend on \(T\) or \(t\).
Here \(I_i(t)\) is a fixed measure constant times
\(\int_0^\infty W_i(y)y^{it}\,dy\).
The product main term is therefore
\(c_{\vartheta,\mathfrak R_*}^2T_1^{1+it}T_2^{1+it}I_1(t)I_2(t)\),
which agrees for the two rectangles when \(T_1T_2\) agrees.
Lemma~\ref{lem:centered-lattice-cancellation} will quantify the remaining
error.  This cancellation is used only on the transformed rows inducing
characters in \(\Theta\); the common coefficient below is retained on
all rows until that later division into cases.

For a remaining slot set \(\mathcal I'\), write
\(\Pi=\prod_{i\in\mathcal I'}p_i\).  A coefficient with a common
character and mask means a coefficient of the form
\begin{equation}
 \begin{split}
 &\left(\prod_{i\in\mathcal I'}
       \nu_i(p_i)W_i^{\rm slot}(q_{p_i}/P_i)\right)
 D_{\mathbf b}(l_1,l_2)\,
 \tau_1(u)q_u^{it}
 1_{(u,\mathfrak R)=1}1_{s\mid u},\\
 &\hspace{40mm}u=\Pi l_1l_2.
 \end{split}
 \label{eq:centered-coefficient-form}
\end{equation}
Here \(\tau_1\) is one zero-extended product of a fixed finite-ray
character and moving residue-symbol factors; \(\mathfrak R\) is a
fixed squarefree extra mask; \(s\) is a fixed squarefree ideal; and
\(t\in\mathbb R\).  The condition \(s\mid u\) is omitted when
\(s=1\).  The slots retain their original disjoint underlying
supports.  They may share primes with either plain variable.

These common data factor multiplicatively on the \emph{full} product.
For any ideals \(v_1,\ldots,v_r\), even with common primes,
\[
 \chi_{\prod v_i}(h)=\prod_i\chi_{v_i}(h),\quad
 q_{\prod v_i}^{it}=\prod_iq_{v_i}^{it},\quad
 1_{(\prod v_i,\mathfrak R)=1}
   =\prod_i1_{(v_i,\mathfrak R)=1}.
\]
The first identity includes all zeros: a present prime whose total
exponent is divisible by six gives the zero-extended principal
factor, not the constant one.  The same multiplicativity holds for
\(\tau_1\).  A fixed-ray character evaluated on a full product also
factors with the same character on every variable.

We fix a support convention that will also control the slot-mesh
quantifier.  Let \(N\) be the original fixed number of slots, and
choose fixed intervals \([a_i,b_i]\) containing the support of each
\(W_i^{\rm slot}\).  Put
\(h_i=\max\{|\log a_i|,|\log b_i|\}\).  Choose a fixed number
\(H_N\) at least \(\sum_i h_i\), enlarged to include the two plain
profile windows, \(\log C_{\rm ref}\), fixed arithmetic normalizations, the relative dyadic
boxes, and the finitely many support enlargements through the \(D\)
operation stages.  The number \(H_N\) may depend on all the fixed
data and on \(N\), but not on \(Z\) or any moving label.  For every
subset \(I\) of live or frozen slots,
\begin{equation}
 \left|\log\prod_{i\in I}\frac{q_{p_i}}{P_i}\right|
 \le\sum_{i\in I}h_i\le H_N,
 \qquad \theta_N:=\frac{H_N}{\log Z}.
 \label{eq:centered-aggregate-support}
\end{equation}
Every residual full product divided by its nominal product scale lies
in \(\exp([-H_N,H_N])\), after increasing \(H_N\) once for the
fixed list of operations.  Both rectangles use that same box.  Extracted
plain prime powers use their exact norms, and a frozen slot contributes
its ratio \(q_p/P_i\) just once.  If a nonempty plain scale below one
is clipped to one, its error is at most the logarithm of that plain
window's fixed endpoint divided by \(\log Z\), once for that plain.
Thus a whole subset of slots or a whole divisor extraction contributes
one aggregate \(O(\theta_N)\) boundary error, not one copy of a
preselected tolerance for each slot or prime.

\subsection{The first Poisson transform and its target bound}\label{sec:plain-first-transform}

We now estimate either the uncentered product or the centered
difference selected above.  Put \(H=Z^m\) and \(X=Z^A\).  Choose a
fixed nonnegative smooth radial function that majorizes the row ball.
The full index product in either rectangle has norm in a fixed
multiple of \(X\).  All comparisons of exponents in this subsection
are first made on fixed dyadic norm intervals.  Their bounded
relative widths change a logarithmic length by \(O(1/\log Z)\);
the final loss discussion includes these changes.

We use the following common localization for both Poisson formulas.
Fix a nonnegative smooth dyadic partition
\(\sum_\lambda\omega_\lambda(q/T_\lambda)=1\) for \(q>0\),
with each weight supported in \(C_d^{-1}\le q/T_\lambda\le1\)
for one fixed \(C_d\).  A
sector fixes dyadic boxes for the finite list of aggregate outer norms,
not a separate box for each slot.  Let \(T_{\rm sec}\) be the
supremum in that sector of the nominal frequency scale.  Its ratio to
the scale at any one set of outer labels is at most a fixed
\(C_{\rm sec}\).  Retain exactly the whole weights whose support
meets
\[
 0<q\le T_{\rm sec}Z^{\xi/2}.
\]
Their union is contained in \(q\le C_dT_{\rm sec}Z^{\xi/2}\);
every discarded weight is supported above \(T_{\rm sec}Z^{\xi/2}\).
The selection depends only on the sector, never on a live column.
After \(N\) and the support data are fixed, take \(Z\) large enough
that
\begin{equation}
 2\theta_N+\frac{\log(C_dC_{\rm sec})}{\log Z}<\frac\xi4.
 \label{eq:centered-localization-threshold}
\end{equation}
One fixed enlargement of \(C_{\rm sec}\) covers all the finitely many
sector endpoint factors.

Here is a direct tail bound that justifies this operation on the
genuine sums.  In both applications the kernel argument is at least
\(q/(e^{2H_N}T_{\rm sec})\).  On every discarded weight it is
therefore at least \(Z^{\xi/4}\), by
Equation~\eqref{eq:centered-localization-threshold}.
Lemma~\ref{lem:kernel-seminorms} gives arbitrary decay
\((1+\text{argument})^{-B}\).  For the first formula use
\(|G(a,h)|\le q_a^{1/2}\), and for the second use the finite
definition \(|F(u,v;j)|\le q_uq_v\).  Absolute ideal counting and
the divisor bounds for the fixed number of factors bound all raw
columns and prefactors in a sector by \(C_NZ^{B_0}\) times a fixed
polynomial in the retained heights, where \(B_0\) depends only on
the bounded total lengths and chosen subsidiary power shares.  A
lattice norm dyad of scale \(T\) has \(O(1+T)\) frequencies.
Summing the radial decay over discarded dyads consequently gives
\(C_{N,B}Z^{B_1-B\xi/4}\) times that height polynomial, for a
bounded \(B_1\) independent of \(B\) and \(N\), and a convergent
geometric sum when \(B>1\).
Choose \(B\) sufficiently large to obtain any prescribed power
saving, including the polynomial number of frozen outer labels.
This is an absolute bound for the original terms, before any
off-coprime extension or Fourier absolutization.  It uses neither
the later \(s\)-radical saving nor a factorized off-coprime identity.
On the retained weights the full smooth kernel is kept; no sharp
condition comparing a row norm with the two live column norms is
inserted.  A retained range below the first nonzero lattice norm is
empty, apart from a bounded boundary dyad covered by the clipping
convention below.

At zero frequency in the first row Poisson formula, a character mean
can be nonzero only if its exponent at every prime is zero modulo
six.  In particular, no prime occurs to total multiplicity one in
the product of the two full index products.  This product is
therefore a powerful ideal.  There are \(O_\epsilon(X^{1+\epsilon})\)
such products of norm \(O(X^2)\): each powerful ideal is a square
times a cube of a squarefree ideal, and summing over the latter gives
the usual \(O(Y^{1/2+\epsilon})\) bound up to norm \(Y\).
Allocations to the fixed number of factors are divisor-bounded.
The central factor is \(X^{-1}\), and the row mean has size at most
\(H\).  Thus the zero frequency is \(O(Z^{m+\epsilon_1})\).

For the nonzero frequencies write the two full products as \(Ca,Db\),
where \(C,D\) contain their complete common prime support and
\[
 (a,b)=1,\qquad (ab,CD)=1.
\]
At each common prime, fix its exact
valuation in each plain variable and whether it is supplied by a
slot.  Dividing out these valuations adds that prime to the mask of
\emph{every} remaining factor on the side.  A slot that supplied the
prime is frozen and removed.  This description is valid even when a
slot and one or both plain variables supplied that prime.  The
allocation is made once for the coefficient, so the same ideals
\(\mathfrak b_i\) occur in the two terms of \(D_{\mathbf b}\).
An impossible allocation is a zero term.  In particular, a scalar
forced to vanish by an old moving zero or a common mask is not
replaced by its absolute upper bound before these support conditions
have been imposed.

Let \(c,d\) be the logarithmic norms of \(C,D\), and let \(p\) be the logarithmic norm of
their common radical.  Let \(\mathfrak r\) be the product of common
primes whose net exponents are nonzero modulo six; its logarithmic
norm is \(R\).  The corresponding character
\(\xi_{\mathfrak r}\) is primitive modulo \(\mathfrak r\).
In the M\"obius expansion of the complementary common row mask,
write \(\mathfrak e\) for the selected divisor and \(E\) for its
logarithmic norm.  In particular
\[
 R\le p,\qquad E\le p-R.
\]

Let \(A_C(a)\) and \(A_D(b)\) be the allocated convolution
coefficients, with the common character \(\tau\) omitted.  They
include all profiles, live slot coefficients, and fixed masks; in
the centered case each contains the entire allocated difference.

Here is a sum definition of the allocated coefficients and their
summation sets.  At the entrance to a stage, a full column tuple is
\(\mathbf n=(l_1,l_2,(p_i)_{i\in I})\), with full product
\(n=l_1l_2\prod_{i\in I}p_i\).  Its coefficient, with the common
character and common norm power removed, is
\[
 D_{\mathbf 1}(l_1,l_2)\prod_{i\in I}
              \nu_i(p_i)W_i^{\rm slot}(q_{p_i}/P_i)
                 1_{(n,\mathfrak R)=1};
\]
for an uncentered input the second rectangle is omitted.
For a pair of full products, let
\(\mathfrak c=\operatorname{rad}((n_1,n_2))\).
An allocation on one side consists of the two exact ideals
\(\mathfrak b_j=\prod_{p\mid\mathfrak c}p^{v_p(l_j)}\)
and every slot prime that belongs to \(\mathfrak c\).
Let \(F\subset I\) be these frozen slots, and put
\[
 C=\mathfrak b_1\mathfrak b_2\prod_{i\in F}p_i,
 \qquad I_C=I\setminus F.
\]
The allocation is allowed only when \(\operatorname{rad}C=\mathfrak c\);
the analogous product on the other side is \(D\), with the same
radical and possibly different valuations and frozen slots.
Fix those frozen data.  The complete residual coefficient is
\begin{equation}\label{own:plain-allocation}
\begin{split}
 A_C(a)={}&
 \left(\prod_{i\in F}\nu_i(p_i)W_i^{\rm slot}(q_{p_i}/P_i)\right)
 1_{(C,\mathfrak R)=1}\\
 &{}\times
 \sum_{\substack{l_1,l_2,\ (p_i)\in\prod_{i\in I_C}\mathcal P_i\\
             l_1l_2\prod_{i\in I_C}p_i=a\\
             (a,\mathfrak c\mathfrak R)=1}}
 D_{(\mathfrak b_1,\mathfrak b_2)}(l_1,l_2)
       \prod_{i\in I_C}\nu_i(p_i)
                    W_i^{\rm slot}(q_{p_i}/P_i).
\end{split}
\end{equation}
The common factor \(\tau(C)\) stays outer and is not divided out
at a zero.  If it vanishes, this allocation has zero coefficient.
This definition includes every original window and full signed
rectangle, even when one rectangle is empty.  A supported original
pair determines exactly one such pair of allocations with
\((a,b)=1\); conversely such an allocated pair reconstructs it.
This proves the expansion before Poisson, without unspecified
``allocated coefficients.''  Previously extracted powers are
multiplied into the \(\mathfrak b_j\); hence the same definition
works at each later stage.  In the second transformation its
existing \(s\)-divisibility is retained literally until the
common radical is required to contain \(s\).

Set
\[
 \tau_C(a)=\tau(a)\chi_a(\mathfrak e\mathfrak r)
                       \xi_{\mathfrak r}(a),\qquad
 \tau_D(b)=\tau(b)\chi_b(\mathfrak e\mathfrak r)
                       \overline{\xi_{\mathfrak r}(b)}.
\]
Residue-class Poisson, with the self-dual lattice measure from
Appendix~\ref{app:conventions}, gives for this allocation the following
nonzero-frequency expression, up to a scalar of bounded modulus in
the frozen labels:
\begin{equation}
 \begin{split}
 \frac{H}{Xq_{\mathfrak e}\sqrt{q_{\mathfrak r}}}
 \sum_{h\ne0}G_{\xi_{\mathfrak r}}(\mathfrak r,h)
 \sum_{(a,b)=1}
 &\frac{A_C(a)\overline{A_D(b)}}{\sqrt{q_aq_b}}\,
 \tau_C(a)\overline{\tau_D(b)}
 \overline{\mathcal R(a,b)}\,
 G(a,h)\overline{G(b,-h)}
 \\
 &\quad\cdot
 \widehat\Phi_1\left(
    \frac{Hq_h}{q_{\mathfrak e}q_{\mathfrak r}q_aq_b}
                    \right).
 \end{split}
 \label{eq:first-poisson-bridge}
\end{equation}
Here \(G_{\xi_{\mathfrak r}}\) is the normalized primitive Gauss
sum and has modulus at most one.  To check the normalization, the
substitution \(k=\mathfrak e k'\) and Poisson modulo
\(\mathfrak r ab\) give \(H/(q_{\mathfrak e}q_{\mathfrak r}q_aq_b)\).
The three unnormalized Gauss sums restore
\(\sqrt{q_{\mathfrak r}q_aq_b}\).  The global squared central
normalization is \(X^{-1}\).  For the phase, CRT for the pairwise coprime moduli
\(a,b,\mathfrak r\) supplies
\[
 \chi_a(b\mathfrak r)\overline{\chi_b(a\mathfrak r)}
       \xi_{\mathfrak r}(ab).
\]
The substitution \(k=\mathfrak e k'\) supplies
\(\chi_a(\mathfrak e)\overline{\chi_b(\mathfrak e)}\), up to a
scalar in the frozen labels.  Reciprocity changes
\(\chi_a(b)\overline{\chi_b(a)}\) into
\(\overline{\mathcal R(a,b)}\); the remaining factors are precisely
\(\tau_C(a)\overline{\tau_D(b)}\).  This proves
Equation~\eqref{eq:first-poisson-bridge}.  Each new character acts on
its whole residual product.  Its new moving primes belong to the
extracted support and puncture every residual factor.  Their full displayed
union is counted even if the factors at \(\mathfrak r\) cancel on units.

The new full displayed moving support has logarithmic norm at most
\(\widetilde q=q+R+E\), including any canceled factors at
\(\mathfrak r\).  Its nominal frequency scale is exactly
\[
 T_1(C,D,\mathfrak e,\mathfrak r)
 =\frac{q_{\mathfrak e}q_{\mathfrak r}X^2}
        {Hq_Cq_D}=Z^{K_0},
 \qquad K_0=2A-c-d+R+E-m.
\]
The product support convention gives
\(q_a/(X/q_C),q_b/(X/q_D)\in\exp([-H_N,H_N])\), so the
argument in Equation~\eqref{eq:first-poisson-bridge} is at least
\(q_h/(e^{2H_N}T_{\rm sec})\) for the supremum of \(T_1\)
in the outer sector.  Apply the preceding whole-dyad localization
to this genuine coprime bridge.  If \(K=\log_ZT_\lambda\) is
the upper length of a retained dyad, then
\begin{equation}
 \widetilde q=q+R+E,\qquad
 K\le K_0+\delta_{{\rm fr},1},\qquad
 0\le\delta_{{\rm fr},1}:=\frac\xi2+
       \frac{\log(C_dC_{\rm sec})}{\log Z}<\xi.
 \label{old-eq:2.3}
\end{equation}
The zero frequency already estimated above was the zero term of the
original common smooth row ball; it is not restored on a truncated
row domain.  The later ledgers retain the possible inequality
\(K_0-K\ge-\delta_{{\rm fr},1}\).

For each retained dyad and genuine common-support allocation, define
\(\mathcal C_1(a,b;h)\) to be the full bridge summand also on noncoprime
residual pairs individually disjoint from \(CD\) and allowed by the old
masks.  Use the fixed bicharacter \(\mathcal R(a,b)\), the full
zero-extended \(G(a,h),G(b,-h)\), the displayed whole-product characters,
and the same dyad weight, kernel, inverse roots, and formal product norms.
CRT identifies this definition with the genuine summand only on
\((a,b)=1\).  On each finite column shell the exact identity is
\[
 \sum_{(a,b)=1}\mathcal C_1(a,b;h)
 =\sum_{a,b}\mathcal C_1(a,b;h)\sum_{s\mid a,b}\mu(s).
\]
The full squarefree M\"obius sum annihilates the artificial noncoprime
pairs.  Insert it before estimating independent factors, and put
\(s_0=\log_Zq_s\).
Separate the fixed-ray phases and all smooth factors in the normalized
row norm and the two whole-product norms.  In these variables the
kernel is \(\widehat\Phi_1(R_{\rm sc}x/(y_1y_2))\) on fixed logarithmic
boxes.  If its aggregate scale ratio \(R_{\rm sc}\) varies over the sector,
include that single normalized outer ratio as another coordinate.
The cutoffs are chosen on the common product annulus of
Equation~\eqref{eq:centered-aggregate-support}, before fixing live slot
labels.  Thus one Fourier coefficient measure is common to the rows and
all live labels.  It supplies a row phase and only one norm power on each
whole column, together with phases in frozen outer norms.  It introduces
no separate powers on the two plain variables or the two rectangles.
The full kernel and inverse roots are kept until this separation;
taking their absolute supremum inside \(D_{\mathbf b}\) would not
preserve the coefficient.  Apply Cauchy--Schwarz in \(h\) only after this separation,
and only then use \(|G_{\xi_{\mathfrak r}}|\le1\).  The resulting positive norm
on the \(C\) side is
\[
 \mathcal N_C=
 \sum_{h\ne0}\omega_K(q_h/Z^K)
 \left|Z^{-(A-c)/2}
   \sum_{s\mid a}A_C(a)\tau_C'(a)q_a^{it}G(a,h)\right|^2,
 \qquad \omega_K\ge0,
\]
and there is an analogous \(D\) norm.  Here \(\tau_C'\) includes
one separated fixed-ray character.  All fixed masks in \(A_C\)
remain present.  Complete extraction and multiplicativity therefore leave
the coefficient in Equation~\eqref{eq:centered-coefficient-form}, now
multiplied by \(G(a,h)\).  In particular the two rectangles retain the
same allocated plain powers, character, puncture, and norm power.

The factors outside these two norms have exponent
\(m-A-R/2-E\).  The absolute number of common-support labels and
M\"obius labels has exponent \(p+s_0+\epsilon_1\).  To justify
the common-support count uniformly, first count its radical, giving
\(O(Z^{p+\epsilon_1})\).  For a fixed radical \(\mathfrak c\) of
polynomial norm and any fixed \(T\), Rankin's bound gives
\[
 \#\{v:\operatorname{rad}(v)\mid\mathfrak c,\ q_v\le Z^T\}
 \le Z^{aT}\prod_{\substack{r\mid\mathfrak c\\r\ {\rm prime}}}
             (1-q_r^{-a})^{-1}
 \ll Z^{aT+\epsilon_1}
\]
for every fixed \(a>0\).  Choose \(a\) small and use the
polynomial-size Euler-product estimate.  This bounds the choices of
the powers in \(C,D\) by an arbitrarily small power.  Their
allocations and the choice of \(\mathfrak e\) have divisor-bounded
multiplicity.  Finally there are \(O(Z^{s_0+\epsilon_1})\)
possible \(s\).

The allowance \(\epsilon_G\) below denotes the error envelope for the
current induction depth, together with its local small-power shares.
These envelopes will be chosen compatibly when the finite induction is
completed, using the same slot mesh throughout.
We will prove the following sufficient bound for the squared \(C\)
norm:
\begin{equation}
 \mathcal N_C\ll
 Z^{A-c+\widetilde q+\mathcal B_c-s_0+\epsilon_G},
 \qquad
 \mathcal B_c=\max\{0,(3c-5d-R)/6\}.
 \label{old-eq:2.5}
\end{equation}
The analogous quantity is
\(\mathcal B_d=\max\{0,(3d-5c-R)/6\}\).
The complete exponent ledger for this implication is
\[
 \begin{aligned}
 &(m-A-R/2-E)+(p+s_0)\\
 &\quad+\tfrac12\{A-c+\widetilde q+\mathcal B_c-s_0
                  +A-d+\widetilde q+\mathcal B_d-s_0\}\\
 &=M+\tfrac12\{\mathcal B_c+\mathcal B_d-(c+d-2p-R)\}.
 \end{aligned}
\]
The last brace is nonpositive:
\begin{equation}
 \mathcal B_c+\mathcal B_d\le c+d-2p-R.
 \label{old-eq:2.6}
\end{equation}
Indeed, at a common prime of multiplicities \(i\ge j\ge1\), put
\(r=1_{6\nmid i-j}\).  Only the \(i\) side can have a positive
local numerator.  Its contribution is at most
\((3i-5j-r)_+/6\), whereas the right side contributes
\(i+j-2-r\).  The latter is nonnegative.  If the former is
positive, six times their difference is
\(3i+11j-12-5r\), which is nonnegative: for \(r=0\) it is at
least \(2\), and for \(r=1\) one has \(i\ge j+1\), giving at
least \(0\).  The positive part of a sum is at most the sum of
positive parts.  Multiplication by each prime's logarithmic norm
and summation proves Equation~\eqref{old-eq:2.6}.  Thus it remains
to establish Equation~\eqref{old-eq:2.5}.

\subsection{Enlarging the Gauss-row norm}\label{sec:plain-gauss-enlargement}

This subsection describes three possible positive norms to which
the second transform will be applied.  Let \(w\) denote a length
removed from the \(C\) column, and let \(w_o\le w\) be the
additional moving-radical length created by that removal.  Put
\(a_0=A-c-w\).  When a squared extraction coefficient of size
\(Z^{-w_o}\) has been removed, Equation~\eqref{old-eq:2.5}
allows the unweighted remaining norm the exponent
\begin{equation}
 \Lambda_c=
 a_0+\widetilde q+w_o+w+\mathcal B_c-s_0+\epsilon_G.
 \label{old-eq:2.7}
\end{equation}
This is just \(A-c+\widetilde q+\mathcal B_c-s_0+\epsilon_G+w_o\).

For the initial norm, with \(w=w_o=0\) and \(a_0=A-c\), the first
transform gives the following conditional summand after fixing \(\Pi\)
and omitting its outer scalar:
\begin{equation}
 Z^{-a_0/2}\sum_{l_1,l_2}
 D_{\mathbf b}(l_1,l_2)\tau_1(l_1l_2)q_{l_1l_2}^{it}
 1_{(\Pi l_1l_2,\mathfrak R)=1}1_{s\mid\Pi l_1l_2}
 G(\Pi l_1l_2,h).
 \label{old-eq:2.18c}
\end{equation}
The slot sums are restored before this polynomial is squared.
The same statement with the second rectangle omitted applies to
uncentered products.  The local calculations below will show that the
extracted terms retain this form at their shortened length \(a_0=A-c-w\).

Define
\begin{equation}
 J=d-c+(K_0-K)-2w+w_o,\qquad J_+=\max(J,0).
 \label{old-eq:2.8}
\end{equation}

In the zero-slot proof take \(w=w_o=0\), and set
\[
 \ell=0,\qquad g=J_++\sigma.
\]
The positive norm \(\mathcal N_C\) is at most the corresponding
norm over a smooth ball of length \(K+g\).  There is no
multiplication of rows in this case.

For positive slots, start with \(w=w_o=0\).  Set
\(\ell_*=\sigma/3\), choose the slot mesh \(\eta<\sigma/6\),
and take the fixed pool
\[
 \mathcal P=\{p: Z^{\ell_*}/2<q_p\le Z^{\ell_*},\
                    p\notin\mathcal S\}.
\]
The prime ideal theorem in the fixed field, equivalently its fixed
ray-class form \cite[Theorem 1.1]{TZ}, gives
\(|\mathcal P|=Z^{\ell_*+o(1)}\).  Because the live slot lengths
are at most \(\eta\) and their relative annular supports are fixed,
the single inequality
\(Z^{\ell_*-\eta}>2\max_i b_i\) makes this pool disjoint from
every live slot window.  Its exponent gap is at least \(\sigma/6\),
independently of \(N\); only the threshold depends on the fixed
windows.  For a row \(h\ne0\), omit pool primes dividing
\(h\), \(s\), or any frozen support.  Each such integer or ideal
has polynomial norm, so \(O(\log Z)\) primes are omitted.
The remaining set \(\mathcal P_h\) has size comparable to
\(|\mathcal P|\), uniformly in the row and frozen labels.

To compare the rows \(h\) and \(hp^6\), write a full modulus as \(p^iu\) with
\((p,u)=1\).  CRT and reciprocity give the exact formula
\[
 G(p^iu,h)
 =\mathcal R(p,u)^i\chi_p(u)^{2i}G(p^i,h)G(u,h).
\]
Also \(G(u,hp^6)=G(u,h)\), since \(p\) is a unit modulo \(u\)
and \(\chi_u(p^6)=1\).
For \(p\nmid h\), Equation~\eqref{eq:gauss-local} shows that
replacing \(h\) by \(hp^6\) changes only \(i=1,6,7\).
The factors \(G(p^i,h)\) that occur are a scalar in the row and
\(p\); the displayed remaining factor is one fixed-ray character
times one residue character on the whole \(u\).  Since \(p\) is
outside every live slot window, \(p^i\) is allocated only to the two
plain variables.  The residual product is punctured at \(p\), and
the same allocation updates \(D_{\mathbf b}\) in both rectangles.
Because \(p\nmid s\), the condition \(s\mid u\) remains.

Here is the precise averaging argument.  Write the normalized
polynomial in \(\mathcal N_C\) as \(\mathcal H(h)\).  For each
eligible \(p\), local Gauss evaluation gives
\(\mathcal H(h)=\mathcal H(hp^6)\) plus the extracted terms with
\(p\)-adic column valuations \(1,6,7\).  There are only boundedly
many allocations of each of these valuations to the two plain
variables.  Jensen's inequality, first in \(p\) and then for this
fixed finite sum, gives
\[
 |\mathcal H(h)|^2\ll
 \frac1{|\mathcal P_h|}\sum_{p\in\mathcal P_h}
 \left\{|\mathcal H(hp^6)|^2+
        \sum_{i=1,6,7}\sum_{\rm allocations}
        |c_{p,i}(h)|^2|\mathcal H_{p,i}(h)|^2\right\}.
\]
The \(\mathcal H_{p,i}\) are normalized at their shortened column
scales.  The local identity and allocation just described are applied to the
whole centered coefficient, so each error retains
Equation~\eqref{eq:centered-coefficient-form}.

On the support of the original row weight, the new row \(hp^6\)
has norm \(O(Z^{K+6\ell_*})\).  A fixed output row has only
boundedly many representations as \(hp^6\) with \(p\in\mathcal P_h\):
every such \(p\) divides that row, and all such primes have norm at
least \(Z^{\ell_*}/2\), while the output norm has bounded
logarithmic length.  Summing the first term over \(h\) therefore
gives the factor \(|\mathcal P|^{-1}=Z^{-\ell_*+o(1)}\)
times a positive norm on the new rows.  For this main term set
\[
 \ell=\ell_*,\qquad g=J_++2\sigma.
\]
Since \(6\ell_*=2\sigma\), a smooth ball of length \(K+g\)
contains all the new rows.

For the extracted terms put \(P=q_p\).  When \(p\nmid h\),
Equation~\eqref{eq:gauss-local} gives the following central
coefficients.  At valuation one, the old Gauss sum has modulus one
and the new one is zero, so extraction gives squared coefficient
\(P^{-1}\).  At valuation six, the new Gauss sum is
\(P^3(1-P^{-1})\); its central factor is \(P^{-3}\), giving
\((1-P^{-1})^2\).  At valuation seven, the new Gauss sum has
modulus \(P^3\), and the central factor \(P^{-7/2}\) gives
\(P^{-1}\).  The row phase is \(\overline{\chi_p(h)}\) for
valuations one and seven, and is one for valuation six.  It
multiplies both rectangles.

Let \(\ell_p=\log_ZP\).  For sufficiently large \(Z\),
\(\ell_*/2\le\ell_p\le\ell_*\).  The removal and new
moving-radical lengths for the three terms are
\begin{equation}
 (w,w_o)=(i\ell_p,e_i\ell_p),\qquad
 i=1,6,7,\quad e_1=e_7=1,\quad e_6=0.
 \label{old-eq:2.9}
\end{equation}
The two squared factors \(P^{-1}\) give exactly \(Z^{-w_o}\).
For the valuation-six term, the exact factor
\(\chi_p(u)^{12}=1_{(u,p)=1}\) is represented by the already
common fixed puncture at \(p\), with its fixed-ray phase \(\mathcal R(p,u)^6\) retained.
At valuations one and seven the displayed local factor remains, and its
prime is counted in the moving radical.  In each
error, freeze \(p\), retain its common
column puncture, and discard its row eligibility restriction only
after taking the positive norm.  The average remains
\(|\mathcal P|^{-1}\sum_{p\in\mathcal P}\), so a bound uniform in
the frozen \(p\) introduces no prime-count factor.  For each error
set
\[
 \ell=0,\qquad g=J_++\sigma
\]
and enlarge its original rows directly to a smooth ball of length
\(K+g\).  Errors are not amplified again.  All three types of
norm satisfy
\begin{equation}
 0\le w_o\le w\le7\sigma/3.
 \label{old-eq:2.10}
\end{equation}

The row zero is added only at the final smooth ball.  Put
\(Y_{\rm col}=e^{H_N}Z^{a_0}=Z^{a_0+\theta_N}\).  A nonempty column
shell satisfies \(Y_{\rm col}\ge1\), and hence \(a_0\ge-\theta_N\).
By Equation~\eqref{eq:gauss-local}, \(G(u,0)=0\) unless \(u\) is a
sixth power, and \(|G(u,0)|\le q_u^{1/2}\le Y_{\rm col}^{1/2}\).
There are \(O(Y_{\rm col}^{1/6})\) sixth powers on this support.
If \(s\mid u\), then \(q_s^6\le q_u\le Y_{\rm col}\), so
\(s_0\le(a_0+\theta_N)/6\).  Including the assigned divisor-bounded
convolution loss, the squared contribution after central normalization is
\[
 \ll Z^{-a_0}\bigl(Y_{\rm col}^{1/6}Y_{\rm col}^{1/2}\bigr)^2
          Z^{\epsilon_1}
 =Z^{a_0/3+4\theta_N/3+\epsilon_1}
 \le Z^{a_0-s_0+2\theta_N+\epsilon_1}.
\]
The last inequality follows from \(a_0\ge-\theta_N\) and the displayed
bound for \(s_0\).  Thus Equation~\eqref{old-eq:2.7} bounds this added
row with the numerical \(2\theta_N\) correction included in the
\(C_*\xi\) stage allowance.  Here \(C_*\ge1\) denotes the common
constant for aggregate support errors in one stage, chosen independently
of the slot count \(N\).  The estimates below establish that one such
choice covers all operations in a stage.  Zero is never multiplied by a
pool prime.

\subsection{The second transform and smaller-width products}\label{sec:plain-second-transform}

Restore all live slot sums before expanding the square.  Let
\(B(u)\) denote the full allocated convolution coefficient in
one of the preceding Gauss polynomials, including its fixed mask
and \(s\)-divisibility condition.  It contains the whole difference
when centering is used.  Its moduli satisfy
\(q_u/Z^{a_0}\in\exp([-H_N,H_N])\).  Let \(\Phi_2\) be the nonnegative
smooth radial function defining the final row ball.  Apart from
the \(Z^{o(1)}\) loss in the prime density, the norm to be bounded is
\[
 Z^{-\ell}\sum_h\Phi_2(q_h/Z^{K+g})
 \left|Z^{-a_0/2}\sum_u B(u)\tau_1(u)q_u^{it}G(u,h)\right|^2.
\]
The Fourier identity in Lemma~\ref{lem:full-correlation} gives its
exact expansion
\[
 \begin{split}
 Z^{K+g-\ell-a_0}\sum_{u,v}
 \frac{B(u)\overline{B(v)}\tau_1(u)\overline{\tau_1(v)}
       q_u^{it}q_v^{-it}}{\sqrt{q_uq_v}}
 \sum_j F(u,v;j)
 \widehat\Phi_2\left(\frac{Z^{K+g}q_j}{q_uq_v}\right).
 \end{split}
\]
The kernel that occurs here is
\begin{equation}
 \widehat\Phi_2\left(\frac{Z^{K+g}q_j}{q_uq_v}\right).
 \label{old-eq:2.18b}
\end{equation}
Both the displayed inverse square roots and this kernel are retained
as functions of full \(u,v\) until their whole-product separation below.  In particular the
formula includes every live prime and every shared-prime multiplicity.

At \(j=0\), Lemma~\ref{lem:full-correlation} leaves only \(u=v\),
with \(F(u,u;0)=\varphi(u)\le q_u\).  The number of supported
moduli divisible by \(s\), on a nonempty shell, is at most
\[
 C\frac{Y_{\rm col}}{q_s}=C Z^{a_0-s_0+\theta_N},
\]
because writing \(u=sv\) gives \(Y_{\rm col}/q_s\ge1\) on that shell.
The divisor-bounded coefficients use their separate \(\epsilon_1\)
share.  Thus the diagonal exponent is
\(K+g-\ell-s_0+\theta_N\), including when \(a_0-s_0\) is slightly
negative.  Subtracting the allowance in Equation~\eqref{old-eq:2.7}
from its nominal part, without \(\epsilon_G\), gives
\begin{equation}
 \begin{split}
 &(K+g-\ell-s_0)
 -(a_0+\widetilde q+w_o+w+\mathcal B_c-s_0)\\
 &\quad=(A-M)-d-(K_0-K)-w_o-\mathcal B_c+g-\ell\\
 &\quad\le A-M+5\sigma/3+\delta_{{\rm fr},1}.
 \end{split}
 \label{old-eq:2.12}
\end{equation}
For the inequality, put \(D_0=d+K_0-K\).  The contribution
\(-D_0-w_o+J_+\) is at most \(\delta_{{\rm fr},1}\): it is at most
\(-c-2w\) when \(J\ge0\), and when \(J<0\) use
\(D_0\ge-\delta_{{\rm fr},1}\).
The largest remaining increment is
\(2\sigma-\ell_*=5\sigma/3\) in the amplified main norm;
the other norms have increment \(\sigma\).  Thus the diagonal
requires only the displayed terminal loss \(5\sigma/3\), plus
the already reserved frequency perturbation \(\delta_{{\rm fr},1}\)
and numerical support correction \(\theta_N\), when \(A\le M\),
and an additional \(\delta\) in the padded zero-slot core.  The
\(\theta_N\) correction is included in the \(C_*\xi\) stage allowance.

For \(j\ne0\), first perform the whole-dyad localization on this
genuine full \((u,v,j)\) sum, before extracting any common support.
The nominal scale is
\(T_2=Z^{2a_0-K-g}\), and the kernel argument is at least
\(q_j/(e^{2H_N}T_{\rm sec})\) for its supremum in the current
Gauss sector.  The tail estimate above applies using
\(|F(u,v;j)|\le q_uq_v\), with the literal conditions
\(s\mid u,v\) still present.  Write
\[
 \Omega_{\rm ret}(q_j)
 :=\sum_{\lambda\ {\rm retained}}\omega_\lambda(q_j/T_\lambda).
\]
This common row weight satisfies \(0\le\Omega_{\rm ret}\le1\).
The full kernel in Equation~\eqref{old-eq:2.18b} remains on every
retained term.  All subsequent signed Fourier separations are performed
one retained \(\lambda\) at a time, so the normalized row variable
stays on a fixed log box; \(\Omega_{\rm ret}\) records their sum and
common domain.  There are \(O(\log Z)\) relevant retained dyads in
the bounded polynomial ranges.  The preceding \(j=0\) term is exactly the diagonal
of the same final \(\Phi_2\) ball, including the previously added
Gauss row zero; neither zero term is restored on a different domain.

Now write \(u=D_2a,\ v=E_2b\), extracting the genuine complete
common support, so that \((a,b)=1\) and \((ab,D_2E_2)=1\).
Allocate all extracted powers to the plain variables and slots as at the
first transform.  Lemma~\ref{lem:complete-support-correlation} applies on
this genuine locus and gives
\[
 F(D_2a,E_2b;j)=F(D_2,E_2;j)\mathcal R(a,E_2)
   \overline{\mathcal R(b,D_2)}\mathcal R(a,b)
   \chi_a(j)\overline{\chi_b(-j)}.
\]
It permits arbitrary prime powers in all four moduli and leaves a row
scalar at the common primes.  Every extracted prime punctures all
remaining factors on its side, and every slot supplying that prime is
frozen.

This also explains explicitly the cases where a slot shares a
prime with a plain variable.  If \(p\) occurs only on one side, with
full multiplicity \(i\), its residual factor is
\(\chi_p(j)^i\).  Splitting \(i\) among the two plain variables and
the possible slot gives precisely the same factor by multiplicativity,
including when \(6\mid i\) and \(p\mid j\).  For example, a plain
\(p^5\) and a slot \(p\) give the zero-extended factor
\(\chi_p(j)^6\).  If \(p\) occurs on both sides, its full powers are
removed.  In the unequal case \(i>j_0\ge1\), the local correlation
is zero unless \(6\mid j_0\) and the frequency is \(p^{j_0}k\)
with \(p\nmid k\); when it is nonzero it equals
\[
 P^{j_0-1}(P-1)\chi_p(k)^{i-j_0},\qquad P=q_p.
\]
This is a row scalar, even if the excess valuation on the first side
came partly from a live slot.  For instance \(i=7,j_0=6\) with a
plain \(p^6\) and slot \(p\) leaves the scalar
\(P^5(P-1)\chi_p(k)\), freezes that slot, and punctures both
remaining plain variables at \(p\).  Equal multiplicities have the
scalar factors in Equation~\eqref{eq:correlation-local} and the
same puncture conclusion.  Complete rather than gcd-only extraction
is what makes these local factors independent of the residual
variables.

Define
\[
 c_2=\log_Zq_{D_2},\quad d_2=\log_Zq_{E_2},\quad
 b_2=(c_2+d_2)/2.
\]
Let \(p_2\) be the logarithmic norm of their common radical,
let \(G_c=(D_2,E_2)\), and put \(g_2=\log_Zq_{G_c}\).
The correlation vanishes unless \(G_c\mid j\).  At a common
prime with equal multiplicity \(i\not\equiv0\pmod6\), call the
divided frequency \(j/G_c\) a unit or nonunit according as that
prime does not or does divide it.  Let \(t_2\) be the total
radical length of these unit primes.  Let \(V_{\rm id}\) be
the product of these nonunit primes and put
\(V=\log_Zq_{V_{\rm id}}\).

Equation~\eqref{eq:correlation-local} gives the following absolute
local bounds:
\[
 \begin{array}{c|c}
 \text{common multiplicities and divided frequency}&
       \text{absolute correlation bound}\\ \hline
 i=i,\ 6\nmid i,\ \text{unit}&P^{i-1}\\
 i=i,\ 6\nmid i,\ \text{nonunit}&P^i\\
 i=i,\ 6\mid i,\ \text{either}&P^i\\
 i>j_0,\ 6\mid j_0,\ \text{unit}&P^{j_0}.
 \end{array}
\]
Every unequal case not in the last line is zero.  Fix the indicated
unit/nonunit partition and write \(j=G_cV_{\rm id}h'\).  The
partitioned scalar
\[
 1_{\rm part}(h')
 \frac{F(D_2,E_2;G_cV_{\rm id}h')}{Z^{g_2-t_2}}
\]
is defined to be zero off that partition and has modulus at most one
everywhere with this definition.  We do not assert the unit bound for
the unpartitioned correlation on other rows.  Every frozen allocation
forced to be zero by the old masks is still discarded before this
absolute bound is used.

The complete-support identity leaves the row factor
\(\chi_n(G_cV_{\rm id}h')\) on each whole residual column.
To count its fixed moving support, put \(v=G_cV_{\rm id}\) and
\(e_p=v_p(v)\bmod 6\).  Products of the canonical primary
generators of these good ideals are primary, so the exact all-input
identity is
\[
 \chi_n(v)=
   \prod_{\substack{p\mid v\\e_p\ne0}}\chi_n(p)^{e_p}
   \prod_{\substack{p\mid v\\e_p=0}}1_{(n,p)=1},
\]
with any unit of the original row retained in \(h'\).  All primes in this identity
already puncture every residual factor, because they are in the
genuine extracted support.  The active \(e_p\ne0\) primes are
contained in the unit set counted by \(t_2\) and the nonunit set
counted by \(V\).  Equal six-divisible and unequal-minimum-six-divisible
primes have \(e_p=0\) and are represented solely by the common
puncture in this fixed factor.  A nonunit prime can also have
\(e_p=0\), in which case counting it in \(V\) only enlarges the
bound.  This factors only \(\chi_n(v)\): the natural row factor
\(\chi_n(h')\), with all its zeros, and all fixed-ray reciprocity
phases remain.  It does not reclassify a cancellation with the
varying row as an extra puncture.  Thus the full new moving support
has length at most \(\widetilde q+w_o+t_2+V\).

Define the nominal row and total widths by
\begin{equation}
 \begin{aligned}
 m'&=2a_0-K-g-g_2-V,&
 q'&=\widetilde q+w_o+t_2+V,\\
 M'&=m'+q'
    =M+J-g-g_2+t_2
    \le M-\sigma.
 \end{aligned}
 \label{old-eq:2.13}
\end{equation}
Here \(g_2\ge t_2\), and the definitions of \(g\) give
\(g-J\ge\sigma\).  The identities follow by substituting
\(a_0=A-c-w\) and the definition of \(K_0\) preceding
Equation~\eqref{old-eq:2.3}.  Put
\[
 0\le\delta_{{\rm fr},2}:=\frac\xi2+
       \frac{\log(C_dC_{\rm sec})}{\log Z}<\xi.
\]
The sector constants here are those for the second transform.
The retained row weight pulls back exactly to
\(\Omega_{\rm ret}(q_{G_cV_{\rm id}}q_{h'})\).  On its support,
\[
 q_{h'}\le
 \frac{C_dT_{\rm sec}Z^{\xi/2}}{q_{G_cV_{\rm id}}}
 \le\frac{C_dC_{\rm sec}Z^{2a_0-K-g+\xi/2}}
          {q_{G_cV_{\rm id}}}
 =Z^{m'+\delta_{{\rm fr},2}}.
\]
Thus it lies in that common enclosing row ball.  On a
nonempty retained range define the declared nonnegative row length
\(m'_{\rm act}=\max\{0,m'+\delta_{{\rm fr},2}\}\) and
\(M'_{\rm act}=m'_{\rm act}+q'\).  Nonemptiness implies
\(m'+\delta_{{\rm fr},2}\ge0\), so
\(0\le M'_{\rm act}-M'\le\delta_{{\rm fr},2}\).
Regard the pulled-back weight as zero on the rest of this enclosing
ball throughout the signed calculation, and do the same for the
partitioned scalar.  This is an exact extension by zero.  They will
be replaced by their absolute bounds only after a nonnegative child
norm or an exceptional absolute product has been formed.
Complete extraction only shortens a plain variable or freezes an
entire slot, so the surviving slot length satisfies \(z'\le z\).

We record every other exponent in this second transformation.  The
condition \(s\mid u,v\) has not been dropped: it implies that every
prime of \(s\) is in the second complete common radical.  Thus after the complete extraction there is no remaining
\(s\)-divisibility condition on the residual columns.  The number of
possible radicals of length \(p_2\), for this fixed
\(s\), is \(O(Z^{p_2-s_0+\epsilon_1})\): write that radical as
\(s\mathfrak r_2\) with \((s,\mathfrak r_2)=1\), and count the
squarefree \(\mathfrak r_2\) of the remaining norm.  If \(p_2<s_0\)
outside the fixed shell boundary, there are none.  The same Rankin
argument used for \(C,D\) bounds their power and allocation
multiplicities.  Choosing the unit/nonunit partition costs at most
\(2^{\omega(\operatorname{rad}(D_2E_2))}\), another divisor-bounded
factor included in \(\epsilon_1\).  The exponents are
\[
 \begin{array}{l|r}
 \text{factor}&\text{exponent}\\ \hline
 \text{squared central normalization}&-a_0\\
 \text{row Poisson factor and normalized inverse roots}&K+g-a_0\\
 \text{prime density, when present}&-\ell\\
 \text{common-support count with fixed }s&p_2-s_0\\
 \text{common correlation}&g_2-t_2\\
 \text{conversion to normalized residual products}&a_0-b_2
 \end{array}
\]
The last line is
\(\{(a_0-c_2)+(a_0-d_2)\}/2\).
Their sum is
\[
 K+g-\ell-a_0+p_2-s_0+g_2-t_2-b_2.
\]
Subtracting this sum from the allowance
Equation~\eqref{old-eq:2.7} leaves the exponent permitted for
the inner plain products.  This is the bound to be supplied either by
smaller-width moments or by the exceptional-row estimate:
\begin{equation}
 M'+\Delta_{\rm child},\qquad
 \Delta_{\rm child}=b_2-p_2+w+\mathcal B_c+\ell
                   \ge w+\ell.
 \label{old-eq:2.14}
\end{equation}
The inequality uses \(b_2\ge p_2\), since both multiplicities
at each common prime are at least one.

We now remove the residual coprimality before forming independent
children.  For the fixed genuine \(D_2,E_2\), define
\[
 \widetilde F_{D_2,E_2}(a,b;j)
 :=F(D_2,E_2;j)\mathcal R(a,E_2)
   \overline{\mathcal R(b,D_2)}\mathcal R(a,b)
   \chi_a(j)\overline{\chi_b(-j)}
\]
on all residual pairs individually coprime to \(D_2E_2\) and allowed by
the old masks.  It equals the genuine correlation only when \((a,b)=1\).
Define the remaining
\(\mathcal C_2(a,b;j)\) on all individually allowed residual pairs
by the allocated convolution formulas and the old common masks,
including \((ab,D_2E_2)=1\).  Keep the formal full-product norms
\(q_{D_2}q_a,q_{E_2}q_b\), both inverse roots, the full smooth
kernel, and the same pulled-back row weight extended by zero on its
enclosing ball.  For each retained row the exact identity is
\[
 \begin{split}
 \sum_{(a,b)=1}\mathcal C_2(a,b;j)F(D_2a,E_2b;j)
 &=\sum_{\mathfrak t\ {\rm squarefree}}\mu(\mathfrak t)
   \sum_{\mathfrak t\mid a,b}\mathcal C_2(a,b;j)
                \widetilde F_{D_2,E_2}(a,b;j).
 \end{split}
\]
This is the full M\"obius indicator, not a truncation.  The individual
column shells and the retained row ball are finite.  Its equality
therefore follows by interchanging finite sums and using the genuine
correlation identity only when \((a,b)=1\).  Extra common primes of
\(\mathfrak t\) are not part of the genuine \(D_2,E_2\) support or
of its frequency restrictions.  Since the residual masks already
exclude that support, every nonzero \(\mathfrak t\) is disjoint from
it and from \(s\).

On each retained dyad, separate the full kernel and inverse roots in
the normalized \(h'\) norm and the two whole-product norms, as for the first transform.
The frozen factor \(G_cV_{\rm id}\) contributes only an outer phase
to the row Fourier power.  The fixed support boxes are chosen before
the current live labels, so the resulting coefficient measure is
common to them, to both rectangles, and to the row.  Then fix \(\mathfrak t\).  For each of its primes \(p\),
use the exact factor-allocation identity
\[
 1_{p\mid\prod_i n_i}
 =\sum_{\varnothing\ne J}(-1)^{|J|+1}
       \prod_{i\in J}1_{p\mid n_i},
\]
where the factors \(n_i\) are the two plain variables and the live
slots.  For a selected plain variable write \(n_i=pl_i\), with no
restriction on \(l_i\).  This extracts only a row scalar and
\(q_p^{it}\), and replaces both \(X_i,Y_i\) by
\(X_i/q_p,Y_i/q_p\).  It introduces no one-variable puncture.
A selected slot is frozen.  If two distinct divisor primes select
the same prime slot, the term is zero.  The quotients and unselected
factors may still contain \(p\), and may overlap the opposite side;
no new coprimality is imposed.  Old common masks remain common, and
the product of the two new plain scales is equal in the two rectangles.
The extracted row scalars, including zero scalars, are retained in
the signed identity.

Only now, for a fixed retained dyad and fixed Fourier parameters,
is each separated summand a product of two independent residual
convolutions, where \(\mathbf J\) records the factor allocations.
Their row characters are
\[
 \tau_1(n)\rho(n)\chi_n(G_cV_{\rm id}h'),\qquad
 \tau_1(n)\rho'(n)\chi_n(-G_cV_{\rm id}h'),
 \quad \rho,\rho'\in\Theta.
\]
They are understood with the fixed-factor support representation
above.  In particular their inducing-character ratio belongs to
\(\Theta\), including the supplementary factor
\(n\mapsto\chi_n(-1)\); this assertion takes no quotient at a zero.
Their natural row zeros are retained.  Extracting
the selected factors contributes row scalars, not new factors of the
residual character and not new support in \(q'\).  These scalars,
including zeros, remain in the signed identity.

\begin{lemma}[Whole-product separation and transport]
\label{lem:whole-product-transport}
Consider a finite column shell with allocated coefficients
\(A_C(a),A_D(b)\) defined by \eqref{own:plain-allocation}.
Keep all arithmetic masks, the complete signed rectangle
differences, and all live prime sums in those coefficients.
A smooth multiplier depending only on the row norm, frozen outer
norms, and the two whole norms \(q_a,q_b\), on the fixed boxes
specified above, has a Fourier separation common to every row and
live tuple.  On either column it introduces exactly one additional
norm power of the whole product.  Its two plain variables and all
live slots therefore acquire that same norm power, not independently
selected powers.

This conclusion persists through both of the following operations:
complete extraction at a fixed genuine common support, and the full
M\"obius allocation of a residual divisibility condition.  The former
leaves one common puncture on every remaining factor; the latter
does not introduce a puncture on any quotient.  In either case the
two formal products of plain scales remain exactly equal.
\end{lemma}
\begin{proof}
Let \(T_a,T_b\) be the fixed whole-column scales in the sector,
and use normalized independent coordinates
\(x=q_a/T_a,\ y=q_b/T_b\) and \(\boldsymbol z\) for the row and
outer norms occurring in the smooth multiplier \(H\).
Multiply \(H\) by larger cutoffs that are one on the \emph{full}
individual supports.  The resulting compact logarithmic profile
has the exact inversion
\[
 H(x,y,\boldsymbol z)
  =(2\pi)^{-d}\int
       \widehat H(t_a,t_b,\boldsymbol t)
       x^{it_a}y^{it_b}\boldsymbol z^{i\boldsymbol t}
             \,dt_a\,dt_b\,d\boldsymbol t .
\]
This density is defined on its ambient box before any arithmetic
indices are evaluated.  All arithmetic relations restrict evaluation
points, not that density.  Its weighted absolute integrals are
bounded by the finite-seminorm estimate in
Appendix~\ref{app:smooth}.  For the kernels here the bounds are uniform
in their positive scale parameter by
Lemma~\ref{lem:kernel-seminorms}; normalized inverse roots are fixed
smooth powers on the same boxes.

In the product \(A_C(a)\overline{A_D(b)}\), put the first mode
inside \(A_C\) and put the \emph{negative} second mode inside the
conjugated polynomial.  The two column factors are then
\[
 A_C(a)(q_a/T_a)^{it_a},\qquad
 \overline{A_D(b)(q_b/T_b)^{-it_b}},
\]
which restore the two modes exactly.  Since \(a=\Pi l_1l_2\),
the only nonconstant factor on its side is
\[
 q_a^{it_a}=q_\Pi^{it_a}q_{l_1}^{it_a}q_{l_2}^{it_a}.
\]
It multiplies the whole difference \(D_{\mathbf b}\), unchanged,
and has the same exponent in both rectangles.  The constants
\(T_a^{-it_a},T_b^{-it_b}\) are outer.  The argument also covers
an inherited norm power by adding its exponent.  Fixed-ray
separation similarly supplies one multiplicative character on the
whole product.  A common coprimality mask factors without any
division, including at a zero.

For genuine complete extraction, fix all extracted prime powers
in each plain variable and all frozen slots.  Write the two
extracted plain ideals as \(d_1,d_2\).  Exactness of the valuations
is equivalent to puncturing \emph{every} remaining factor at the
extracted support; this is precisely the mask in
\eqref{own:plain-allocation}.  Its two pairs of scales become
\[
 (X_1/q_{d_1},X_2/q_{d_2}),\qquad
 (Y_1/q_{d_1},Y_2/q_{d_2}),
\]
whose products are equal.  The extracted character and norm powers
are common scalars, still including any zero.  An empty rectangle
has value zero but remains in this formal identity.

For a residual divisor prime, the inclusion--exclusion identity
\[
 1_{p\mid\prod n_i}
   =1-\prod_i(1-1_{p\mid n_i})
\]
holds pointwise even if several factors contain \(p\).
In a chosen nonempty subset, replacing a selected plain variable
by \(pl\) imposes \emph{no} restriction \(p\nmid l\).
It extracts the same scalar from both rectangles and divides
their corresponding scale by the same exact \(q_p\).
A selected prime slot is frozen; disjoint underlying slot supports
make inconsistent multiple selections identically zero.
The old common masks remain unchanged.  Multiplying these finite
identities over the squarefree divisor, before taking an absolute
value, proves the last assertion.  Extra overlaps introduced by
the full M\"obius sum are thus allowed, rather than mistakenly
treated as a new primitive conductor.
\end{proof}

\begin{lemma}[Common coefficient under complete extraction]
\label{lem:centered-coefficient-invariant}
For a centered input, the two transforms and the intervening Gauss-row
enlargement just constructed preserve the following coefficient data.
Each Gauss polynomial and each amplifier error retains
Equation~\eqref{eq:centered-coefficient-form}, multiplied by \(G(u,h)\),
apart from bounded frozen scalars and row scalars of modulus at most one.
For an uncentered input, omit the second rectangle throughout.
After the second transform and the full M\"obius factor allocation, each
separated child side has that coefficient form without \(G\) and without
\(s\mid u\).  Every surviving slot has the child's one whole-product
character and its original coefficient \(\nu_i\), with no Gauss coefficient.
Within a side the plain variables have the same character, puncture mask,
and norm power in both rectangles.  The two sides of one squared Gauss norm
may have different norm powers; their inducing characters differ by a
member of \(\Theta\).

If the new inducing character belongs to \(\Theta\), fixing the live slot
labels leaves the plain coefficient
\[
 \vartheta(l_1)\vartheta(l_2)
 1_{(l_1l_2,\mathfrak R_*)=1}q_{l_1l_2}^{it}D_{\mathbf b}(l_1,l_2)
\]
for one \(\vartheta\in\Theta\), one squarefree mask
\(\mathfrak R_*\) of polynomial norm, and one real \(t\).
The mask may depend on the frozen row, but is common to both variables
and both rectangles.  All coefficients are understood with their retained
frozen scalars, including zeros, and the exact normalizations recorded below.
\end{lemma}

\begin{proof}
The first-transform calculation proved the Gauss coefficient form; the
local \(p\)-power calculation proved it for each amplifier error.
Lemma~\ref{lem:whole-product-transport} supplies the common Fourier
density and the exact preservation of the signed rectangle through
both extractions.  In between, the genuine complete-support
correlation identity removes the Gauss factors and supplies one
character of the whole residual product.  All primes of that
genuine common support have already been punctured.  The full
M\"obius allocation therefore proves the stated child form, with
one common character and norm power on each side.  Thus ``common'' concerns one separated side;
the displayed character ratio is the relation between the two sides.
When the inducing character belongs to \(\Theta\), it equals some
\(\vartheta\in\Theta\) on units.  Its redundant natural zeros and
the fixed punctures combine into one \(\mathfrak R_*\) of polynomial
norm.  For each fixed \(\Pi\), multiplicativity factors its value,
mask, and norm power on \(\Pi l_1l_2\) into a scalar in \(\Pi\)
times the three factors stated in the conclusion.  This proves the exceptional assertion.
\end{proof}

We next record the exact normalization of the two residual convolutions.  Put
\[
 \alpha_1=a_0-c_2,\qquad \alpha_2=a_0-d_2,\qquad
 \frac{\alpha_1+\alpha_2}{2}=a_0-b_2.
\]
At the entrance to the current two-transform stage, let
\(\mathcal I_{\rm in}\) be its live slot set and let \(X_1,X_2\)
be the formal scales of its first rectangle.  Their exact convention is
\(X_1X_2\prod_{i\in\mathcal I_{\rm in}}P_i=Z^A\).
For side \(j\in\{1,2\}\), let \(\mathcal F_j\) be the slots from this
entrance set frozen by the first and second genuine common-support
extractions, let \(\mathcal I_j\) be the slots still live before the
final \(\mathfrak t\)-allocation, and let \(Q_{{\rm plain},j}\) be
the product of the exact norms of the plain powers extracted at those
steps and at the amplifier.  Slots removed in an ancestor stage are not
included in \(\mathcal F_j\).  Let \(T_j\) be the common formal
pre-\(\mathfrak t\) product of the two plain scales in its rectangles;
for an uncentered input use its one formal plain product.  Before using
slot ratios, discard a frozen-slot profile-zero term, which is identically
zero by its frozen data independently of the row and live labels.
Complete extraction gives the exact identities
\[
 \begin{aligned}
 Z^{A-\alpha_j}
   &=Q_{{\rm plain},j}\prod_{i\in\mathcal F_j}q_{p_i},\\
 T_j\prod_{i\in\mathcal I_j}P_i
   &=\frac{Z^A}{Q_{{\rm plain},j}\prod_{i\in\mathcal F_j}P_i}
     =Z^{\alpha_j+e_j},\\
 e_j&=\sum_{i\in\mathcal F_j}\log_Z(q_{p_i}/P_i),
 \qquad |e_j|\le\theta_N.
 \end{aligned}
\]
The amplifier contributes only exact plain powers to this calculation.
Each earlier frozen slot occurs once, even if it shared its prime with
a plain.  The already separated inverse roots, kernel, and fixed
normalization constants remain outside \(e_j\).

For side \(j\), let \(d_{j,i}\) be the product of final
\(\mathfrak t\)-primes selected in plain \(i\), and let
\(\mathcal J_j\subseteq\mathcal I_j\) be the newly frozen slots.  Put
\[
 a_{j,i}=\log_Zq_{d_{j,i}},\qquad
 \widetilde r_j=a_{j,1}+a_{j,2}+\sum_{i\in\mathcal J_j}z_i\ge0,
 \qquad
 \omega_j=\sum_{i\in\mathcal J_j}\log_Z(q_{p_i}/P_i).
\]
The corresponding formal scale is divided by \(q_{d_{j,i}}\) in both
rectangles.  An assignment selecting one prime slot at two distinct
\(\mathfrak t\)-primes, or a newly frozen slot with zero profile value,
is identically zero from the frozen data and is discarded before these
ratio bounds are used.  All other extracted zeros remain until the
nonnegative or absolute estimate described below.

Use lower-endpoint divisor dyads \(T\le q_{\mathfrak t}<2T\),
\(T\ge1\), and put \(t_-:=\log_ZT\).  Each dyad contains
\(O(Z^{t_-})\) ideals.  The actual selected product on side \(j\)
is divisible by \(\mathfrak t\), so
\begin{equation}
 \begin{aligned}
 \widetilde r_j+\omega_j
 &=\log_Z\left(q_{d_{j,1}d_{j,2}}
                  \prod_{i\in\mathcal J_j}q_{p_i}\right)
 \ge\log_Zq_{\mathfrak t}\ge t_-,\\
 |\omega_j|&\le\theta_N,\qquad |e_j+\omega_j|\le\theta_N
 \quad(j=1,2).
 \end{aligned}
 \label{eq:centered-divisor-boundary}
\end{equation}
The last inequality uses the disjoint subsets \(\mathcal F_j\) and
\(\mathcal J_j\) of the current stage entrance slots.  It does not
include ancestor slots or fixed separation constants.

Writing \(T'_j=T_j/q_{d_{j,1}d_{j,2}}\), the post formal factor product
and the nominal raw child scale are
\begin{equation}
 T'_j\prod_{i\in\mathcal I_j\setminus\mathcal J_j}P_i
   =Z^{\alpha_j+e_j-\widetilde r_j},\qquad
 N_{{\rm post},j}=Z^{\alpha_j-\widetilde r_j}.
 \label{eq:centered-raw-normalization}
\end{equation}
Define \(P_{C,\mathfrak t,\mathbf J}(h')\) and
\(P_{D,\mathfrak t,\mathbf J}(h')\) to be their residual convolutions
multiplied by \(N_{{\rm post},j}^{-1/2}\).  Relative to the nominal
pre-normalizer \(Z^{-\alpha_j/2}\), the exact extraction coefficient
on side \(j\) is \(Z^{-\widetilde r_j/2}\), apart from the frozen slot
amplitudes, retained extracted row scalars, and already separated outer
factors.  These raw children retain the formal signed rectangles and are
not moment-lemma invocations; their formal logarithmic lengths may be
negative.  No clipping has occurred.

The factor assignments are divisor-bounded for fixed \(N\); their
total is bounded using \(C_N^{\omega(\mathfrak t)}
\ll_{N,a}q_{\mathfrak t}^a\) for any fixed \(a>0\), not by
spending a fixed exponent at each prime.  The \(O(\log Z)\) divisor
dyads in the bounded column range use their existing logarithmic share.
These discrete losses are separate from the numerical \(\theta_N\) terms.

For precision, the other common row scalar after extracting
\(Z^{g_2-t_2}\) is
\[
 v_{{\rm part},\lambda}(h')=
 \omega_\lambda(q_{G_cV_{\rm id}}q_{h'}/T_\lambda)
 1_{\rm part}(h')
 \frac{F(D_2,E_2;G_cV_{\rm id}h')}{Z^{g_2-t_2}}\zeta(h'),
 \qquad |\zeta(h')|\le1,
\]
where \(\zeta\) contains the separated row phases for this dyad and
these Fourier parameters.  It is zero off the retained partition and
enclosing ball, and \(|v_{{\rm part},\lambda}|\le1\).
The subsequent estimates are made dyad by dyad and then summed with
the already allowed \(O(\log Z)\) mass.  Split the already factorized row sum according to membership of the
inducing character in \(\Theta\).  On rows outside \(\Theta\),
Cauchy--Schwarz yields two nonnegative child norms; on rows in
\(\Theta\), take a pointwise absolute product.  Only at these steps
may the absolute values of the common row weight, partition scalar,
or extracted row scalars be bounded by one and the rows enlarged to
the common \(m'_{\rm act}\) ball.  This can admit extra exceptional
rows that violate a former unit restriction, but the character-only
count below includes them.  Both sides use the same eligibility class
and compare with Equation~\eqref{old-eq:2.14}.

Consider first rows whose child inducing character is outside
\(\Theta\).  The two sides have the same eligibility condition
because their characters differ by a member of \(\Theta\).
On such rows, first bound a centered child norm by the sum of the norms
of its two rectangles; for an uncentered child there is one rectangle.
For one fixed side \(j\) and rectangle \(\varrho\), write its pre and
post formal plain scales as \(S_{\varrho,i}\) and
\(S'_{\varrho,i}=S_{\varrho,i}/q_{d_{j,i}}\).  Choose a fixed upper
support endpoint \(B_i\) for each current plain profile type before
the current row, divisor, and live labels.  Omit a rectangle only when
\(S'_{\varrho,i}B_i<1\) for some \(i\); then that plain factor is zero
on every nonzero integral ideal.  Retain equality and every other profile,
mask, row-scalar, or arithmetic zero.  This test is independent of the
row and live labels.

For a retained rectangle put
\(x_i=\log_ZS_{\varrho,i}\) and \(y_i=x_i-a_{j,i}\).  The fixed
endpoint test and \(a_{j,i}\ge0\), with \(H_N\) enlarged to contain
the two positive parts of the upper endpoint logs, give the exact clipped identities
\begin{equation}
 \begin{aligned}
 \pi_{0,j,\varrho}&=\sum_{i=1}^2(-x_i)_+,\qquad
 \pi_{j,\varrho}=\sum_{i=1}^2(-y_i)_+,
 \qquad 0\le\pi_{0,j,\varrho}\le\pi_{j,\varrho}\le\theta_N,\\
 r_{{\rm clip},j,\varrho}
 &=\sum_{i\in\mathcal J_j}z_i+
       \sum_{i=1}^2\bigl((x_i)_+-(y_i)_+\bigr)
 =\widetilde r_j-(\pi_{j,\varrho}-\pi_{0,j,\varrho})\ge0,\\
 A_{{\rm clip},j,\varrho}
 &=\sum_{i=1}^2(y_i)_++\sum_{i\in\mathcal I_j\setminus\mathcal J_j}z_i
 =\alpha_j+e_j-\widetilde r_j+\pi_{j,\varrho}\\
 &=\alpha_j+e_j+\pi_{0,j,\varrho}-r_{{\rm clip},j,\varrho}.
 \end{aligned}
 \label{eq:centered-clipped-rectangle}
\end{equation}
In particular the clipped total decreases by \(r_{{\rm clip},j,\varrho}\)
from its own pre-clipped total \(\alpha_j+e_j+\pi_{0,j,\varrho}\).
Different rectangles can have different clipped totals.  A retained
subunit scale lies in the fixed interval \([1/B_i,1]\), so its dilation
to scale one preserves finite seminorm bounds.  The row character, common
mask, and surviving slot weights are unchanged.  The exact coefficient
from the pre-normalizer \(Z^{-\alpha_j/2}\) to this standard clipped
product is
\[
 Z^{(A_{{\rm clip},j,\varrho}-\alpha_j)/2}
 =Z^{(-r_{{\rm clip},j,\varrho}+e_j+\pi_{0,j,\varrho})/2}.
\]
For a fixed pair of retained rectangles, suppress their rectangle indices.
Equation~\eqref{eq:centered-divisor-boundary} then bounds the paired
divisor count and these coefficients by
\begin{equation}
 \begin{split}
 t_- -\frac{r_{{\rm clip},1}+r_{{\rm clip},2}}2
     +\frac{e_1+\pi_{0,1}+e_2+\pi_{0,2}}2
 &\le\frac{(e_1+\omega_1)+\pi_1+(e_2+\omega_2)+\pi_2}{2}\\
 &\le2\theta_N.
 \end{split}
 \label{eq:centered-clipped-shell}
\end{equation}
There are at most two rectangles per side, hence at most four such pairs;
their triangle factor is fixed.  No common clipped reduction is used for
the signed difference.

On the \(c_2\) side, before the last coprimality extraction, the
nominal total length is \(A'=\alpha_1=A-c-w-c_2\).  For the positive-slot parameters,
\begin{equation}
 \begin{split}
 (A'-M')-(A-M)
 &=g+w-d-c_2-(K_0-K)-w_o+g_2-t_2\\
 &\le6(w+\ell)+\delta_{{\rm fr},1}.
 \end{split}
 \label{old-eq:3.14}
\end{equation}
To prove the inequality, first use \(g_2-t_2\le c_2\).
In the amplified main norm, \(w=w_o=0\) and
\[
 g=J_++2\sigma
 \le d+(K_0-K)+6\ell+\delta_{{\rm fr},1}.
\]
For an error,
\(g\le d+(K_0-K)+\sigma+\delta_{{\rm fr},1}\), so the excess is at most
\(w-w_o+\sigma+\delta_{{\rm fr},1}\).  These inequalities use
\(d+K_0-K\ge-\delta_{{\rm fr},1}\) and the
\(1\)-Lipschitz property of the positive part.  The three choices in
Equation~\eqref{old-eq:2.9}, together with
\(\ell_p\ge\sigma/6\), give
\(w-w_o+\sigma\le6w\).  This proves
Equation~\eqref{old-eq:3.14}.

For each retained clipped rectangle, if no slot survives, apply the
completed unrestricted zero-slot assertion at the declared width
\(M'_{\rm act}\).  If slots survive, first apply the algebraic mask deletion in
Equations~\eqref{old-eq:2.1a}--\eqref{old-eq:2.1b}.
It expresses this standard clipped product as natural products and only
decreases its nonnegative affine expression.  For each actual product, after both frequency
enclosures, the \(\mathfrak t\)-allocation, and mask deletion, let
\(A_{\rm act},z_{\rm act}\) be its declared total and slot lengths
and put
\[
 F_{\rm act}
 =\bigl(A_{\rm act}-M'_{\rm act}+(6\kappa-1)z_{\rm act}\bigr)_+.
\]
The parent satisfies Equation~\eqref{old-eq:3.9}.  Before this mask
deletion, let \(z'\le z\) be the surviving slot length of the clipped
rectangle.  Equations~\eqref{old-eq:3.14} and
\eqref{eq:centered-clipped-rectangle}, together with \(M'_{\rm act}\ge M'\), give
\[
 \begin{split}
 A_{{\rm clip},j,\varrho}-M'_{\rm act}+(6\kappa-1)z'
 &\le \alpha_j-M'+(6\kappa-1)z
       +e_j+\pi_{j,\varrho}-\widetilde r_j-(M'_{\rm act}-M')\\
 &\le6(w+\ell)+\delta_{{\rm fr},1}+e_j+\pi_{j,\varrho}.
 \end{split}
\]
The same calculation holds on the other side with \(c_2,d_2\) exchanged.
Taking the positive part and then deleting the fixed mask therefore gives
\[
 F_{\rm act}\le6(w+\ell)+\delta_{{\rm fr},1}
                 +(e_j+\pi_{j,\varrho})_+
 \le6(w+\ell)+\delta_{{\rm fr},1}+2\theta_N.
\]
Here \(z_{\rm act}\le z'\le z\).  All the displayed ledgers use the fixed list
of aggregate lengths
\[
 A,z,c,d,p,R,E,s_0,K,w,w_o,g,\ell,c_2,d_2,p_2,g_2,t_2,V.
\]
Their affine coefficients and positive-part Lipschitz constants are
numerical and independent of \(N\).  The displayed frequency, raw
normalization, and rectangle-clipping bounds are therefore included in
\(C_*\xi\), for a fixed \(C_*\ge1\) independent of \(N\), once the
aggregate threshold has been imposed.  The existence of this fixed error bound uses the single
\(\theta_N=H_N/\log Z\) before comparing lengths, not \(N\) separate
copies of \(\xi\).  We enlarge \(C_*\) below to cover the fixed number
of operations in a stage.

If \(F_{\rm act}>0\), remove whole slots from this product until the
remaining product satisfies Equation~\eqref{old-eq:3.9} or
until no slot remains.  Removal here means applying the pointwise
bound Equation~\eqref{old-eq:3.5} to the entire selected prime
polynomial; no prime label is frozen.  A removed slot of length
\(d\) decreases \(A_{\rm act}+(6\kappa-1)z_{\rm act}\) by \(6\kappa d\)
and contributes \(\kappa d\) to the squared \(Z\)-exponent, with the
fixed seminorm and height factor retained separately.  Since each
slot has length at most \(\eta\), order the positive live lengths
and take the first prefix reaching \(F_{\rm act}/(6\kappa)\).
Its preceding prefix is smaller than that threshold and its final
slot has length at most \(\eta\).  If no prefix reaches the threshold,
remove all slots, whose total is smaller.  Thus exactly one slot can
cause an overshoot, and the total removed length \(d_z\) satisfies
\begin{equation}
 \begin{split}
 d_z&\le\min\left\{z_{\rm act},
                     \frac{F_{\rm act}}{6\kappa}+\eta\right\},\\
 \kappa d_z&\le\frac{F_{\rm act}}6+\kappa\eta
 \le\Delta_{\rm child}+\frac{\delta_{{\rm fr},1}}6
                       +\frac{\theta_N}{3}+\eta.
 \end{split}
 \label{old-eq:3.16}
\end{equation}
If the slots disappear before the affine boundary is reached, the
remaining plain lengths may be arbitrary; this is exactly why the
completed smaller-width zero-slot assertion includes all lengths.
Otherwise apply the positive-slot induction to the remaining
product.  Its slots still have their original coefficient class,
and its inducing rows remain outside \(\Theta\).  Cauchy--Schwarz
averages the errors of the two separately clipped children, so their
possibly different rectangles and slot sets do not double \(\eta\).
Combining Equation~\eqref{old-eq:3.16} with the paired coefficient
bound in Equation~\eqref{eq:centered-clipped-shell} and
\(M'_{\rm act}-M'\le\delta_{{\rm fr},2}\), the strict-edge scale
cost beyond the child envelope is at most
\begin{equation}
 \delta_{{\rm fr},2}+\frac{\delta_{{\rm fr},1}}6
        +\eta+\frac{7\theta_N}{3}.
 \label{eq:centered-child-edge-cost}
\end{equation}
If no slot survives, the same bound holds without needing a greedy cost.
The numerical \(\theta_N\) terms and the two frequency corrections
give the single \(O(\xi+\eta)\) edge loss in the \(\epsilon_G\)
reserve inherited from Equation~\eqref{old-eq:2.7} when using
Equation~\eqref{old-eq:2.14}; the existing local small-power and fixed
analytic factors remain separate.  The pointwise estimate is requested
once for the entire removed product; its internal small-power shares
may depend on \(N\).  This one greedy operation occurs only after
all boundary defects of the actual child have been included in
\(F_{\rm act}\).  The argument for the \(d_2\)
side is the same with \(c_2,d_2\) exchanged.

We have now bounded all child rows whose inducing characters lie outside
\(\Theta\) using only smaller widths.  For the remaining rows, the next
count and volume bound handle the uncentered range; the centered range
uses the subtraction retained in the coefficient lemma.

\subsection{Rows with inducing characters in \texorpdfstring{\(\Theta\)}{Theta}}

For this subsection, call a child row \emph{exceptional} when its
inducing character belongs to \(\Theta\).  By the coefficient
lemma, either both separated sides are exceptional or neither is.
Every prime in the existing full displayed moving support is a common
column zero, as is every extra common puncture.  This is true initially,
remains true for \(\mathfrak e,\mathfrak r\) in the first transform,
and remains true for an amplifier prime.  For a nonzero genuine
second allocation, the factors \(\tau_1(D_2),\tau_1(E_2)\) and the
old common masks therefore force its complete radical to be disjoint
from all those supports.  This conclusion is made before replacing
any frozen scalar by an upper bound.  The artificial residual
extension keeps \(D_2,E_2\) and those masks fixed and cannot revive
an impossible allocation.

We make explicit the finite-ray reduction for row factors at the fixed
primes.  As in Equation~\eqref{own:row-fixed-ray}, write a nonzero row as
\(h'=\varepsilon h_{\mathcal S}h_{\rm good}\), with
\(\varepsilon\) a unit and the two other factors supported on
\(\mathcal S\) and its complement.  For a primary element \(n\)
prime to \(\mathcal S\), reciprocity gives, with every zero retained,
\[
 \chi_n(h')=\chi_n(\varepsilon h_{\mathcal S})
   \mathcal R(n,h_{\rm good})
   \prod_{p\mid h_{\rm good}}\chi_p(n)^{v_p(h')}.
\]
By the fixed-numerator ray argument in Appendix~\ref{app:conventions}, the first factor belongs to
a finite family of ray characters supported on \(\mathcal S\) after
the unit and the \(\mathcal S\)-valuations modulo six are fixed.
The second belongs to the fixed reciprocity family after fixing the
good ray sector.  We do not absorb any moving good prime into this
fixed family.  The family supplied by the first factor need not be
contained in \(\Theta\); there are simply finitely many such choices.
Every remaining displayed good-prime factor has its actual local
sextic exponent, with nontrivial exponents ramified at that prime.

Let \(v_1\) be the radical length of the primes counted in \(V\)
whose equal multiplicity is one, and put \(f=2v_1\).  At such a
prime the fixed factor \(G_cV_{\rm id}\) in the row has valuation
two.  There is no existing moving character at that prime, and
every character of \(\Theta\) is unramified there.  Sextic
reciprocity therefore shows that exceptional induction requires
\[
 v_p(h')+2\equiv0\pmod6,\qquad\text{that is,}\qquad
 v_p(h')\equiv4\pmod6.
\]
More generally, at every prime outside \(\mathcal S\), exceptional
induction prescribes a single residue class modulo six for
\(v_p(h')\), determined by the frozen moving character and
\(G_cV_{\rm id}\).  Outside their supports that residue is zero.
For each of the finitely many choices at \(\mathcal S\) and of
unit and fixed-ray data, the ideal of \(h'\) consequently has a
unique form
\[
 (h')=\mathfrak h_0\mathfrak v^6,
\]
where \(\mathfrak h_0\) is fixed and sixth-power-free.  The
displayed valuation-four conditions give
\(q_{\mathfrak h_0}\ge Z^{4v_1}=Z^{2f}\).
Ideal counting up to \(q_{h'}\ll Z^{m'_{\rm act}}\) now gives the
stronger bound \(O(Z^{(m'_{\rm act}-2f)/6+\epsilon_1})\) on every
nonempty range, with the usual bounded-scale convention.  We use
only the weaker bound
\begin{equation}
 \#\{h':q_{h'}\ll Z^{m'_{\rm act}},\ \text{exceptional}\}
 \ll Z^{(m'_{\rm act}-f)/6+\epsilon_1}
 \le Z^{(m'-f)/6+\delta_{{\rm fr},2}/6+\epsilon_1}.
 \label{eq:exceptional-row-count}
\end{equation}
If the exponent would describe a scale below one, the forced ideal
\(\mathfrak h_0\) makes the range empty except at the same
bounded-scale boundary.  This argument includes the principal
character and every other member of \(\Theta\).  Additional
conditions at old moving primes can only reduce the count.
It also covers rows admitted when the partition scalar was bounded
after positivity.  At a unit prime, such an added row can be
exceptional even though it was absent from the original partition;
no extra forcing saving from the unit set \(t_2\) is used.

On exceptional rows take the absolute product for each already
allocated pair of children and multiply by
Equation~\eqref{eq:exceptional-row-count}.  Let
\(c_{\mathfrak t,\mathbf J}(h')\) denote its central extraction
coefficient and bounded extracted scalars.  For an allocation with raw
reductions \(\widetilde r_1,\widetilde r_2\),
\(|c_{\mathfrak t,\mathbf J}(h')|\ll_N
 Z^{-(\widetilde r_1+\widetilde r_2)/2}\) after taking absolute values.
Using absolute volume rather than cancellation, the unnormalized plain
difference on side \(j\) is \(O(T'_j)\) at every positive formal scale,
including subunit scales, and the unnormalized live slots contribute
\(O_N(\prod_{i\in\mathcal I_j\setminus\mathcal J_j}P_i)\).
Equation~\eqref{eq:centered-raw-normalization} therefore gives the raw
normalized volume
\[
 |P_{j,\mathfrak t,\mathbf J}(h')|
 \ll_N N_{{\rm post},j}^{-1/2}T'_j
             \prod_{i\in\mathcal I_j\setminus\mathcal J_j}P_i
 =Z^{(\alpha_j-\widetilde r_j)/2+e_j}
\]
up to the fixed seminorm and height factors, where \(j\) denotes its
\(C\) or \(D\) side.  Thus the exponent for an allocated pair and one
lower-endpoint divisor dyad, including its count, is
\[
 a_0-b_2+e_1+e_2+t_- -\widetilde r_1-\widetilde r_2
 \le a_0-b_2+e_1+(e_2+\omega_2)
 \le a_0-b_2+2\theta_N.
\]
Here Equation~\eqref{eq:centered-divisor-boundary} gives
\(t_- -\widetilde r_2\le\omega_2\), and \(\widetilde r_1\ge0\).
Consequently
\[
 \sum_{\mathfrak t,\mathbf J}
 |c_{\mathfrak t,\mathbf J}(h')|
 |P_{C,\mathfrak t,\mathbf J}(h')
       P_{D,\mathfrak t,\mathbf J}(h')|
 \ll Z^{a_0-b_2+2\theta_N+\epsilon_1}.
\]
The displayed sum is over a fixed lower-endpoint divisor dyad; its
allocations have already been included in
\(\epsilon_1\).  It does not assert an independent product before
M\"obius allocation.  Subtracting the inner allowance from the
reference volume and the nominal exceptional count gives the exact
algebraic identity
\begin{equation}
 \frac{m'-f}{6}+a_0-b_2-(M'+\Delta_{\rm child})
 =A-\frac56M-F_1-F_2,
 \label{old-eq:2.15}
\end{equation}
where
\begin{equation}
 \begin{aligned}
 F_1&=c/6+5\{d+(K_0-K)\}/6+w/3+\widetilde q/6
       +w_o+\mathcal B_c-5g/6+\ell,\\
 F_2&=2b_2-\tfrac56g_2-p_2+t_2+V/6+f/6.
 \end{aligned}
 \label{old-eq:2.16}
\end{equation}
The actual excess is bounded by the nominal expression in
Equation~\eqref{old-eq:2.15} plus
\(\delta_{{\rm fr},2}/6+2\theta_N+\epsilon_1\).
For an explicit check, the left side of
Equation~\eqref{old-eq:2.15} first equals
\[
 a_0-2b_2+p_2-w-\mathcal B_c-\ell
       -\tfrac56m'-q'-f/6.
\]
Substitution of Equation~\eqref{old-eq:2.13} and
\(K=2A-c-d+R+E-m-(K_0-K)\) gives
Equations~\eqref{old-eq:2.15}--\eqref{old-eq:2.16}.

The following lower bounds are the reason complete common supports
do not consume the available exponent:
\begin{equation}
 F_1\ge\tfrac23(c+w)-3\sigma-\tfrac56\delta_{{\rm fr},1},\qquad
 F_2\ge\tfrac23b_2.
 \label{old-eq:2.17}
\end{equation}
To prove the first, put \(D_0=d+K_0-K\), initially set
\(g=J_+\) and \(\ell=0\), and recall
\(J=D_0-c-2w+w_o\).  When \(J\ge0\), direct simplification
gives
\[
 F_1=c+2w+\widetilde q/6+w_o/6+\mathcal B_c
       \ge c+2w.
\]
When \(J<0\), use \(\widetilde q\ge R\) and
\[
 \mathcal B_c+\widetilde q/6
 \ge (3c-5D_0)_+/6-\tfrac56\delta_{{\rm fr},1}.
\]
Indeed, the left side is at least
\(\{(3c-5d-R)_++R\}/6\), which is at least
\((3c-5d)_+/6\), while
\(D_0\ge d-\delta_{{\rm fr},1}\).  The positive part is
\(1\)-Lipschitz, giving the displayed correction.
It follows that
\[
 F_1\ge
 \frac{c+5D_0+(3c-5D_0)_+}{6}+\frac w3
       -\frac56\delta_{{\rm fr},1}
 \ge\frac23(c+w)-\frac w3-\frac56\delta_{{\rm fr},1}.
\]
The actual choices have \(g\le J_++2\sigma\) and
\(\ell\ge0\).  Equation~\eqref{old-eq:2.10} therefore bounds
the additional loss from \(2(c+w)/3\) by
\[
 w/3+5\sigma/3+\tfrac56\delta_{{\rm fr},1}
 \le22\sigma/9+\tfrac56\delta_{{\rm fr},1}
 <3\sigma+\tfrac56\delta_{{\rm fr},1}.
\]

For the second inequality in Equation~\eqref{old-eq:2.17}, compute
prime by prime from the definition of \(F_2\).  In units of the
prime's logarithmic norm, the contributions are
\[
 \begin{array}{c|c|c}
 \text{common case}&F_2&b_2\\ \hline
 i=i,\ 6\nmid i,\ \text{unit}&7i/6&i\\
 i=i,\ 6\nmid i,\ \text{nonunit}
       &(7i-5)/6+\tfrac13 1_{i=1}&i\\
 i=i,\ 6\mid i,\ \text{either}&7i/6-1&i\\
 i>j_0,\ 6\mid j_0,\ \text{unit}
       &i+j_0/6-1&(i+j_0)/2 .
 \end{array}
\]
The extra \(1/3\) in the second line is \(f/6\).  The first
line exceeds \(2b_2/3\).  In the second line equality holds at
\(i=1\), and for \(i\ge2\) the difference is
\((3i-5)/6\ge0\).  In the third line the difference is
\((i-2)/2\ge0\), because \(i\ge6\).  In the last line it is
\((4i-j_0-6)/6\ge0\), because \(j_0\ge6\) and
\(i\ge j_0+1\).  Summing proves the claim.

For the uncentered range \(A\le5M/6\),
Equations~\eqref{old-eq:2.15}--\eqref{old-eq:2.17}
bound the nominal exceptional excess by
\(3\sigma+5\delta_{{\rm fr},1}/6\).  The actual-count and aggregate
support corrections recorded above are further \(O(\xi)\) terms
in the fixed stage allowance.  Together with the
diagonal estimate and the smaller-width estimates, this proves the
uncentered stage at the current width with the reserved terminal
loss.  The comparisons constructed earlier may consequently use
that stage.  It remains to estimate the centered exceptional terms
when \(A>5M/6\).

For one separated side of a fixed exceptional row, fix the live slot
labels.  Lemma~\ref{lem:centered-coefficient-invariant} then gives the same
character, puncture, and norm power in the two plain variables and in both
rectangle terms.  The next lemma shows that the coefficient of \(X^{1+it}\)
in each plain sum is independent of \(X\), so the two product main terms
agree at equal products of scales.

\begin{lemma}[Masked rectangle cancellation]
\label{lem:centered-lattice-cancellation}
Let \(\vartheta\) range over a fixed finite set of finite-order
ray characters with conductor primes in \(\mathcal S\).  Let
\(\mathfrak R_*\) be squarefree with \(q_{\mathfrak R_*}\le Z^{B_*}\)
for a fixed \(B_*\), and let \(W_1,W_2\) be smooth profiles on
fixed annuli.  For \(X>0\) and \(t\in\mathbb R\), define
\[
 \mathcal L_i(X,t):=
 \sum_l\vartheta(l)1_{(l,\mathfrak R_*)=1}
           q_l^{it}W_i(q_l/X).
\]
Then
\begin{align}
 \mathcal L_i(X,t)
 &=c_{\vartheta,\mathfrak R_*}X^{1+it}I_i(t)
   +O\left(Z^{\epsilon_1}(1+|t|)^J\right),
 \label{old-eq:2.18d}\\
 |\mathcal L_i(X,t)|&\ll X.
 \label{old-eq:2.18e}
\end{align}
Here \(J\) is fixed, the constants use finitely many seminorms,
\(c_{\vartheta,\mathfrak R_*}\) is independent of \(X,t\), and
\(I_i(t)\) is a fixed measure constant times
\(\int_0^\infty W_i(y)y^{it}\,dy\).
Both estimates are uniform for \(X>0\), for the stated masks, and
for the finite character set.

Let \(D_{\mathbf b}\) be as in
Equation~\eqref{old-eq:2.18a}, and put
\[
 U_i=X_i/q_{\mathfrak b_i},\qquad
 V_i=Y_i/q_{\mathfrak b_i},\qquad
 T=U_1U_2=V_1V_2.
\]
If \(U_i,V_i\ge Z^r\) for \(i=1,2\) and some \(r\ge0\),
then
\begin{equation}
 \begin{split}
 T^{-1/2}\left|
 \sum_{l_1,l_2}
 \vartheta(l_1)\vartheta(l_2)
 1_{(l_1l_2,\mathfrak R_*)=1}
 q_{l_1l_2}^{it}D_{\mathbf b}(l_1,l_2)
 \right|
 \ll Z^{\epsilon_1}(1+|t|)^{2J}T^{1/2}Z^{-r}.
 \end{split}
 \label{old-eq:2.18f}
\end{equation}
Without a nonnegative lower length, the same expression is
\(O(T^{1/2})\).
\end{lemma}

\begin{proof}
First omit the mask.  The primary generators representing ideals
outside \(\mathcal S\), with the fixed character weight
\(\vartheta\), are a finite weighted collection of residue
classes in a fixed lattice.  Apply Poisson on that lattice to
\[
 f_{X,t}(z)=q_z^{it}W_i(q_z/X)
           =X^{it}f_t(z/\sqrt X),\qquad
 f_t(z)=q_z^{it}W_i(q_z).
\]
The Fourier transform is
\(X^{1+it}\widehat f_t(\sqrt X\,\xi)\).
Because \(f_t\) is supported on a fixed annulus, integration by
parts for any fixed \(J_0>2\) gives
\[
 |\widehat f_t(\xi)|
 \ll_{J_0} (1+|t|)^{J_0}(1+|\xi|)^{-J_0},
\]
using finitely many seminorms of \(W_i\).  For \(X\ge1\), the
sum over nonzero points of the fixed dual lattice is therefore
\[
 \ll (1+|t|)^{J_0} X
       \sum_{\xi\ne0}(1+\sqrt X|\xi|)^{-J_0}
 \ll (1+|t|)^{J_0}
\]
after increasing \(J_0\) if necessary.  For \(0<X<1\), both the
lattice sum and its zero-frequency main term are \(O(1)\) by
absolute counting on the fixed annulus.  The zero frequency is
\(c_\vartheta X^{1+it}I_i(t)\), where \(c_\vartheta\) is the
fixed weighted mean of the residue classes.  This proves the
unmasked version of Equation~\eqref{old-eq:2.18d}.

Now insert the mask by inclusion--exclusion.  First remove from
\(\mathfrak R_*\) its primes in \(\mathcal S\), since every
summation ideal already avoids them.  Thus all divisors used below
are outside \(\mathcal S\).  With \(\mathcal L_i^0\) denoting
the unmasked sum, multiplicativity gives
\[
 \mathcal L_i(X,t)
 =\sum_{d\mid\mathfrak R_*}
       \mu(d)\vartheta(d)q_d^{it}\mathcal L_i^0(X/q_d,t).
\]
The main coefficient is
\[
 c_{\vartheta,\mathfrak R_*}
 =c_\vartheta\sum_{d\mid\mathfrak R_*}
             \frac{\mu(d)\vartheta(d)}{q_d}.
\]
The \(q_d^{it}\) from extraction has canceled the
\(q_d^{-it}\) in the main term at \(X/q_d\).  This proves that
the coefficient is independent of both \(X\) and \(t\).
The errors are multiplied by at most the number of divisors
\(\#\{d:d\mid\mathfrak R_*\}\ll_{\epsilon_1,B_*}Z^{\epsilon_1}\).
The coefficient itself is also \(Z^{\epsilon_1}\)-bounded by
the polynomial-size Euler-product estimate.  This proves
Equation~\eqref{old-eq:2.18d}.

For Equation~\eqref{old-eq:2.18e}, take absolute values in the
original sum.  It is empty below a fixed positive scale; whenever
it is nonempty, ideal counting on the fixed annulus gives \(O(X)\)
points.  This is uniform in \(t\) and in the mask.

The masked sum in Equation~\eqref{old-eq:2.18f} is exactly
\[
 \mathcal L_1(U_1,t)\mathcal L_2(U_2,t)
 -\mathcal L_1(V_1,t)\mathcal L_2(V_2,t).
\]
The two product main terms in
Equation~\eqref{old-eq:2.18d} are both
\(c_{\vartheta,\mathfrak R_*}^2T^{1+it}I_1(t)I_2(t)\).
They cancel.  Each cross term with one error is bounded by
\(Z^{\epsilon_1}(1+|t|)^{2J}\) times one of
\(U_1,U_2,V_1,V_2\), and the product of errors has the same
bound after reducing the subsidiary power loss.  If all four
scales are at least \(Z^r\), each is at most \(TZ^{-r}\),
and \(1\le TZ^{-r}\) because \(T\ge Z^{2r}\).
The difference is therefore
\(O(Z^{\epsilon_1}(1+|t|)^{2J}TZ^{-r})\).
Division by \(T^{1/2}\) proves
Equation~\eqref{old-eq:2.18f}.  When no nonnegative lower length
is available, Equation~\eqref{old-eq:2.18e} bounds each product
by \(O(T)\), giving the last assertion.
\end{proof}

Apply this lemma to each already allocated exceptional child.  The
hypotheses on its coefficient and common mask follow from
Lemma~\ref{lem:centered-coefficient-invariant}; in particular the
main terms cancel for the same row, mask, and whole-product norm
power before any absolute value.  Initially all four plain lengths
are at least \(L\).  Before the selected \(\mathfrak t\)-factors
are removed, the \(c_2\) side has lower plain length at least
\(L-c-w-c_2\) and nominal total length \(\alpha_1=a_0-c_2\):
no one plain loses more than the total \(c+w+c_2\) extracted by
the two genuine common supports and the amplifier.  The other side
has the analogous bounds with \(d_2\).  Put
\[
 r=(L-c-w-\min(c_2,d_2))_+,
\]
and call the side with this reference saving the first side, relabeling
its associated raw data together.  If \(\widetilde r_1<r\), each
selected plain length satisfies \(a_{1,i}\le\widetilde r_1\), so all
four formal post plain scales have logarithmic length at least
\(r-a_{1,i}\ge r-\widetilde r_1>0\).  Apply
Equation~\eqref{old-eq:2.18f} on these exact equal-product formal scales.
If \(\widetilde r_1\ge r\), including equality, or if \(r=0\), use
the all-scale absolute-volume fallback in the same lemma.  No formal
centered scale is clipped in this branch, and an identically zero
rectangle remains in the formal difference.  The resulting saving is
\((r-\widetilde r_1)_+\).  With absolute volume on the other side,
this gives the aggregate bound on each lower-endpoint divisor dyad,
\begin{equation}
 \begin{split}
 &\sum_{\mathfrak t,\mathbf J}
 |c_{\mathfrak t,\mathbf J}(h')|
 |P_{C,\mathfrak t,\mathbf J}(h')
       P_{D,\mathfrak t,\mathbf J}(h')|\\
 &\hspace{15mm}\ll
 Z^{a_0-b_2-r+2\theta_N+\epsilon_1}.
 \end{split}
 \label{old-eq:2.18}
\end{equation}
Here and below a fixed polynomial in the separated heights is
understood.  To check the bound, the raw coefficients and volumes give
the reference exponent \(a_0-b_2+e_1+e_2\) and reduction
\(-\widetilde r_1-\widetilde r_2\), while the divisor count gives
\(t_-\).  Equation~\eqref{eq:centered-divisor-boundary} gives the
one-line inequality
\begin{equation}
 \begin{split}
 t_- -\widetilde r_1-\widetilde r_2-(r-\widetilde r_1)_+
 &=(t_- -\widetilde r_2)-\max(r,\widetilde r_1)\\
 &\le\omega_2-r.
 \end{split}
 \label{old-eq:2.18h}
\end{equation}
Adding \(e_1+e_2\) leaves at most
\(e_1+(e_2+\omega_2)-r\le2\theta_N-r\).
The divisor-bounded allocations contribute only \(\epsilon_1\),
proving Equation~\eqref{old-eq:2.18}.
This applies to the artificial terms even when quotients retain
\(\mathfrak t\)-primes, because the coefficient lemma preserves the
common mask, equal product scales, and one norm power within each
rectangle difference.

The live slots are summed only after the finite transform has removed
their Gauss coefficients.  Their unnormalized absolute sums contribute
\(O_N(\prod_{i\ {\rm live}}P_i)\).  Combining this with the nominal
raw normalizer gives the full \(e_j\) volume factor in
Equation~\eqref{eq:centered-raw-normalization}, as included above; it
is not replaced by a factor-product normalizer.  The Fourier measure was fixed before
their labels, so conditioning on those labels changes neither the
common mask nor the centered saving.  A zero conditional slot scalar
is discarded only in the present absolute bound.

Set \(v=c+w+\min(c_2,d_2)\).  By
Equation~\eqref{old-eq:2.17}, \(F_1+F_2\) is at least
\(2v/3-3\sigma-5\delta_{{\rm fr},1}/6\), since
\(b_2\ge\min(c_2,d_2)\).
The centered saving in Equation~\eqref{old-eq:2.18} now bounds
the exceptional deficit, apart from \(3\sigma\) and the recorded
frequency and aggregate support perturbations, by
\begin{equation}
 A-\tfrac56M-\tfrac23v-(L-v)_+
 \le (A-M)_+.
 \label{old-eq:2.19}
\end{equation}
Indeed, for \(0\le v\le L\) the left side is
\(A-5M/6-L+v/3\), which increases with \(v\); for
\(v\ge L\) it is \(A-5M/6-2v/3\), which decreases.
Its maximum is attained at \(v=L=M/4\) and equals \(A-M\).
For positive slots, Equation~\eqref{old-eq:3.9} implies
\(A\le M\).  For the zero-slot core,
Equation~\eqref{old-eq:2.1h} gives \(A-M\le\delta\).
Thus the centered exceptional terms require only the terminal
loss \(3\sigma\), or \(\delta+3\sigma\) in the padded core.
The additional \(\delta_{{\rm fr},1}\), \(\delta_{{\rm fr},2}\),
and \(\theta_N\) terms are included in the same \(C_*\xi\)
stage allowance.
Together with Equation~\eqref{old-eq:2.12}, this completes the
centered stage at the current width.

\subsection{Completion of the finite induction}

We finish by verifying the quantifier order and the finite induction.
Every nonterminal call is to a child with nominal width \(M'\) in
Equation~\eqref{old-eq:2.13}, at least \(\sigma\) below its parent.
The declared row width is \(M'_{\rm act}\), and all nonexceptional
moment-input lengths are taken after the specified support and clipping
operations.  Exceptional raw scales remain formal and are not clipped.
The two frequency enclosures, the aggregate support errors, and these
clippings alter the width comparison by at most \(C_*\xi\).  Choose
this below \(\sigma/2\).  Every nonempty child then has nonnegative
width and leaves its band of length \(\sigma/4\).  An empty
nonzero-frequency range is discarded; a nonempty formal range just
below scale one has already been included by its enclosing length and
clipping error.

At each smaller width, Equations
\eqref{old-eq:2.1a}--\eqref{old-eq:2.1b} delete the fixed extra masks
before the natural positive-slot estimate.  The unrestricted zero-slot
assertion follows once from the padded core by reflecting at most its
two plains.  A centered comparison calls only the earlier uncentered
stage in its band, using its strict margin.  Each Gauss norm has at
most one amplification, whose errors go directly to the second
transform.  Equation~\eqref{old-eq:3.16} is invoked once for each
actual natural child, after all its boundary defects have been
included in \(F_{\rm act}\).  Thus there is no same-band cycle.
The integer \(D\) defined above bounds the strict calls by \(D-2\),
because each drops the width by at least \(\sigma/2\) from an initial
width at most \(M_{\max}\).

Here is why \(C_*\) and the mesh can be chosen independently of the
fixed slot count.  Equation~\eqref{eq:centered-aggregate-support}
bounds the entire logarithmic error of any slot subset by
\(\theta_N=H_N/\log Z\).  Each extracted plain power is measured
exactly, and each of the two possible plain clippings in one
nonexceptional rectangle, after triangle inequality, has one fixed
endpoint error.  The exceptional calculation uses raw reductions and
the disjoint-subset bounds without clipping.  The first and second frequency enclosures use only
their one aggregate outer scale and one common row dyad.  All the
local ledgers are affine or positive parts of affine expressions in
the fixed list of aggregate lengths displayed above.  Their numerical
Lipschitz constants do not depend on \(N\); the local \(F_2\)
inequalities and the radical counts use exact prime norms.  In
particular the change in the positive-slot affine expression is at most
\begin{equation}
 |\Delta M|+|\Delta A|+5|\Delta z|,
 \label{old-eq:2.1j}
\end{equation}
because \(0\le6\kappa-1\le5\).  The total change of \(z\) is
aggregated before this inequality is used.  These observations give
one numerical \(C_*\), enlarged for the fixed number of operations
per stage, that covers all frequency, normalization, and boundary
errors by \(C_*\xi\).  The single greedy prefix contributes at
most \(\eta\), not one \(\eta\) per removed slot.

The analytic separations introduce no derivative-order multiple of
\(\xi\).  On a fixed full-product logarithmic box a radial kernel is
\(\widehat\Phi_i(R_{\rm sc}\exp(L(\mathbf x)))\), where \(L\) is a fixed
linear form in the row and whole-column log norms.  Every derivative
is a fixed combination of Euler derivatives of \(\widehat\Phi_i\).
Lemma~\ref{lem:kernel-seminorms} bounds these uniformly for all
\(R_{\rm sc}>0\), with any prescribed radial decay and derivative orders.
One does not use the weaker bound \(R_{\rm sc}^J\) for
\(R_{\rm sc}\le Z^\xi\).
After its displayed central power has been extracted, an inverse root
is \((q_u/Z^{a_0})^{-1/2}\) on the fixed box; its derivatives cost
only fixed constants depending on \(H_N\) and their order.
the smooth calculus of Appendix~\ref{app:smooth} uses the full weighted Fourier
integral, not a supremum over heights below \(Z^\xi\).  A required
height degree \(J\), which may depend on \(N\), raises an input
seminorm order and a fixed constant, not a \(Z^{J\xi}\) exponent.
Reflection similarly acts on at most two plains with fixed gamma
shape; larger derivative orders raise finite seminorm and height
orders, while the upper-annulus subshare \(\xi/2\) remains fixed.
The fixed polynomial height factor in Equation~\eqref{old-eq:3.5}
is integrated by these weighted Fourier measures at its required finite
order, without changing the previously chosen exponent of \(Z\).

The discrete label counts are also used only once.  After extracting
the displayed exponent for a family of common supports, divisors, or
allocations, divide its absolute weighted counting measure by that
total mass before applying a separated tuple seminorm.  This leaves a
measure of mass at most one; its already extracted exponent is not
counted again in the smooth norm.  Each Fourier measure is common to
the live labels by the whole-product construction.

\paragraph{Choices before specifying the slot count.}
Choose the parameters in the following order.  First, from the bounded
real ranges and \(\epsilon\), choose \(\rho,\sigma,\delta>0\), with
\(\delta\) small compared to \(\rho,\sigma\), so that
\[
 T_{\rm term}:=\rho+\delta+\rho/6+3\sigma+5\sigma/3<\epsilon/4.
\]
This bounds a terminal width-floor, diagonal, or exceptional loss;
the displayed frequency and numerical support corrections belong instead to the per-stage
\(\xi\) allowance.  With \(D\) and \(C_*\ge1\) as above, choose,
still before specifying \(N\) or its fixed profiles,
\begin{equation}
 \begin{aligned}
 \xi&<\min\left\{\frac\delta2,\frac\rho{30},
                  \frac{\sigma}{4C_*},
                  \frac{\epsilon}{16C_*D}\right\},\\
 \eta&<\min\left\{\frac\sigma6,
                   \frac{\epsilon}{16C_*D}\right\},\qquad
 \epsilon_0<\frac{\epsilon}{16C_*D}.
 \end{aligned}
 \label{old-eq:2.1i}
\end{equation}
These choices make the strict drop at least \(\sigma/2\), preserve
the comparison margins and padded core, and are uniform for
\(\kappa\in[13/18,1]\).

\paragraph{Choices for a fixed slot system.}
Now fix any \(Z\)-independent \(N\), arithmetic data, and profile
windows.  Form \(H_N\) and the finite full support boxes through
depth \(D\).  Their fixed factors \(\exp(O(H_N))\) belong to the
constants.  After imposing an eventual threshold
\(\log Z\ge C H_N/\xi\), the threshold in
Equation~\eqref{eq:centered-reflection-length}, and
Equation~\eqref{eq:centered-localization-threshold}, all occurring
column, radical, and mask norms have bounded logarithmic ranges
independent of moving labels.  In those ranges choose every
arbitrarily small power estimate with local shares whose sum in a
stage is at most \(\epsilon_0\).  These shares may depend on \(N\).
For example, the two divisor masses and the at most \(N\) frozen-slot
masses in mask deletion may each use a share
\(\epsilon_0/(10(N+2))\).  In the prime estimate, the principal
squared exponent is exactly \(\kappa z\), and the contour displacement
cost is \(2e z\), with bounded total \(z\), not \(N\eta\).  Choose
\(e\) for its assigned share; the fixed character expansions and
powers of \(\log Z\) are absorbed after an \(N\)-dependent threshold.
The same reasoning applies as individual \(z_i\) approach zero,
since \(P_i\ge1\).  For divisor allocations use the global bounds
\(\tau_{N+2}(v)^C\ll_{N,C,a}q_v^a\) and
\(C_N^{\omega(v)}\ll_{N,a}q_v^a\) with a sufficiently small
\(a>0\), not a fixed loss at every prime.  Rankin's fixed-radical
bound is treated with the same local shares.  The one normalized pool
average has only its single density loss in a stage.

The induction is simultaneous for every surviving subset of these
original \(N\) slots, with the same mesh.  To make this precise, for
\(0\le d\le D-2\), let
\[
 \mathcal E_d
 :=T_{\rm term}+(d+1)C_*(\eta+\xi+\epsilon_0)
\]
be the permitted error when at most \(d\) strict calls remain.
A terminal group, including its bounded reflection or comparison
preprocessing, has error at most \(\mathcal E_0\).  A strict edge
and its fixed number of same-band operations use the child envelope
\(\mathcal E_{d-1}\) plus at most
\(C_*(\eta+\xi+\epsilon_0)\), giving \(\mathcal E_d\).
The Gauss allowance \(\epsilon_G\) in
Equation~\eqref{old-eq:2.5} denotes this appropriate envelope and
the current local shares; it is not a request to reapply the moment
lemma with an \(N\)-dependent loss and a new mesh.  Triangle
inequalities take the maximum error up to the already counted mass,
and Cauchy--Schwarz averages the errors of the two children.
Consequently one terminal loss occurs along a branch, not at every
ancestor.  The largest path error is at most
\(T_{\rm term}+C_*D(\eta+\xi+\epsilon_0)<\epsilon\) by
Equation~\eqref{old-eq:2.1i}.

Finally choose the required finite kernel, reflection, prime, and
terminal seminorm and polynomial-height orders backwards through
these \(D\) stages.  The full weighted Fourier estimates and uniform
Euler-kernel estimates above make every such choice finite.  Choose
one final lower threshold for the aggregate support inequalities,
pool separation and density, the raw frequency tails, and all fixed
logarithmic losses.  The orders, constants, and threshold may grow
with \(N\) and the fixed data, but \(\eta,\xi,D,C_*\) in the
exponent comparison do not.

Each new extra mask is supported on frozen columns or a frozen
amplifier prime, and its logarithmic norm increases by a bounded
total length at a stage.  It is fixed within every child row sum.
The only row-dependent zeros there are the natural zeros of the row;
the declared moving zeros remain in the fixed twist and remain counted.
Over depth \(D\) the masks remain of polynomial norm.
On an exceptional row, the redundant part of the full moving union
and the row radical, together with these extra masks, forms the
common polynomial-size \(\mathfrak R_*\); the lattice cancellation
is uniform for it.  The fixed-numerator ray lemma controls the
finite choices at \(\mathcal S\), without absorbing a moving good
prime into a fixed modulus.

When this estimate and the inverse moment are applied together, all internal
orders in both arguments, including the reflection-kernel orders,
are fixed first.  The later external Fourier-tail order in
the smooth calculus of Appendix~\ref{app:smooth} lies outside this internal
propagation: it raises only the external input seminorm and does
not change a previously fixed internal height order.  This proves
the asserted finite-order uniformity.  The finite induction proves
the natural assertion, and the initial mask deletion proves
Lemma~\ref{lem:plain} in its full stated form.

For the later physical application, \(q=0\) and the row \(u\)
is sixth-power-free.  If a prime outside \(\mathcal S\) has
valuation \(1,\ldots,5\) in \(u\), its local character has order
\(6/\gcd(6,v_p(u))>1\), so the inducing character cannot belong
to \(\Theta\).  The exceptional physical rows are therefore
supported on \(\mathcal S\); sixth-power-freeness makes this a
finite set, up to the finite unit group.  With abstract moving
twists, additional exceptional rows can arise by cancellation.
They were included in Equation~\eqref{eq:exceptional-row-count}
and the fixed-character volume argument.

\paragraph{Endpoint verification.}
At \(\kappa=13/18\) the inequality \(6\kappa-1\ge10/3\)
still gives the strict \(z<M/20\) whenever \(A>5M/6\).
The two comparison margins and the finite-width induction are therefore
valid at this endpoint itself.  All other uses of \(\kappa\) are
\(0<\kappa\le1\), the prime estimate, the affine premise, or
\(\kappa d_z\le F_{\rm act}/6+\eta\).
There is no limiting argument in a strict moment constant.
This completes Proposition~\ref{prop:arithmetic}.

\section{Source map, exact checks, and scope}\label{app:checks}
The pinned source is \cite{QRH}, revision
\nolinkurl{fd4aeeb2ee4fc729c18d98444fed42fd0529eeeb}.
The arithmetic construction and the longer inverse and plain
inductions are adapted from it, with attribution, rather than
claimed as new constructions.  The correspondence is:
\begin{itemize}
\item Sections 6, 7, 12, and 16 of that source give the original
physical probe and its local compensation.  Section~\ref{sec:common-probe}
here includes the defining sum, fixed data, both exact identities,
the local coefficient calculation and selected factor.
\item Source Sections 4.4 and 8 supply the growth argument and
detector.  Section~\ref{sec:detector-complete} here proves their
buffered and saturated forms at the full range \(B\le1\).
\item Source Proposition 16.1 is the dynamic local calculation.
Section~\ref{sec:dynamic-errors} here includes the full local table
and proves the explicit aggregate \(12e\) loss with no per-slot
real exponent loss.
\item Source Lemma 14.3 and Proposition 15.2 correspond to
Appendix~\ref{app:energy}.  The completed reflection used there
is derived from the unconditional cusp formulas of \cite{DR} in
Appendix~\ref{app:reflection}.
\item Source Section 17, physical lines 9218--12342, is the
inverse-moment construction.  Appendix~\ref{app:inverse} here
contains the first raw transformed sum
\eqref{own:inverse-first-raw}, both complete separated sums
\eqref{eq:inverse-first-separated} and
\eqref{eq:inverse-xi-identity}, their positive majorants, and the
initializing identity \eqref{eq:initial-separated-identity}.
The support projections are taken before any independent positive
child enlargement.
\item Source Section 18, physical lines 12477--14984, is the
plain induction.  Appendix~\ref{app:plain} here includes its
complete transformations, with the allocated coefficients defined
by \eqref{own:plain-allocation}.  Its endpoint is extended from
\(3/4\) to \(13/18\) by the displayed centered-comparison proof.
\end{itemize}
These references identify provenance.  None of these essential
arguments is delegated to an unavailable earlier manuscript.

\subsection{Input receipts and source provenance}
The three read-only inputs actually present in this workspace were
read in full and hashed locally:
\begin{itemize}
\item \texttt{answer.tex}: 478018 bytes, SHA256
\nolinkurl{701de0a583717adfc46ca7e9e0156e1512ad78c4c52f82bf3c39b29094167fba}.
\item \texttt{research\_notes.tex}: 253762 bytes, SHA256
\nolinkurl{5b9cd16e44e360c1e88109868b0344e62d37cb2dd491a6c68ae1789574814210}.
\item \texttt{references.bib}: 11730 bytes, SHA256
\nolinkurl{06fd736687606b9661f2316ebb356ca5b9799564ab9b923b4e628f1dc23b0f3b}.
\end{itemize}
The predecessor retains earlier supplied-archive receipts.  They are
preserved below as historical byte checks, not remeasured in this
workspace, where the reference archive itself is absent.
The seed answer has 141293 bytes and recorded SHA256
\nolinkurl{73b570cd6bbb93c7667fdff8253f9c011f3f60bd0a76794d141fcfc15211c6ab}.
The four recorded upstream pins are:
\begin{itemize}
\item Raw seven-eighths TeX: 766316 bytes, 16677 physical lines,
\nolinkurl{42a5ee0febca59fd1def55cfd6c6808c322ef4237d7726303ea52f711deac6a3}.
\item Seven-eighths PDF: 1588969 bytes,
\nolinkurl{8fe93046f8cf5ef1ba5969c89addc02d76311adc4ee907509ff9cd96f7ec99e7}.
\item Eleven-twelfths companion PDF: 784503 bytes,
\nolinkurl{f919b57829b178c8e60e7c17b018cf773e7907cf642ef5a3347d8a826e8dbf18}.
\item Optional Landau--Siegel PDF: 404812 bytes,
\nolinkurl{fab9f6e84ece2a37d834b2c031d8eb14e81413d9f6270791e7b5468b295214ef}.
\end{itemize}
The pinned raw TeX was also consulted through its revision-specific
public URL on 9 October 2026.  A browser-rendered text check is not
a fresh byte-hash measurement.  Byte authentication does not validate
mathematics.  Neither companion zero-free conclusion, either archived
workflow report, nor a claimed formalization is an input to this proof.

\subsection{Reproducing the complete document}
The accompanying research notes include a version-tied source capsule
and a short extractor for an environment in which only that notes file
is mounted.  The capsule contains this entire \texttt{answer.tex},
the bibliography, and all four certificate scripts, with an explicit
member list and source hashes.  It contains no external proof premise.
The included \texttt{verify\_submission.py} runs the four scripts and
then runs \texttt{pdflatex} twice on the complete document, checks
the contract, and records the PDF page count and the final log.
This is a way to perform an actual full-document build, not a
compilation of a preamble or an isolated certificate fragment.
A successful build or a matching hash does not certify the analytic
proof.

\subsection{Checks of the published input hypotheses}
The external inputs were compared with the following precise versions
and locations on 9 October 2026:
\begin{itemize}
\item Goldmakher--Louvel, arXiv v2, Definition~1 and Theorem~1.1,
pp.\ 1--2.  The explicit primitive quadratic family, its unit
normalization, and its common-class conductor property are verified
before \eqref{eq:quad-sieve-proof}; none is an extra sieve assumption.
\item Heath-Brown's author manuscript, Theorem~2, pp.\ 3--4,
and its restatement in \cite[Theorem 9.1]{DR}.
The original author-uploaded text was inspected when the ORA PDF endpoint
was unavailable to the browser.  The indexing is squarefree
Eisenstein elements, not squarefree rational norms.
\item Dunn--Radziwi{\l}{\l}, arXiv v3 dated 14 May 2024,
pp.\ 16--19 and 78.  The three cusp representatives and the
conjugation \(\overline{d_j(-\mu)}\) match
\eqref{eq:theta-cusp}.  Infinity coefficients, the two auxiliary
cusp supports, and their sizes give
\eqref{eq:theta-support}; the full-frequency extension is proved
here, not inferred from a unit-frequency reflection theorem.
\item Poonen, Sections~4.7 and 5.11, Lemma~4.20 and Theorem~5.22.
The factor \(2\) in \(\Gamma_{\mathbb C}\) cancels from the
functional equation, yielding exactly the completion used in
Lemma~\ref{lem:hecke-strip-growth}.
\item Thorner--Zaman, arXiv v3, Theorem~1.1, p.\ 2.
For the fixed ray extension, the possible term
\(\operatorname{Li}(x^{\beta_1})\), with fixed \(\beta_1<1\),
is \(o(\operatorname{Li}(x))\).  This suffices for the slot
asymptotic; no uniform moving-conductor prime asymptotic is invoked.
\item Milne, version~4.03, VIII~(5.3), (5.5), p.\ 245,
for the fixed-numerator ray character; NIST DLMF
(25.15.4), for the primitive/imprimitive identity.
\end{itemize}
All essential changes in parameter ranges are justified locally in
this document.  These checks identify the imported standard formulas,
not an independent computer verification of the complete analytic proof.

\subsection{Certificate scope and ledger}
The rational program in Appendix~\ref{app:certificate-code} has
been executed with exact arithmetic.  It certifies the displayed
geometry and continuous exponent inequalities, not an analytic
moment theorem.  The latter proofs are in this document.
In addition to coefficient inequalities, the script checks every
full dual right side, the displayed continuous interval bounds, and
the root-uniqueness signs.  Its normalized UTF-8 source, with one
terminal newline, has SHA256
\nolinkurl{af406d0f43d651838a65c39da9eb69c3837fba72f2ddad7c15adb4dc9bc7711f}.
An additional exact check of the affine local errors over
\(51/100\le a\le1,\ 0\le e\le1/1000\) gives maxima
for the off-row exponents
\[
 (-51/100,-507/250,-59/100,-257/250),
\]
and for the ramified boundary exponents
\[
 (-497/500,-13/25,-51/100,-53/100,-13/25,-11/20).
\]
Since these functions are affine, the four rational vertices
certify the entire rectangle; the geometric tails and simultaneous
conductor allocation have the separate analytic proofs above.
The additional local script also verifies the five rational identities
for the quotient-free factors before numerical exponent comparisons.
Its normalized source has SHA256
\nolinkurl{77d9833579f17d533ef6c19d3dc9fbc68243e24b5383ebd158e656c39086fe45}.
Both scripts were executed successfully with exact arithmetic.

For a separate finite arithmetic sanity check, unnormalized Gauss
identities and the full finite reflection sums were evaluated exactly
in cyclotomic rings at norms \(7,13,19,25\).
All \(20\) Gauss identities and \(768\) full-frequency identities
passed, including the zero-frequency and zero-extended principal
branches and the inert prime of norm \(25\).  The complete script
and output are retained in the research record.  These finite tests
do not replace the all-prime proofs in
Appendices~\ref{app:conventions}--\ref{app:reflection}.

A separate coefficient-transport script checks the proved finite-ray
quadratic refinement, both relative Gauss phases, and every value in the
zero-extended sixth-root semigroup for
\eqref{eq:inverse-moving-characters}.  It also checks the child
squarefreeness and mask supports and the full factor-allocation identity.
There are \(117649=7^6\) zero-pattern tests of the moving-character
identity, not samples of primes.  The all-modulus CRT proof is
\eqref{eq:inverse-first-arithmetic-identity}; the preservation of
centered coefficients is proved in
Lemma~\ref{lem:whole-product-transport}.
The script \texttt{moment\_transport\_certificate.py} and its output
are preserved in the research record.  Its normalized source hash is
\nolinkurl{06f0b04e5c41c2451de6b3ff5822ca3d30fb1f34b319de3b7efb2119defd5c1f}.
These tests certify finite algebra, not the analytic norm estimates;
those still require the displayed signed identities, positive
fibre bounds, strict decrease, and finite-order tail arguments.

The ledger distinguishes the following objects.
\begin{itemize}
\item The theorem's explicit uniform rational and its improvement are
\[
 \theta=\frac{874957019421}{10^{12}},\qquad
 \frac{3499999}{4000000}-\theta=\frac{42730579}{10^{12}}.
\]
\item The algebraic obstruction
\(0.8749570194200989\ldots\) is a lower bound only for the
restricted exponent program of Proposition~\ref{prop:restricted-barrier}.
It is not the metadata rational.
\item A preliminary finite grid suggested
\(0.874956994782\ldots\).  An off-grid high exponent was positive
by more than \(10^{-6}\), so this candidate was rejected.
\item Smaller values obtained by replacing the inverse
amplification slope by \(2B-1\) remain heuristic.  The current
sixth-power amplification does not give that replacement.
\end{itemize}
All starting scales and constants may depend on the target.
Only the zero-free boundary and the positive comparison savings
are common to the family; that is what
Lemma~\ref{lem:criterion} requires.  No height cutoff appears
in the conclusion.

\section{Reproducible exact certificate}\label{app:certificate-code}
The following script uses only rational arithmetic and symbolic
polynomial identities.  It verifies all geometry checks, every
coefficient inequality in the continuous certificate, the tail
and endpoint margins, the six low-energy dual vectors, and the
root isolation used for the restricted optimization bound.
It does not purport to verify an analytic theorem by computation.
It is also supplied as \texttt{exact\_certificate.py}.
For the recorded digest, strip trailing blank lines from the following
block and append one terminal newline.
\begin{verbatim}
"""Exact stage-1 certificates. Python 3 + SymPy. No network or floats.
The analytic statements are proved in answer.tex, not by this script.
"""
import sympy as s
from math import comb
Q = s.Rational
delta, x, y, L, b, c = s.symbols("delta x y L b c")
theta = Q(874957019421, 10**12)
incumbent = Q(3499999, 4000000)
ell = Q(11, 3) - 4*theta
kappa = 2*theta - 1
bv = -(3*ell+1)*(36*ell**2-75*ell+19)
bv /= 3*(114*ell**2-159*ell-7)
M = 1-ell
lx, ly = (M-bv)/2, (M+bv)/2
h = 1-lx+ell
z0, zeta, dmin = Q(17,50), Q(1,10**14), Q(1,100)
assert Q(1,2) < theta < incumbent
assert Q(1,6) < ell < Q(167,1000) < Q(1,5)
assert Q(123,1000) < bv < Q(124,1000)
assert M-2*ell > Q(49,100)
assert lx-ell > Q(9,50)
assert ly-ell-Q(11,6)*bv > Q(2,25)
assert Q(35,100) < lx < Q(36,100)
assert Q(47,100) < ly < Q(49,100)
assert Q(81,100) < h < Q(82,100)
assert ell/(h+zeta) > Q(1,5)
assert Q(43,50) < theta
assert Q(13,18) < Q(37,50) < kappa < Q(3,4)
assert theta == 1-h/6+bv/12
assert lx/2-1+h/6 == -Q(2,3)-bv/6
assert Q(30,169) < Q(1,4)
alpha = Q(5,6)
D = 3-(1+2*c)*x
P = (2-2*c*x)*(1-x)
J = (alpha-delta)*D+delta*P
R = 1-delta+(alpha-delta)*delta*P/(2*J)
hg, lyg = (1+3*L+b)/2, (1-L+b)/2
thg = Q(11,12)-L/4
a = (1+delta)/2
Eh = a-thg-hg/6-a*lyg-L/2+x*delta*L
Eh += hg*(R+delta/2)
F = s.Poly(s.cancel(-J*Eh), x, delta).as_expr()
F = s.expand(F.subs(x,Q(1,2)-y))
assert s.degree(F, delta) <= 2 and s.degree(F, y) <= 3
bL = -(3*L+1)*(36*L**2-75*L+19)
bL /= 3*(114*L**2-159*L-7)
cL = 2/(5-3*L)
dv = (5-9*L)/(6*(3*L+1))
pol = 657*L**3-954*L**2+21*L+20
Av = (3*L+1)*(414*L**3-1191*L**2+878*L-29)
Av /= 4*(3*L-5)*(114*L**2-159*L-7)
Dv = -(153*L**2-201*L-20)*pol
Dv /= 36*(3*L-5)*(3*L+1)*(114*L**2-159*L-7)
Fzero = F.subs({y:0,b:bL,c:cL})
assert s.cancel(Fzero-Av*(delta-dv)**2-Dv) == 0
assert pol.subs(L,ell) == Q(
    40963882847150664198386369,46875*10**30)
assert pol.subs(L,ell) > Q(8,10**10)
assert Av.subs(L,ell) > Q(1,5)
assert (Dv/pol).subs(L,ell) > Q(1,160)
assert Dv.subs(L,ell) > Q(5,10**12)
# Continuous interval bounds, not just the fixed rational ell.
lo, hi = Q(1,6), Q(167,1000)
numD = -153*hi**2+201*lo+20
denD = 36*(5-3*lo)*(3*hi+1)*(-114*lo**2+159*hi+7)
numA = (3*lo+1)*(414*lo**3-1191*hi**2+878*lo-29)
denA = 4*(5-3*lo)*(-114*lo**2+159*hi+7)
assert 5-3*hi > 0 and -114*hi**2+159*lo+7 > 0
assert numD == Q(49232983,10**6)
assert denD == Q(3694400793,500000)
assert numD > 0 and numD/denD > Q(1,160)
assert numA > 0 and numA/denA > Q(1,5)
Fa = s.expand(F.subs({L:ell,b:bv,c:1/(3*kappa)}))
u = s.symbols("u")
bounds_shift = [(3,112,605),(41,269,434),(16,98,148)]
for i, row in enumerate(bounds_shift, start=1):
    gi = s.Poly(Fa.coeff(y,i).subs(delta,Q(1,3)+u),u)
    for j, bound in enumerate(row):
        assert gi.coeff_monomial(u**j) >= Q(bound,1000)

def bernstein(lo, hi):
    t, v = s.symbols("t v")
    g = s.Poly(Fa.subs({delta:lo+(hi-lo)*t,y:v/2}),t,v)
    ans = []
    for i in range(3):
        row = []
        for j in range(4):
            val = sum(
                g.coeff_monomial(t**k*v**l)
                * Q(comb(i,k),comb(2,k))
                * Q(comb(j,l),comb(3,l))
                for k in range(i+1) for l in range(j+1))
            row.append(val)
        ans.append(row)
    return ans

matrices = [
 [[35772,41252,46731,52210],
  [12757,12166,11370,10371],
  [4549,4195,5696,10209]],
 [[4549,4195,5696,10209],
  [1813,1539,3804,10155],
  [723,1228,5211,14732]]]
for interval, lower in zip([(Q(0),Q(1,4)),
                             (Q(1,4),Q(1,3))],matrices):
    bb = bernstein(*interval)
    for i in range(3):
        for j in range(4):
            assert bb[i][j] >= Q(lower[i][j],10**6)
# J<=5/2 and all regions F>5/10^12 certify E<-2/10^12.
assert Q(723,10**6) > Q(5,10**12)
assert Q(5,10**12)/Q(5,2) == Q(2,10**12)
assert Q(2,10**12)-2*zeta > Q(1,10**12)

rr, dd = s.symbols("rr dd")
E = a-theta+h*(z0-Q(1,6))-a*ly-ell/2
E += x*delta*ell+dd*(rr+delta/2-z0)
floor = E.subs({dd:h,rr:1,x:Q(1,2),delta:Q(1,50)})
short = E.subs({dd:Q(1,2),rr:Q(76,75)-2*delta/3,x:Q(1,2)})
large = E.subs({dd:h,rr:1-delta,x:Q(1,2)})
assert floor < -Q(1,200)
assert s.diff(short,delta) > 0
assert short.subs(delta,alpha) < -Q(1,500)
assert s.diff(large,delta) < 0
assert large.subs(delta,alpha) < -Q(7,100)
assert h*(z0-Q(1,6))-ly/2+2*dmin < -Q(7,100)
eps, zinfty = Q(1,10**14), 4*10**14
B0 = lx/2+1+ly
tail = B0+(h+zeta)*(1+eps)-zeta*zinfty+eps
assert tail < -1

# Exact six-branch dual certificate for the relaxed low energy.
O,H,A,N,G,z,v,t,e = s.symbols("O H A N G z v t e")
vs = [O,H,A,N,G,z,v,t,e]
m0,l0 = s.symbols("m0 l0")
Am = s.Matrix([[1,1,0,0,0,0,0,0,0],
 [-1,0,2,0,0,0,0,0,0],[0,0,-1,1,0,0,0,0,0],
 [0,0,0,0,0,1,0,0,0],[0,-2,-2,1,3,-2,1,3,1]])
rhs = s.Matrix([m0,0,0,l0,-1-l0])
duals = [(Q(1,2),0,Q(1,2),Q(1,2),Q(1,4)),
 (1,0,0,0,Q(1,2)),(Q(1,2),0,Q(1,2),Q(1,6),Q(1,4)),
 (Q(1,2),0,0,1,Q(1,2)),(2,1,0,2,Q(3,2)),
 (Q(4,3),Q(2,3),0,2,Q(7,6))]
T = 2*H+2*A+2*z-1-l0-N-3*G
objs = []
for mx in [H,v+t]:
    for mn in [v,z,(v+z)/3]:
        objs.append(s.expand(O/2+mx-G+z-mn-t-2*e/3
                             -(T-v-3*t-e)/2))
dual_bounds = [m0/2+Q(1,4)+3*l0/4, m0,
    m0/2+Q(1,4)+5*l0/12, m0/2+l0,
    2*m0+l0-1, 4*(m0+l0)/3-Q(2,3)]
for obj, dvect, bound in zip(objs,duals,dual_bounds):
    row = s.Matrix([[obj.coeff(var) for var in vs]])
    dual = s.Matrix([dvect])
    constant = obj.subs(dict.fromkeys(vs,0))
    assert all(num >= 0 for num in dvect)
    assert all(num >= 0 for num in list(dual*Am-row))
    assert s.expand(obj-(row*s.Matrix(vs))[0]-constant) == 0
    assert s.expand((dual*rhs)[0]+constant-bound) == 0

mv,lv,th = s.symbols("mv lv th")
T1 = Q(5,6)+mv/12-lv/6
T2 = Q(23,24)-mv/6+5*lv/24
T3 = Q(1,3)+7*mv/12+lv/3
assert s.expand(Q(17,30)*T1+Q(2,5)*T2
                +Q(1,30)*T3-Q(13,15)) == 0
H0 = (1+delta)*(1-mv/2)+delta*lv
combo = (6*delta+Q(4,3))*(T1-th)+Q(2,3)*(T3-th)
lower = -(6*delta+2)*th+6*delta+Q(7,3)
assert s.expand(H0-lower+combo) == 0
Rcoupled = R.subs({x:Q(1,2),c:cL,delta:dv})
obstruction = pol/(2*(3*L+1)*(153*L**2-201*L-20))
assert s.cancel(Rcoupled-Q(2,3)-obstruction) == 0
# Root uniqueness on [1/6,1/5]: p' is decreasing and negative.
assert pol.subs(L,Q(1,6)) == Q(1,24)
assert pol.subs(L,Q(1,5)) == -Q(1088,125)
assert s.diff(pol,L).subs(L,Q(1,6)) == -Q(969,4)
assert s.diff(pol,L,2).subs(L,Q(1,5)) == -Q(5598,5)
d0 = (9*th-7)/(18*(1-th))
assert s.cancel(lower.subs({delta:d0})-th) == 0
theta_lo = Q(8749570194200989,10**16)
theta_hi = Q(8749570194200990,10**16)
assert pol.subs(L,Q(11,3)-4*theta_lo) < 0
assert pol.subs(L,Q(11,3)-4*theta_hi) > 0
assert theta > theta_hi and theta-theta_lo < Q(1,10**12)
# Compare with Liu, arXiv:2610.12234v1, without a decimal radical.
prior_lo, prior_hi = Q(5037,10**11), Q(5038,10**11)
for gap in (prior_lo, prior_hi):
    assert 1507-1653*(theta+gap) > 0
assert (1507-1653*(theta+prior_lo))**2 > 4*921
assert (1507-1653*(theta+prior_hi))**2 < 4*921
print("All exact assertions passed.")
print("theta =",theta)
print("improvement =",incumbent-theta)
print("ell =",ell)
print("reference kappa =",kappa)
print("continuous endpoint saving >",Q(2,10**12))
print("post-extension saving >",Q(1,10**12))
print("restricted threshold in",theta_lo,theta_hi)
print("prior-bound gap in",prior_lo,prior_hi)

\end{verbatim}
Its output is
\begin{verbatim}
All exact assertions passed.
theta = 874957019421/1000000000000
improvement = 42730579/1000000000000
ell = 125128941737/750000000000
reference kappa = 374957019421/500000000000
continuous endpoint saving > 1/500000000000
post-extension saving > 1/1000000000000
\end{verbatim}
The rational root endpoints printed by the script are exactly
those in \eqref{eq:root-isolation}.  It also prints the exact
comparison
\begin{verbatim}
prior-bound gap in 5037/100000000000 2519/50000000000
\end{verbatim}
for \(B_{\rm Liu}-\theta\).

\section{Remaining optimization questions}
There is no unresolved lemma being assumed in
Theorem~\ref{thm:main}.
Its rational half-plane, the unrestricted Hecke auxiliary statements,
and the transfer to every modulus are the proved claims.

The following questions concern stronger results, not gaps in that
theorem.  To pass below the restricted threshold of
Proposition~\ref{prop:restricted-barrier}, at least one bound in
its explicitly defined exponent program must change.
The simultaneous geometry variation and the six-branch energy
relaxation were optimized above; adding the same marked amplifier
was examined and does not improve its long-row envelope.
A stronger inverse moment (in particular, a justified replacement
for its \(5/6\) slope), a mixed witness/prime moment, or a sharper
arithmetic energy estimate would remove an actual restriction of
that program.  None follows from the old workflow approvals.
No claim that such improvements are impossible, or that this
rational is an absolute optimal zero-free constant, is made.
\begin{thebibliography}{99}

\bibitem{Seed}
\emph{An explicit uniform descent below seven eighths},
supplied 48-page research manuscript, dated 8 October 2026,
with its notes and two AI workflow reviews.
Read-only answer SHA256:
\nolinkurl{73b570cd6bbb93c7667fdff8253f9c011f3f60bd0a76794d141fcfc15211c6ab}.
The arithmetic construction and retained appendix proofs are
attributed to this seed and its cited sources; the reviews are
not axioms.

\bibitem{QRH}
OpenAI, \emph{The Quasi-Riemann Hypothesis: A Zero-Free Half-Plane
\(\Rea s>7/8\)}, preprint dated 30 September 2026.
\href{https://raw.githubusercontent.com/openai/math/fd4aeeb2ee4fc729c18d98444fed42fd0529eeeb/preprints/The-Quasi-Riemann-Hypothesis-September-30-2026/build/paper.tex}{Pinned source}
at the revision in Section~\ref{sec:result}.
The historical supplied-byte receipts and the present
revision-specific source check are distinguished in Appendix~\ref{app:checks}.
\bibitem{DLMF}
NIST Digital Library of Mathematical Functions,
\href{https://dlmf.nist.gov/25.15}{Section 25.15(i)},
especially (25.15.4).
\bibitem{Poonen}
B. Poonen, \emph{Tate's thesis}, MIT 18.786 notes (2015),
\href{https://math.mit.edu/~poonen/786/notes.pdf}{Course notes},
Sections 4.7 and 5.11, Lemma 4.20 and Theorem 5.22(b)--(c).
\bibitem{DR}
A. Dunn and M. Radziwi\l\l,
\emph{Bias in cubic Gauss sums: Patterson's conjecture},
Annals of Mathematics \textbf{200} (2024), 967--1057.
\href{https://arxiv.org/pdf/2109.07463v3}{Version 3},
Equations (5.1), (5.4), (5.7)--(5.8), (5.9), (5.12)--(5.16),
Appendix A rows 1, 10, 19; no GRH-conditional prime conclusion.
\bibitem{Exploration}
tomoto0, \emph{Independent verification of OpenAI's preprint},
\href{https://github.com/tomoto0/quasi-riemann-hypothesis-7-8-verification/blob/main/reports/improvement_exploration.md}{improvement exploration}
and
\href{https://github.com/tomoto0/quasi-riemann-hypothesis-7-8-verification/blob/main/reports/barrier_analysis.md}{barrier analysis}.
Revisited 9 October 2026; another revision and heuristic parameter
claims, not an imported proof of improvement.

\bibitem{Argonaut}
Argonaut Math,
\emph{Beyond the Seven-Eighths Barrier: A Sharper Quasi-Riemann Boundary},
\href{https://github.com/Argonaut-Math/argonaut-math-quasi-riemann-boundary/releases/tag/v0.1.8}{release v0.1.8},
8 October 2026; its TeX header is dated 7 October 2026.
\href{https://raw.githubusercontent.com/Argonaut-Math/argonaut-math-quasi-riemann-boundary/v0.1.8/preprints/Beyond_the_Seven_Eighths_Barrier.tex}{Release-tagged source}
consulted 9 October 2026 for attribution only.
Its stated threshold is \(3499999/4000000\); neither its proof nor
its formal-verification claims are used here.

\bibitem{Liu}
B. Liu, \emph{Slightly improved zero-free half-planes for the
quasi-Riemann hypothesis},
\href{https://arxiv.org/pdf/2610.12234v1}{arXiv:2610.12234v1},
8 October 2026, Theorem 21.1 and Sections 23--25.
Version-specific abstract and PDF inspected on 9 October 2026;
no fresh byte-hash receipt and no theorem imported.

\bibitem{HB}
D. R. Heath-Brown,
\emph{Kummer's conjecture for cubic Gauss sums},
Israel Journal of Mathematics \textbf{120} (2000), 97--124.
\href{https://ora.ox.ac.uk/objects/uuid:16c3966e-61cb-499d-839c-559b3794fcff/files/m502494cff282fae8a90f44a466f2b941}{Author manuscript},
Theorem 2, PDF pp.\ 3--4.
\bibitem{GL}
L. Goldmakher and B. Louvel,
\emph{A quadratic large sieve inequality over number fields},
Math.\ Proc.\ Cambridge Philos.\ Soc.\ \textbf{154} (2013), 193--212.
\href{https://arxiv.org/pdf/1112.1642v2}{Version 2},
Definition 1 and Theorem 1.1.
\bibitem{TZ}
J. Thorner and A. Zaman,
\emph{A unified and improved Chebotarev density theorem},
Algebra \& Number Theory \textbf{13} (2019), 1039--1068.
\href{https://arxiv.org/pdf/1803.02823v3}{Version 3},
Theorem 1.1.  Only the fixed-extension consequence is used.

\bibitem{MilneCFT}
J. S. Milne,
\emph{Class Field Theory}, version 4.03, 6 August 2020.
\href{https://www.jmilne.org/math/CourseNotes/CFT.pdf}{Course notes},
Chapter VIII, (5.3) and (5.5).
\end{thebibliography}
\end{document}
